\newcommand{\NWtarget}[2]{#2}
\newcommand{\NWlink}[2]{#2}
\newcommand{\NWtxtMacroDefBy}{Fragment defined by}
\newcommand{\NWtxtMacroRefIn}{Fragment referenced in}
\newcommand{\NWtxtMacroNoRef}{Fragment never referenced}
\newcommand{\NWtxtDefBy}{Defined by}
\newcommand{\NWtxtRefIn}{Referenced in}
\newcommand{\NWtxtNoRef}{Not referenced}
\newcommand{\NWtxtFileDefBy}{File defined by}
\newcommand{\NWtxtIdentsUsed}{Uses:}
\newcommand{\NWtxtIdentsNotUsed}{Never used}
\newcommand{\NWtxtIdentsDefed}{Defines:}
\newcommand{\NWsep}{${\diamond}$}
\newcommand{\NWnotglobal}{(not defined globally)}
\newcommand{\NWuseHyperlinks}{}
% kcc.weft — Master file for kcc: A Formal C17 Compiler in ARM64 Assembly
%
% This is a literate program in the tradition of Donald Knuth.
% The .weft source files are the ONLY source of truth.
% All assembly (.s) files are generated by tangling.
%
% Build:  just build        (tangle + assemble + link)
% Weave:  just doc          (generate LaTeX documentation)


\documentclass[11pt,a4paper,twoside,openright,DIV=15]{scrbook}

% Page layout: DIV=15 gives a wider/taller type area suited to
% code-heavy literate programs.  KOMA-Script computes margins from
% DIV; 15 keeps readable margins while accommodating long code scraps.

% Engine-agnostic encoding and fonts
\usepackage{iftex}
\ifPDFTeX
  \usepackage[utf8]{inputenc}
  \usepackage[T1]{fontenc}
  \usepackage{lmodern}
\else
  \usepackage{fontspec}
\fi

% Mathematics
\usepackage{amsmath,amssymb,amsthm}

% Code listings
\usepackage{listings}
\usepackage{xcolor}

% Cross-references and links
\usepackage{hyperref}
\hypersetup{
  colorlinks=true,
  linkcolor=blue!60!black,
  citecolor=green!50!black,
  urlcolor=blue!70!black,
  pdftitle={kcc --- A Formal C17 Compiler in ARM64 Assembly},
  pdfauthor={kcc project},
}

% Section styling
\usepackage{kcc-sections}

% weft weave emits old font commands that KOMA-Script disables by default:
\DeclareOldFontCommand{\tt}{\normalfont\ttfamily}{\mathtt}
\DeclareOldFontCommand{\bf}{\normalfont\bfseries}{\mathbf}
\DeclareOldFontCommand{\rm}{\normalfont\rmfamily}{\mathrm}
\DeclareOldFontCommand{\it}{\normalfont\itshape}{\mathit}
\DeclareOldFontCommand{\sf}{\normalfont\sffamily}{\mathsf}
\DeclareOldFontCommand{\sc}{\normalfont\scshape}{\mathsc}

% Listings configuration for ARM64 assembly
\lstdefinelanguage{aarch64}{
  morekeywords={mov,add,sub,stp,ldp,ldr,str,ldrb,strb,adrp,adr,
    bl,b,ret,svc,cbz,cbnz,cmp,tst,and,orr,eor,lsl,lsr,asr,
    mul,udiv,sdiv,madd,msub,neg,mvn,bic,
    b.eq,b.ne,b.lt,b.le,b.gt,b.ge,b.lo,b.hi,b.hs,b.ls,b.cc,b.cs,
    sxtw,uxtw,sxth,uxth,sxtb,uxtb,
    nop,brk,mrs,msr},
  morekeywords=[2]{x0,x1,x2,x3,x4,x5,x6,x7,x8,x9,x10,x11,x12,
    x13,x14,x15,x16,x17,x18,x19,x20,x21,x22,x23,x24,x25,x26,
    x27,x28,x29,x30,sp,xzr,wzr,
    w0,w1,w2,w3,w4,w5,w6,w7,w8,w9,w10,w11,w12,w13,w14,w15,
    w16,w17,w18,w19,w20,w21,w22,w23,w24,w25,w26,w27,w28,w29,w30},
  morekeywords=[3]{.globl,.p2align,.text,.data,.bss,.equ,.asciz,
    .ascii,.quad,.long,.short,.byte,.space,.zerofill,.macro,.endm,
    .section,.comm,.if,.else,.endif},
  sensitive=true,
  morecomment=[l]{//},
  morestring=[b]",
}
\lstset{
  language=aarch64,
  basicstyle=\ttfamily\small,
  keywordstyle=\color{blue!70!black}\bfseries,
  keywordstyle=[2]\color{red!60!black},
  keywordstyle=[3]\color{green!40!black}\bfseries,
  commentstyle=\color{gray}\itshape,
  stringstyle=\color{orange!80!black},
  numbers=left,
  numberstyle=\tiny\color{gray},
  numbersep=8pt,
  tabsize=4,
  showstringspaces=false,
  breaklines=true,
  frame=single,
  rulecolor=\color{gray!30},
  xleftmargin=1.5em,
  framexleftmargin=1.5em,
}

% Theorem environments for formal semantics
\theoremstyle{definition}
\newtheorem{definition}{Definition}[chapter]
\newtheorem{requirement}{Requirement}[chapter]
\theoremstyle{plain}
\newtheorem{theorem}{Theorem}[chapter]
\newtheorem{lemma}[theorem]{Lemma}
\theoremstyle{remark}
\newtheorem{remark}{Remark}[chapter]

% Custom commands
\newcommand{\iso}[1]{\textsc{iso\,#1}}
\newcommand{\ub}{\textsc{ub}}
\newcommand{\idb}{\textsc{ib}}
\newcommand{\stdref}[1]{\S\,#1}
\newcommand{\kcc}{\textsf{kcc}}
\newcommand{\reg}[1]{\texttt{#1}}

% Title
\title{\kcc\\[0.5em]
  \Large A Formal C17 Compiler in ARM64 Assembly\\[0.3em]
  \large A Literate Program}
\author{The \kcc{} Project}
\date{\today}

\begin{document}
\maketitle
\tableofcontents

% Tolerate uneven page fills — weft scraps are unbreakable blocks
% that frequently leave partial pages. Without this, LaTeX forces
% equal-height columns and produces overfull/underfull vbox warnings.
\raggedbottom

% Forward-reference placeholders for chapters not yet implemented.
% Remove each line when the real chapter is added.
\makeatletter
\newcommand{\forwardlabel}[1]{\immediate\write\@auxout{\string\newlabel{#1}{{??}{??}{}{}{}}}}
\makeatother
% ═══════════════════════════════════════════════════════════════════
% Part I: Foundations
% ═══════════════════════════════════════════════════════════════════
\part{Foundations}

% 01-normative-corpus.weft — Standards and Citation Conventions
%
% This chapter establishes the normative corpus that governs every
% decision in kcc.  No code is emitted here; this is pure prose
% that defines the language of obligation used throughout.

\chapter{The Normative Corpus}
\label{ch:normative-corpus}

\section{Primary Standard}

The sole governing document for the C language accepted by \kcc{} is:

\begin{quote}
\textbf{ISO/IEC 9899:2018} --- \emph{Information technology ---
Programming languages --- C}.  Fourth edition, 2018-06-01.
\end{quote}

This document, hereafter cited as \textbf{C17} or simply ``the Standard,''
supersedes all prior editions (C89, C99, C11).  Where the Standard is
silent, \kcc{} treats the behavior as implementation-defined and documents
it in Appendix~A (\stdref{J.3}).

\section{Secondary References}

The following documents inform the implementation but do not override
the Standard:

\begin{enumerate}
\item \textbf{ARM Architecture Reference Manual (ARM ARM), DDI 0487}
      --- defines the A64 instruction set, encoding rules, and system
      registers.  All instruction mnemonics and encodings cite this
      document by section number.

\item \textbf{Procedure Call Standard for the Arm 64-bit Architecture
      (AAPCS64)}, ARM IHI 0055 --- defines the calling convention,
      register usage, and stack layout.

\item \textbf{Apple ARM64 ABI documentation} --- describes
      platform-specific deviations from AAPCS64: mandatory frame
      pointer, reserved \reg{x18}, no red zone, Mach-O symbol
      conventions.

\item \textbf{Mach-O file format reference} --- governs the binary
      format emitted by the backend.

\item \textbf{macOS system call interface} --- the \texttt{svc \#0x80}
      convention with syscall number in \reg{x16}.
\end{enumerate}

\section{Citation Conventions}

Throughout this document, we use the following shorthand:

\begin{itemize}
\item \stdref{6.5.3.1} refers to ISO 9899:2018 section 6.5.3.1.
\item \textbf{ARM ARM C6.2.5} refers to the ARM Architecture Reference
      Manual, section C6.2.5.
\item \textbf{AAPCS64 \S5.4} refers to the Procedure Call Standard,
      section 5.4.
\end{itemize}

The keywords \textsc{shall}, \textsc{shall not}, \textsc{may}, and
\textsc{undefined behavior} carry their ISO meanings.  We abbreviate
undefined behavior as \ub{} and implementation-defined behavior
as \idb{}.
% 02-conformance-model.weft — Conformance Obligations
%
% Establishes the conformance model: what kcc must do, what it may
% choose, and how it handles undefined behavior.  No code emitted.

\chapter{Conformance Model}
\label{ch:conformance}

\section{Hosted vs.\ Freestanding}

The Standard (\stdref{5.1.2}) distinguishes two execution environments:

\begin{itemize}
\item \textbf{Hosted} (\stdref{5.1.2.2}) --- the full C library is
      available; program entry is \texttt{main}.
\item \textbf{Freestanding} (\stdref{5.1.2.1}) --- a minimal library
      subset; entry point is implementation-defined.
\end{itemize}

\kcc{} targets a \emph{hosted} environment on macOS/ARM64.
The program entry point is \texttt{\_main} per Mach-O convention,
which corresponds to \texttt{main} as defined in \stdref{5.1.2.2.1}.

\section{Translation Limits}

Section \stdref{5.2.4.1} specifies minimum translation limits that
every conforming implementation must support.  \kcc{} meets or exceeds
all limits listed therein.  Notable limits:

\begin{itemize}
\item 127 nesting levels of blocks
\item 63 nesting levels of conditional inclusion
\item 4095 characters in a logical source line (after line splicing)
\item 4095 characters in a string literal
\item 511 identifiers with block scope in one block
\item 127 parameters in one function definition
\item 63 levels of nested structure or union definitions
\end{itemize}

These limits are enforced through the diagnostic engine
(\S\ref{ch:diagnostics}).

\section{Implementation-Defined Behavior}

The Standard identifies approximately 200 implementation-defined
behaviors (\stdref{J.3}).  Each is documented in Appendix~A of this
document with:

\begin{enumerate}
\item The \stdref{J.3} clause number.
\item The choice made by \kcc{}.
\item The rationale for that choice.
\end{enumerate}

\section{Undefined Behavior Policy}

\kcc{} adopts a \emph{diagnostic-first} policy for undefined behavior:

\begin{requirement}[UB Policy]
Where \kcc{} can detect undefined behavior at compile time with
reasonable cost, it shall emit a diagnostic.  Where detection is
infeasible or unreasonably expensive, \kcc{} documents the treatment
in Appendix~B (UB Policy Matrix) and proceeds with the most natural
ARM64 behavior.
\end{requirement}

This policy balances strict conformance with the practical reality
of writing a compiler in assembly language.  The UB matrix is the
authoritative reference for each specific case.

% ═══════════════════════════════════════════════════════════════════
% Part II: Machine
% ═══════════════════════════════════════════════════════════════════
\part{The Machine}

% 07-arm64-conventions.weft — ARM64 Register Map, Apple ABI, Macro Definitions
%
% This chapter establishes the assembly-level conventions used throughout
% kcc.  Every function in the compiler follows these rules.

\chapter{ARM64 Conventions}
\label{ch:arm64-conventions}

\section{Register Assignment}

The AAPCS64 (ARM IHI 0055) divides the 31 general-purpose registers
into well-defined roles.  \kcc{} further refines the callee-saved
registers for persistent compiler state:

\begin{center}
\begin{tabular}{lll}
\hline
Register(s) & AAPCS64 Role & \kcc{} Usage \\
\hline
\reg{x0}--\reg{x7}   & Arguments/results      & Per AAPCS64 \\
\reg{x8}              & Indirect result         & Scratch \\
\reg{x9}--\reg{x15}   & Temporaries             & Scratch (caller-saved) \\
\reg{x16}--\reg{x17}  & Intra-procedure scratch & Syscall number, linker veneers \\
\reg{x18}             & Platform register       & \textbf{Reserved} (Apple) \\
\reg{x19}             & Callee-saved            & Arena pointer (global) \\
\reg{x20}             & Callee-saved            & String table pointer (global) \\
\reg{x21}             & Callee-saved            & Current source buffer \\
\reg{x22}             & Callee-saved            & Current token pointer \\
\reg{x23}--\reg{x28}  & Callee-saved            & Available for local use \\
\reg{x29}             & Frame pointer           & \textbf{Mandatory} (Apple) \\
\reg{x30}             & Link register           & Return address \\
\reg{sp}              & Stack pointer           & 16-byte aligned always \\
\hline
\end{tabular}
\end{center}

\section{Apple Platform Constraints}

Apple ARM64 deviates from vanilla AAPCS64 in several critical ways:

\begin{enumerate}
\item \textbf{Frame pointer is mandatory.}  Every function must
      establish a frame record: \texttt{stp x29, x30, [sp, \#-N]!}
      followed by \texttt{mov x29, sp} (or \texttt{add x29, sp, \#off}).

\item \textbf{\reg{x18} is reserved.}  The platform register must
      never be read or written by user code.

\item \textbf{No red zone.}  Unlike x86-64, there is no guaranteed
      scratch space below \reg{sp}.  All local storage must be
      explicitly allocated by decrementing \reg{sp}.

\item \textbf{Stack alignment.}  \reg{sp} must be 16-byte aligned
      at every instruction that accesses memory through \reg{sp}.

\item \textbf{Mach-O symbol prefix.}  All external symbols carry
      a leading underscore: \texttt{\_main}, \texttt{\_kcc\_arena\_init}.

\item \textbf{Syscall convention.}  System calls use \texttt{svc \#0x80}
      with the syscall number in \reg{x16}, arguments in \reg{x0}--\reg{x5},
      and the return value in \reg{x0}.  The carry flag indicates error.
\end{enumerate}

\section{Common Macros}

We define three macros used throughout the codebase.  Every tangled
\texttt{.s} file includes these at the top.

\texttt{FUNC} establishes a standard function prologue with the
mandatory frame pointer.  \texttt{ENDFUNC} tears down the frame
and returns.  \texttt{SYSCALL} invokes a macOS system call.

\begin{flushleft} \small
\begin{minipage}{\linewidth}\label{scrap1}\raggedright\small
\NWtarget{weft14}{} $\langle\,${\itshape ARM64: common macros}\nobreak\ {\footnotesize {14}}$\,\rangle\equiv$
\vspace{-1ex}
\begin{list}{}{} \item
\mbox{}\verb@@\\
\mbox{}\verb@// -- kcc common macros ----------------------------------------------@\\
\mbox{}\verb@// FUNC name — standard function prologue (Apple ARM64 ABI)@\\
\mbox{}\verb@// Saves frame pointer (x29) and link register (x30).@\\
\mbox{}\verb@// The function is exported with the Mach-O underscore prefix.@\\
\mbox{}\verb@.macro FUNC name@\\
\mbox{}\verb@    .globl _\name@\\
\mbox{}\verb@    .p2align 2@\\
\mbox{}\verb@_\name:@\\
\mbox{}\verb@    stp x29, x30, [sp, #-16]!@\\
\mbox{}\verb@    mov x29, sp@\\
\mbox{}\verb@.endm@\\
\mbox{}\verb@@\\
\mbox{}\verb@// ENDFUNC — standard function epilogue@\\
\mbox{}\verb@// Restores x29/x30 and returns to caller.@\\
\mbox{}\verb@.macro ENDFUNC@\\
\mbox{}\verb@    ldp x29, x30, [sp], #16@\\
\mbox{}\verb@    ret@\\
\mbox{}\verb@.endm@\\
\mbox{}\verb@@\\
\mbox{}\verb@// SYSCALL num — invoke macOS syscall@\\
\mbox{}\verb@// Places the syscall number in x16 and traps to kernel.@\\
\mbox{}\verb@// Arguments must already be in x0–x5 per convention.@\\
\mbox{}\verb@// Returns result in x0; carry flag set on error.@\\
\mbox{}\verb@.macro SYSCALL num@\\
\mbox{}\verb@    mov x16, #\num@\\
\mbox{}\verb@    svc #0x80@\\
\mbox{}\verb@.endm@\\
\mbox{}\verb@@{\NWsep}
\end{list}
\vspace{-1.5ex}
\footnotesize
\begin{list}{}{\setlength{\itemsep}{-\parsep}\setlength{\itemindent}{-\leftmargin}}
\item \NWtxtMacroRefIn\ \NWlink{weft?}{?}\NWlink{weft?}{, ?}\NWlink{weft?}{, ?}\NWlink{weft?}{, ?}\NWlink{weft?}{, ?}\NWlink{weft?}{, ?}\NWlink{weft?}{, ?}\NWlink{weft?}{, ?}\NWlink{weft?}{, ?}\NWlink{weft?}{, ?}\NWlink{weft?}{, ?}\NWlink{weft?}{, ?}\NWlink{weft?}{, ?}\NWlink{weft?}{, ?}.

\item{}
\end{list}
\end{minipage}\vspace{4ex}
\end{flushleft}
% 08-syscalls.weft — macOS System Call Wrappers
%
% kcc uses no C library.  Every interaction with the operating system
% goes through the macOS ARM64 syscall interface: place the syscall
% number in x16, arguments in x0–x5, and execute svc #0x80.
%
% This chapter provides thin wrappers that establish proper frame
% records and present a clean calling convention to the rest of kcc.

\chapter{System Calls}
\label{ch:syscalls}

\section{Syscall Number Constants}

These constants come from the macOS/XNU kernel headers
(\texttt{<sys/syscall.h>}).  On ARM64, syscall numbers are placed
in \reg{x16} before the \texttt{svc} instruction.

\begin{flushleft} \small
\begin{minipage}{\linewidth}\label{scrap2}\raggedright\small
\NWtarget{weft15a}{} $\langle\,${\itshape Syscall: constants}\nobreak\ {\footnotesize {15a}}$\,\rangle\equiv$
\vspace{-1ex}
\begin{list}{}{} \item
\mbox{}\verb@@\\
\mbox{}\verb@// -- macOS ARM64 syscall numbers ------------------------------------@\\
\mbox{}\verb@// Source: XNU bsd/kern/syscalls.master@\\
\mbox{}\verb@.equ SYS_exit,     1@\\
\mbox{}\verb@.equ SYS_read,     3@\\
\mbox{}\verb@.equ SYS_write,    4@\\
\mbox{}\verb@.equ SYS_open,     5@\\
\mbox{}\verb@.equ SYS_close,    6@\\
\mbox{}\verb@.equ SYS_mmap,   197          // 0xC5@\\
\mbox{}\verb@.equ SYS_munmap,  73          // 0x49@\\
\mbox{}\verb@.equ SYS_lseek,  199          // 0xC7@\\
\mbox{}\verb@@\\
\mbox{}\verb@// lseek whence constants@\\
\mbox{}\verb@.equ SEEK_SET,     0@\\
\mbox{}\verb@.equ SEEK_CUR,     1@\\
\mbox{}\verb@.equ SEEK_END,     2@\\
\mbox{}\verb@@\\
\mbox{}\verb@// mmap constants for file mapping@\\
\mbox{}\verb@.equ PROT_READ,    0x1       // PROT_READ only@\\
\mbox{}\verb@.equ MAP_PRIVATE,  0x2       // MAP_PRIVATE (no MAP_ANON)@\\
\mbox{}\verb@@{\NWsep}
\end{list}
\vspace{-1.5ex}
\footnotesize
\begin{list}{}{\setlength{\itemsep}{-\parsep}\setlength{\itemindent}{-\leftmargin}}
\item \NWtxtMacroRefIn\ \NWlink{weft?}{?}\NWlink{weft?}{, ?}\NWlink{weft?}{, ?}.

\item{}
\end{list}
\end{minipage}\vspace{4ex}
\end{flushleft}
\section{Process Control}

The simplest syscall: terminate the process.  \texttt{SYS\_exit}
takes a single argument in \reg{x0} (the exit status) and does
not return.

\begin{flushleft} \small
\begin{minipage}{\linewidth}\label{scrap3}\raggedright\small
\NWtarget{weft15b}{} $\langle\,${\itshape Syscall: exit}\nobreak\ {\footnotesize {15b}}$\,\rangle\equiv$
\vspace{-1ex}
\begin{list}{}{} \item
\mbox{}\verb@@\\
\mbox{}\verb@// -- SYS_exit ---------------------------------------------------@\\
\mbox{}\verb@// void kcc_sys_exit(int status)@\\
\mbox{}\verb@//   x0 = exit status@\\
\mbox{}\verb@// Does not return.@\\
\mbox{}\verb@FUNC kcc_sys_exit@\\
\mbox{}\verb@    SYSCALL SYS_exit@\\
\mbox{}\verb@    brk #1                    // unreachable; trap if kernel returns@\\
\mbox{}\verb@ENDFUNC@\\
\mbox{}\verb@@{\NWsep}
\end{list}
\vspace{-1.5ex}
\footnotesize
\begin{list}{}{\setlength{\itemsep}{-\parsep}\setlength{\itemindent}{-\leftmargin}}
\item \NWtxtMacroRefIn\ \NWlink{weft?}{?}.

\item{}
\end{list}
\end{minipage}\vspace{4ex}
\end{flushleft}
\section{File I/O}

The core I/O syscalls.  Each wrapper preserves the AAPCS64 calling
convention: arguments arrive in \reg{x0}--\reg{x5}, the result is
returned in \reg{x0}, and the carry flag indicates an error (in which
case \reg{x0} holds the \texttt{errno} value).

\subsection{write}

Write bytes to a file descriptor.  This is the fundamental output
primitive; all formatted output in \kcc{} ultimately reduces to
calls to \texttt{kcc\_sys\_write}.

\begin{flushleft} \small
\begin{minipage}{\linewidth}\label{scrap4}\raggedright\small
\NWtarget{weft16a}{} $\langle\,${\itshape Syscall: write}\nobreak\ {\footnotesize {16a}}$\,\rangle\equiv$
\vspace{-1ex}
\begin{list}{}{} \item
\mbox{}\verb@@\\
\mbox{}\verb@// -- SYS_write --------------------------------------------------@\\
\mbox{}\verb@// ssize_t kcc_sys_write(int fd, const void *buf, size_t len)@\\
\mbox{}\verb@//   x0 = file descriptor@\\
\mbox{}\verb@//   x1 = buffer pointer@\\
\mbox{}\verb@//   x2 = byte count@\\
\mbox{}\verb@// Returns bytes written in x0, or (uint64_t)-1 with carry set.@\\
\mbox{}\verb@FUNC kcc_sys_write@\\
\mbox{}\verb@    SYSCALL SYS_write@\\
\mbox{}\verb@ENDFUNC@\\
\mbox{}\verb@@{\NWsep}
\end{list}
\vspace{-1.5ex}
\footnotesize
\begin{list}{}{\setlength{\itemsep}{-\parsep}\setlength{\itemindent}{-\leftmargin}}
\item \NWtxtMacroRefIn\ \NWlink{weft?}{?}.

\item{}
\end{list}
\end{minipage}\vspace{4ex}
\end{flushleft}
\subsection{read}

Read bytes from a file descriptor.  Used by the source manager to read source files when
\texttt{mmap} is unsuitable.

\begin{flushleft} \small
\begin{minipage}{\linewidth}\label{scrap5}\raggedright\small
\NWtarget{weft16b}{} $\langle\,${\itshape Syscall: read}\nobreak\ {\footnotesize {16b}}$\,\rangle\equiv$
\vspace{-1ex}
\begin{list}{}{} \item
\mbox{}\verb@@\\
\mbox{}\verb@// -- SYS_read ---------------------------------------------------@\\
\mbox{}\verb@// ssize_t kcc_sys_read(int fd, void *buf, size_t len)@\\
\mbox{}\verb@//   x0 = file descriptor@\\
\mbox{}\verb@//   x1 = buffer pointer@\\
\mbox{}\verb@//   x2 = maximum bytes to read@\\
\mbox{}\verb@// Returns bytes read in x0, 0 at EOF, or -1 on error.@\\
\mbox{}\verb@FUNC kcc_sys_read@\\
\mbox{}\verb@    SYSCALL SYS_read@\\
\mbox{}\verb@ENDFUNC@\\
\mbox{}\verb@@{\NWsep}
\end{list}
\vspace{-1.5ex}
\footnotesize
\begin{list}{}{\setlength{\itemsep}{-\parsep}\setlength{\itemindent}{-\leftmargin}}
\item \NWtxtMacroRefIn\ \NWlink{weft?}{?}.

\item{}
\end{list}
\end{minipage}\vspace{4ex}
\end{flushleft}
\subsection{open}

Open a file by pathname.  Returns a file descriptor or $-1$ on error.

\begin{flushleft} \small
\begin{minipage}{\linewidth}\label{scrap6}\raggedright\small
\NWtarget{weft17a}{} $\langle\,${\itshape Syscall: open}\nobreak\ {\footnotesize {17a}}$\,\rangle\equiv$
\vspace{-1ex}
\begin{list}{}{} \item
\mbox{}\verb@@\\
\mbox{}\verb@// -- SYS_open ---------------------------------------------------@\\
\mbox{}\verb@// int kcc_sys_open(const char *path, int flags, mode_t mode)@\\
\mbox{}\verb@//   x0 = path (null-terminated)@\\
\mbox{}\verb@//   x1 = flags (O_RDONLY=0, O_WRONLY=1, O_RDWR=2, ...)@\\
\mbox{}\verb@//   x2 = mode (permissions, used with O_CREAT)@\\
\mbox{}\verb@// Returns fd in x0, or -1 on error.@\\
\mbox{}\verb@FUNC kcc_sys_open@\\
\mbox{}\verb@    SYSCALL SYS_open@\\
\mbox{}\verb@ENDFUNC@\\
\mbox{}\verb@@{\NWsep}
\end{list}
\vspace{-1.5ex}
\footnotesize
\begin{list}{}{\setlength{\itemsep}{-\parsep}\setlength{\itemindent}{-\leftmargin}}
\item \NWtxtMacroRefIn\ \NWlink{weft?}{?}.

\item{}
\end{list}
\end{minipage}\vspace{4ex}
\end{flushleft}
\subsection{close}

Close a file descriptor.

\begin{flushleft} \small
\begin{minipage}{\linewidth}\label{scrap7}\raggedright\small
\NWtarget{weft17b}{} $\langle\,${\itshape Syscall: close}\nobreak\ {\footnotesize {17b}}$\,\rangle\equiv$
\vspace{-1ex}
\begin{list}{}{} \item
\mbox{}\verb@@\\
\mbox{}\verb@// -- SYS_close --------------------------------------------------@\\
\mbox{}\verb@// int kcc_sys_close(int fd)@\\
\mbox{}\verb@//   x0 = file descriptor@\\
\mbox{}\verb@// Returns 0 on success, -1 on error.@\\
\mbox{}\verb@FUNC kcc_sys_close@\\
\mbox{}\verb@    SYSCALL SYS_close@\\
\mbox{}\verb@ENDFUNC@\\
\mbox{}\verb@@{\NWsep}
\end{list}
\vspace{-1.5ex}
\footnotesize
\begin{list}{}{\setlength{\itemsep}{-\parsep}\setlength{\itemindent}{-\leftmargin}}
\item \NWtxtMacroRefIn\ \NWlink{weft?}{?}.

\item{}
\end{list}
\end{minipage}\vspace{4ex}
\end{flushleft}
\section{Memory Management}

\subsection{mmap}

Map pages of memory.  This is how \kcc{} obtains all dynamic memory:
the arena allocator (Chapter~\ref{ch:arena}) calls \texttt{mmap}
with \texttt{MAP\_PRIVATE | MAP\_ANON} to get fresh pages from the
kernel.

\begin{flushleft} \small
\begin{minipage}{\linewidth}\label{scrap8}\raggedright\small
\NWtarget{weft18}{} $\langle\,${\itshape Syscall: mmap}\nobreak\ {\footnotesize {18}}$\,\rangle\equiv$
\vspace{-1ex}
\begin{list}{}{} \item
\mbox{}\verb@@\\
\mbox{}\verb@// -- SYS_mmap ---------------------------------------------------@\\
\mbox{}\verb@// void *kcc_sys_mmap(void *addr, size_t len, int prot,@\\
\mbox{}\verb@//                    int flags, int fd, off_t offset)@\\
\mbox{}\verb@//   x0 = requested address (0 for kernel choice)@\\
\mbox{}\verb@//   x1 = length in bytes@\\
\mbox{}\verb@//   x2 = protection (PROT_READ|PROT_WRITE = 0x3)@\\
\mbox{}\verb@//   x3 = flags (MAP_PRIVATE|MAP_ANON = 0x1002)@\\
\mbox{}\verb@//   x4 = file descriptor (-1 for anonymous)@\\
\mbox{}\verb@//   x5 = offset@\\
\mbox{}\verb@// Returns mapped address in x0, or MAP_FAILED ((void *)-1).@\\
\mbox{}\verb@FUNC kcc_sys_mmap@\\
\mbox{}\verb@    SYSCALL SYS_mmap@\\
\mbox{}\verb@ENDFUNC@\\
\mbox{}\verb@@{\NWsep}
\end{list}
\vspace{-1.5ex}
\footnotesize
\begin{list}{}{\setlength{\itemsep}{-\parsep}\setlength{\itemindent}{-\leftmargin}}
\item \NWtxtMacroRefIn\ \NWlink{weft?}{?}.

\item{}
\end{list}
\end{minipage}\vspace{4ex}
\end{flushleft}
\subsection{lseek}

Reposition the file offset.  Used by the source manager to determine
file size before \texttt{mmap}: \texttt{lseek(fd, 0, SEEK\_END)}
returns the file length.

\begin{flushleft} \small
\begin{minipage}{\linewidth}\label{scrap9}\raggedright\small
\NWtarget{weft19}{} $\langle\,${\itshape Syscall: lseek}\nobreak\ {\footnotesize {19}}$\,\rangle\equiv$
\vspace{-1ex}
\begin{list}{}{} \item
\mbox{}\verb@@\\
\mbox{}\verb@// -- SYS_lseek --------------------------------------------------@\\
\mbox{}\verb@// off_t kcc_sys_lseek(int fd, off_t offset, int whence)@\\
\mbox{}\verb@//   x0 = file descriptor@\\
\mbox{}\verb@//   x1 = offset@\\
\mbox{}\verb@//   x2 = whence (SEEK_SET=0, SEEK_CUR=1, SEEK_END=2)@\\
\mbox{}\verb@// Returns new file offset in x0, or -1 on error.@\\
\mbox{}\verb@FUNC kcc_sys_lseek@\\
\mbox{}\verb@    SYSCALL SYS_lseek@\\
\mbox{}\verb@ENDFUNC@\\
\mbox{}\verb@@{\NWsep}
\end{list}
\vspace{-1.5ex}
\footnotesize
\begin{list}{}{\setlength{\itemsep}{-\parsep}\setlength{\itemindent}{-\leftmargin}}
\item \NWtxtMacroRefIn\ \NWlink{weft?}{?}.

\item{}
\end{list}
\end{minipage}\vspace{4ex}
\end{flushleft}
% 09-memory-management.weft — Arena Allocator in ARM64 Assembly
%
% kcc manages all dynamic memory through a single arena allocator.
% The arena obtains large chunks from the kernel via mmap and
% sub-allocates from them with a simple bump pointer.  There is
% no individual free operation --- the entire arena is released at
% once when compilation completes.
%
% This strategy is ideal for a compiler: allocations are numerous,
% small, and all share the same lifetime (one translation unit).

\chapter{Memory Management: The Arena}
\label{ch:arena}

\section{Design Rationale}

A compiler allocates thousands of small objects (tokens, AST nodes,
type descriptors, symbol entries) that all live until the end of
compilation.  A general-purpose allocator (malloc/free) would waste
time tracking individual lifetimes.  An arena allocator reduces
allocation to a pointer bump --- a single \texttt{add} instruction ---
and deallocation to a single \texttt{munmap} of the entire arena.

\section{Chunk Layout}

Each arena chunk is a contiguous region obtained from \texttt{mmap}.
The first 24 bytes form a header:

\begin{center}
\begin{tabular}{rll}
\hline
Offset & Field & Description \\
\hline
0  & \texttt{next}  & Pointer to the next chunk (linked list) \\
8  & \texttt{limit} & One past the last usable byte \\
16 & \texttt{avail} & Next free byte (bump pointer) \\
\hline
\end{tabular}
\end{center}

\begin{flushleft} \small
\begin{minipage}{\linewidth}\label{scrap10}\raggedright\small
\NWtarget{weft20a}{} $\langle\,${\itshape Arena: constants}\nobreak\ {\footnotesize {20a}}$\,\rangle\equiv$
\vspace{-1ex}
\begin{list}{}{} \item
\mbox{}\verb@@\\
\mbox{}\verb@// -- Arena chunk layout ---------------------------------------------@\\
\mbox{}\verb@.equ CHUNK_NEXT,    0         // pointer to next chunk@\\
\mbox{}\verb@.equ CHUNK_LIMIT,   8         // end of usable region@\\
\mbox{}\verb@.equ CHUNK_AVAIL,  16         // next free byte (bump pointer)@\\
\mbox{}\verb@.equ CHUNK_HEADER, 24         // total header size@\\
\mbox{}\verb@@\\
\mbox{}\verb@// Default chunk size: 64 KiB (16 pages of 4 KiB)@\\
\mbox{}\verb@.equ ARENA_CHUNK_SIZE, 65536@\\
\mbox{}\verb@@\\
\mbox{}\verb@// mmap constants (macOS)@\\
\mbox{}\verb@.equ PROT_RW,       0x3      // PROT_READ | PROT_WRITE@\\
\mbox{}\verb@.equ MAP_PRIV_ANON, 0x1002   // MAP_PRIVATE | MAP_ANON@\\
\mbox{}\verb@.equ MAP_FAILED,    -1@\\
\mbox{}\verb@@{\NWsep}
\end{list}
\vspace{-1.5ex}
\footnotesize
\begin{list}{}{\setlength{\itemsep}{-\parsep}\setlength{\itemindent}{-\leftmargin}}
\item \NWtxtMacroRefIn\ \NWlink{weft?}{?}.

\item{}
\end{list}
\end{minipage}\vspace{4ex}
\end{flushleft}
\section{Initialization}

\texttt{kcc\_arena\_init} obtains the first chunk from the kernel
and initializes the global arena state.  The arena pointer is stored
in a global variable; in later phases, the main processing loop
will cache it in \reg{x19}.

\begin{flushleft} \small
\begin{minipage}{\linewidth}\label{scrap11}\raggedright\small
\NWtarget{weft20b}{} $\langle\,${\itshape Arena: init}\nobreak\ {\footnotesize {20b}}$\,\rangle\equiv$
\vspace{-1ex}
\begin{list}{}{} \item
\mbox{}\verb@@\\
\mbox{}\verb@// -- kcc_arena_init -------------------------------------------------@\\
\mbox{}\verb@// void kcc_arena_init(void)@\\
\mbox{}\verb@// Allocates the initial arena chunk via mmap.@\\
\mbox{}\verb@// On failure (mmap returns MAP_FAILED), exits with status 1.@\\
\mbox{}\verb@.globl _kcc_arena_init@\\
\mbox{}\verb@.p2align 2@\\
\mbox{}\verb@_kcc_arena_init:@\\
\mbox{}\verb@    stp x29, x30, [sp, #-16]!@\\
\mbox{}\verb@    mov x29, sp@\\
\mbox{}\verb@@\hbox{$\langle\,${\itshape Arena: init mmap call}\nobreak\ {\footnotesize \NWlink{weft21}{21}}$\,\rangle$}\verb@@\\
\mbox{}\verb@@\hbox{$\langle\,${\itshape Arena: init chunk header}\nobreak\ {\footnotesize \NWlink{weft22}{22}}$\,\rangle$}\verb@@\\
\mbox{}\verb@    ldp x29, x30, [sp], #16@\\
\mbox{}\verb@    ret@\\
\mbox{}\verb@@\\
\mbox{}\verb@.Larena_init_fail:@\\
\mbox{}\verb@    mov x0, #1                    // exit status = 1@\\
\mbox{}\verb@    bl _kcc_sys_exit@\\
\mbox{}\verb@@{\NWsep}
\end{list}
\vspace{-1.5ex}
\footnotesize
\begin{list}{}{\setlength{\itemsep}{-\parsep}\setlength{\itemindent}{-\leftmargin}}
\item \NWtxtMacroRefIn\ \NWlink{weft?}{?}.

\item{}
\end{list}
\end{minipage}\vspace{4ex}
\end{flushleft}
The initialization proceeds in two steps.  First, we request a
64\,KiB anonymous mapping from the kernel:

\begin{flushleft} \small
\begin{minipage}{\linewidth}\label{scrap12}\raggedright\small
\NWtarget{weft21}{} $\langle\,${\itshape Arena: init mmap call}\nobreak\ {\footnotesize {21}}$\,\rangle\equiv$
\vspace{-1ex}
\begin{list}{}{} \item
\mbox{}\verb@@\\
\mbox{}\verb@    // mmap(NULL, ARENA_CHUNK_SIZE, PROT_RW, MAP_PRIV_ANON, -1, 0)@\\
\mbox{}\verb@    mov x0, #0                    // addr = NULL (kernel chooses)@\\
\mbox{}\verb@    mov x1, #ARENA_CHUNK_SIZE     // len  = 64 KiB@\\
\mbox{}\verb@    mov x2, #PROT_RW              // prot = read|write@\\
\mbox{}\verb@    mov x3, #MAP_PRIV_ANON        // flags = private|anonymous@\\
\mbox{}\verb@    mov x4, #-1                   // fd   = -1 (no file)@\\
\mbox{}\verb@    mov x5, #0                    // offset = 0@\\
\mbox{}\verb@    bl _kcc_sys_mmap@\\
\mbox{}\verb@@\\
\mbox{}\verb@    // Check for MAP_FAILED@\\
\mbox{}\verb@    cmn x0, #1                    // compare with (uint64_t)-1@\\
\mbox{}\verb@    b.eq .Larena_init_fail@\\
\mbox{}\verb@@{\NWsep}
\end{list}
\vspace{-1.5ex}
\footnotesize
\begin{list}{}{\setlength{\itemsep}{-\parsep}\setlength{\itemindent}{-\leftmargin}}
\item \NWtxtMacroRefIn\ \NWlink{weft20b}{20b}.

\item{}
\end{list}
\end{minipage}\vspace{4ex}
\end{flushleft}
Then we write the three header fields and record the chunk as current:

\begin{flushleft} \small
\begin{minipage}{\linewidth}\label{scrap13}\raggedright\small
\NWtarget{weft22}{} $\langle\,${\itshape Arena: init chunk header}\nobreak\ {\footnotesize {22}}$\,\rangle\equiv$
\vspace{-1ex}
\begin{list}{}{} \item
\mbox{}\verb@@\\
\mbox{}\verb@    // Initialize chunk header (x0 = base address)@\\
\mbox{}\verb@    str xzr, [x0, #CHUNK_NEXT]   // next = NULL@\\
\mbox{}\verb@@\\
\mbox{}\verb@    mov x1, x0@\\
\mbox{}\verb@    add x1, x1, #ARENA_CHUNK_SIZE@\\
\mbox{}\verb@    str x1, [x0, #CHUNK_LIMIT]   // limit = base + size@\\
\mbox{}\verb@@\\
\mbox{}\verb@    mov x1, x0@\\
\mbox{}\verb@    add x1, x1, #CHUNK_HEADER@\\
\mbox{}\verb@    str x1, [x0, #CHUNK_AVAIL]   // avail = base + header@\\
\mbox{}\verb@@\\
\mbox{}\verb@    // Store as the current chunk@\\
\mbox{}\verb@    adrp x1, _kcc_arena_current@{\tt @}\verb@PAGE@\\
\mbox{}\verb@    add  x1, x1, _kcc_arena_current@{\tt @}\verb@PAGEOFF@\\
\mbox{}\verb@    str  x0, [x1]@\\
\mbox{}\verb@@{\NWsep}
\end{list}
\vspace{-1.5ex}
\footnotesize
\begin{list}{}{\setlength{\itemsep}{-\parsep}\setlength{\itemindent}{-\leftmargin}}
\item \NWtxtMacroRefIn\ \NWlink{weft20b}{20b}.

\item{}
\end{list}
\end{minipage}\vspace{4ex}
\end{flushleft}
\section{Allocation: Fast Path}

\texttt{kcc\_arena\_alloc} is the primary allocation function.  It
takes a size in bytes (in \reg{x0}), rounds it up to 8-byte alignment,
and returns a pointer to the allocated region.

The fast path is a single comparison and pointer bump --- three
instructions after the alignment arithmetic:

\begin{flushleft} \small
\begin{minipage}{\linewidth}\label{scrap14}\raggedright\small
\NWtarget{weft23}{} $\langle\,${\itshape Arena: alloc fast path}\nobreak\ {\footnotesize {23}}$\,\rangle\equiv$
\vspace{-1ex}
\begin{list}{}{} \item
\mbox{}\verb@@\\
\mbox{}\verb@// -- kcc_arena_alloc ------------------------------------------------@\\
\mbox{}\verb@// void *kcc_arena_alloc(size_t size)@\\
\mbox{}\verb@//   x0 = number of bytes to allocate@\\
\mbox{}\verb@// Returns pointer to allocated memory in x0.@\\
\mbox{}\verb@// All allocations are 8-byte aligned.@\\
\mbox{}\verb@.globl _kcc_arena_alloc@\\
\mbox{}\verb@.p2align 2@\\
\mbox{}\verb@_kcc_arena_alloc:@\\
\mbox{}\verb@    stp x29, x30, [sp, #-32]!@\\
\mbox{}\verb@    mov x29, sp@\\
\mbox{}\verb@    str x19, [sp, #16]@\\
\mbox{}\verb@@\\
\mbox{}\verb@    // Round size up to 8-byte alignment: (size + 7) & ~7@\\
\mbox{}\verb@    add x0, x0, #7@\\
\mbox{}\verb@    and x0, x0, #~7@\\
\mbox{}\verb@    mov x19, x0                   // x19 = aligned size@\\
\mbox{}\verb@@\\
\mbox{}\verb@    // Load current chunk pointer@\\
\mbox{}\verb@    adrp x1, _kcc_arena_current@{\tt @}\verb@PAGE@\\
\mbox{}\verb@    add  x1, x1, _kcc_arena_current@{\tt @}\verb@PAGEOFF@\\
\mbox{}\verb@    ldr  x2, [x1]                 // x2 = current chunk@\\
\mbox{}\verb@@\\
\mbox{}\verb@    // Load avail and limit@\\
\mbox{}\verb@    ldr x3, [x2, #CHUNK_AVAIL]   // x3 = avail@\\
\mbox{}\verb@    ldr x4, [x2, #CHUNK_LIMIT]   // x4 = limit@\\
\mbox{}\verb@@\\
\mbox{}\verb@    // Check: avail + size <= limit?@\\
\mbox{}\verb@    add x5, x3, x19@\\
\mbox{}\verb@    cmp x5, x4@\\
\mbox{}\verb@    b.hi .Larena_alloc_slow       // overflow: need new chunk@\\
\mbox{}\verb@@\\
\mbox{}\verb@    // Fast path: bump the pointer@\\
\mbox{}\verb@    str x5, [x2, #CHUNK_AVAIL]   // avail += size@\\
\mbox{}\verb@    mov x0, x3                    // return old avail@\\
\mbox{}\verb@@\\
\mbox{}\verb@    ldr x19, [sp, #16]@\\
\mbox{}\verb@    ldp x29, x30, [sp], #32@\\
\mbox{}\verb@    ret@\\
\mbox{}\verb@@{\NWsep}
\end{list}
\vspace{-1.5ex}
\footnotesize
\begin{list}{}{\setlength{\itemsep}{-\parsep}\setlength{\itemindent}{-\leftmargin}}
\item \NWtxtMacroRefIn\ \NWlink{weft?}{?}.

\item{}
\end{list}
\end{minipage}\vspace{4ex}
\end{flushleft}
\section{Allocation: Slow Path}

When the current chunk is exhausted, we allocate a new one via
\texttt{mmap}, link it to the chain, and fulfill the request from
the fresh chunk.

\begin{flushleft} \small
\begin{minipage}{\linewidth}\label{scrap15}\raggedright\small
\NWtarget{weft24a}{} $\langle\,${\itshape Arena: alloc slow path}\nobreak\ {\footnotesize {24a}}$\,\rangle\equiv$
\vspace{-1ex}
\begin{list}{}{} \item
\mbox{}\verb@@\\
\mbox{}\verb@.Larena_alloc_slow:@\\
\mbox{}\verb@    // Chunk size = max(ARENA_CHUNK_SIZE, CHUNK_HEADER + size)@\\
\mbox{}\verb@    add x0, x19, #CHUNK_HEADER@\\
\mbox{}\verb@    mov x1, #ARENA_CHUNK_SIZE@\\
\mbox{}\verb@    cmp x0, x1@\\
\mbox{}\verb@    csel x1, x0, x1, hi          // x1 = max(needed, default)@\\
\mbox{}\verb@@\\
\mbox{}\verb@    // mmap the new chunk@\\
\mbox{}\verb@    mov x0, #0@\\
\mbox{}\verb@    mov x2, #PROT_RW@\\
\mbox{}\verb@    mov x3, #MAP_PRIV_ANON@\\
\mbox{}\verb@    mov x4, #-1@\\
\mbox{}\verb@    mov x5, #0@\\
\mbox{}\verb@    bl _kcc_sys_mmap@\\
\mbox{}\verb@@\\
\mbox{}\verb@    cmn x0, #1@\\
\mbox{}\verb@    b.eq .Larena_alloc_oom@\\
\mbox{}\verb@@\hbox{$\langle\,${\itshape Arena: alloc link new chunk}\nobreak\ {\footnotesize \NWlink{weft24b}{24b}}$\,\rangle$}\verb@@\\
\mbox{}\verb@@\\
\mbox{}\verb@.Larena_alloc_oom:@\\
\mbox{}\verb@    mov x0, #2                   // exit status = 2@\\
\mbox{}\verb@    bl _kcc_sys_exit@\\
\mbox{}\verb@@{\NWsep}
\end{list}
\vspace{-1.5ex}
\footnotesize
\begin{list}{}{\setlength{\itemsep}{-\parsep}\setlength{\itemindent}{-\leftmargin}}
\item \NWtxtMacroRefIn\ \NWlink{weft?}{?}.

\item{}
\end{list}
\end{minipage}\vspace{4ex}
\end{flushleft}
Once the new chunk is mapped, we link it into the chain, initialize
its header, and return the first allocation from it:

\begin{flushleft} \small
\begin{minipage}{\linewidth}\label{scrap16}\raggedright\small
\NWtarget{weft24b}{} $\langle\,${\itshape Arena: alloc link new chunk}\nobreak\ {\footnotesize {24b}}$\,\rangle\equiv$
\vspace{-1ex}
\begin{list}{}{} \item
\mbox{}\verb@@\\
\mbox{}\verb@    // x0 = new chunk base; link to previous chunk@\\
\mbox{}\verb@    adrp x1, _kcc_arena_current@{\tt @}\verb@PAGE@\\
\mbox{}\verb@    add  x1, x1, _kcc_arena_current@{\tt @}\verb@PAGEOFF@\\
\mbox{}\verb@    ldr  x2, [x1]                 // x2 = old chunk@\\
\mbox{}\verb@@\\
\mbox{}\verb@    str x2, [x0, #CHUNK_NEXT]    // new->next = old@\\
\mbox{}\verb@@\\
\mbox{}\verb@    mov x3, x0@\\
\mbox{}\verb@    add x3, x3, #ARENA_CHUNK_SIZE@\\
\mbox{}\verb@    str x3, [x0, #CHUNK_LIMIT]   // limit (conservative)@\\
\mbox{}\verb@@\\
\mbox{}\verb@    mov x3, x0@\\
\mbox{}\verb@    add x3, x3, #CHUNK_HEADER@\\
\mbox{}\verb@    add x4, x3, x19              // new avail after allocation@\\
\mbox{}\verb@    str x4, [x0, #CHUNK_AVAIL]@\\
\mbox{}\verb@@\\
\mbox{}\verb@    str x0, [x1]                  // update arena_current@\\
\mbox{}\verb@    mov x0, x3                    // return = first byte after header@\\
\mbox{}\verb@@\\
\mbox{}\verb@    ldr x19, [sp, #16]@\\
\mbox{}\verb@    ldp x29, x30, [sp], #32@\\
\mbox{}\verb@    ret@\\
\mbox{}\verb@@{\NWsep}
\end{list}
\vspace{-1.5ex}
\footnotesize
\begin{list}{}{\setlength{\itemsep}{-\parsep}\setlength{\itemindent}{-\leftmargin}}
\item \NWtxtMacroRefIn\ \NWlink{weft24a}{24a}.

\item{}
\end{list}
\end{minipage}\vspace{4ex}
\end{flushleft}
\section{Global State}

The arena's mutable state is a single pointer to the current (most
recently allocated) chunk.  We store it in the \texttt{\_\_DATA} segment.

\begin{flushleft} \small
\begin{minipage}{\linewidth}\label{scrap17}\raggedright\small
\NWtarget{weft25}{} $\langle\,${\itshape Arena: data}\nobreak\ {\footnotesize {25}}$\,\rangle\equiv$
\vspace{-1ex}
\begin{list}{}{} \item
\mbox{}\verb@@\\
\mbox{}\verb@// -- Arena global state ---------------------------------------------@\\
\mbox{}\verb@.data@\\
\mbox{}\verb@.p2align 3@\\
\mbox{}\verb@_kcc_arena_current: .quad 0   // pointer to current arena chunk@\\
\mbox{}\verb@@{\NWsep}
\end{list}
\vspace{-1.5ex}
\footnotesize
\begin{list}{}{\setlength{\itemsep}{-\parsep}\setlength{\itemindent}{-\leftmargin}}
\item \NWtxtMacroRefIn\ \NWlink{weft?}{?}.

\item{}
\end{list}
\end{minipage}\vspace{4ex}
\end{flushleft}
% 10-data-structures.weft — Hash Table, Vector, and String Interning
%
% A compiler needs three fundamental data structures: a hash table
% for O(1) lookups (keyword recognition, symbol tables), a dynamic
% vector for growable sequences (token streams, AST children), and
% a string intern table for canonicalizing identifiers.
%
% All three are arena-allocated, meaning they never individually
% free memory --- the entire arena is released at program end.

\chapter{Data Structures}
\label{ch:data-structures}

\section{FNV-1a Hash Function}

FNV-1a (Fowler/Noll/Vo, variant ``a'') is a non-cryptographic hash
function with excellent distribution for short strings --- exactly the
workload a compiler presents (identifiers average 6--12 bytes).  The
algorithm is simple: start from an offset basis, and for each byte
\texttt{XOR} it into the hash then multiply by the FNV prime.

The 64-bit constants are:
\begin{itemize}
\item Offset basis: \texttt{0xCBF29CE484222325} ($= 14695981039346656037$)
\item Prime: \texttt{0x00000100000001B3} ($= 1099511628211$)
\end{itemize}

\begin{flushleft} \small
\begin{minipage}{\linewidth}\label{scrap18}\raggedright\small
\NWtarget{weft26}{} $\langle\,${\itshape Hash: FNV-1a}\nobreak\ {\footnotesize {26}}$\,\rangle\equiv$
\vspace{-1ex}
\begin{list}{}{} \item
\mbox{}\verb@@\\
\mbox{}\verb@// -- kcc_fnv1a ------------------------------------------------------@\\
\mbox{}\verb@// uint64_t kcc_fnv1a(const void *data, size_t len)@\\
\mbox{}\verb@//   x0 = pointer to data@\\
\mbox{}\verb@//   x1 = byte count@\\
\mbox{}\verb@// Returns 64-bit FNV-1a hash in x0.@\\
\mbox{}\verb@.globl _kcc_fnv1a@\\
\mbox{}\verb@.p2align 2@\\
\mbox{}\verb@_kcc_fnv1a:@\\
\mbox{}\verb@    stp x29, x30, [sp, #-16]!@\\
\mbox{}\verb@    mov x29, sp@\\
\mbox{}\verb@@\\
\mbox{}\verb@    // Load FNV-1a offset basis: 0xCBF29CE484222325@\\
\mbox{}\verb@    movz x2, #0xCBF2, lsl #48@\\
\mbox{}\verb@    movk x2, #0x9CE4, lsl #32@\\
\mbox{}\verb@    movk x2, #0x8422, lsl #16@\\
\mbox{}\verb@    movk x2, #0x2325@\\
\mbox{}\verb@@\\
\mbox{}\verb@    // Load FNV-1a prime: 0x00000100000001B3@\\
\mbox{}\verb@    movz x3, #0x0100, lsl #32@\\
\mbox{}\verb@    movk x3, #0x01B3@\\
\mbox{}\verb@@\\
\mbox{}\verb@    cbz x1, .Lfnv_done@\\
\mbox{}\verb@@\\
\mbox{}\verb@.Lfnv_loop:@\\
\mbox{}\verb@    ldrb w4, [x0], #1           // load byte@\\
\mbox{}\verb@    eor  x2, x2, x4             // hash ^= byte@\\
\mbox{}\verb@    mul  x2, x2, x3             // hash *= prime@\\
\mbox{}\verb@    subs x1, x1, #1@\\
\mbox{}\verb@    b.ne .Lfnv_loop@\\
\mbox{}\verb@@\\
\mbox{}\verb@.Lfnv_done:@\\
\mbox{}\verb@    mov x0, x2                  // return hash@\\
\mbox{}\verb@    ldp x29, x30, [sp], #16@\\
\mbox{}\verb@    ret@\\
\mbox{}\verb@@{\NWsep}
\end{list}
\vspace{-1.5ex}
\footnotesize
\begin{list}{}{\setlength{\itemsep}{-\parsep}\setlength{\itemindent}{-\leftmargin}}
\item \NWtxtMacroRefIn\ \NWlink{weft?}{?}.

\item{}
\end{list}
\end{minipage}\vspace{4ex}
\end{flushleft}
\section{Hash Table}

The hash table uses open addressing with linear probing.  Each entry
stores the cached hash, a pointer to the key, the key length, and a
64-bit value.  Empty slots are indicated by a null key pointer.

\begin{center}
\begin{tabular}{rll}
\hline
Offset & Field & Description \\
\hline
0  & \texttt{hash}   & Cached FNV-1a hash \\
8  & \texttt{key}    & Pointer to key bytes \\
16 & \texttt{keylen} & Key length in bytes \\
24 & \texttt{value}  & Associated 64-bit value \\
\hline
\end{tabular}
\end{center}

\begin{flushleft} \small
\begin{minipage}{\linewidth}\label{scrap19}\raggedright\small
\NWtarget{weft27a}{} $\langle\,${\itshape Hash: constants}\nobreak\ {\footnotesize {27a}}$\,\rangle\equiv$
\vspace{-1ex}
\begin{list}{}{} \item
\mbox{}\verb@@\\
\mbox{}\verb@// -- Hash table layout ----------------------------------------------@\\
\mbox{}\verb@.equ HT_ENTRIES,       0     // pointer to entries array@\\
\mbox{}\verb@.equ HT_CAP,           8     // capacity (power of 2)@\\
\mbox{}\verb@.equ HT_COUNT,        16     // number of occupied slots@\\
\mbox{}\verb@.equ HT_SIZE,         24     // total header size@\\
\mbox{}\verb@@\\
\mbox{}\verb@.equ HTE_HASH,         0     // cached hash value@\\
\mbox{}\verb@.equ HTE_KEY,          8     // pointer to key@\\
\mbox{}\verb@.equ HTE_KEYLEN,      16     // key length@\\
\mbox{}\verb@.equ HTE_VALUE,       24     // associated value@\\
\mbox{}\verb@.equ HTE_SIZE,        32     // entry size@\\
\mbox{}\verb@@{\NWsep}
\end{list}
\vspace{-1.5ex}
\footnotesize
\begin{list}{}{\setlength{\itemsep}{-\parsep}\setlength{\itemindent}{-\leftmargin}}
\item \NWtxtMacroRefIn\ \NWlink{weft?}{?}.

\item{}
\end{list}
\end{minipage}\vspace{4ex}
\end{flushleft}
\subsection{Initialization}

\texttt{kcc\_ht\_init} creates a hash table with the given capacity
(which must be a power of~2).  The entries array is arena-allocated
and zero-filled, so all key pointers start as \texttt{NULL}.

\begin{flushleft} \small
\begin{minipage}{\linewidth}\label{scrap20}\raggedright\small
\NWtarget{weft27b}{} $\langle\,${\itshape Hash: init}\nobreak\ {\footnotesize {27b}}$\,\rangle\equiv$
\vspace{-1ex}
\begin{list}{}{} \item
\mbox{}\verb@@\\
\mbox{}\verb@// -- kcc_ht_init ----------------------------------------------------@\\
\mbox{}\verb@// void *kcc_ht_init(uint64_t capacity)@\\
\mbox{}\verb@//   x0 = capacity (must be power of 2)@\\
\mbox{}\verb@// Returns pointer to hash table header.@\\
\mbox{}\verb@.globl _kcc_ht_init@\\
\mbox{}\verb@.p2align 2@\\
\mbox{}\verb@_kcc_ht_init:@\\
\mbox{}\verb@    stp x29, x30, [sp, #-48]!@\\
\mbox{}\verb@    mov x29, sp@\\
\mbox{}\verb@    stp x19, x20, [sp, #16]@\\
\mbox{}\verb@    str x21, [sp, #32]@\\
\mbox{}\verb@@\\
\mbox{}\verb@    mov x19, x0                  // x19 = capacity@\\
\mbox{}\verb@@\\
\mbox{}\verb@    // Allocate the header (HT_SIZE = 24 bytes)@\\
\mbox{}\verb@    mov x0, #HT_SIZE@\\
\mbox{}\verb@    bl _kcc_arena_alloc@\\
\mbox{}\verb@    mov x20, x0                  // x20 = header pointer@\\
\mbox{}\verb@@\\
\mbox{}\verb@    // Allocate entries array (capacity * HTE_SIZE)@\\
\mbox{}\verb@    mov x0, x19@\\
\mbox{}\verb@    mov x1, #HTE_SIZE@\\
\mbox{}\verb@    mul x0, x0, x1@\\
\mbox{}\verb@    mov x21, x0                  // x21 = total bytes@\\
\mbox{}\verb@    bl _kcc_arena_alloc@\\
\mbox{}\verb@    // x0 = entries pointer@\\
\mbox{}\verb@@\\
\mbox{}\verb@    // Zero-fill entries (arena doesn't guarantee zeroed memory)@\\
\mbox{}\verb@    mov x1, #0                   // fill byte@\\
\mbox{}\verb@    mov x2, x21                  // byte count@\\
\mbox{}\verb@    // x0 already = entries ptr, save it first@\\
\mbox{}\verb@    str x0, [x20, #HT_ENTRIES]  // store entries pointer@\\
\mbox{}\verb@    mov x3, x0                   // save entries ptr@\\
\mbox{}\verb@    bl _kcc_memset@\\
\mbox{}\verb@@\\
\mbox{}\verb@    // Fill header@\\
\mbox{}\verb@    str x19, [x20, #HT_CAP]     // capacity@\\
\mbox{}\verb@    str xzr, [x20, #HT_COUNT]   // count = 0@\\
\mbox{}\verb@@\\
\mbox{}\verb@    mov x0, x20                  // return header pointer@\\
\mbox{}\verb@@\\
\mbox{}\verb@    ldr x21, [sp, #32]@\\
\mbox{}\verb@    ldp x19, x20, [sp, #16]@\\
\mbox{}\verb@    ldp x29, x30, [sp], #48@\\
\mbox{}\verb@    ret@\\
\mbox{}\verb@@{\NWsep}
\end{list}
\vspace{-1.5ex}
\footnotesize
\begin{list}{}{\setlength{\itemsep}{-\parsep}\setlength{\itemindent}{-\leftmargin}}
\item \NWtxtMacroRefIn\ \NWlink{weft?}{?}.

\item{}
\end{list}
\end{minipage}\vspace{4ex}
\end{flushleft}
\subsection{Lookup}

\texttt{kcc\_ht\_get} searches for a key in the table.  It returns
the associated value if found, or 0 if not.  Since 0 is a valid
``not found'' sentinel for our use cases (keyword token types start
at~10), this is sufficient.

\begin{flushleft} \small
\begin{minipage}{\linewidth}\label{scrap21}\raggedright\small
\NWtarget{weft29}{} $\langle\,${\itshape Hash: get}\nobreak\ {\footnotesize {29}}$\,\rangle\equiv$
\vspace{-1ex}
\begin{list}{}{} \item
\mbox{}\verb@@\\
\mbox{}\verb@// -- kcc_ht_get -----------------------------------------------------@\\
\mbox{}\verb@// uint64_t kcc_ht_get(void *ht, const void *key, uint64_t keylen)@\\
\mbox{}\verb@//   x0 = hash table pointer@\\
\mbox{}\verb@//   x1 = key pointer@\\
\mbox{}\verb@//   x2 = key length@\\
\mbox{}\verb@// Returns value if found, 0 if not found.@\\
\mbox{}\verb@.globl _kcc_ht_get@\\
\mbox{}\verb@.p2align 2@\\
\mbox{}\verb@_kcc_ht_get:@\\
\mbox{}\verb@    stp x29, x30, [sp, #-80]!@\\
\mbox{}\verb@    mov x29, sp@\\
\mbox{}\verb@    stp x19, x20, [sp, #16]@\\
\mbox{}\verb@    stp x21, x22, [sp, #32]@\\
\mbox{}\verb@    stp x23, x24, [sp, #48]@\\
\mbox{}\verb@    str x25, [sp, #64]@\\
\mbox{}\verb@@\\
\mbox{}\verb@    mov x19, x0                  // x19 = ht@\\
\mbox{}\verb@    mov x20, x1                  // x20 = key@\\
\mbox{}\verb@    mov x21, x2                  // x21 = keylen@\\
\mbox{}\verb@@\\
\mbox{}\verb@    // Hash the key@\\
\mbox{}\verb@    mov x0, x1@\\
\mbox{}\verb@    mov x1, x2@\\
\mbox{}\verb@    bl _kcc_fnv1a@\\
\mbox{}\verb@    mov x22, x0                  // x22 = hash@\\
\mbox{}\verb@@\\
\mbox{}\verb@    // Load table fields@\\
\mbox{}\verb@    ldr x23, [x19, #HT_ENTRIES] // x23 = entries base@\\
\mbox{}\verb@    ldr x24, [x19, #HT_CAP]    // x24 = capacity@\\
\mbox{}\verb@    sub x3, x24, #1             // x3 = mask (cap - 1)@\\
\mbox{}\verb@    and x25, x22, x3            // x25 = index (callee-saved!)@\\
\mbox{}\verb@@\\
\mbox{}\verb@.Lht_get_probe:@\\
\mbox{}\verb@    // Compute entry address: entries + index * HTE_SIZE@\\
\mbox{}\verb@    mov x5, #HTE_SIZE@\\
\mbox{}\verb@    mul x6, x25, x5@\\
\mbox{}\verb@    add x6, x23, x6             // x6 = entry pointer@\\
\mbox{}\verb@@\\
\mbox{}\verb@    // Check if slot is empty (key == NULL)@\\
\mbox{}\verb@    ldr x7, [x6, #HTE_KEY]@\\
\mbox{}\verb@    cbz x7, .Lht_get_not_found@\\
\mbox{}\verb@@\\
\mbox{}\verb@    // Check hash match first (fast reject)@\\
\mbox{}\verb@    ldr x8, [x6, #HTE_HASH]@\\
\mbox{}\verb@    cmp x8, x22@\\
\mbox{}\verb@    b.ne .Lht_get_next@\\
\mbox{}\verb@@\\
\mbox{}\verb@    // Check key length@\\
\mbox{}\verb@    ldr x8, [x6, #HTE_KEYLEN]@\\
\mbox{}\verb@    cmp x8, x21@\\
\mbox{}\verb@    b.ne .Lht_get_next@\\
\mbox{}\verb@@\\
\mbox{}\verb@    // Compare key bytes@\\
\mbox{}\verb@    mov x0, x7                   // entry key@\\
\mbox{}\verb@    mov x1, x20                  // search key@\\
\mbox{}\verb@    mov x2, x21                  // length@\\
\mbox{}\verb@    bl _kcc_memcmp@\\
\mbox{}\verb@    cbnz x0, .Lht_get_next@\\
\mbox{}\verb@@\\
\mbox{}\verb@    // Found! Return value (x25 = index, callee-saved across memcmp)@\\
\mbox{}\verb@    mov x5, #HTE_SIZE@\\
\mbox{}\verb@    mul x6, x25, x5@\\
\mbox{}\verb@    add x6, x23, x6@\\
\mbox{}\verb@    ldr x0, [x6, #HTE_VALUE]@\\
\mbox{}\verb@    b .Lht_get_done@\\
\mbox{}\verb@@\\
\mbox{}\verb@.Lht_get_next:@\\
\mbox{}\verb@    add x25, x25, #1@\\
\mbox{}\verb@    sub x3, x24, #1@\\
\mbox{}\verb@    and x25, x25, x3            // wrap around@\\
\mbox{}\verb@    b .Lht_get_probe@\\
\mbox{}\verb@@\\
\mbox{}\verb@.Lht_get_not_found:@\\
\mbox{}\verb@    mov x0, #0@\\
\mbox{}\verb@@\\
\mbox{}\verb@.Lht_get_done:@\\
\mbox{}\verb@    ldr x25, [sp, #64]@\\
\mbox{}\verb@    ldp x23, x24, [sp, #48]@\\
\mbox{}\verb@    ldp x21, x22, [sp, #32]@\\
\mbox{}\verb@    ldp x19, x20, [sp, #16]@\\
\mbox{}\verb@    ldp x29, x30, [sp], #80@\\
\mbox{}\verb@    ret@\\
\mbox{}\verb@@{\NWsep}
\end{list}
\vspace{-1.5ex}
\footnotesize
\begin{list}{}{\setlength{\itemsep}{-\parsep}\setlength{\itemindent}{-\leftmargin}}
\item \NWtxtMacroRefIn\ \NWlink{weft?}{?}.

\item{}
\end{list}
\end{minipage}\vspace{4ex}
\end{flushleft}
\subsection{Insertion}

\texttt{kcc\_ht\_put} inserts a key-value pair.  If the key already
exists, the value is updated.  No resizing is performed --- the caller
must ensure the table has sufficient capacity.  For the keyword table
(44~entries in a 128-slot table, 34\% load factor), this is always safe.

\begin{flushleft} \small
\begin{minipage}{\linewidth}\label{scrap22}\raggedright\small
\NWtarget{weft30}{} $\langle\,${\itshape Hash: put}\nobreak\ {\footnotesize {30}}$\,\rangle\equiv$
\vspace{-1ex}
\begin{list}{}{} \item
\mbox{}\verb@@\\
\mbox{}\verb@// -- kcc_ht_put -----------------------------------------------------@\\
\mbox{}\verb@// void kcc_ht_put(void *ht, const void *key, uint64_t keylen, uint64_t value)@\\
\mbox{}\verb@//   x0 = hash table pointer@\\
\mbox{}\verb@//   x1 = key pointer@\\
\mbox{}\verb@//   x2 = key length@\\
\mbox{}\verb@//   x3 = value to store@\\
\mbox{}\verb@// Inserts or updates the key-value pair.@\\
\mbox{}\verb@.globl _kcc_ht_put@\\
\mbox{}\verb@.p2align 2@\\
\mbox{}\verb@_kcc_ht_put:@\\
\mbox{}\verb@    stp x29, x30, [sp, #-96]!@\\
\mbox{}\verb@    mov x29, sp@\\
\mbox{}\verb@    stp x19, x20, [sp, #16]@\\
\mbox{}\verb@    stp x21, x22, [sp, #32]@\\
\mbox{}\verb@    stp x23, x24, [sp, #48]@\\
\mbox{}\verb@    stp x25, x26, [sp, #64]@\\
\mbox{}\verb@    str x27, [sp, #80]@\\
\mbox{}\verb@@\\
\mbox{}\verb@    mov x19, x0                  // x19 = ht@\\
\mbox{}\verb@    mov x20, x1                  // x20 = key@\\
\mbox{}\verb@    mov x21, x2                  // x21 = keylen@\\
\mbox{}\verb@    mov x25, x3                  // x25 = value@\\
\mbox{}\verb@@\\
\mbox{}\verb@    // Hash the key@\\
\mbox{}\verb@    mov x0, x1@\\
\mbox{}\verb@    mov x1, x2@\\
\mbox{}\verb@    bl _kcc_fnv1a@\\
\mbox{}\verb@    mov x22, x0                  // x22 = hash@\\
\mbox{}\verb@@\\
\mbox{}\verb@    // Load table fields@\\
\mbox{}\verb@    ldr x23, [x19, #HT_ENTRIES] // x23 = entries base@\\
\mbox{}\verb@    ldr x24, [x19, #HT_CAP]    // x24 = capacity@\\
\mbox{}\verb@    sub x3, x24, #1             // x3 = mask@\\
\mbox{}\verb@    and x27, x22, x3            // x27 = index (callee-saved!)@\\
\mbox{}\verb@@\\
\mbox{}\verb@.Lht_put_probe:@\\
\mbox{}\verb@    mov x5, #HTE_SIZE@\\
\mbox{}\verb@    mul x6, x27, x5@\\
\mbox{}\verb@    add x6, x23, x6             // x6 = entry pointer@\\
\mbox{}\verb@@\\
\mbox{}\verb@    // Empty slot? Insert here@\\
\mbox{}\verb@    ldr x7, [x6, #HTE_KEY]@\\
\mbox{}\verb@    cbz x7, .Lht_put_insert@\\
\mbox{}\verb@@\\
\mbox{}\verb@    // Check if same key (update case)@\\
\mbox{}\verb@    ldr x8, [x6, #HTE_HASH]@\\
\mbox{}\verb@    cmp x8, x22@\\
\mbox{}\verb@    b.ne .Lht_put_next@\\
\mbox{}\verb@@\\
\mbox{}\verb@    ldr x8, [x6, #HTE_KEYLEN]@\\
\mbox{}\verb@    cmp x8, x21@\\
\mbox{}\verb@    b.ne .Lht_put_next@\\
\mbox{}\verb@@\\
\mbox{}\verb@    mov x26, x6                 // save entry ptr (callee-saved)@\\
\mbox{}\verb@    mov x0, x7@\\
\mbox{}\verb@    mov x1, x20@\\
\mbox{}\verb@    mov x2, x21@\\
\mbox{}\verb@    bl _kcc_memcmp@\\
\mbox{}\verb@    cbnz x0, .Lht_put_next@\\
\mbox{}\verb@@\\
\mbox{}\verb@    // Key exists: update value@\\
\mbox{}\verb@    str x25, [x26, #HTE_VALUE]@\\
\mbox{}\verb@    b .Lht_put_done@\\
\mbox{}\verb@@\\
\mbox{}\verb@.Lht_put_next:@\\
\mbox{}\verb@    add x27, x27, #1@\\
\mbox{}\verb@    sub x3, x24, #1@\\
\mbox{}\verb@    and x27, x27, x3@\\
\mbox{}\verb@    b .Lht_put_probe@\\
\mbox{}\verb@@\\
\mbox{}\verb@.Lht_put_insert:@\\
\mbox{}\verb@    // x6 = empty entry pointer (freshly computed, no call since)@\\
\mbox{}\verb@    str x22, [x6, #HTE_HASH]@\\
\mbox{}\verb@    str x20, [x6, #HTE_KEY]@\\
\mbox{}\verb@    str x21, [x6, #HTE_KEYLEN]@\\
\mbox{}\verb@    str x25, [x6, #HTE_VALUE]@\\
\mbox{}\verb@@\\
\mbox{}\verb@    // Increment count@\\
\mbox{}\verb@    ldr x0, [x19, #HT_COUNT]@\\
\mbox{}\verb@    add x0, x0, #1@\\
\mbox{}\verb@    str x0, [x19, #HT_COUNT]@\\
\mbox{}\verb@@\\
\mbox{}\verb@.Lht_put_done:@\\
\mbox{}\verb@    ldr x27, [sp, #80]@\\
\mbox{}\verb@    ldp x25, x26, [sp, #64]@\\
\mbox{}\verb@    ldp x23, x24, [sp, #48]@\\
\mbox{}\verb@    ldp x21, x22, [sp, #32]@\\
\mbox{}\verb@    ldp x19, x20, [sp, #16]@\\
\mbox{}\verb@    ldp x29, x30, [sp], #96@\\
\mbox{}\verb@    ret@\\
\mbox{}\verb@@{\NWsep}
\end{list}
\vspace{-1.5ex}
\footnotesize
\begin{list}{}{\setlength{\itemsep}{-\parsep}\setlength{\itemindent}{-\leftmargin}}
\item \NWtxtMacroRefIn\ \NWlink{weft?}{?}.

\item{}
\end{list}
\end{minipage}\vspace{4ex}
\end{flushleft}
\section{Dynamic Vector}

A vector is a growable array of fixed-size elements.  When the vector
is full, it doubles in capacity by allocating a new array from the
arena and copying the old contents (the old array is wasted, as is
typical with arena allocators).

\begin{center}
\begin{tabular}{rll}
\hline
Offset & Field & Description \\
\hline
0  & \texttt{data}      & Pointer to element array \\
8  & \texttt{count}     & Number of elements stored \\
16 & \texttt{cap}       & Capacity (number of elements) \\
24 & \texttt{elem\_size} & Size of each element in bytes \\
\hline
\end{tabular}
\end{center}

\begin{flushleft} \small
\begin{minipage}{\linewidth}\label{scrap23}\raggedright\small
\NWtarget{weft31}{} $\langle\,${\itshape Vec: constants}\nobreak\ {\footnotesize {31}}$\,\rangle\equiv$
\vspace{-1ex}
\begin{list}{}{} \item
\mbox{}\verb@@\\
\mbox{}\verb@// -- Vector layout --------------------------------------------------@\\
\mbox{}\verb@.equ VEC_DATA,      0         // pointer to element array@\\
\mbox{}\verb@.equ VEC_COUNT,     8         // number of elements@\\
\mbox{}\verb@.equ VEC_CAP,      16         // capacity@\\
\mbox{}\verb@.equ VEC_ELEMSIZE, 24         // element size in bytes@\\
\mbox{}\verb@.equ VEC_SIZE,     32         // total header size@\\
\mbox{}\verb@@{\NWsep}
\end{list}
\vspace{-1.5ex}
\footnotesize
\begin{list}{}{\setlength{\itemsep}{-\parsep}\setlength{\itemindent}{-\leftmargin}}
\item \NWtxtMacroRefIn\ \NWlink{weft?}{?}.

\item{}
\end{list}
\end{minipage}\vspace{4ex}
\end{flushleft}
\subsection{Initialization}

\begin{flushleft} \small
\begin{minipage}{\linewidth}\label{scrap24}\raggedright\small
\NWtarget{weft32}{} $\langle\,${\itshape Vec: init}\nobreak\ {\footnotesize {32}}$\,\rangle\equiv$
\vspace{-1ex}
\begin{list}{}{} \item
\mbox{}\verb@@\\
\mbox{}\verb@// -- kcc_vec_init ---------------------------------------------------@\\
\mbox{}\verb@// void *kcc_vec_init(uint64_t elem_size)@\\
\mbox{}\verb@//   x0 = element size in bytes@\\
\mbox{}\verb@// Returns pointer to vector header.  Initial capacity is 16.@\\
\mbox{}\verb@.globl _kcc_vec_init@\\
\mbox{}\verb@.p2align 2@\\
\mbox{}\verb@_kcc_vec_init:@\\
\mbox{}\verb@    stp x29, x30, [sp, #-48]!@\\
\mbox{}\verb@    mov x29, sp@\\
\mbox{}\verb@    stp x19, x20, [sp, #16]@\\
\mbox{}\verb@    str x21, [sp, #32]@\\
\mbox{}\verb@@\\
\mbox{}\verb@    mov x19, x0                  // x19 = elem_size@\\
\mbox{}\verb@@\\
\mbox{}\verb@    // Allocate header@\\
\mbox{}\verb@    mov x0, #VEC_SIZE@\\
\mbox{}\verb@    bl _kcc_arena_alloc@\\
\mbox{}\verb@    mov x20, x0                  // x20 = header@\\
\mbox{}\verb@@\\
\mbox{}\verb@    // Allocate initial data (16 * elem_size)@\\
\mbox{}\verb@    mov x21, #16                 // initial capacity@\\
\mbox{}\verb@    mov x0, x21@\\
\mbox{}\verb@    mul x0, x0, x19@\\
\mbox{}\verb@    bl _kcc_arena_alloc@\\
\mbox{}\verb@@\\
\mbox{}\verb@    // Fill header@\\
\mbox{}\verb@    str  x0, [x20, #VEC_DATA]@\\
\mbox{}\verb@    str xzr, [x20, #VEC_COUNT]@\\
\mbox{}\verb@    str x21, [x20, #VEC_CAP]@\\
\mbox{}\verb@    str x19, [x20, #VEC_ELEMSIZE]@\\
\mbox{}\verb@@\\
\mbox{}\verb@    mov x0, x20                  // return header@\\
\mbox{}\verb@@\\
\mbox{}\verb@    ldr x21, [sp, #32]@\\
\mbox{}\verb@    ldp x19, x20, [sp, #16]@\\
\mbox{}\verb@    ldp x29, x30, [sp], #48@\\
\mbox{}\verb@    ret@\\
\mbox{}\verb@@{\NWsep}
\end{list}
\vspace{-1.5ex}
\footnotesize
\begin{list}{}{\setlength{\itemsep}{-\parsep}\setlength{\itemindent}{-\leftmargin}}
\item \NWtxtMacroRefIn\ \NWlink{weft?}{?}.

\item{}
\end{list}
\end{minipage}\vspace{4ex}
\end{flushleft}
\subsection{Push}

Append an element.  If the vector is full, a new array of double
capacity is allocated from the arena, and the old contents are copied.

\begin{flushleft} \small
\begin{minipage}{\linewidth}\label{scrap25}\raggedright\small
\NWtarget{weft36}{} $\langle\,${\itshape Vec: push}\nobreak\ {\footnotesize {36}}$\,\rangle\equiv$
\vspace{-1ex}
\begin{list}{}{} \item
\mbox{}\verb@@\\
\mbox{}\verb@// -- kcc_vec_push ---------------------------------------------------@\\
\mbox{}\verb@// void kcc_vec_push(void *vec, const void *elem)@\\
\mbox{}\verb@//   x0 = vector pointer@\\
\mbox{}\verb@//   x1 = pointer to element to copy@\\
\mbox{}\verb@// Appends a copy of elem to the vector.@\\
\mbox{}\verb@.globl _kcc_vec_push@\\
\mbox{}\verb@.p2align 2@\\
\mbox{}\verb@_kcc_vec_push:@\\
\mbox{}\verb@    stp x29, x30, [sp, #-48]!@\\
\mbox{}\verb@    mov x29, sp@\\
\mbox{}\verb@    stp x19, x20, [sp, #16]@\\
\mbox{}\verb@    stp x21, x22, [sp, #32]@\\
\mbox{}\verb@@\\
\mbox{}\verb@    mov x19, x0                  // x19 = vec@\\
\mbox{}\verb@    mov x20, x1                  // x20 = elem@\\
\mbox{}\verb@@\\
\mbox{}\verb@    ldr x21, [x19, #VEC_COUNT]  // x21 = count@\\
\mbox{}\verb@    ldr x22, [x19, #VEC_CAP]   // x22 = cap@\\
\mbox{}\verb@@\\
\mbox{}\verb@    cmp x21, x22@\\
\mbox{}\verb@    b.lo .Lvec_push_store@\\
\mbox{}\verb@@\\
\mbox{}\verb@    // Grow: double capacity@\\
\mbox{}\verb@    lsl x22, x22, #1            // new_cap = cap * 2@\\
\mbox{}\verb@    ldr x0, [x19, #VEC_ELEMSIZE]@\\
\mbox{}\verb@    mul x0, x22, x0             // new_cap * elem_size@\\
\mbox{}\verb@    bl _kcc_arena_alloc@\\
\mbox{}\verb@    // x0 = new array@\\
\mbox{}\verb@@\\
\mbox{}\verb@    // Copy old data@\\
\mbox{}\verb@    mov x3, x0                  // save new array@\\
\mbox{}\verb@    ldr x1, [x19, #VEC_DATA]   // src = old array@\\
\mbox{}\verb@    mov x4, x1                  // save old for memcpy dst/src swap@\\
\mbox{}\verb@    ldr x2, [x19, #VEC_ELEMSIZE]@\\
\mbox{}\verb@    mul x2, x21, x2             // old count * elem_size@\\
\mbox{}\verb@    mov x0, x3                  // dst = new array@\\
\mbox{}\verb@    mov x1, x4                  // src = old array@\\
\mbox{}\verb@    bl _kcc_memcpy@\\
\mbox{}\verb@@\\
\mbox{}\verb@    // Update vec header@\\
\mbox{}\verb@    str x3, [x19, #VEC_DATA]@\\
\mbox{}\verb@    str x22, [x19, #VEC_CAP]@\\
\mbox{}\verb@@\\
\mbox{}\verb@.Lvec_push_store:@\\
\mbox{}\verb@    // Compute dest: data + count * elem_size@\\
\mbox{}\verb@    ldr x0, [x19, #VEC_DATA]@\\
\mbox{}\verb@    ldr x2, [x19, #VEC_ELEMSIZE]@\\
\mbox{}\verb@    mul x3, x21, x2@\\
\mbox{}\verb@    add x0, x0, x3              // dst@\\
\mbox{}\verb@@\\
\mbox{}\verb@    mov x1, x20                 // src = elem@\\
\mbox{}\verb@    // x2 = elem_size (still loaded)@\\
\mbox{}\verb@    bl _kcc_memcpy@\\
\mbox{}\verb@@\\
\mbox{}\verb@    // Increment count@\\
\mbox{}\verb@    add x21, x21, #1@\\
\mbox{}\verb@    str x21, [x19, #VEC_COUNT]@\\
\mbox{}\verb@@\\
\mbox{}\verb@    ldp x21, x22, [sp, #32]@\\
\mbox{}\verb@    ldp x19, x20, [sp, #16]@\\
\mbox{}\verb@    ldp x29, x30, [sp], #48@\\
\mbox{}\verb@    ret@\\
\mbox{}\verb@@{\NWsep}
\end{list}
\vspace{-1.5ex}
\footnotesize
\begin{list}{}{\setlength{\itemsep}{-\parsep}\setlength{\itemindent}{-\leftmargin}}
\item \NWtxtMacroRefIn\ \NWlink{weft?}{?}.

\item{}
\end{list}
\end{minipage}\vspace{4ex}
\end{flushleft}
\subsection{Element Access}

\begin{flushleft} \small
\begin{minipage}{\linewidth}\label{scrap26}\raggedright\small
\NWtarget{weft37}{} $\langle\,${\itshape Vec: get}\nobreak\ {\footnotesize {37}}$\,\rangle\equiv$
\vspace{-1ex}
\begin{list}{}{} \item
\mbox{}\verb@@\\
\mbox{}\verb@// -- kcc_vec_get ----------------------------------------------------@\\
\mbox{}\verb@// void *kcc_vec_get(void *vec, uint64_t index)@\\
\mbox{}\verb@//   x0 = vector pointer@\\
\mbox{}\verb@//   x1 = element index@\\
\mbox{}\verb@// Returns pointer to element at the given index.@\\
\mbox{}\verb@// No bounds checking (caller's responsibility).@\\
\mbox{}\verb@.globl _kcc_vec_get@\\
\mbox{}\verb@.p2align 2@\\
\mbox{}\verb@_kcc_vec_get:@\\
\mbox{}\verb@    stp x29, x30, [sp, #-16]!@\\
\mbox{}\verb@    mov x29, sp@\\
\mbox{}\verb@@\\
\mbox{}\verb@    ldr x2, [x0, #VEC_ELEMSIZE]@\\
\mbox{}\verb@    ldr x3, [x0, #VEC_DATA]@\\
\mbox{}\verb@    mul x1, x1, x2@\\
\mbox{}\verb@    add x0, x3, x1              // data + index * elem_size@\\
\mbox{}\verb@@\\
\mbox{}\verb@    ldp x29, x30, [sp], #16@\\
\mbox{}\verb@    ret@\\
\mbox{}\verb@@{\NWsep}
\end{list}
\vspace{-1.5ex}
\footnotesize
\begin{list}{}{\setlength{\itemsep}{-\parsep}\setlength{\itemindent}{-\leftmargin}}
\item \NWtxtMacroRefIn\ \NWlink{weft?}{?}.

\item{}
\end{list}
\end{minipage}\vspace{4ex}
\end{flushleft}
\section{String Interning}

String interning ensures that each unique string is stored exactly
once in the arena.  Two interned strings can be compared for equality
by pointer comparison rather than byte-by-byte --- a critical
optimization for the symbol table.

The intern table is a hash table mapping string contents to their
canonical arena-allocated copy.

\begin{flushleft} \small
\begin{minipage}{\linewidth}\label{scrap27}\raggedright\small
\NWtarget{weft39a}{} $\langle\,${\itshape Strtab: init}\nobreak\ {\footnotesize {39a}}$\,\rangle\equiv$
\vspace{-1ex}
\begin{list}{}{} \item
\mbox{}\verb@@\\
\mbox{}\verb@// -- kcc_strtab_init ------------------------------------------------@\\
\mbox{}\verb@// void kcc_strtab_init(void)@\\
\mbox{}\verb@// Creates the global string intern table.@\\
\mbox{}\verb@.globl _kcc_strtab_init@\\
\mbox{}\verb@.p2align 2@\\
\mbox{}\verb@_kcc_strtab_init:@\\
\mbox{}\verb@    stp x29, x30, [sp, #-16]!@\\
\mbox{}\verb@    mov x29, sp@\\
\mbox{}\verb@@\\
\mbox{}\verb@    mov x0, #256                 // initial capacity@\\
\mbox{}\verb@    bl _kcc_ht_init@\\
\mbox{}\verb@@\\
\mbox{}\verb@    adrp x1, _kcc_strtab@{\tt @}\verb@PAGE@\\
\mbox{}\verb@    add  x1, x1, _kcc_strtab@{\tt @}\verb@PAGEOFF@\\
\mbox{}\verb@    str  x0, [x1]               // store table pointer@\\
\mbox{}\verb@@\\
\mbox{}\verb@    ldp x29, x30, [sp], #16@\\
\mbox{}\verb@    ret@\\
\mbox{}\verb@@{\NWsep}
\end{list}
\vspace{-1.5ex}
\footnotesize
\begin{list}{}{\setlength{\itemsep}{-\parsep}\setlength{\itemindent}{-\leftmargin}}
\item \NWtxtMacroRefIn\ \NWlink{weft?}{?}.

\item{}
\end{list}
\end{minipage}\vspace{4ex}
\end{flushleft}
\texttt{kcc\_strtab\_intern} looks up a string in the table.  If found,
it returns the existing interned pointer.  If not, it copies the string
into the arena, inserts it, and returns the new pointer.

\begin{flushleft} \small
\begin{minipage}{\linewidth}\label{scrap28}\raggedright\small
\NWtarget{weft39b}{} $\langle\,${\itshape Strtab: intern}\nobreak\ {\footnotesize {39b}}$\,\rangle\equiv$
\vspace{-1ex}
\begin{list}{}{} \item
\mbox{}\verb@@\\
\mbox{}\verb@// -- kcc_strtab_intern ----------------------------------------------@\\
\mbox{}\verb@// const char *kcc_strtab_intern(const char *s, uint64_t len)@\\
\mbox{}\verb@//   x0 = pointer to string (not necessarily null-terminated)@\\
\mbox{}\verb@//   x1 = string length@\\
\mbox{}\verb@// Returns pointer to the interned (null-terminated) copy.@\\
\mbox{}\verb@.globl _kcc_strtab_intern@\\
\mbox{}\verb@.p2align 2@\\
\mbox{}\verb@_kcc_strtab_intern:@\\
\mbox{}\verb@    stp x29, x30, [sp, #-48]!@\\
\mbox{}\verb@    mov x29, sp@\\
\mbox{}\verb@    stp x19, x20, [sp, #16]@\\
\mbox{}\verb@    str x21, [sp, #32]@\\
\mbox{}\verb@@\\
\mbox{}\verb@    mov x19, x0                  // x19 = string@\\
\mbox{}\verb@    mov x20, x1                  // x20 = length@\\
\mbox{}\verb@@\\
\mbox{}\verb@    // Look up in the table@\\
\mbox{}\verb@    adrp x0, _kcc_strtab@{\tt @}\verb@PAGE@\\
\mbox{}\verb@    add  x0, x0, _kcc_strtab@{\tt @}\verb@PAGEOFF@\\
\mbox{}\verb@    ldr  x0, [x0]               // x0 = table pointer@\\
\mbox{}\verb@    mov  x1, x19@\\
\mbox{}\verb@    mov  x2, x20@\\
\mbox{}\verb@    bl _kcc_ht_get@\\
\mbox{}\verb@    cbnz x0, .Lstrtab_found     // already interned@\\
\mbox{}\verb@@\\
\mbox{}\verb@    // Not found: allocate len+1 bytes and copy@\\
\mbox{}\verb@    add x0, x20, #1@\\
\mbox{}\verb@    bl _kcc_arena_alloc@\\
\mbox{}\verb@    mov x21, x0                  // x21 = new buffer@\\
\mbox{}\verb@@\\
\mbox{}\verb@    // Copy string bytes@\\
\mbox{}\verb@    mov x1, x19                  // src@\\
\mbox{}\verb@    mov x2, x20                  // count@\\
\mbox{}\verb@    bl _kcc_memcpy@\\
\mbox{}\verb@@\\
\mbox{}\verb@    // Null-terminate@\\
\mbox{}\verb@    strb wzr, [x21, x20]@\\
\mbox{}\verb@@\\
\mbox{}\verb@    // Insert into table@\\
\mbox{}\verb@    adrp x0, _kcc_strtab@{\tt @}\verb@PAGE@\\
\mbox{}\verb@    add  x0, x0, _kcc_strtab@{\tt @}\verb@PAGEOFF@\\
\mbox{}\verb@    ldr  x0, [x0]@\\
\mbox{}\verb@    mov  x1, x21                 // key = interned copy@\\
\mbox{}\verb@    mov  x2, x20                 // keylen@\\
\mbox{}\verb@    mov  x3, x21                 // value = same pointer@\\
\mbox{}\verb@    bl _kcc_ht_put@\\
\mbox{}\verb@@\\
\mbox{}\verb@    mov x0, x21                  // return interned pointer@\\
\mbox{}\verb@@\\
\mbox{}\verb@.Lstrtab_found:@\\
\mbox{}\verb@    ldr x21, [sp, #32]@\\
\mbox{}\verb@    ldp x19, x20, [sp, #16]@\\
\mbox{}\verb@    ldp x29, x30, [sp], #48@\\
\mbox{}\verb@    ret@\\
\mbox{}\verb@@{\NWsep}
\end{list}
\vspace{-1.5ex}
\footnotesize
\begin{list}{}{\setlength{\itemsep}{-\parsep}\setlength{\itemindent}{-\leftmargin}}
\item \NWtxtMacroRefIn\ \NWlink{weft?}{?}.

\item{}
\end{list}
\end{minipage}\vspace{4ex}
\end{flushleft}
\subsection{Global State}

\begin{flushleft} \small
\begin{minipage}{\linewidth}\label{scrap29}\raggedright\small
\NWtarget{weft39c}{} $\langle\,${\itshape Strtab: data}\nobreak\ {\footnotesize {39c}}$\,\rangle\equiv$
\vspace{-1ex}
\begin{list}{}{} \item
\mbox{}\verb@@\\
\mbox{}\verb@// -- String intern table global state -------------------------------@\\
\mbox{}\verb@.data@\\
\mbox{}\verb@.p2align 3@\\
\mbox{}\verb@_kcc_strtab: .quad 0             // pointer to hash table@\\
\mbox{}\verb@@{\NWsep}
\end{list}
\vspace{-1.5ex}
\footnotesize
\begin{list}{}{\setlength{\itemsep}{-\parsep}\setlength{\itemindent}{-\leftmargin}}
\item \NWtxtMacroRefIn\ \NWlink{weft?}{?}.

\item{}
\end{list}
\end{minipage}\vspace{4ex}
\end{flushleft}
% 11-string-operations.weft — String Operations in ARM64 Assembly
%
% Since kcc uses no C library, we provide our own implementations
% of the fundamental string and memory operations.  These are used
% throughout the compiler: by the lexer to compare identifiers,
% by the arena to zero-fill memory, by the I/O layer to copy
% output into buffers.
%
% Each function follows AAPCS64 conventions and includes the
% mandatory Apple frame pointer.

\chapter{String Operations}
\label{ch:strings}

\section{String Length}

Compute the length of a null-terminated string by scanning for the
zero byte.  This is the moral equivalent of \texttt{strlen(3)}.

\begin{flushleft} \small
\begin{minipage}{\linewidth}\label{scrap30}\raggedright\small
\NWtarget{weft40a}{} $\langle\,${\itshape String: strlen}\nobreak\ {\footnotesize {40a}}$\,\rangle\equiv$
\vspace{-1ex}
\begin{list}{}{} \item
\mbox{}\verb@@\\
\mbox{}\verb@// -- kcc_strlen -----------------------------------------------------@\\
\mbox{}\verb@// size_t kcc_strlen(const char *s)@\\
\mbox{}\verb@//   x0 = pointer to null-terminated string@\\
\mbox{}\verb@// Returns length (not counting the null terminator) in x0.@\\
\mbox{}\verb@FUNC kcc_strlen@\\
\mbox{}\verb@    mov x1, x0                   // x1 = scan pointer@\\
\mbox{}\verb@.Lstrlen_loop:@\\
\mbox{}\verb@    ldrb w2, [x1], #1            // load byte, advance@\\
\mbox{}\verb@    cbnz w2, .Lstrlen_loop       // loop until null@\\
\mbox{}\verb@    sub x0, x1, x0               // length = end - start@\\
\mbox{}\verb@    sub x0, x0, #1               // exclude the null byte itself@\\
\mbox{}\verb@ENDFUNC@\\
\mbox{}\verb@@{\NWsep}
\end{list}
\vspace{-1.5ex}
\footnotesize
\begin{list}{}{\setlength{\itemsep}{-\parsep}\setlength{\itemindent}{-\leftmargin}}
\item \NWtxtMacroRefIn\ \NWlink{weft?}{?}.

\item{}
\end{list}
\end{minipage}\vspace{4ex}
\end{flushleft}
\section{Memory Copy}

Copy \texttt{n} bytes from source to destination.  Overlapping
regions produce undefined results (as with \texttt{memcpy}, not
\texttt{memmove}).

\begin{flushleft} \small
\begin{minipage}{\linewidth}\label{scrap31}\raggedright\small
\NWtarget{weft40b}{} $\langle\,${\itshape String: memcpy}\nobreak\ {\footnotesize {40b}}$\,\rangle\equiv$
\vspace{-1ex}
\begin{list}{}{} \item
\mbox{}\verb@@\\
\mbox{}\verb@// -- kcc_memcpy -----------------------------------------------------@\\
\mbox{}\verb@// void *kcc_memcpy(void *dst, const void *src, size_t n)@\\
\mbox{}\verb@//   x0 = destination@\\
\mbox{}\verb@//   x1 = source@\\
\mbox{}\verb@//   x2 = byte count@\\
\mbox{}\verb@// Returns dst in x0.@\\
\mbox{}\verb@FUNC kcc_memcpy@\\
\mbox{}\verb@    mov x3, x0                   // save dst for return value@\\
\mbox{}\verb@    cbz x2, .Lmemcpy_done        // nothing to copy@\\
\mbox{}\verb@.Lmemcpy_loop:@\\
\mbox{}\verb@    ldrb w4, [x1], #1            // load from src@\\
\mbox{}\verb@    strb w4, [x0], #1            // store to dst@\\
\mbox{}\verb@    subs x2, x2, #1@\\
\mbox{}\verb@    b.ne .Lmemcpy_loop@\\
\mbox{}\verb@.Lmemcpy_done:@\\
\mbox{}\verb@    mov x0, x3                   // return original dst@\\
\mbox{}\verb@ENDFUNC@\\
\mbox{}\verb@@{\NWsep}
\end{list}
\vspace{-1.5ex}
\footnotesize
\begin{list}{}{\setlength{\itemsep}{-\parsep}\setlength{\itemindent}{-\leftmargin}}
\item \NWtxtMacroRefIn\ \NWlink{weft?}{?}.

\item{}
\end{list}
\end{minipage}\vspace{4ex}
\end{flushleft}
\section{Memory Compare}

Compare \texttt{n} bytes.  Returns 0 if equal, negative if the first
differing byte in \texttt{s1} is less than in \texttt{s2}, positive
otherwise.

\begin{flushleft} \small
\begin{minipage}{\linewidth}\label{scrap32}\raggedright\small
\NWtarget{weft?}{} $\langle\,${\itshape String: memcmp}\nobreak\ {\footnotesize {?}}$\,\rangle\equiv$
\vspace{-1ex}
\begin{list}{}{} \item
\mbox{}\verb@@\\
\mbox{}\verb@// -- kcc_memcmp -----------------------------------------------------@\\
\mbox{}\verb@// int kcc_memcmp(const void *s1, const void *s2, size_t n)@\\
\mbox{}\verb@//   x0 = first buffer@\\
\mbox{}\verb@//   x1 = second buffer@\\
\mbox{}\verb@//   x2 = byte count@\\
\mbox{}\verb@// Returns comparison result in x0.@\\
\mbox{}\verb@FUNC kcc_memcmp@\\
\mbox{}\verb@    cbz x2, .Lmemcmp_equal       // zero length → equal@\\
\mbox{}\verb@.Lmemcmp_loop:@\\
\mbox{}\verb@    ldrb w3, [x0], #1@\\
\mbox{}\verb@    ldrb w4, [x1], #1@\\
\mbox{}\verb@    cmp w3, w4@\\
\mbox{}\verb@    b.ne .Lmemcmp_diff@\\
\mbox{}\verb@    subs x2, x2, #1@\\
\mbox{}\verb@    b.ne .Lmemcmp_loop@\\
\mbox{}\verb@.Lmemcmp_equal:@\\
\mbox{}\verb@    mov x0, #0@\\
\mbox{}\verb@    b .Lmemcmp_ret@\\
\mbox{}\verb@.Lmemcmp_diff:@\\
\mbox{}\verb@    sub x0, x3, x4               // s1[i] - s2[i]@\\
\mbox{}\verb@.Lmemcmp_ret:@\\
\mbox{}\verb@ENDFUNC@\\
\mbox{}\verb@@{\NWsep}
\end{list}
\vspace{-1.5ex}
\footnotesize
\begin{list}{}{\setlength{\itemsep}{-\parsep}\setlength{\itemindent}{-\leftmargin}}
\item \NWtxtMacroRefIn\ \NWlink{weft?}{?}.

\item{}
\end{list}
\end{minipage}\vspace{4ex}
\end{flushleft}
\section{Memory Set}

Fill \texttt{n} bytes of memory with the value \texttt{c}.

\begin{flushleft} \small
\begin{minipage}{\linewidth}\label{scrap33}\raggedright\small
\NWtarget{weft?}{} $\langle\,${\itshape String: memset}\nobreak\ {\footnotesize {?}}$\,\rangle\equiv$
\vspace{-1ex}
\begin{list}{}{} \item
\mbox{}\verb@@\\
\mbox{}\verb@// -- kcc_memset -----------------------------------------------------@\\
\mbox{}\verb@// void *kcc_memset(void *s, int c, size_t n)@\\
\mbox{}\verb@//   x0 = destination@\\
\mbox{}\verb@//   x1 = fill byte (only low 8 bits used)@\\
\mbox{}\verb@//   x2 = byte count@\\
\mbox{}\verb@// Returns s in x0.@\\
\mbox{}\verb@FUNC kcc_memset@\\
\mbox{}\verb@    mov x3, x0                   // save dst for return value@\\
\mbox{}\verb@    cbz x2, .Lmemset_done@\\
\mbox{}\verb@.Lmemset_loop:@\\
\mbox{}\verb@    strb w1, [x0], #1@\\
\mbox{}\verb@    subs x2, x2, #1@\\
\mbox{}\verb@    b.ne .Lmemset_loop@\\
\mbox{}\verb@.Lmemset_done:@\\
\mbox{}\verb@    mov x0, x3                   // return original dst@\\
\mbox{}\verb@ENDFUNC@\\
\mbox{}\verb@@{\NWsep}
\end{list}
\vspace{-1.5ex}
\footnotesize
\begin{list}{}{\setlength{\itemsep}{-\parsep}\setlength{\itemindent}{-\leftmargin}}
\item \NWtxtMacroRefIn\ \NWlink{weft?}{?}.

\item{}
\end{list}
\end{minipage}\vspace{4ex}
\end{flushleft}
\section{String Compare}

Compare two null-terminated strings byte by byte.  Returns 0 if equal,
a negative value if the first differing byte in \texttt{s1} is less than
in \texttt{s2}, a positive value otherwise --- the same contract as
\texttt{strcmp(3)}.

\begin{flushleft} \small
\begin{minipage}{\linewidth}\label{scrap34}\raggedright\small
\NWtarget{weft?}{} $\langle\,${\itshape String: strcmp}\nobreak\ {\footnotesize {?}}$\,\rangle\equiv$
\vspace{-1ex}
\begin{list}{}{} \item
\mbox{}\verb@@\\
\mbox{}\verb@// -- kcc_strcmp ------------------------------------------------------@\\
\mbox{}\verb@// int kcc_strcmp(const char *s1, const char *s2)@\\
\mbox{}\verb@//   x0 = first string (null-terminated)@\\
\mbox{}\verb@//   x1 = second string (null-terminated)@\\
\mbox{}\verb@// Returns comparison result in x0.@\\
\mbox{}\verb@FUNC kcc_strcmp@\\
\mbox{}\verb@.Lstrcmp_loop:@\\
\mbox{}\verb@    ldrb w2, [x0], #1@\\
\mbox{}\verb@    ldrb w3, [x1], #1@\\
\mbox{}\verb@    cmp w2, w3@\\
\mbox{}\verb@    b.ne .Lstrcmp_diff@\\
\mbox{}\verb@    cbz w2, .Lstrcmp_equal        // both reached null -> equal@\\
\mbox{}\verb@    b .Lstrcmp_loop@\\
\mbox{}\verb@.Lstrcmp_equal:@\\
\mbox{}\verb@    mov x0, #0@\\
\mbox{}\verb@    b .Lstrcmp_ret@\\
\mbox{}\verb@.Lstrcmp_diff:@\\
\mbox{}\verb@    sub x0, x2, x3                // s1[i] - s2[i]@\\
\mbox{}\verb@.Lstrcmp_ret:@\\
\mbox{}\verb@ENDFUNC@\\
\mbox{}\verb@@{\NWsep}
\end{list}
\vspace{-1.5ex}
\footnotesize
\begin{list}{}{\setlength{\itemsep}{-\parsep}\setlength{\itemindent}{-\leftmargin}}
\item \NWtxtMacroRefIn\ \NWlink{weft?}{?}.

\item{}
\end{list}
\end{minipage}\vspace{4ex}
\end{flushleft}
% 12-io.weft — Buffered File I/O
%
% All output from kcc (diagnostics, generated code, debug traces)
% passes through this buffered I/O layer.  Rather than issuing a
% syscall for every character, we accumulate output in a 4 KiB buffer
% and flush it in bulk.  This dramatically reduces syscall overhead.
%
% For Phase 0, the buffer is statically allocated in .bss.  In later
% phases, we may switch to arena-allocated buffers to support multiple
% output streams simultaneously.

\chapter{Buffered I/O}
\label{ch:io}

\section{Buffer Layout}

The I/O state consists of four values: the file descriptor, the
buffer base address, the current write position, and the capacity.
For Phase~0, we use a single static buffer writing to \texttt{stdout}.

\begin{flushleft} \small
\begin{minipage}{\linewidth}\label{scrap35}\raggedright\small
\NWtarget{weft?}{} $\langle\,${\itshape IO: data}\nobreak\ {\footnotesize {?}}$\,\rangle\equiv$
\vspace{-1ex}
\begin{list}{}{} \item
\mbox{}\verb@@\\
\mbox{}\verb@// -- I/O global state -----------------------------------------------@\\
\mbox{}\verb@.data@\\
\mbox{}\verb@.p2align 3@\\
\mbox{}\verb@_kcc_io_fd:   .quad 1            // file descriptor (1 = stdout)@\\
\mbox{}\verb@_kcc_io_pos:  .quad 0            // current write position@\\
\mbox{}\verb@_kcc_io_cap:  .quad 4096         // buffer capacity@\\
\mbox{}\verb@@\\
\mbox{}\verb@.bss@\\
\mbox{}\verb@.p2align 4@\\
\mbox{}\verb@_kcc_io_buf:  .space 4096        // output buffer@\\
\mbox{}\verb@@{\NWsep}
\end{list}
\vspace{-1.5ex}
\footnotesize
\begin{list}{}{\setlength{\itemsep}{-\parsep}\setlength{\itemindent}{-\leftmargin}}
\item \NWtxtMacroRefIn\ \NWlink{weft?}{?}.

\item{}
\end{list}
\end{minipage}\vspace{4ex}
\end{flushleft}
\section{Flush}

\texttt{kcc\_io\_flush} writes the buffered contents to the file
descriptor and resets the position to zero.  This is called
automatically when the buffer fills, and explicitly before program
exit.

\begin{flushleft} \small
\begin{minipage}{\linewidth}\label{scrap36}\raggedright\small
\NWtarget{weft?}{} $\langle\,${\itshape IO: flush}\nobreak\ {\footnotesize {?}}$\,\rangle\equiv$
\vspace{-1ex}
\begin{list}{}{} \item
\mbox{}\verb@@\\
\mbox{}\verb@// -- kcc_io_flush ---------------------------------------------------@\\
\mbox{}\verb@// void kcc_io_flush(void)@\\
\mbox{}\verb@// Writes buffered data to the file descriptor, resets position.@\\
\mbox{}\verb@.globl _kcc_io_flush@\\
\mbox{}\verb@.p2align 2@\\
\mbox{}\verb@_kcc_io_flush:@\\
\mbox{}\verb@    stp x29, x30, [sp, #-32]!@\\
\mbox{}\verb@    mov x29, sp@\\
\mbox{}\verb@    str x19, [sp, #16]@\\
\mbox{}\verb@@\\
\mbox{}\verb@    // Load position — if zero, nothing to flush@\\
\mbox{}\verb@    adrp x19, _kcc_io_pos@{\tt @}\verb@PAGE@\\
\mbox{}\verb@    add  x19, x19, _kcc_io_pos@{\tt @}\verb@PAGEOFF@\\
\mbox{}\verb@    ldr  x2, [x19]               // x2 = bytes to write@\\
\mbox{}\verb@    cbz  x2, .Lflush_done@\\
\mbox{}\verb@@\\
\mbox{}\verb@    // Load fd@\\
\mbox{}\verb@    adrp x0, _kcc_io_fd@{\tt @}\verb@PAGE@\\
\mbox{}\verb@    add  x0, x0, _kcc_io_fd@{\tt @}\verb@PAGEOFF@\\
\mbox{}\verb@    ldr  x0, [x0]                // x0 = fd@\\
\mbox{}\verb@@\\
\mbox{}\verb@    // Load buffer base@\\
\mbox{}\verb@    adrp x1, _kcc_io_buf@{\tt @}\verb@PAGE@\\
\mbox{}\verb@    add  x1, x1, _kcc_io_buf@{\tt @}\verb@PAGEOFF  // x1 = buf@\\
\mbox{}\verb@@\\
\mbox{}\verb@    // x2 already has the byte count@\\
\mbox{}\verb@    bl _kcc_sys_write@\\
\mbox{}\verb@@\\
\mbox{}\verb@    // Reset position to 0@\\
\mbox{}\verb@    str xzr, [x19]@\\
\mbox{}\verb@@\\
\mbox{}\verb@.Lflush_done:@\\
\mbox{}\verb@    ldr x19, [sp, #16]@\\
\mbox{}\verb@    ldp x29, x30, [sp], #32@\\
\mbox{}\verb@    ret@\\
\mbox{}\verb@@{\NWsep}
\end{list}
\vspace{-1.5ex}
\footnotesize
\begin{list}{}{\setlength{\itemsep}{-\parsep}\setlength{\itemindent}{-\leftmargin}}
\item \NWtxtMacroRefIn\ \NWlink{weft?}{?}.

\item{}
\end{list}
\end{minipage}\vspace{4ex}
\end{flushleft}
\section{Put Character}

\texttt{kcc\_io\_putc} appends a single byte to the buffer.  If
the buffer is full, it flushes first.

\begin{flushleft} \small
\begin{minipage}{\linewidth}\label{scrap37}\raggedright\small
\NWtarget{weft?}{} $\langle\,${\itshape IO: putc}\nobreak\ {\footnotesize {?}}$\,\rangle\equiv$
\vspace{-1ex}
\begin{list}{}{} \item
\mbox{}\verb@@\\
\mbox{}\verb@// -- kcc_io_putc ----------------------------------------------------@\\
\mbox{}\verb@// void kcc_io_putc(int c)@\\
\mbox{}\verb@//   x0 = character to write (low 8 bits)@\\
\mbox{}\verb@FUNC kcc_io_putc@\\
\mbox{}\verb@    // Load buffer state@\\
\mbox{}\verb@    adrp x1, _kcc_io_pos@{\tt @}\verb@PAGE@\\
\mbox{}\verb@    add  x1, x1, _kcc_io_pos@{\tt @}\verb@PAGEOFF@\\
\mbox{}\verb@    ldr  x2, [x1]                // x2 = current position@\\
\mbox{}\verb@@\\
\mbox{}\verb@    adrp x3, _kcc_io_cap@{\tt @}\verb@PAGE@\\
\mbox{}\verb@    add  x3, x3, _kcc_io_cap@{\tt @}\verb@PAGEOFF@\\
\mbox{}\verb@    ldr  x3, [x3]                // x3 = capacity@\\
\mbox{}\verb@@\\
\mbox{}\verb@    // Flush if buffer is full@\\
\mbox{}\verb@    cmp x2, x3@\\
\mbox{}\verb@    b.lo .Lputc_store@\\
\mbox{}\verb@@\\
\mbox{}\verb@    // Save the character across the flush call@\\
\mbox{}\verb@    str x0, [sp, #-16]!@\\
\mbox{}\verb@    bl _kcc_io_flush@\\
\mbox{}\verb@    ldr x0, [sp], #16@\\
\mbox{}\verb@@\\
\mbox{}\verb@    // Reload position (should be 0 after flush)@\\
\mbox{}\verb@    adrp x1, _kcc_io_pos@{\tt @}\verb@PAGE@\\
\mbox{}\verb@    add  x1, x1, _kcc_io_pos@{\tt @}\verb@PAGEOFF@\\
\mbox{}\verb@    ldr  x2, [x1]@\\
\mbox{}\verb@@\\
\mbox{}\verb@.Lputc_store:@\\
\mbox{}\verb@    // Store the byte@\\
\mbox{}\verb@    adrp x3, _kcc_io_buf@{\tt @}\verb@PAGE@\\
\mbox{}\verb@    add  x3, x3, _kcc_io_buf@{\tt @}\verb@PAGEOFF@\\
\mbox{}\verb@    strb w0, [x3, x2]@\\
\mbox{}\verb@@\\
\mbox{}\verb@    // Increment position@\\
\mbox{}\verb@    add x2, x2, #1@\\
\mbox{}\verb@    str x2, [x1]@\\
\mbox{}\verb@ENDFUNC@\\
\mbox{}\verb@@{\NWsep}
\end{list}
\vspace{-1.5ex}
\footnotesize
\begin{list}{}{\setlength{\itemsep}{-\parsep}\setlength{\itemindent}{-\leftmargin}}
\item \NWtxtMacroRefIn\ \NWlink{weft?}{?}.

\item{}
\end{list}
\end{minipage}\vspace{4ex}
\end{flushleft}
\section{Put String}

\texttt{kcc\_io\_puts} writes a null-terminated string to the buffer.
It processes the string byte by byte, flushing the buffer whenever
it fills.  The null terminator is not written.

This is the primary output function used by the compiler driver,
the diagnostic engine, and (eventually) the code emitter.

\begin{flushleft} \small
\begin{minipage}{\linewidth}\label{scrap38}\raggedright\small
\NWtarget{weft?}{} $\langle\,${\itshape IO: puts}\nobreak\ {\footnotesize {?}}$\,\rangle\equiv$
\vspace{-1ex}
\begin{list}{}{} \item
\mbox{}\verb@@\\
\mbox{}\verb@// -- kcc_io_puts ----------------------------------------------------@\\
\mbox{}\verb@// void kcc_io_puts(const char *s)@\\
\mbox{}\verb@//   x0 = pointer to null-terminated string@\\
\mbox{}\verb@// Writes all bytes (excluding the null terminator) to the buffer.@\\
\mbox{}\verb@.globl _kcc_io_puts@\\
\mbox{}\verb@.p2align 2@\\
\mbox{}\verb@_kcc_io_puts:@\\
\mbox{}\verb@    stp x29, x30, [sp, #-48]!@\\
\mbox{}\verb@    mov x29, sp@\\
\mbox{}\verb@    stp x19, x20, [sp, #16]@\\
\mbox{}\verb@    stp x21, x22, [sp, #32]@\\
\mbox{}\verb@@\\
\mbox{}\verb@    mov x19, x0                  // x19 = string pointer@\\
\mbox{}\verb@@\\
\mbox{}\verb@    // Cache buffer state in callee-saved registers@\\
\mbox{}\verb@    adrp x20, _kcc_io_pos@{\tt @}\verb@PAGE@\\
\mbox{}\verb@    add  x20, x20, _kcc_io_pos@{\tt @}\verb@PAGEOFF@\\
\mbox{}\verb@    ldr  x21, [x20]              // x21 = current position@\\
\mbox{}\verb@@\\
\mbox{}\verb@    adrp x22, _kcc_io_buf@{\tt @}\verb@PAGE@\\
\mbox{}\verb@    add  x22, x22, _kcc_io_buf@{\tt @}\verb@PAGEOFF  // x22 = buffer base@\\
\mbox{}\verb@@\\
\mbox{}\verb@.Lputs_loop:@\\
\mbox{}\verb@    ldrb w0, [x19], #1           // load byte, advance string ptr@\\
\mbox{}\verb@    cbz w0, .Lputs_done          // null terminator → stop@\\
\mbox{}\verb@@\\
\mbox{}\verb@    // Store byte in buffer@\\
\mbox{}\verb@    strb w0, [x22, x21]@\\
\mbox{}\verb@    add x21, x21, #1@\\
\mbox{}\verb@@\\
\mbox{}\verb@    // Check if buffer is full (4096 bytes)@\\
\mbox{}\verb@    cmp x21, #4096@\\
\mbox{}\verb@    b.lo .Lputs_loop@\\
\mbox{}\verb@@\\
\mbox{}\verb@    // Buffer full — flush@\\
\mbox{}\verb@    str x21, [x20]               // save pos before flush@\\
\mbox{}\verb@    bl _kcc_io_flush@\\
\mbox{}\verb@    ldr x21, [x20]               // reload pos (should be 0)@\\
\mbox{}\verb@    b .Lputs_loop@\\
\mbox{}\verb@@\\
\mbox{}\verb@.Lputs_done:@\\
\mbox{}\verb@    str x21, [x20]               // save final position@\\
\mbox{}\verb@@\\
\mbox{}\verb@    ldp x21, x22, [sp, #32]@\\
\mbox{}\verb@    ldp x19, x20, [sp, #16]@\\
\mbox{}\verb@    ldp x29, x30, [sp], #48@\\
\mbox{}\verb@    ret@\\
\mbox{}\verb@@{\NWsep}
\end{list}
\vspace{-1.5ex}
\footnotesize
\begin{list}{}{\setlength{\itemsep}{-\parsep}\setlength{\itemindent}{-\leftmargin}}
\item \NWtxtMacroRefIn\ \NWlink{weft?}{?}.

\item{}
\end{list}
\end{minipage}\vspace{4ex}
\end{flushleft}
\section{Write N Bytes}

\texttt{kcc\_io\_write} writes exactly \texttt{n} bytes from a pointer
to the output buffer.  Unlike \texttt{puts}, it does not look for a null
terminator --- the caller supplies the byte count.  This is essential for
printing token text, which is a \texttt{(pointer, length)} pair into the
source buffer rather than a null-terminated string.

\begin{flushleft} \small
\begin{minipage}{\linewidth}\label{scrap39}\raggedright\small
\NWtarget{weft?}{} $\langle\,${\itshape IO: write}\nobreak\ {\footnotesize {?}}$\,\rangle\equiv$
\vspace{-1ex}
\begin{list}{}{} \item
\mbox{}\verb@@\\
\mbox{}\verb@// -- kcc_io_write ---------------------------------------------------@\\
\mbox{}\verb@// void kcc_io_write(const void *ptr, size_t len)@\\
\mbox{}\verb@//   x0 = pointer to data@\\
\mbox{}\verb@//   x1 = byte count@\\
\mbox{}\verb@// Writes exactly len bytes to the output buffer.@\\
\mbox{}\verb@.globl _kcc_io_write@\\
\mbox{}\verb@.p2align 2@\\
\mbox{}\verb@_kcc_io_write:@\\
\mbox{}\verb@    stp x29, x30, [sp, #-48]!@\\
\mbox{}\verb@    mov x29, sp@\\
\mbox{}\verb@    stp x19, x20, [sp, #16]@\\
\mbox{}\verb@    stp x21, x22, [sp, #32]@\\
\mbox{}\verb@@\\
\mbox{}\verb@    mov x19, x0                  // x19 = source pointer@\\
\mbox{}\verb@    mov x20, x1                  // x20 = remaining count@\\
\mbox{}\verb@@\\
\mbox{}\verb@    // Cache buffer state@\\
\mbox{}\verb@    adrp x21, _kcc_io_pos@{\tt @}\verb@PAGE@\\
\mbox{}\verb@    add  x21, x21, _kcc_io_pos@{\tt @}\verb@PAGEOFF@\\
\mbox{}\verb@@\\
\mbox{}\verb@    adrp x22, _kcc_io_buf@{\tt @}\verb@PAGE@\\
\mbox{}\verb@    add  x22, x22, _kcc_io_buf@{\tt @}\verb@PAGEOFF@\\
\mbox{}\verb@@\\
\mbox{}\verb@    cbz x20, .Lwrite_done@\\
\mbox{}\verb@@\\
\mbox{}\verb@.Lwrite_loop:@\\
\mbox{}\verb@    ldr  x1, [x21]              // x1 = current pos@\\
\mbox{}\verb@    ldrb w2, [x19], #1          // load byte from source@\\
\mbox{}\verb@    strb w2, [x22, x1]          // store in buffer@\\
\mbox{}\verb@    add  x1, x1, #1@\\
\mbox{}\verb@    str  x1, [x21]              // update pos@\\
\mbox{}\verb@@\\
\mbox{}\verb@    cmp x1, #4096@\\
\mbox{}\verb@    b.lo .Lwrite_no_flush@\\
\mbox{}\verb@    bl _kcc_io_flush@\\
\mbox{}\verb@.Lwrite_no_flush:@\\
\mbox{}\verb@    subs x20, x20, #1@\\
\mbox{}\verb@    b.ne .Lwrite_loop@\\
\mbox{}\verb@@\\
\mbox{}\verb@.Lwrite_done:@\\
\mbox{}\verb@    ldp x21, x22, [sp, #32]@\\
\mbox{}\verb@    ldp x19, x20, [sp, #16]@\\
\mbox{}\verb@    ldp x29, x30, [sp], #48@\\
\mbox{}\verb@    ret@\\
\mbox{}\verb@@{\NWsep}
\end{list}
\vspace{-1.5ex}
\footnotesize
\begin{list}{}{\setlength{\itemsep}{-\parsep}\setlength{\itemindent}{-\leftmargin}}
\item \NWtxtMacroRefIn\ \NWlink{weft?}{?}.

\item{}
\end{list}
\end{minipage}\vspace{4ex}
\end{flushleft}
\section{Print Unsigned Integer}

\texttt{kcc\_io\_putu} converts an unsigned 64-bit integer to decimal
ASCII and writes it to the output buffer.  This is used by
\texttt{--dump-tokens} to print line and column numbers.

The algorithm divides repeatedly by~10, storing digits in reverse order
in a small stack buffer, then writes them out with \texttt{kcc\_io\_write}.

\begin{flushleft} \small
\begin{minipage}{\linewidth}\label{scrap40}\raggedright\small
\NWtarget{weft?}{} $\langle\,${\itshape IO: putu}\nobreak\ {\footnotesize {?}}$\,\rangle\equiv$
\vspace{-1ex}
\begin{list}{}{} \item
\mbox{}\verb@@\\
\mbox{}\verb@// -- kcc_io_putu ----------------------------------------------------@\\
\mbox{}\verb@// void kcc_io_putu(uint64_t n)@\\
\mbox{}\verb@//   x0 = unsigned integer to print@\\
\mbox{}\verb@// Outputs the decimal representation to the buffered output.@\\
\mbox{}\verb@.globl _kcc_io_putu@\\
\mbox{}\verb@.p2align 2@\\
\mbox{}\verb@_kcc_io_putu:@\\
\mbox{}\verb@    stp x29, x30, [sp, #-64]!@\\
\mbox{}\verb@    mov x29, sp@\\
\mbox{}\verb@    str x19, [sp, #16]@\\
\mbox{}\verb@    // digit buffer: sp+24 through sp+47 (24 bytes, room for 20 digits)@\\
\mbox{}\verb@@\\
\mbox{}\verb@    mov x19, x0                  // x19 = number@\\
\mbox{}\verb@    add x3, sp, #48             // x3 = one past end of buffer@\\
\mbox{}\verb@    mov x4, x3                  // x4 = saved end pointer@\\
\mbox{}\verb@@\\
\mbox{}\verb@    // Handle zero specially@\\
\mbox{}\verb@    cbnz x19, .Lputu_nonzero@\\
\mbox{}\verb@    mov w1, #'0'@\\
\mbox{}\verb@    sub x3, x3, #1@\\
\mbox{}\verb@    strb w1, [x3]@\\
\mbox{}\verb@    b .Lputu_print@\\
\mbox{}\verb@@\\
\mbox{}\verb@.Lputu_nonzero:@\\
\mbox{}\verb@    mov x5, #10@\\
\mbox{}\verb@.Lputu_loop:@\\
\mbox{}\verb@    udiv x1, x19, x5            // x1 = n / 10@\\
\mbox{}\verb@    msub x2, x1, x5, x19        // x2 = n % 10@\\
\mbox{}\verb@    add  w2, w2, #'0'           // convert to ASCII@\\
\mbox{}\verb@    sub  x3, x3, #1@\\
\mbox{}\verb@    strb w2, [x3]@\\
\mbox{}\verb@    mov  x19, x1@\\
\mbox{}\verb@    cbnz x19, .Lputu_loop@\\
\mbox{}\verb@@\\
\mbox{}\verb@.Lputu_print:@\\
\mbox{}\verb@    mov x0, x3                  // pointer to first digit@\\
\mbox{}\verb@    sub x1, x4, x3             // length@\\
\mbox{}\verb@    bl _kcc_io_write@\\
\mbox{}\verb@@\\
\mbox{}\verb@    ldr x19, [sp, #16]@\\
\mbox{}\verb@    ldp x29, x30, [sp], #64@\\
\mbox{}\verb@    ret@\\
\mbox{}\verb@@{\NWsep}
\end{list}
\vspace{-1.5ex}
\footnotesize
\begin{list}{}{\setlength{\itemsep}{-\parsep}\setlength{\itemindent}{-\leftmargin}}
\item \NWtxtMacroRefIn\ \NWlink{weft?}{?}.

\item{}
\end{list}
\end{minipage}\vspace{4ex}
\end{flushleft}
\section{Stderr Output}

\texttt{kcc\_io\_err} writes a null-terminated string directly to
\texttt{stderr} (file descriptor~2) via a raw syscall.  No buffering
is used --- diagnostic messages should appear immediately, even if
the program crashes afterwards.

\begin{flushleft} \small
\begin{minipage}{\linewidth}\label{scrap41}\raggedright\small
\NWtarget{weft?}{} $\langle\,${\itshape IO: err}\nobreak\ {\footnotesize {?}}$\,\rangle\equiv$
\vspace{-1ex}
\begin{list}{}{} \item
\mbox{}\verb@@\\
\mbox{}\verb@// -- kcc_io_err -----------------------------------------------------@\\
\mbox{}\verb@// void kcc_io_err(const char *s)@\\
\mbox{}\verb@//   x0 = null-terminated string@\\
\mbox{}\verb@// Writes directly to stderr (fd=2), unbuffered.@\\
\mbox{}\verb@.globl _kcc_io_err@\\
\mbox{}\verb@.p2align 2@\\
\mbox{}\verb@_kcc_io_err:@\\
\mbox{}\verb@    stp x29, x30, [sp, #-32]!@\\
\mbox{}\verb@    mov x29, sp@\\
\mbox{}\verb@    str x19, [sp, #16]@\\
\mbox{}\verb@@\\
\mbox{}\verb@    mov x19, x0                  // save string pointer@\\
\mbox{}\verb@    bl _kcc_strlen               // x0 = length@\\
\mbox{}\verb@@\\
\mbox{}\verb@    mov x2, x0                  // x2 = byte count@\\
\mbox{}\verb@    mov x1, x19                 // x1 = buffer@\\
\mbox{}\verb@    mov x0, #2                  // x0 = fd (stderr)@\\
\mbox{}\verb@    bl _kcc_sys_write@\\
\mbox{}\verb@@\\
\mbox{}\verb@    ldr x19, [sp, #16]@\\
\mbox{}\verb@    ldp x29, x30, [sp], #32@\\
\mbox{}\verb@    ret@\\
\mbox{}\verb@@{\NWsep}
\end{list}
\vspace{-1.5ex}
\footnotesize
\begin{list}{}{\setlength{\itemsep}{-\parsep}\setlength{\itemindent}{-\leftmargin}}
\item \NWtxtMacroRefIn\ \NWlink{weft?}{?}.

\item{}
\end{list}
\end{minipage}\vspace{4ex}
\end{flushleft}
% 13-diagnostics.weft — Diagnostic Engine
%
% ISO 9899:2018 §5.1.1.3 requires that a conforming implementation
% produce at least one diagnostic message if a preprocessing translation
% unit or translation unit contains a violation of any syntax rule or
% constraint.
%
% The diagnostic engine provides structured error reporting in the
% format universally adopted by compilers:
%
%     file:line:col: severity: message
%
% Three severity levels are supported: error (constraint violation),
% warning (suspicious but permitted construct), and note (additional
% context for a preceding diagnostic).  An error count is maintained,
% and compilation aborts if it exceeds a configurable fatal limit.

\chapter{Diagnostic Engine}
\label{ch:diagnostics}

\section{Severity Levels}

The C standard distinguishes only ``diagnostic message'' from other
output, but modern practice recognizes three levels: errors that prevent
translation, warnings that flag suspicious code, and notes that provide
supplementary context.  We encode these as integer constants.

The fatal error limit prevents the compiler from drowning the user
in cascading errors.  After \texttt{DIAG\_MAX\_ERRORS} errors, the
engine prints a summary and calls \texttt{sys\_exit(1)}.

\begin{flushleft} \small
\begin{minipage}{\linewidth}\label{scrap42}\raggedright\small
\NWtarget{weft?}{} $\langle\,${\itshape Diag: constants}\nobreak\ {\footnotesize {?}}$\,\rangle\equiv$
\vspace{-1ex}
\begin{list}{}{} \item
\mbox{}\verb@@\\
\mbox{}\verb@// -- Diagnostic severity codes ------------------------------------------@\\
\mbox{}\verb@.equ DIAG_ERROR,       0@\\
\mbox{}\verb@.equ DIAG_WARNING,     1@\\
\mbox{}\verb@.equ DIAG_NOTE,        2@\\
\mbox{}\verb@@\\
\mbox{}\verb@// Maximum errors before aborting@\\
\mbox{}\verb@.equ DIAG_MAX_ERRORS, 20@\\
\mbox{}\verb@@{\NWsep}
\end{list}
\vspace{-1.5ex}
\footnotesize
\begin{list}{}{\setlength{\itemsep}{-\parsep}\setlength{\itemindent}{-\leftmargin}}
\item \NWtxtMacroRefIn\ \NWlink{weft?}{?}.

\item{}
\end{list}
\end{minipage}\vspace{4ex}
\end{flushleft}
\section{Stderr Integer Output}

The diagnostic engine needs to print line and column numbers to stderr.
This helper converts a \texttt{uint64\_t} to decimal ASCII and writes
it directly to file descriptor~2 via \texttt{sys\_write}.  The algorithm
mirrors \texttt{kcc\_io\_putu} (Chapter~\ref{ch:io}) but bypasses the
stdout buffer.

\begin{flushleft} \small
\begin{minipage}{\linewidth}\label{scrap43}\raggedright\small
\NWtarget{weft?}{} $\langle\,${\itshape Diag: putu stderr}\nobreak\ {\footnotesize {?}}$\,\rangle\equiv$
\vspace{-1ex}
\begin{list}{}{} \item
\mbox{}\verb@@\\
\mbox{}\verb@// -- _kcc_diag_putu -----------------------------------------------------@\\
\mbox{}\verb@// void _kcc_diag_putu(uint64_t n)@\\
\mbox{}\verb@//   x0 = unsigned integer to print@\\
\mbox{}\verb@// Outputs the decimal representation directly to stderr (fd=2).@\\
\mbox{}\verb@.p2align 2@\\
\mbox{}\verb@_kcc_diag_putu:@\\
\mbox{}\verb@    stp x29, x30, [sp, #-64]!@\\
\mbox{}\verb@    mov x29, sp@\\
\mbox{}\verb@    str x19, [sp, #16]@\\
\mbox{}\verb@@\\
\mbox{}\verb@    mov x19, x0                  // x19 = number@\\
\mbox{}\verb@    add x3, sp, #48              // x3 = one past end of digit buffer@\\
\mbox{}\verb@    mov x4, x3                   // x4 = saved end pointer@\\
\mbox{}\verb@@\\
\mbox{}\verb@    // Handle zero specially@\\
\mbox{}\verb@    cbnz x19, .Ldiag_putu_nonzero@\\
\mbox{}\verb@    mov w1, #'0'@\\
\mbox{}\verb@    sub x3, x3, #1@\\
\mbox{}\verb@    strb w1, [x3]@\\
\mbox{}\verb@    b .Ldiag_putu_print@\\
\mbox{}\verb@@\\
\mbox{}\verb@.Ldiag_putu_nonzero:@\\
\mbox{}\verb@    mov x5, #10@\\
\mbox{}\verb@.Ldiag_putu_loop:@\\
\mbox{}\verb@    udiv x1, x19, x5            // x1 = n / 10@\\
\mbox{}\verb@    msub x2, x1, x5, x19        // x2 = n % 10@\\
\mbox{}\verb@    add  w2, w2, #'0'           // convert to ASCII digit@\\
\mbox{}\verb@    sub  x3, x3, #1@\\
\mbox{}\verb@    strb w2, [x3]@\\
\mbox{}\verb@    mov  x19, x1@\\
\mbox{}\verb@    cbnz x19, .Ldiag_putu_loop@\\
\mbox{}\verb@@\\
\mbox{}\verb@.Ldiag_putu_print:@\\
\mbox{}\verb@    // sys_write(2, digits, count)@\\
\mbox{}\verb@    mov x0, #2                  // fd = stderr@\\
\mbox{}\verb@    mov x1, x3                  // buffer@\\
\mbox{}\verb@    sub x2, x4, x3              // length@\\
\mbox{}\verb@    bl _kcc_sys_write@\\
\mbox{}\verb@@\\
\mbox{}\verb@    ldr x19, [sp, #16]@\\
\mbox{}\verb@    ldp x29, x30, [sp], #64@\\
\mbox{}\verb@    ret@\\
\mbox{}\verb@@{\NWsep}
\end{list}
\vspace{-1.5ex}
\footnotesize
\begin{list}{}{\setlength{\itemsep}{-\parsep}\setlength{\itemindent}{-\leftmargin}}
\item \NWtxtMacroRefIn\ \NWlink{weft?}{?}.

\item{}
\end{list}
\end{minipage}\vspace{4ex}
\end{flushleft}
\section{Core Emission}

The heart of the diagnostic engine is \texttt{kcc\_diag\_emit}.  Given
a severity level, source location, and message string, it writes the
full diagnostic to stderr in the standard format.

Each component (file path, colon, number, severity label, message) is
written as a separate string to stderr.  This avoids the need for
\texttt{sprintf}-style formatting.  Stderr writes are unbuffered, so
diagnostics appear immediately even if the program crashes.

After an error diagnostic, the error counter is incremented.  If it
reaches \texttt{DIAG\_MAX\_ERRORS}, a ``too many errors'' message is
printed and the process exits with status~1.

\begin{flushleft} \small
\begin{minipage}{\linewidth}\label{scrap44}\raggedright\small
\NWtarget{weft?}{} $\langle\,${\itshape Diag: emit}\nobreak\ {\footnotesize {?}}$\,\rangle\equiv$
\vspace{-1ex}
\begin{list}{}{} \item
\mbox{}\verb@@\\
\mbox{}\verb@// -- kcc_diag_emit ------------------------------------------------------@\\
\mbox{}\verb@// void kcc_diag_emit(int severity, uint64_t line, uint64_t col,@\\
\mbox{}\verb@//                    const char *msg)@\\
\mbox{}\verb@//   x0 = severity (DIAG_ERROR/DIAG_WARNING/DIAG_NOTE)@\\
\mbox{}\verb@//   x1 = line number (1-based)@\\
\mbox{}\verb@//   x2 = column number (1-based)@\\
\mbox{}\verb@//   x3 = pointer to null-terminated message string@\\
\mbox{}\verb@// Writes: file:line:col: severity: message\n   to stderr.@\\
\mbox{}\verb@.globl _kcc_diag_emit@\\
\mbox{}\verb@.p2align 2@\\
\mbox{}\verb@_kcc_diag_emit:@\\
\mbox{}\verb@    stp x29, x30, [sp, #-64]!@\\
\mbox{}\verb@    mov x29, sp@\\
\mbox{}\verb@    stp x19, x20, [sp, #16]@\\
\mbox{}\verb@    stp x21, x22, [sp, #32]@\\
\mbox{}\verb@    str x23, [sp, #48]@\\
\mbox{}\verb@@\\
\mbox{}\verb@    mov x19, x0                  // x19 = severity@\\
\mbox{}\verb@    mov x20, x1                  // x20 = line@\\
\mbox{}\verb@    mov x21, x2                  // x21 = column@\\
\mbox{}\verb@    mov x22, x3                  // x22 = message@\\
\mbox{}\verb@@\\
\mbox{}\verb@    // Print file path (from source manager)@\\
\mbox{}\verb@    adrp x0, _kcc_src_path@{\tt @}\verb@PAGE@\\
\mbox{}\verb@    add  x0, x0, _kcc_src_path@{\tt @}\verb@PAGEOFF@\\
\mbox{}\verb@    ldr  x0, [x0]@\\
\mbox{}\verb@    cbz  x0, .Ldiag_no_path@\\
\mbox{}\verb@    bl _kcc_io_err@\\
\mbox{}\verb@    b .Ldiag_colon1@\\
\mbox{}\verb@.Ldiag_no_path:@\\
\mbox{}\verb@    adrp x0, _kcc_diag_unknown_file@{\tt @}\verb@PAGE@\\
\mbox{}\verb@    add  x0, x0, _kcc_diag_unknown_file@{\tt @}\verb@PAGEOFF@\\
\mbox{}\verb@    bl _kcc_io_err@\\
\mbox{}\verb@@\\
\mbox{}\verb@.Ldiag_colon1:@\\
\mbox{}\verb@    // Print ":"@\\
\mbox{}\verb@    adrp x0, _kcc_diag_colon@{\tt @}\verb@PAGE@\\
\mbox{}\verb@    add  x0, x0, _kcc_diag_colon@{\tt @}\verb@PAGEOFF@\\
\mbox{}\verb@    bl _kcc_io_err@\\
\mbox{}\verb@@\\
\mbox{}\verb@    // Print line number@\\
\mbox{}\verb@    mov x0, x20@\\
\mbox{}\verb@    bl _kcc_diag_putu@\\
\mbox{}\verb@@\\
\mbox{}\verb@    // Print ":"@\\
\mbox{}\verb@    adrp x0, _kcc_diag_colon@{\tt @}\verb@PAGE@\\
\mbox{}\verb@    add  x0, x0, _kcc_diag_colon@{\tt @}\verb@PAGEOFF@\\
\mbox{}\verb@    bl _kcc_io_err@\\
\mbox{}\verb@@\\
\mbox{}\verb@    // Print column number@\\
\mbox{}\verb@    mov x0, x21@\\
\mbox{}\verb@    bl _kcc_diag_putu@\\
\mbox{}\verb@@\\
\mbox{}\verb@    // Print ": "@\\
\mbox{}\verb@    adrp x0, _kcc_diag_colon_space@{\tt @}\verb@PAGE@\\
\mbox{}\verb@    add  x0, x0, _kcc_diag_colon_space@{\tt @}\verb@PAGEOFF@\\
\mbox{}\verb@    bl _kcc_io_err@\\
\mbox{}\verb@@\\
\mbox{}\verb@    // Print severity string based on level@\\
\mbox{}\verb@    cmp x19, #DIAG_WARNING@\\
\mbox{}\verb@    b.eq .Ldiag_sev_warn@\\
\mbox{}\verb@    cmp x19, #DIAG_NOTE@\\
\mbox{}\verb@    b.eq .Ldiag_sev_note@\\
\mbox{}\verb@    // Default: error@\\
\mbox{}\verb@    adrp x0, _kcc_diag_str_error@{\tt @}\verb@PAGE@\\
\mbox{}\verb@    add  x0, x0, _kcc_diag_str_error@{\tt @}\verb@PAGEOFF@\\
\mbox{}\verb@    b .Ldiag_sev_done@\\
\mbox{}\verb@.Ldiag_sev_warn:@\\
\mbox{}\verb@    adrp x0, _kcc_diag_str_warning@{\tt @}\verb@PAGE@\\
\mbox{}\verb@    add  x0, x0, _kcc_diag_str_warning@{\tt @}\verb@PAGEOFF@\\
\mbox{}\verb@    b .Ldiag_sev_done@\\
\mbox{}\verb@.Ldiag_sev_note:@\\
\mbox{}\verb@    adrp x0, _kcc_diag_str_note@{\tt @}\verb@PAGE@\\
\mbox{}\verb@    add  x0, x0, _kcc_diag_str_note@{\tt @}\verb@PAGEOFF@\\
\mbox{}\verb@.Ldiag_sev_done:@\\
\mbox{}\verb@    bl _kcc_io_err@\\
\mbox{}\verb@@\\
\mbox{}\verb@    // Print ": "@\\
\mbox{}\verb@    adrp x0, _kcc_diag_colon_space@{\tt @}\verb@PAGE@\\
\mbox{}\verb@    add  x0, x0, _kcc_diag_colon_space@{\tt @}\verb@PAGEOFF@\\
\mbox{}\verb@    bl _kcc_io_err@\\
\mbox{}\verb@@\\
\mbox{}\verb@    // Print message@\\
\mbox{}\verb@    mov x0, x22@\\
\mbox{}\verb@    bl _kcc_io_err@\\
\mbox{}\verb@@\\
\mbox{}\verb@    // Print newline@\\
\mbox{}\verb@    adrp x0, _kcc_diag_newline@{\tt @}\verb@PAGE@\\
\mbox{}\verb@    add  x0, x0, _kcc_diag_newline@{\tt @}\verb@PAGEOFF@\\
\mbox{}\verb@    bl _kcc_io_err@\\
\mbox{}\verb@@\\
\mbox{}\verb@    // Update counters based on severity@\\
\mbox{}\verb@    cmp x19, #DIAG_ERROR@\\
\mbox{}\verb@    b.ne .Ldiag_check_warn@\\
\mbox{}\verb@@\\
\mbox{}\verb@    // Increment error count@\\
\mbox{}\verb@    adrp x0, _kcc_diag_error_count@{\tt @}\verb@PAGE@\\
\mbox{}\verb@    add  x0, x0, _kcc_diag_error_count@{\tt @}\verb@PAGEOFF@\\
\mbox{}\verb@    ldr  x1, [x0]@\\
\mbox{}\verb@    add  x1, x1, #1@\\
\mbox{}\verb@    str  x1, [x0]@\\
\mbox{}\verb@@\\
\mbox{}\verb@    // Check fatal error limit@\\
\mbox{}\verb@    cmp x1, #DIAG_MAX_ERRORS@\\
\mbox{}\verb@    b.lt .Ldiag_done@\\
\mbox{}\verb@@\\
\mbox{}\verb@    // Too many errors — abort@\\
\mbox{}\verb@    adrp x0, _kcc_diag_too_many@{\tt @}\verb@PAGE@\\
\mbox{}\verb@    add  x0, x0, _kcc_diag_too_many@{\tt @}\verb@PAGEOFF@\\
\mbox{}\verb@    bl _kcc_io_err@\\
\mbox{}\verb@    mov x0, #1@\\
\mbox{}\verb@    bl _kcc_sys_exit@\\
\mbox{}\verb@@\\
\mbox{}\verb@.Ldiag_check_warn:@\\
\mbox{}\verb@    cmp x19, #DIAG_WARNING@\\
\mbox{}\verb@    b.ne .Ldiag_done@\\
\mbox{}\verb@@\\
\mbox{}\verb@    // Increment warning count@\\
\mbox{}\verb@    adrp x0, _kcc_diag_warn_count@{\tt @}\verb@PAGE@\\
\mbox{}\verb@    add  x0, x0, _kcc_diag_warn_count@{\tt @}\verb@PAGEOFF@\\
\mbox{}\verb@    ldr  x1, [x0]@\\
\mbox{}\verb@    add  x1, x1, #1@\\
\mbox{}\verb@    str  x1, [x0]@\\
\mbox{}\verb@@\\
\mbox{}\verb@.Ldiag_done:@\\
\mbox{}\verb@    ldr x23, [sp, #48]@\\
\mbox{}\verb@    ldp x21, x22, [sp, #32]@\\
\mbox{}\verb@    ldp x19, x20, [sp, #16]@\\
\mbox{}\verb@    ldp x29, x30, [sp], #64@\\
\mbox{}\verb@    ret@\\
\mbox{}\verb@@{\NWsep}
\end{list}
\vspace{-1.5ex}
\footnotesize
\begin{list}{}{\setlength{\itemsep}{-\parsep}\setlength{\itemindent}{-\leftmargin}}
\item \NWtxtMacroRefIn\ \NWlink{weft?}{?}.

\item{}
\end{list}
\end{minipage}\vspace{4ex}
\end{flushleft}
\section{Convenience Wrappers}

Three wrappers provide a simpler calling convention for the common
case where the caller already has line, column, and message:

\begin{center}
\begin{tabular}{ll}
\texttt{kcc\_diag\_error(line, col, msg)} & emits an error \\
\texttt{kcc\_diag\_warning(line, col, msg)} & emits a warning \\
\texttt{kcc\_diag\_note(line, col, msg)} & emits a note \\
\end{tabular}
\end{center}

These are tail-call wrappers: they rearrange the arguments and branch
(not \texttt{bl}) to \texttt{kcc\_diag\_emit}, so the return address
in \reg{x30} passes through untouched.

\begin{flushleft} \small
\begin{minipage}{\linewidth}\label{scrap45}\raggedright\small
\NWtarget{weft?}{} $\langle\,${\itshape Diag: error}\nobreak\ {\footnotesize {?}}$\,\rangle\equiv$
\vspace{-1ex}
\begin{list}{}{} \item
\mbox{}\verb@@\\
\mbox{}\verb@// -- kcc_diag_error -----------------------------------------------------@\\
\mbox{}\verb@// void kcc_diag_error(uint64_t line, uint64_t col, const char *msg)@\\
\mbox{}\verb@//   x0 = line, x1 = col, x2 = msg@\\
\mbox{}\verb@.globl _kcc_diag_error@\\
\mbox{}\verb@.p2align 2@\\
\mbox{}\verb@_kcc_diag_error:@\\
\mbox{}\verb@    mov x3, x2                   // msg → x3@\\
\mbox{}\verb@    mov x2, x1                   // col → x2@\\
\mbox{}\verb@    mov x1, x0                   // line → x1@\\
\mbox{}\verb@    mov x0, #DIAG_ERROR          // severity → x0@\\
\mbox{}\verb@    b _kcc_diag_emit             // tail call@\\
\mbox{}\verb@@{\NWsep}
\end{list}
\vspace{-1.5ex}
\footnotesize
\begin{list}{}{\setlength{\itemsep}{-\parsep}\setlength{\itemindent}{-\leftmargin}}
\item \NWtxtMacroRefIn\ \NWlink{weft?}{?}.

\item{}
\end{list}
\end{minipage}\vspace{4ex}
\end{flushleft}
\begin{flushleft} \small
\begin{minipage}{\linewidth}\label{scrap46}\raggedright\small
\NWtarget{weft?}{} $\langle\,${\itshape Diag: warning}\nobreak\ {\footnotesize {?}}$\,\rangle\equiv$
\vspace{-1ex}
\begin{list}{}{} \item
\mbox{}\verb@@\\
\mbox{}\verb@// -- kcc_diag_warning ---------------------------------------------------@\\
\mbox{}\verb@// void kcc_diag_warning(uint64_t line, uint64_t col, const char *msg)@\\
\mbox{}\verb@.globl _kcc_diag_warning@\\
\mbox{}\verb@.p2align 2@\\
\mbox{}\verb@_kcc_diag_warning:@\\
\mbox{}\verb@    mov x3, x2@\\
\mbox{}\verb@    mov x2, x1@\\
\mbox{}\verb@    mov x1, x0@\\
\mbox{}\verb@    mov x0, #DIAG_WARNING@\\
\mbox{}\verb@    b _kcc_diag_emit@\\
\mbox{}\verb@@{\NWsep}
\end{list}
\vspace{-1.5ex}
\footnotesize
\begin{list}{}{\setlength{\itemsep}{-\parsep}\setlength{\itemindent}{-\leftmargin}}
\item \NWtxtMacroRefIn\ \NWlink{weft?}{?}.

\item{}
\end{list}
\end{minipage}\vspace{4ex}
\end{flushleft}
\begin{flushleft} \small
\begin{minipage}{\linewidth}\label{scrap47}\raggedright\small
\NWtarget{weft?}{} $\langle\,${\itshape Diag: note}\nobreak\ {\footnotesize {?}}$\,\rangle\equiv$
\vspace{-1ex}
\begin{list}{}{} \item
\mbox{}\verb@@\\
\mbox{}\verb@// -- kcc_diag_note ------------------------------------------------------@\\
\mbox{}\verb@// void kcc_diag_note(uint64_t line, uint64_t col, const char *msg)@\\
\mbox{}\verb@.globl _kcc_diag_note@\\
\mbox{}\verb@.p2align 2@\\
\mbox{}\verb@_kcc_diag_note:@\\
\mbox{}\verb@    mov x3, x2@\\
\mbox{}\verb@    mov x2, x1@\\
\mbox{}\verb@    mov x1, x0@\\
\mbox{}\verb@    mov x0, #DIAG_NOTE@\\
\mbox{}\verb@    b _kcc_diag_emit@\\
\mbox{}\verb@@{\NWsep}
\end{list}
\vspace{-1.5ex}
\footnotesize
\begin{list}{}{\setlength{\itemsep}{-\parsep}\setlength{\itemindent}{-\leftmargin}}
\item \NWtxtMacroRefIn\ \NWlink{weft?}{?}.

\item{}
\end{list}
\end{minipage}\vspace{4ex}
\end{flushleft}
\section{Error Count Query}

\texttt{kcc\_diag\_had\_errors} returns the current error count.
The driver uses this after tokenization to decide whether to continue
with subsequent phases or exit with a failure status.

\begin{flushleft} \small
\begin{minipage}{\linewidth}\label{scrap48}\raggedright\small
\NWtarget{weft?}{} $\langle\,${\itshape Diag: had errors}\nobreak\ {\footnotesize {?}}$\,\rangle\equiv$
\vspace{-1ex}
\begin{list}{}{} \item
\mbox{}\verb@@\\
\mbox{}\verb@// -- kcc_diag_had_errors ------------------------------------------------@\\
\mbox{}\verb@// uint64_t kcc_diag_had_errors(void)@\\
\mbox{}\verb@// Returns the error count in x0.@\\
\mbox{}\verb@.globl _kcc_diag_had_errors@\\
\mbox{}\verb@.p2align 2@\\
\mbox{}\verb@_kcc_diag_had_errors:@\\
\mbox{}\verb@    adrp x0, _kcc_diag_error_count@{\tt @}\verb@PAGE@\\
\mbox{}\verb@    add  x0, x0, _kcc_diag_error_count@{\tt @}\verb@PAGEOFF@\\
\mbox{}\verb@    ldr  x0, [x0]@\\
\mbox{}\verb@    ret@\\
\mbox{}\verb@@{\NWsep}
\end{list}
\vspace{-1.5ex}
\footnotesize
\begin{list}{}{\setlength{\itemsep}{-\parsep}\setlength{\itemindent}{-\leftmargin}}
\item \NWtxtMacroRefIn\ \NWlink{weft?}{?}.

\item{}
\end{list}
\end{minipage}\vspace{4ex}
\end{flushleft}
\section{Reset}

\texttt{kcc\_diag\_reset} clears both error and warning counters.
This is used by the test harness between individual tests to prevent
counters from one test affecting the next.

\begin{flushleft} \small
\begin{minipage}{\linewidth}\label{scrap49}\raggedright\small
\NWtarget{weft?}{} $\langle\,${\itshape Diag: reset}\nobreak\ {\footnotesize {?}}$\,\rangle\equiv$
\vspace{-1ex}
\begin{list}{}{} \item
\mbox{}\verb@@\\
\mbox{}\verb@// -- kcc_diag_reset -----------------------------------------------------@\\
\mbox{}\verb@// void kcc_diag_reset(void)@\\
\mbox{}\verb@// Resets error and warning counts to zero.@\\
\mbox{}\verb@.globl _kcc_diag_reset@\\
\mbox{}\verb@.p2align 2@\\
\mbox{}\verb@_kcc_diag_reset:@\\
\mbox{}\verb@    adrp x0, _kcc_diag_error_count@{\tt @}\verb@PAGE@\\
\mbox{}\verb@    add  x0, x0, _kcc_diag_error_count@{\tt @}\verb@PAGEOFF@\\
\mbox{}\verb@    str  xzr, [x0]@\\
\mbox{}\verb@    adrp x0, _kcc_diag_warn_count@{\tt @}\verb@PAGE@\\
\mbox{}\verb@    add  x0, x0, _kcc_diag_warn_count@{\tt @}\verb@PAGEOFF@\\
\mbox{}\verb@    str  xzr, [x0]@\\
\mbox{}\verb@    ret@\\
\mbox{}\verb@@{\NWsep}
\end{list}
\vspace{-1.5ex}
\footnotesize
\begin{list}{}{\setlength{\itemsep}{-\parsep}\setlength{\itemindent}{-\leftmargin}}
\item \NWtxtMacroRefIn\ \NWlink{weft?}{?}.

\item{}
\end{list}
\end{minipage}\vspace{4ex}
\end{flushleft}
\section{Global State}

The diagnostic engine maintains two counters and a set of format
strings.  The counters are initialized to zero by the \texttt{.data}
section.

\begin{flushleft} \small
\begin{minipage}{\linewidth}\label{scrap50}\raggedright\small
\NWtarget{weft?}{} $\langle\,${\itshape Diag: data}\nobreak\ {\footnotesize {?}}$\,\rangle\equiv$
\vspace{-1ex}
\begin{list}{}{} \item
\mbox{}\verb@@\\
\mbox{}\verb@// -- Diagnostic global state --------------------------------------------@\\
\mbox{}\verb@.data@\\
\mbox{}\verb@.p2align 3@\\
\mbox{}\verb@_kcc_diag_error_count:  .quad 0    // running error count@\\
\mbox{}\verb@_kcc_diag_warn_count:   .quad 0    // running warning count@\\
\mbox{}\verb@@\\
\mbox{}\verb@// Format strings (read-only)@\\
\mbox{}\verb@.section __TEXT,__cstring,cstring_literals@\\
\mbox{}\verb@_kcc_diag_colon:         .asciz ":"@\\
\mbox{}\verb@_kcc_diag_colon_space:   .asciz ": "@\\
\mbox{}\verb@_kcc_diag_newline:       .asciz "\n"@\\
\mbox{}\verb@_kcc_diag_str_error:     .asciz "error"@\\
\mbox{}\verb@_kcc_diag_str_warning:   .asciz "warning"@\\
\mbox{}\verb@_kcc_diag_str_note:      .asciz "note"@\\
\mbox{}\verb@_kcc_diag_unknown_file:  .asciz "<input>"@\\
\mbox{}\verb@_kcc_diag_too_many:      .asciz "kcc: fatal: too many errors emitted, stopping\n"@\\
\mbox{}\verb@@{\NWsep}
\end{list}
\vspace{-1.5ex}
\footnotesize
\begin{list}{}{\setlength{\itemsep}{-\parsep}\setlength{\itemindent}{-\leftmargin}}
\item \NWtxtMacroRefIn\ \NWlink{weft?}{?}.

\item{}
\end{list}
\end{minipage}\vspace{4ex}
\end{flushleft}
% 14-source-manager.weft — Source File Management
%
% The source manager handles reading C source files from disk.
% Files are memory-mapped with mmap for zero-copy access --- the
% kernel pages in file contents on demand, and the lexer can scan
% directly through the mapped region.
%
% Source location tracking is trivial: since the lexer walks a
% contiguous buffer, the current line and column are maintained
% as simple counters during scanning.

\chapter{Source File Management}
\label{ch:source}

\section{Design}

A compiler must read source files efficiently.  The simplest approach
is to \texttt{mmap} the entire file: the kernel provides a read-only
view of the file's contents, paged in on demand.  The lexer then scans
directly through the mapped memory without any copying.

For each source file we track:
\begin{itemize}
\item The base pointer (start of mapped region)
\item The file size (length of mapped region)
\item The file path (for diagnostic messages)
\end{itemize}

\begin{center}
\begin{tabular}{rll}
\hline
Offset & Field & Description \\
\hline
0  & \texttt{base} & Pointer to start of mapped file \\
8  & \texttt{size} & File size in bytes \\
16 & \texttt{path} & Pointer to file path string \\
\hline
\end{tabular}
\end{center}

\begin{flushleft} \small
\begin{minipage}{\linewidth}\label{scrap51}\raggedright\small
\NWtarget{weft?}{} $\langle\,${\itshape Source: constants}\nobreak\ {\footnotesize {?}}$\,\rangle\equiv$
\vspace{-1ex}
\begin{list}{}{} \item
\mbox{}\verb@@\\
\mbox{}\verb@// -- Source buffer layout --------------------------------------------@\\
\mbox{}\verb@.equ SRC_BASE,      0         // pointer to mmap'd file data@\\
\mbox{}\verb@.equ SRC_SIZE,      8         // file size in bytes@\\
\mbox{}\verb@.equ SRC_PATH,     16         // pointer to file path (null-terminated)@\\
\mbox{}\verb@.equ SRC_BUFSIZE,  24         // structure size@\\
\mbox{}\verb@@{\NWsep}
\end{list}
\vspace{-1.5ex}
\footnotesize
\begin{list}{}{\setlength{\itemsep}{-\parsep}\setlength{\itemindent}{-\leftmargin}}
\item \NWtxtMacroRefIn\ \NWlink{weft?}{?}.

\item{}
\end{list}
\end{minipage}\vspace{4ex}
\end{flushleft}
\section{Opening a Source File}

\texttt{kcc\_source\_open} opens a C source file by path.  The
sequence is:
\begin{enumerate}
\item \texttt{open} the file read-only.
\item \texttt{lseek(fd, 0, SEEK\_END)} to determine the file size.
\item \texttt{mmap} the file with \texttt{PROT\_READ | MAP\_PRIVATE}.
\item \texttt{close} the file descriptor (the mapping persists).
\item Store the base pointer, size, and path in the global source state.
\end{enumerate}

If any step fails, the function prints an error message to stderr
and exits with status~1.

\begin{flushleft} \small
\begin{minipage}{\linewidth}\label{scrap52}\raggedright\small
\NWtarget{weft?}{} $\langle\,${\itshape Source: open file}\nobreak\ {\footnotesize {?}}$\,\rangle\equiv$
\vspace{-1ex}
\begin{list}{}{} \item
\mbox{}\verb@@\\
\mbox{}\verb@// -- kcc_source_open ------------------------------------------------@\\
\mbox{}\verb@// void kcc_source_open(const char *path)@\\
\mbox{}\verb@//   x0 = null-terminated file path@\\
\mbox{}\verb@// Opens and mmaps the file, stores result in global state.@\\
\mbox{}\verb@// On failure, prints error to stderr and exits.@\\
\mbox{}\verb@.globl _kcc_source_open@\\
\mbox{}\verb@.p2align 2@\\
\mbox{}\verb@_kcc_source_open:@\\
\mbox{}\verb@    stp x29, x30, [sp, #-48]!@\\
\mbox{}\verb@    mov x29, sp@\\
\mbox{}\verb@    stp x19, x20, [sp, #16]@\\
\mbox{}\verb@    str x21, [sp, #32]@\\
\mbox{}\verb@@\\
\mbox{}\verb@    mov x19, x0                  // x19 = path@\\
\mbox{}\verb@@\\
\mbox{}\verb@    // Store path in global state@\\
\mbox{}\verb@    adrp x1, _kcc_src_path@{\tt @}\verb@PAGE@\\
\mbox{}\verb@    add  x1, x1, _kcc_src_path@{\tt @}\verb@PAGEOFF@\\
\mbox{}\verb@    str  x19, [x1]@\\
\mbox{}\verb@@\\
\mbox{}\verb@    // open(path, O_RDONLY, 0)@\\
\mbox{}\verb@    mov x0, x19                  // path@\\
\mbox{}\verb@    mov x1, #0                   // O_RDONLY = 0@\\
\mbox{}\verb@    mov x2, #0                   // mode (unused for O_RDONLY)@\\
\mbox{}\verb@    bl _kcc_sys_open@\\
\mbox{}\verb@@\\
\mbox{}\verb@    // Check for error@\\
\mbox{}\verb@    cmn x0, #1@\\
\mbox{}\verb@    b.eq .Lsource_open_fail@\\
\mbox{}\verb@    mov x20, x0                  // x20 = fd@\\
\mbox{}\verb@@\\
\mbox{}\verb@    // lseek(fd, 0, SEEK_END) to get file size@\\
\mbox{}\verb@    mov x0, x20@\\
\mbox{}\verb@    mov x1, #0@\\
\mbox{}\verb@    mov x2, #SEEK_END@\\
\mbox{}\verb@    bl _kcc_sys_lseek@\\
\mbox{}\verb@@\\
\mbox{}\verb@    cmn x0, #1@\\
\mbox{}\verb@    b.eq .Lsource_seek_fail@\\
\mbox{}\verb@    mov x21, x0                  // x21 = file size@\\
\mbox{}\verb@@\\
\mbox{}\verb@    // Reject empty files@\\
\mbox{}\verb@    cbz x21, .Lsource_empty@\\
\mbox{}\verb@@\\
\mbox{}\verb@    // mmap(NULL, size, PROT_READ, MAP_PRIVATE, fd, 0)@\\
\mbox{}\verb@    mov x0, #0                   // addr = NULL@\\
\mbox{}\verb@    mov x1, x21                  // len = file size@\\
\mbox{}\verb@    mov x2, #PROT_READ           // prot = read only@\\
\mbox{}\verb@    mov x3, #MAP_PRIVATE         // flags = private@\\
\mbox{}\verb@    mov x4, x20                  // fd@\\
\mbox{}\verb@    mov x5, #0                   // offset = 0@\\
\mbox{}\verb@    bl _kcc_sys_mmap@\\
\mbox{}\verb@@\\
\mbox{}\verb@    cmn x0, #1@\\
\mbox{}\verb@    b.eq .Lsource_mmap_fail@\\
\mbox{}\verb@@\\
\mbox{}\verb@    // Store base and size in global state@\\
\mbox{}\verb@    adrp x1, _kcc_src_base@{\tt @}\verb@PAGE@\\
\mbox{}\verb@    add  x1, x1, _kcc_src_base@{\tt @}\verb@PAGEOFF@\\
\mbox{}\verb@    str  x0, [x1]@\\
\mbox{}\verb@@\\
\mbox{}\verb@    adrp x1, _kcc_src_size@{\tt @}\verb@PAGE@\\
\mbox{}\verb@    add  x1, x1, _kcc_src_size@{\tt @}\verb@PAGEOFF@\\
\mbox{}\verb@    str  x21, [x1]@\\
\mbox{}\verb@@\\
\mbox{}\verb@    // Close the file descriptor (mapping persists)@\\
\mbox{}\verb@    mov x0, x20@\\
\mbox{}\verb@    bl _kcc_sys_close@\\
\mbox{}\verb@@\\
\mbox{}\verb@    ldr x21, [sp, #32]@\\
\mbox{}\verb@    ldp x19, x20, [sp, #16]@\\
\mbox{}\verb@    ldp x29, x30, [sp], #48@\\
\mbox{}\verb@    ret@\\
\mbox{}\verb@@\\
\mbox{}\verb@.Lsource_open_fail:@\\
\mbox{}\verb@    adrp x0, _kcc_err_open@{\tt @}\verb@PAGE@\\
\mbox{}\verb@    add  x0, x0, _kcc_err_open@{\tt @}\verb@PAGEOFF@\\
\mbox{}\verb@    bl _kcc_io_err@\\
\mbox{}\verb@    mov x0, x19@\\
\mbox{}\verb@    bl _kcc_io_err@\\
\mbox{}\verb@    adrp x0, _kcc_src_newline@{\tt @}\verb@PAGE@\\
\mbox{}\verb@    add  x0, x0, _kcc_src_newline@{\tt @}\verb@PAGEOFF@\\
\mbox{}\verb@    bl _kcc_io_err@\\
\mbox{}\verb@    mov x0, #1@\\
\mbox{}\verb@    bl _kcc_sys_exit@\\
\mbox{}\verb@@\\
\mbox{}\verb@.Lsource_seek_fail:@\\
\mbox{}\verb@    mov x0, x20@\\
\mbox{}\verb@    bl _kcc_sys_close@\\
\mbox{}\verb@    adrp x0, _kcc_err_read@{\tt @}\verb@PAGE@\\
\mbox{}\verb@    add  x0, x0, _kcc_err_read@{\tt @}\verb@PAGEOFF@\\
\mbox{}\verb@    bl _kcc_io_err@\\
\mbox{}\verb@    mov x0, #1@\\
\mbox{}\verb@    bl _kcc_sys_exit@\\
\mbox{}\verb@@\\
\mbox{}\verb@.Lsource_empty:@\\
\mbox{}\verb@    mov x0, x20@\\
\mbox{}\verb@    bl _kcc_sys_close@\\
\mbox{}\verb@    adrp x0, _kcc_err_empty@{\tt @}\verb@PAGE@\\
\mbox{}\verb@    add  x0, x0, _kcc_err_empty@{\tt @}\verb@PAGEOFF@\\
\mbox{}\verb@    bl _kcc_io_err@\\
\mbox{}\verb@    mov x0, #1@\\
\mbox{}\verb@    bl _kcc_sys_exit@\\
\mbox{}\verb@@\\
\mbox{}\verb@.Lsource_mmap_fail:@\\
\mbox{}\verb@    mov x0, x20@\\
\mbox{}\verb@    bl _kcc_sys_close@\\
\mbox{}\verb@    adrp x0, _kcc_err_mmap@{\tt @}\verb@PAGE@\\
\mbox{}\verb@    add  x0, x0, _kcc_err_mmap@{\tt @}\verb@PAGEOFF@\\
\mbox{}\verb@    bl _kcc_io_err@\\
\mbox{}\verb@    mov x0, #1@\\
\mbox{}\verb@    bl _kcc_sys_exit@\\
\mbox{}\verb@@{\NWsep}
\end{list}
\vspace{-1.5ex}
\footnotesize
\begin{list}{}{\setlength{\itemsep}{-\parsep}\setlength{\itemindent}{-\leftmargin}}
\item \NWtxtMacroRefIn\ \NWlink{weft?}{?}.

\item{}
\end{list}
\end{minipage}\vspace{4ex}
\end{flushleft}
\section{Global State}

The source manager's state consists of the mapped buffer, its size,
and the file path.  These are accessed by the lexer when scanning.

\begin{flushleft} \small
\begin{minipage}{\linewidth}\label{scrap53}\raggedright\small
\NWtarget{weft?}{} $\langle\,${\itshape Source: data}\nobreak\ {\footnotesize {?}}$\,\rangle\equiv$
\vspace{-1ex}
\begin{list}{}{} \item
\mbox{}\verb@@\\
\mbox{}\verb@// -- Source manager global state ------------------------------------@\\
\mbox{}\verb@.data@\\
\mbox{}\verb@.p2align 3@\\
\mbox{}\verb@.globl _kcc_src_base@\\
\mbox{}\verb@_kcc_src_base:  .quad 0       // pointer to mmap'd file contents@\\
\mbox{}\verb@.globl _kcc_src_size@\\
\mbox{}\verb@_kcc_src_size:  .quad 0       // file size in bytes@\\
\mbox{}\verb@.globl _kcc_src_path@\\
\mbox{}\verb@_kcc_src_path:  .quad 0       // pointer to file path string@\\
\mbox{}\verb@@\\
\mbox{}\verb@.section __TEXT,__cstring,cstring_literals@\\
\mbox{}\verb@_kcc_err_open:   .asciz "kcc: error: cannot open file: "@\\
\mbox{}\verb@_kcc_err_read:   .asciz "kcc: error: cannot read file\n"@\\
\mbox{}\verb@_kcc_err_empty:  .asciz "kcc: error: source file is empty\n"@\\
\mbox{}\verb@_kcc_err_mmap:   .asciz "kcc: error: cannot map file\n"@\\
\mbox{}\verb@_kcc_src_newline: .asciz "\n"@\\
\mbox{}\verb@@{\NWsep}
\end{list}
\vspace{-1.5ex}
\footnotesize
\begin{list}{}{\setlength{\itemsep}{-\parsep}\setlength{\itemindent}{-\leftmargin}}
\item \NWtxtMacroRefIn\ \NWlink{weft?}{?}.

\item{}
\end{list}
\end{minipage}\vspace{4ex}
\end{flushleft}
% ═══════════════════════════════════════════════════════════════════
% Part III: The Frontend
% ═══════════════════════════════════════════════════════════════════
\part{The Frontend}

% 15-source-encoding.weft — Translation Phase 1: Trigraph Replacement
%
% ISO 9899:2018 §5.1.1.2 paragraph 1 specifies the first phase of
% translation: "Physical source file multibyte characters are mapped,
% in an implementation-defined manner, to the source character set
% [...] if necessary.  Trigraph sequences (§5.2.1.1) are replaced by
% corresponding single-character internal representations."
%
% Trigraphs are a historical artifact from an era when some keyboards
% lacked certain characters.  Each trigraph is a three-character sequence
% beginning with ?? that maps to a single replacement character.  While
% rare in modern code, conformance requires handling them.

\chapter{Source Encoding --- Translation Phase 1}
\label{ch:source-encoding}

\section{Trigraph Sequences}

\stdref{5.2.1.1} defines exactly nine trigraph sequences:

\begin{center}
\begin{tabular}{cl@{\qquad}cl@{\qquad}cl}
\hline
Trigraph & Replacement & Trigraph & Replacement & Trigraph & Replacement \\
\hline
\texttt{??=} & \texttt{\#} & \texttt{??(} & \texttt{[} & \texttt{??/} & \texttt{\textbackslash} \\
\texttt{??)} & \texttt{]} & \texttt{??'} & \texttt{\^{}} & \texttt{??<} & \texttt{\{} \\
\texttt{??>} & \texttt{\}} & \texttt{??!} & \texttt{|} & \texttt{??-} & \texttt{\~{}} \\
\hline
\end{tabular}
\end{center}

The replacement must occur before any other processing.  In particular,
\texttt{??/} produces a backslash that participates in line splicing
(Phase~2), so Phase~1 must complete before Phase~2 begins.

\section{Implementation}

\texttt{kcc\_source\_encode} performs trigraph replacement on the
source buffer.  Because the mmap'd file is read-only
(\texttt{PROT\_READ}), we always copy the source into a writable
arena-allocated buffer.  This copy also prepares the buffer for
in-place line splicing in Phase~2.

The function scans the source with a read pointer and writes to the
new buffer.  When \texttt{??} is encountered, the third character is
checked against the nine trigraph mappings.  If a match is found, the
single replacement character is written and the read pointer advances
by~3.  Otherwise, a single \texttt{?} is written and the read pointer
advances by~1.

\begin{flushleft} \small
\begin{minipage}{\linewidth}\label{scrap54}\raggedright\small
\NWtarget{weft?}{} $\langle\,${\itshape Encoding: translate}\nobreak\ {\footnotesize {?}}$\,\rangle\equiv$
\vspace{-1ex}
\begin{list}{}{} \item
\mbox{}\verb@@\\
\mbox{}\verb@// -- kcc_source_encode --------------------------------------------------@\\
\mbox{}\verb@// void kcc_source_encode(void)@\\
\mbox{}\verb@// Translation Phase 1: replaces trigraphs in the source buffer.@\\
\mbox{}\verb@// Allocates a writable arena copy of the mmap'd source, performs@\\
\mbox{}\verb@// trigraph replacement during the copy, and updates _kcc_src_base@\\
\mbox{}\verb@// and _kcc_src_size to point to the new buffer.@\\
\mbox{}\verb@.globl _kcc_source_encode@\\
\mbox{}\verb@.p2align 2@\\
\mbox{}\verb@_kcc_source_encode:@\\
\mbox{}\verb@    stp x29, x30, [sp, #-64]!@\\
\mbox{}\verb@    mov x29, sp@\\
\mbox{}\verb@    stp x19, x20, [sp, #16]@\\
\mbox{}\verb@    stp x21, x22, [sp, #32]@\\
\mbox{}\verb@    stp x23, x24, [sp, #48]@\\
\mbox{}\verb@@\\
\mbox{}\verb@    // Load source base and size@\\
\mbox{}\verb@    adrp x19, _kcc_src_base@{\tt @}\verb@PAGE@\\
\mbox{}\verb@    add  x19, x19, _kcc_src_base@{\tt @}\verb@PAGEOFF@\\
\mbox{}\verb@    ldr  x20, [x19]              // x20 = read pointer (old base)@\\
\mbox{}\verb@@\\
\mbox{}\verb@    adrp x21, _kcc_src_size@{\tt @}\verb@PAGE@\\
\mbox{}\verb@    add  x21, x21, _kcc_src_size@{\tt @}\verb@PAGEOFF@\\
\mbox{}\verb@    ldr  x22, [x21]              // x22 = source size@\\
\mbox{}\verb@@\\
\mbox{}\verb@    // Allocate writable buffer from arena (same size as original)@\\
\mbox{}\verb@    mov x0, x22@\\
\mbox{}\verb@    bl _kcc_arena_alloc@\\
\mbox{}\verb@    mov x23, x0                  // x23 = write base (new buffer)@\\
\mbox{}\verb@    mov x24, x0                  // x24 = write pointer@\\
\mbox{}\verb@@\\
\mbox{}\verb@    // Compute end of source@\\
\mbox{}\verb@    add x22, x20, x22            // x22 = read end@\\
\mbox{}\verb@@\\
\mbox{}\verb@.Lenc_loop:@\\
\mbox{}\verb@    cmp x20, x22@\\
\mbox{}\verb@    b.hs .Lenc_done@\\
\mbox{}\verb@@\\
\mbox{}\verb@    ldrb w0, [x20]@\\
\mbox{}\verb@@\\
\mbox{}\verb@    // Check for '?' — potential trigraph start@\\
\mbox{}\verb@    cmp w0, #'?'@\\
\mbox{}\verb@    b.ne .Lenc_copy_byte@\\
\mbox{}\verb@@\\
\mbox{}\verb@    // Need at least 2 more characters for a trigraph@\\
\mbox{}\verb@    add x1, x20, #2@\\
\mbox{}\verb@    cmp x1, x22@\\
\mbox{}\verb@    b.hs .Lenc_copy_byte@\\
\mbox{}\verb@@\\
\mbox{}\verb@    // Check second '?'@\\
\mbox{}\verb@    ldrb w1, [x20, #1]@\\
\mbox{}\verb@    cmp w1, #'?'@\\
\mbox{}\verb@    b.ne .Lenc_copy_byte@\\
\mbox{}\verb@@\\
\mbox{}\verb@    // Check third character for trigraph match@\\
\mbox{}\verb@    ldrb w1, [x20, #2]@\\
\mbox{}\verb@@\\
\mbox{}\verb@    cmp w1, #'='@\\
\mbox{}\verb@    b.eq .Lenc_trig_hash       // ??= → #@\\
\mbox{}\verb@    cmp w1, #'('@\\
\mbox{}\verb@    b.eq .Lenc_trig_lbracket   // ??( → [@\\
\mbox{}\verb@    cmp w1, #'/'@\\
\mbox{}\verb@    b.eq .Lenc_trig_backslash  // ??/ → backslash@\\
\mbox{}\verb@    cmp w1, #')'@\\
\mbox{}\verb@    b.eq .Lenc_trig_rbracket   // ??) → ]@\\
\mbox{}\verb@    cmp w1, #'\''@\\
\mbox{}\verb@    b.eq .Lenc_trig_caret      // ??' → ^@\\
\mbox{}\verb@    cmp w1, #'<'@\\
\mbox{}\verb@    b.eq .Lenc_trig_lbrace     // ??< → {@\\
\mbox{}\verb@    cmp w1, #'!'@\\
\mbox{}\verb@    b.eq .Lenc_trig_pipe       // ??! → |@\\
\mbox{}\verb@    cmp w1, #'>'@\\
\mbox{}\verb@    b.eq .Lenc_trig_rbrace     // ??> → }@\\
\mbox{}\verb@    cmp w1, #'-'@\\
\mbox{}\verb@    b.eq .Lenc_trig_tilde      // ??- → ~@\\
\mbox{}\verb@@\\
\mbox{}\verb@    // Not a valid trigraph — copy the first '?' and advance by 1@\\
\mbox{}\verb@    b .Lenc_copy_byte@\\
\mbox{}\verb@@\\
\mbox{}\verb@.Lenc_trig_hash:@\\
\mbox{}\verb@    mov w0, #'#'@\\
\mbox{}\verb@    b .Lenc_write_trig@\\
\mbox{}\verb@.Lenc_trig_lbracket:@\\
\mbox{}\verb@    mov w0, #'['@\\
\mbox{}\verb@    b .Lenc_write_trig@\\
\mbox{}\verb@.Lenc_trig_backslash:@\\
\mbox{}\verb@    mov w0, #'\\'@\\
\mbox{}\verb@    b .Lenc_write_trig@\\
\mbox{}\verb@.Lenc_trig_rbracket:@\\
\mbox{}\verb@    mov w0, #']'@\\
\mbox{}\verb@    b .Lenc_write_trig@\\
\mbox{}\verb@.Lenc_trig_caret:@\\
\mbox{}\verb@    mov w0, #'^'@\\
\mbox{}\verb@    b .Lenc_write_trig@\\
\mbox{}\verb@.Lenc_trig_lbrace:@\\
\mbox{}\verb@    mov w0, #'{'@\\
\mbox{}\verb@    b .Lenc_write_trig@\\
\mbox{}\verb@.Lenc_trig_pipe:@\\
\mbox{}\verb@    mov w0, #'|'@\\
\mbox{}\verb@    b .Lenc_write_trig@\\
\mbox{}\verb@.Lenc_trig_rbrace:@\\
\mbox{}\verb@    mov w0, #'}'@\\
\mbox{}\verb@    b .Lenc_write_trig@\\
\mbox{}\verb@.Lenc_trig_tilde:@\\
\mbox{}\verb@    mov w0, #'~'@\\
\mbox{}\verb@    // fall through to write_trig@\\
\mbox{}\verb@@\\
\mbox{}\verb@.Lenc_write_trig:@\\
\mbox{}\verb@    strb w0, [x24], #1          // write replacement character@\\
\mbox{}\verb@    add x20, x20, #3            // skip 3 source characters@\\
\mbox{}\verb@    b .Lenc_loop@\\
\mbox{}\verb@@\\
\mbox{}\verb@.Lenc_copy_byte:@\\
\mbox{}\verb@    ldrb w0, [x20], #1          // read and advance@\\
\mbox{}\verb@    strb w0, [x24], #1          // write and advance@\\
\mbox{}\verb@    b .Lenc_loop@\\
\mbox{}\verb@@\\
\mbox{}\verb@.Lenc_done:@\\
\mbox{}\verb@    // Update global source state@\\
\mbox{}\verb@    str x23, [x19]               // _kcc_src_base = new buffer@\\
\mbox{}\verb@    sub x24, x24, x23           // compute actual output size@\\
\mbox{}\verb@    str x24, [x21]               // _kcc_src_size = output size@\\
\mbox{}\verb@@\\
\mbox{}\verb@    ldp x23, x24, [sp, #48]@\\
\mbox{}\verb@    ldp x21, x22, [sp, #32]@\\
\mbox{}\verb@    ldp x19, x20, [sp, #16]@\\
\mbox{}\verb@    ldp x29, x30, [sp], #64@\\
\mbox{}\verb@    ret@\\
\mbox{}\verb@@{\NWsep}
\end{list}
\vspace{-1.5ex}
\footnotesize
\begin{list}{}{\setlength{\itemsep}{-\parsep}\setlength{\itemindent}{-\leftmargin}}
\item \NWtxtMacroRefIn\ \NWlink{weft?}{?}.

\item{}
\end{list}
\end{minipage}\vspace{4ex}
\end{flushleft}
% 16-line-splicing.weft — Translation Phase 2: Line Splicing
%
% ISO 9899:2018 §5.1.1.2 paragraph 2 specifies the second phase
% of translation: "Each instance of a backslash character (\)
% immediately followed by a new-line character is deleted, splicing
% physical source lines to form logical source lines."
%
% This enables macros and string literals to span multiple physical
% lines in the source file.  Phase 2 must occur after Phase 1
% (trigraph replacement), because the trigraph ??/ produces a
% backslash that participates in line splicing.

\chapter{Line Splicing --- Translation Phase 2}
\label{ch:line-splicing}

\section{Design}

Line splicing removes every \texttt{\textbackslash}\texttt{\\n} pair
from the source.  Since this operation can only \emph{shrink} the
buffer (deleting 2 characters per splice), it can be performed
in-place on the writable arena buffer produced by Phase~1.

The function uses a read pointer and a write pointer, both starting
at \texttt{src\_base}.  When the read pointer encounters
\texttt{\textbackslash} followed by \texttt{\\n} (0x0A), both
characters are skipped.  Otherwise, the byte is copied from read
to write position.  The final buffer size is the write pointer's
offset from the base.

\begin{flushleft} \small
\begin{minipage}{\linewidth}\label{scrap55}\raggedright\small
\NWtarget{weft?}{} $\langle\,${\itshape Splice: translate}\nobreak\ {\footnotesize {?}}$\,\rangle\equiv$
\vspace{-1ex}
\begin{list}{}{} \item
\mbox{}\verb@@\\
\mbox{}\verb@// -- kcc_source_splice --------------------------------------------------@\\
\mbox{}\verb@// void kcc_source_splice(void)@\\
\mbox{}\verb@// Translation Phase 2: removes backslash-newline sequences in-place.@\\
\mbox{}\verb@// Requires that _kcc_src_base points to a writable buffer (as@\\
\mbox{}\verb@// guaranteed by Phase 1's arena copy).@\\
\mbox{}\verb@// Updates _kcc_src_size to reflect the (possibly shorter) buffer.@\\
\mbox{}\verb@.globl _kcc_source_splice@\\
\mbox{}\verb@.p2align 2@\\
\mbox{}\verb@_kcc_source_splice:@\\
\mbox{}\verb@    stp x29, x30, [sp, #-48]!@\\
\mbox{}\verb@    mov x29, sp@\\
\mbox{}\verb@    stp x19, x20, [sp, #16]@\\
\mbox{}\verb@    stp x21, x22, [sp, #32]@\\
\mbox{}\verb@@\\
\mbox{}\verb@    // Load source base and size@\\
\mbox{}\verb@    adrp x19, _kcc_src_base@{\tt @}\verb@PAGE@\\
\mbox{}\verb@    add  x19, x19, _kcc_src_base@{\tt @}\verb@PAGEOFF@\\
\mbox{}\verb@    ldr  x20, [x19]              // x20 = base (read pointer)@\\
\mbox{}\verb@@\\
\mbox{}\verb@    adrp x21, _kcc_src_size@{\tt @}\verb@PAGE@\\
\mbox{}\verb@    add  x21, x21, _kcc_src_size@{\tt @}\verb@PAGEOFF@\\
\mbox{}\verb@    ldr  x22, [x21]              // x22 = size@\\
\mbox{}\verb@@\\
\mbox{}\verb@    // Compute read end@\\
\mbox{}\verb@    add x1, x20, x22             // x1 = read end@\\
\mbox{}\verb@    mov x2, x20                  // x2 = write pointer@\\
\mbox{}\verb@    mov x3, x20                  // x3 = read pointer@\\
\mbox{}\verb@@\\
\mbox{}\verb@.Lspl_loop:@\\
\mbox{}\verb@    cmp x3, x1@\\
\mbox{}\verb@    b.hs .Lspl_done@\\
\mbox{}\verb@@\\
\mbox{}\verb@    ldrb w4, [x3]@\\
\mbox{}\verb@@\\
\mbox{}\verb@    // Check for backslash@\\
\mbox{}\verb@    cmp w4, #'\\'@\\
\mbox{}\verb@    b.ne .Lspl_copy@\\
\mbox{}\verb@@\\
\mbox{}\verb@    // Check if next character is newline@\\
\mbox{}\verb@    add x5, x3, #1@\\
\mbox{}\verb@    cmp x5, x1@\\
\mbox{}\verb@    b.hs .Lspl_copy              // backslash at very end — keep it@\\
\mbox{}\verb@@\\
\mbox{}\verb@    ldrb w6, [x5]@\\
\mbox{}\verb@    cmp w6, #0x0A                // newline?@\\
\mbox{}\verb@    b.ne .Lspl_copy@\\
\mbox{}\verb@@\\
\mbox{}\verb@    // Skip the backslash-newline pair@\\
\mbox{}\verb@    add x3, x3, #2@\\
\mbox{}\verb@    b .Lspl_loop@\\
\mbox{}\verb@@\\
\mbox{}\verb@.Lspl_copy:@\\
\mbox{}\verb@    strb w4, [x2], #1           // copy byte, advance write pointer@\\
\mbox{}\verb@    add x3, x3, #1              // advance read pointer@\\
\mbox{}\verb@    b .Lspl_loop@\\
\mbox{}\verb@@\\
\mbox{}\verb@.Lspl_done:@\\
\mbox{}\verb@    // Update size = write pointer - base@\\
\mbox{}\verb@    sub x2, x2, x20@\\
\mbox{}\verb@    str x2, [x21]                // _kcc_src_size = new size@\\
\mbox{}\verb@@\\
\mbox{}\verb@    ldp x21, x22, [sp, #32]@\\
\mbox{}\verb@    ldp x19, x20, [sp, #16]@\\
\mbox{}\verb@    ldp x29, x30, [sp], #48@\\
\mbox{}\verb@    ret@\\
\mbox{}\verb@@{\NWsep}
\end{list}
\vspace{-1.5ex}
\footnotesize
\begin{list}{}{\setlength{\itemsep}{-\parsep}\setlength{\itemindent}{-\leftmargin}}
\item \NWtxtMacroRefIn\ \NWlink{weft?}{?}.

\item{}
\end{list}
\end{minipage}\vspace{4ex}
\end{flushleft}
% 17-lexer.weft — Lexical Analysis (C17 Tokenizer)
%
% The lexer converts a stream of source characters into a stream of
% tokens according to ISO 9899:2018 section 6.4.  This is Phase 3 of
% translation (section 5.1.1.2 paragraph 3): each preprocessing token
% is the longest sequence of characters matching a token pattern.
%
% The lexer operates directly on the mmap'd source buffer provided
% by the source manager.  It maintains a cursor (current position),
% tracks line and column numbers, and produces one token at a time.
%
% Token value parsing (converting "42" to the integer 42) is deferred
% to Phase 5 (section 5.1.1.2 paragraph 5).

\chapter{Lexical Analysis}
\label{ch:lexer}

\section{Token Types}

Every C17 token falls into one of the categories defined in \stdref{6.4}:
keywords, identifiers, constants, string literals, and punctuators.
We assign a unique integer code to each.  Codes 0--9 are structural,
10--53 are the 44 keywords, and 60+ are punctuators.

\begin{flushleft} \small
\begin{minipage}{\linewidth}\label{scrap56}\raggedright\small
\NWtarget{weft?}{} $\langle\,${\itshape Lexer: token constants}\nobreak\ {\footnotesize {?}}$\,\rangle\equiv$
\vspace{-1ex}
\begin{list}{}{} \item
\mbox{}\verb@@\\
\mbox{}\verb@// -- Token type codes -----------------------------------------------@\\
\mbox{}\verb@// Structural tokens@\\
\mbox{}\verb@.equ TOK_EOF,          0@\\
\mbox{}\verb@.equ TOK_IDENT,        1@\\
\mbox{}\verb@.equ TOK_NUMBER,       2      // pp-number (integer or float)@\\
\mbox{}\verb@.equ TOK_CHAR_LIT,     3      // character constant@\\
\mbox{}\verb@.equ TOK_STRING_LIT,   4      // string literal@\\
\mbox{}\verb@.equ TOK_UNKNOWN,      5      // unrecognized character@\\
\mbox{}\verb@@\\
\mbox{}\verb@// C17 keywords (section 6.4.1) — 44 keywords, codes 10–53@\\
\mbox{}\verb@.equ TOK_AUTO,        10@\\
\mbox{}\verb@.equ TOK_BREAK,       11@\\
\mbox{}\verb@.equ TOK_CASE,        12@\\
\mbox{}\verb@.equ TOK_CHAR,        13@\\
\mbox{}\verb@.equ TOK_CONST,       14@\\
\mbox{}\verb@.equ TOK_CONTINUE,    15@\\
\mbox{}\verb@.equ TOK_DEFAULT,     16@\\
\mbox{}\verb@.equ TOK_DO,          17@\\
\mbox{}\verb@.equ TOK_DOUBLE,      18@\\
\mbox{}\verb@.equ TOK_ELSE,        19@\\
\mbox{}\verb@.equ TOK_ENUM,        20@\\
\mbox{}\verb@.equ TOK_EXTERN,      21@\\
\mbox{}\verb@.equ TOK_FLOAT,       22@\\
\mbox{}\verb@.equ TOK_FOR,         23@\\
\mbox{}\verb@.equ TOK_GOTO,        24@\\
\mbox{}\verb@.equ TOK_IF,          25@\\
\mbox{}\verb@.equ TOK_INLINE,      26@\\
\mbox{}\verb@.equ TOK_INT,         27@\\
\mbox{}\verb@.equ TOK_LONG,        28@\\
\mbox{}\verb@.equ TOK_REGISTER,    29@\\
\mbox{}\verb@.equ TOK_RESTRICT,    30@\\
\mbox{}\verb@.equ TOK_RETURN,      31@\\
\mbox{}\verb@.equ TOK_SHORT,       32@\\
\mbox{}\verb@.equ TOK_SIGNED,      33@\\
\mbox{}\verb@.equ TOK_SIZEOF,      34@\\
\mbox{}\verb@.equ TOK_STATIC,      35@\\
\mbox{}\verb@.equ TOK_STRUCT,      36@\\
\mbox{}\verb@.equ TOK_SWITCH,      37@\\
\mbox{}\verb@.equ TOK_TYPEDEF,     38@\\
\mbox{}\verb@.equ TOK_UNION,       39@\\
\mbox{}\verb@.equ TOK_UNSIGNED,    40@\\
\mbox{}\verb@.equ TOK_VOID,        41@\\
\mbox{}\verb@.equ TOK_VOLATILE,    42@\\
\mbox{}\verb@.equ TOK_WHILE,       43@\\
\mbox{}\verb@.equ TOK__ALIGNAS,    44@\\
\mbox{}\verb@.equ TOK__ALIGNOF,    45@\\
\mbox{}\verb@.equ TOK__ATOMIC,     46@\\
\mbox{}\verb@.equ TOK__BOOL,       47@\\
\mbox{}\verb@.equ TOK__COMPLEX,    48@\\
\mbox{}\verb@.equ TOK__GENERIC,    49@\\
\mbox{}\verb@.equ TOK__IMAGINARY,  50@\\
\mbox{}\verb@.equ TOK__NORETURN,   51@\\
\mbox{}\verb@.equ TOK__STATIC_ASSERT, 52@\\
\mbox{}\verb@.equ TOK__THREAD_LOCAL,  53@\\
\mbox{}\verb@@\\
\mbox{}\verb@// Punctuators (section 6.4.6) — codes 60+@\\
\mbox{}\verb@.equ TOK_LPAREN,      60      // (@\\
\mbox{}\verb@.equ TOK_RPAREN,      61      // )@\\
\mbox{}\verb@.equ TOK_LBRACE,      62      // {@\\
\mbox{}\verb@.equ TOK_RBRACE,      63      // }@\\
\mbox{}\verb@.equ TOK_LBRACKET,    64      // [@\\
\mbox{}\verb@.equ TOK_RBRACKET,    65      // ]@\\
\mbox{}\verb@.equ TOK_SEMI,        66      // ;@\\
\mbox{}\verb@.equ TOK_COMMA,       67      // ,@\\
\mbox{}\verb@.equ TOK_DOT,         68      // .@\\
\mbox{}\verb@.equ TOK_ARROW,       69      // ->@\\
\mbox{}\verb@.equ TOK_INC,         70      // ++@\\
\mbox{}\verb@.equ TOK_DEC,         71      // --@\\
\mbox{}\verb@.equ TOK_AMP,         72      // &@\\
\mbox{}\verb@.equ TOK_STAR,        73      // *@\\
\mbox{}\verb@.equ TOK_PLUS,        74      // +@\\
\mbox{}\verb@.equ TOK_MINUS,       75      // -@\\
\mbox{}\verb@.equ TOK_TILDE,       76      // ~@\\
\mbox{}\verb@.equ TOK_BANG,        77      // !@\\
\mbox{}\verb@.equ TOK_SLASH,       78      // /@\\
\mbox{}\verb@.equ TOK_PERCENT,     79      // %@\\
\mbox{}\verb@.equ TOK_SHL,         80      // <<@\\
\mbox{}\verb@.equ TOK_SHR,         81      // >>@\\
\mbox{}\verb@.equ TOK_LT,          82      // <@\\
\mbox{}\verb@.equ TOK_GT,          83      // >@\\
\mbox{}\verb@.equ TOK_LE,          84      // <=@\\
\mbox{}\verb@.equ TOK_GE,          85      // >=@\\
\mbox{}\verb@.equ TOK_EQ,          86      // ==@\\
\mbox{}\verb@.equ TOK_NE,          87      // !=@\\
\mbox{}\verb@.equ TOK_CARET,       88      // ^@\\
\mbox{}\verb@.equ TOK_PIPE,        89      // |@\\
\mbox{}\verb@.equ TOK_AND,         90      // &&@\\
\mbox{}\verb@.equ TOK_OR,          91      // ||@\\
\mbox{}\verb@.equ TOK_QUESTION,    92      // ?@\\
\mbox{}\verb@.equ TOK_COLON,       93      // :@\\
\mbox{}\verb@.equ TOK_ASSIGN,      94      // =@\\
\mbox{}\verb@.equ TOK_MUL_ASSIGN,  95      // *=@\\
\mbox{}\verb@.equ TOK_DIV_ASSIGN,  96      // /=@\\
\mbox{}\verb@.equ TOK_MOD_ASSIGN,  97      // %=@\\
\mbox{}\verb@.equ TOK_ADD_ASSIGN,  98      // +=@\\
\mbox{}\verb@.equ TOK_SUB_ASSIGN,  99      // -=@\\
\mbox{}\verb@.equ TOK_SHL_ASSIGN, 100      // <<=@\\
\mbox{}\verb@.equ TOK_SHR_ASSIGN, 101      // >>=@\\
\mbox{}\verb@.equ TOK_AND_ASSIGN, 102      // &=@\\
\mbox{}\verb@.equ TOK_XOR_ASSIGN, 103      // ^=@\\
\mbox{}\verb@.equ TOK_OR_ASSIGN,  104      // |=@\\
\mbox{}\verb@.equ TOK_HASH,       105      // #@\\
\mbox{}\verb@.equ TOK_DHASH,      106      // ##@\\
\mbox{}\verb@.equ TOK_ELLIPSIS,   107      // ...@\\
\mbox{}\verb@@{\NWsep}
\end{list}
\vspace{-1.5ex}
\footnotesize
\begin{list}{}{\setlength{\itemsep}{-\parsep}\setlength{\itemindent}{-\leftmargin}}
\item \NWtxtMacroRefIn\ \NWlink{weft?}{?}\NWlink{weft?}{, ?}.

\item{}
\end{list}
\end{minipage}\vspace{4ex}
\end{flushleft}
\section{Lexer Initialization}

\texttt{kcc\_lex\_init} prepares the lexer to scan a source buffer.
It sets the cursor to the start of the buffer and initializes the
line counter to~1.

\begin{flushleft} \small
\begin{minipage}{\linewidth}\label{scrap57}\raggedright\small
\NWtarget{weft?}{} $\langle\,${\itshape Lexer: init}\nobreak\ {\footnotesize {?}}$\,\rangle\equiv$
\vspace{-1ex}
\begin{list}{}{} \item
\mbox{}\verb@@\\
\mbox{}\verb@// -- kcc_lex_init ---------------------------------------------------@\\
\mbox{}\verb@// void kcc_lex_init(const char *base, uint64_t size)@\\
\mbox{}\verb@//   x0 = pointer to source buffer@\\
\mbox{}\verb@//   x1 = buffer size in bytes@\\
\mbox{}\verb@// Initializes lexer state to scan the given buffer.@\\
\mbox{}\verb@.globl _kcc_lex_init@\\
\mbox{}\verb@.p2align 2@\\
\mbox{}\verb@_kcc_lex_init:@\\
\mbox{}\verb@    stp x29, x30, [sp, #-16]!@\\
\mbox{}\verb@    mov x29, sp@\\
\mbox{}\verb@@\\
\mbox{}\verb@    // Set cursor and end@\\
\mbox{}\verb@    adrp x2, _kcc_lex_cursor@{\tt @}\verb@PAGE@\\
\mbox{}\verb@    add  x2, x2, _kcc_lex_cursor@{\tt @}\verb@PAGEOFF@\\
\mbox{}\verb@    str  x0, [x2]@\\
\mbox{}\verb@@\\
\mbox{}\verb@    add  x1, x0, x1             // end = base + size@\\
\mbox{}\verb@    adrp x2, _kcc_lex_end@{\tt @}\verb@PAGE@\\
\mbox{}\verb@    add  x2, x2, _kcc_lex_end@{\tt @}\verb@PAGEOFF@\\
\mbox{}\verb@    str  x1, [x2]@\\
\mbox{}\verb@@\\
\mbox{}\verb@    // Line starts at 1, line_start = base@\\
\mbox{}\verb@    mov  x1, #1@\\
\mbox{}\verb@    adrp x2, _kcc_lex_line@{\tt @}\verb@PAGE@\\
\mbox{}\verb@    add  x2, x2, _kcc_lex_line@{\tt @}\verb@PAGEOFF@\\
\mbox{}\verb@    str  x1, [x2]@\\
\mbox{}\verb@@\\
\mbox{}\verb@    adrp x2, _kcc_lex_line_start@{\tt @}\verb@PAGE@\\
\mbox{}\verb@    add  x2, x2, _kcc_lex_line_start@{\tt @}\verb@PAGEOFF@\\
\mbox{}\verb@    str  x0, [x2]@\\
\mbox{}\verb@@\\
\mbox{}\verb@    ldp x29, x30, [sp], #16@\\
\mbox{}\verb@    ret@\\
\mbox{}\verb@@{\NWsep}
\end{list}
\vspace{-1.5ex}
\footnotesize
\begin{list}{}{\setlength{\itemsep}{-\parsep}\setlength{\itemindent}{-\leftmargin}}
\item \NWtxtMacroRefIn\ \NWlink{weft?}{?}.

\item{}
\end{list}
\end{minipage}\vspace{4ex}
\end{flushleft}
\section{Next Token}

\texttt{kcc\_lex\_next} is the main entry point.  It skips whitespace
and comments, then dispatches to the appropriate scanner based on the
first character.  The result is stored in the global token variables.

\begin{flushleft} \small
\begin{minipage}{\linewidth}\label{scrap58}\raggedright\small
\NWtarget{weft?}{} $\langle\,${\itshape Lexer: next token}\nobreak\ {\footnotesize {?}}$\,\rangle\equiv$
\vspace{-1ex}
\begin{list}{}{} \item
\mbox{}\verb@@\\
\mbox{}\verb@// -- kcc_lex_next ---------------------------------------------------@\\
\mbox{}\verb@// void kcc_lex_next(void)@\\
\mbox{}\verb@// Advances to the next token.  Stores result in _kcc_tok_* globals.@\\
\mbox{}\verb@.globl _kcc_lex_next@\\
\mbox{}\verb@.p2align 2@\\
\mbox{}\verb@_kcc_lex_next:@\\
\mbox{}\verb@    stp x29, x30, [sp, #-48]!@\\
\mbox{}\verb@    mov x29, sp@\\
\mbox{}\verb@    stp x19, x20, [sp, #16]@\\
\mbox{}\verb@    stp x21, x22, [sp, #32]@\\
\mbox{}\verb@@\\
\mbox{}\verb@    // Load cursor and end@\\
\mbox{}\verb@    adrp x19, _kcc_lex_cursor@{\tt @}\verb@PAGE@\\
\mbox{}\verb@    add  x19, x19, _kcc_lex_cursor@{\tt @}\verb@PAGEOFF@\\
\mbox{}\verb@    ldr  x20, [x19]             // x20 = cursor@\\
\mbox{}\verb@@\\
\mbox{}\verb@    adrp x1, _kcc_lex_end@{\tt @}\verb@PAGE@\\
\mbox{}\verb@    add  x1, x1, _kcc_lex_end@{\tt @}\verb@PAGEOFF@\\
\mbox{}\verb@    ldr  x21, [x1]              // x21 = end@\\
\mbox{}\verb@@\\
\mbox{}\verb@.Llex_skip_ws:@\\
\mbox{}\verb@    // Check EOF@\\
\mbox{}\verb@    cmp x20, x21@\\
\mbox{}\verb@    b.hs .Llex_eof@\\
\mbox{}\verb@@\\
\mbox{}\verb@    ldrb w0, [x20]@\\
\mbox{}\verb@@\\
\mbox{}\verb@    // Check whitespace (space, tab, \r, \v, \f)@\\
\mbox{}\verb@    cmp w0, #' '@\\
\mbox{}\verb@    b.eq .Llex_ws_advance@\\
\mbox{}\verb@    cmp w0, #0x09               // \t@\\
\mbox{}\verb@    b.eq .Llex_ws_advance@\\
\mbox{}\verb@    cmp w0, #0x0D               // \r@\\
\mbox{}\verb@    b.eq .Llex_ws_advance@\\
\mbox{}\verb@    cmp w0, #0x0B               // \v@\\
\mbox{}\verb@    b.eq .Llex_ws_advance@\\
\mbox{}\verb@    cmp w0, #0x0C               // \f@\\
\mbox{}\verb@    b.eq .Llex_ws_advance@\\
\mbox{}\verb@@\\
\mbox{}\verb@    // Check newline@\\
\mbox{}\verb@    cmp w0, #0x0A               // \n@\\
\mbox{}\verb@    b.eq .Llex_newline@\\
\mbox{}\verb@@\\
\mbox{}\verb@    // Check for comments: // or /*@\\
\mbox{}\verb@    cmp w0, #'/'@\\
\mbox{}\verb@    b.eq .Llex_maybe_comment@\\
\mbox{}\verb@@\\
\mbox{}\verb@    // Not whitespace; proceed to character dispatch@\\
\mbox{}\verb@    b .Llex_dispatch@\\
\mbox{}\verb@@\\
\mbox{}\verb@.Llex_ws_advance:@\\
\mbox{}\verb@    add x20, x20, #1@\\
\mbox{}\verb@    b .Llex_skip_ws@\\
\mbox{}\verb@@\\
\mbox{}\verb@.Llex_newline:@\\
\mbox{}\verb@    add x20, x20, #1@\\
\mbox{}\verb@    // Increment line counter, update line_start@\\
\mbox{}\verb@    adrp x1, _kcc_lex_line@{\tt @}\verb@PAGE@\\
\mbox{}\verb@    add  x1, x1, _kcc_lex_line@{\tt @}\verb@PAGEOFF@\\
\mbox{}\verb@    ldr  x2, [x1]@\\
\mbox{}\verb@    add  x2, x2, #1@\\
\mbox{}\verb@    str  x2, [x1]@\\
\mbox{}\verb@    adrp x1, _kcc_lex_line_start@{\tt @}\verb@PAGE@\\
\mbox{}\verb@    add  x1, x1, _kcc_lex_line_start@{\tt @}\verb@PAGEOFF@\\
\mbox{}\verb@    str  x20, [x1]@\\
\mbox{}\verb@    b .Llex_skip_ws@\\
\mbox{}\verb@@\\
\mbox{}\verb@.Llex_maybe_comment:@\\
\mbox{}\verb@    // Check next char@\\
\mbox{}\verb@    add x1, x20, #1@\\
\mbox{}\verb@    cmp x1, x21@\\
\mbox{}\verb@    b.hs .Llex_dispatch          // just a '/' at EOF@\\
\mbox{}\verb@    ldrb w2, [x1]@\\
\mbox{}\verb@    cmp w2, #'/'@\\
\mbox{}\verb@    b.eq .Llex_line_comment@\\
\mbox{}\verb@    cmp w2, #'*'@\\
\mbox{}\verb@    b.eq .Llex_block_comment@\\
\mbox{}\verb@    b .Llex_dispatch             // just a '/' punctuator@\\
\mbox{}\verb@@\\
\mbox{}\verb@.Llex_line_comment:@\\
\mbox{}\verb@    // Skip until end of line@\\
\mbox{}\verb@    add x20, x20, #2            // skip "//"@\\
\mbox{}\verb@.Llex_lc_loop:@\\
\mbox{}\verb@    cmp x20, x21@\\
\mbox{}\verb@    b.hs .Llex_skip_ws@\\
\mbox{}\verb@    ldrb w0, [x20]@\\
\mbox{}\verb@    cmp w0, #0x0A@\\
\mbox{}\verb@    b.eq .Llex_skip_ws           // newline ends comment (don't skip it)@\\
\mbox{}\verb@    add x20, x20, #1@\\
\mbox{}\verb@    b .Llex_lc_loop@\\
\mbox{}\verb@@\\
\mbox{}\verb@.Llex_block_comment:@\\
\mbox{}\verb@    // Skip until */@\\
\mbox{}\verb@    add x20, x20, #2            // skip "/*"@\\
\mbox{}\verb@.Llex_bc_loop:@\\
\mbox{}\verb@    cmp x20, x21@\\
\mbox{}\verb@    b.hs .Llex_bc_unterminated@\\
\mbox{}\verb@    ldrb w0, [x20]@\\
\mbox{}\verb@    cmp w0, #0x0A@\\
\mbox{}\verb@    b.eq .Llex_bc_newline@\\
\mbox{}\verb@    cmp w0, #'*'@\\
\mbox{}\verb@    b.ne .Llex_bc_advance@\\
\mbox{}\verb@    // Check for '/' after '*'@\\
\mbox{}\verb@    add x1, x20, #1@\\
\mbox{}\verb@    cmp x1, x21@\\
\mbox{}\verb@    b.hs .Llex_bc_advance@\\
\mbox{}\verb@    ldrb w2, [x1]@\\
\mbox{}\verb@    cmp w2, #'/'@\\
\mbox{}\verb@    b.ne .Llex_bc_advance@\\
\mbox{}\verb@    add x20, x1, #1             // skip past "*/"@\\
\mbox{}\verb@    b .Llex_skip_ws@\\
\mbox{}\verb@@\\
\mbox{}\verb@.Llex_bc_newline:@\\
\mbox{}\verb@    add x20, x20, #1@\\
\mbox{}\verb@    adrp x1, _kcc_lex_line@{\tt @}\verb@PAGE@\\
\mbox{}\verb@    add  x1, x1, _kcc_lex_line@{\tt @}\verb@PAGEOFF@\\
\mbox{}\verb@    ldr  x2, [x1]@\\
\mbox{}\verb@    add  x2, x2, #1@\\
\mbox{}\verb@    str  x2, [x1]@\\
\mbox{}\verb@    adrp x1, _kcc_lex_line_start@{\tt @}\verb@PAGE@\\
\mbox{}\verb@    add  x1, x1, _kcc_lex_line_start@{\tt @}\verb@PAGEOFF@\\
\mbox{}\verb@    str  x20, [x1]@\\
\mbox{}\verb@    b .Llex_bc_loop@\\
\mbox{}\verb@@\\
\mbox{}\verb@.Llex_bc_advance:@\\
\mbox{}\verb@    add x20, x20, #1@\\
\mbox{}\verb@    b .Llex_bc_loop@\\
\mbox{}\verb@@\\
\mbox{}\verb@.Llex_bc_unterminated:@\\
\mbox{}\verb@    // Report error at current position via diagnostic engine@\\
\mbox{}\verb@    adrp x0, _kcc_lex_line@{\tt @}\verb@PAGE@\\
\mbox{}\verb@    add  x0, x0, _kcc_lex_line@{\tt @}\verb@PAGEOFF@\\
\mbox{}\verb@    ldr  x0, [x0]                // line@\\
\mbox{}\verb@    mov  x1, #1                  // col (approximate — comment start unknown)@\\
\mbox{}\verb@    adrp x2, _kcc_err_unterm_comment@{\tt @}\verb@PAGE@\\
\mbox{}\verb@    add  x2, x2, _kcc_err_unterm_comment@{\tt @}\verb@PAGEOFF@\\
\mbox{}\verb@    bl _kcc_diag_error@\\
\mbox{}\verb@    // Treat as EOF@\\
\mbox{}\verb@    b .Llex_eof@\\
\mbox{}\verb@@\\
\mbox{}\verb@.Llex_eof:@\\
\mbox{}\verb@    // Record token location before setting type@\\
\mbox{}\verb@    bl _kcc_lex_record_loc@\\
\mbox{}\verb@    mov w0, #TOK_EOF@\\
\mbox{}\verb@    adrp x1, _kcc_tok_type@{\tt @}\verb@PAGE@\\
\mbox{}\verb@    add  x1, x1, _kcc_tok_type@{\tt @}\verb@PAGEOFF@\\
\mbox{}\verb@    str  x0, [x1]@\\
\mbox{}\verb@    str  xzr, [x19]             // cursor stays at end@\\
\mbox{}\verb@    // tok_start = cursor, tok_len = 0@\\
\mbox{}\verb@    adrp x1, _kcc_tok_start@{\tt @}\verb@PAGE@\\
\mbox{}\verb@    add  x1, x1, _kcc_tok_start@{\tt @}\verb@PAGEOFF@\\
\mbox{}\verb@    str  x20, [x1]@\\
\mbox{}\verb@    adrp x1, _kcc_tok_len@{\tt @}\verb@PAGE@\\
\mbox{}\verb@    add  x1, x1, _kcc_tok_len@{\tt @}\verb@PAGEOFF@\\
\mbox{}\verb@    str  xzr, [x1]@\\
\mbox{}\verb@    b .Llex_done@\\
\mbox{}\verb@@\\
\mbox{}\verb@.Llex_dispatch:@\\
\mbox{}\verb@    // Save cursor state@\\
\mbox{}\verb@    str x20, [x19]@\\
\mbox{}\verb@@\\
\mbox{}\verb@    // Record location (line:col) for this token@\\
\mbox{}\verb@    bl _kcc_lex_record_loc@\\
\mbox{}\verb@@\\
\mbox{}\verb@    // Reload cursor@\\
\mbox{}\verb@    ldr x20, [x19]@\\
\mbox{}\verb@    ldrb w0, [x20]@\\
\mbox{}\verb@@\\
\mbox{}\verb@    // Dispatch based on character class@\\
\mbox{}\verb@    // Identifier start: a-z, A-Z, _@\\
\mbox{}\verb@    cmp w0, #'_'@\\
\mbox{}\verb@    b.eq .Llex_go_ident@\\
\mbox{}\verb@    cmp w0, #'a'@\\
\mbox{}\verb@    b.hs .Llex_check_lower@\\
\mbox{}\verb@    cmp w0, #'A'@\\
\mbox{}\verb@    b.hs .Llex_check_upper@\\
\mbox{}\verb@    b .Llex_check_digit@\\
\mbox{}\verb@@\\
\mbox{}\verb@.Llex_check_lower:@\\
\mbox{}\verb@    cmp w0, #'z'@\\
\mbox{}\verb@    b.ls .Llex_go_ident@\\
\mbox{}\verb@    b .Llex_check_digit@\\
\mbox{}\verb@@\\
\mbox{}\verb@.Llex_check_upper:@\\
\mbox{}\verb@    cmp w0, #'Z'@\\
\mbox{}\verb@    b.ls .Llex_go_ident@\\
\mbox{}\verb@    b .Llex_check_digit@\\
\mbox{}\verb@@\\
\mbox{}\verb@.Llex_go_ident:@\\
\mbox{}\verb@    bl _kcc_lex_scan_ident@\\
\mbox{}\verb@    b .Llex_done@\\
\mbox{}\verb@@\\
\mbox{}\verb@.Llex_check_digit:@\\
\mbox{}\verb@    cmp w0, #'0'@\\
\mbox{}\verb@    b.lo .Llex_check_string@\\
\mbox{}\verb@    cmp w0, #'9'@\\
\mbox{}\verb@    b.ls .Llex_go_number@\\
\mbox{}\verb@@\\
\mbox{}\verb@    // Check for '.' followed by digit (float like .5)@\\
\mbox{}\verb@    b .Llex_check_string@\\
\mbox{}\verb@@\\
\mbox{}\verb@.Llex_go_number:@\\
\mbox{}\verb@    bl _kcc_lex_scan_number@\\
\mbox{}\verb@    b .Llex_done@\\
\mbox{}\verb@@\\
\mbox{}\verb@.Llex_check_string:@\\
\mbox{}\verb@    cmp w0, #'"'@\\
\mbox{}\verb@    b.eq .Llex_go_string@\\
\mbox{}\verb@    cmp w0, #'\''@\\
\mbox{}\verb@    b.eq .Llex_go_char@\\
\mbox{}\verb@    cmp w0, #'.'@\\
\mbox{}\verb@    b.eq .Llex_check_dot_number@\\
\mbox{}\verb@    b .Llex_go_punct@\\
\mbox{}\verb@@\\
\mbox{}\verb@.Llex_check_dot_number:@\\
\mbox{}\verb@    // '.' might start a float (.5) or be a punctuator@\\
\mbox{}\verb@    add x1, x20, #1@\\
\mbox{}\verb@    cmp x1, x21@\\
\mbox{}\verb@    b.hs .Llex_go_punct          // '.' at EOF@\\
\mbox{}\verb@    ldrb w2, [x1]@\\
\mbox{}\verb@    cmp w2, #'0'@\\
\mbox{}\verb@    b.lo .Llex_go_punct@\\
\mbox{}\verb@    cmp w2, #'9'@\\
\mbox{}\verb@    b.ls .Llex_go_number         // '.digit' -> float@\\
\mbox{}\verb@    b .Llex_go_punct@\\
\mbox{}\verb@@\\
\mbox{}\verb@.Llex_go_string:@\\
\mbox{}\verb@    bl _kcc_lex_scan_string@\\
\mbox{}\verb@    b .Llex_done@\\
\mbox{}\verb@@\\
\mbox{}\verb@.Llex_go_char:@\\
\mbox{}\verb@    bl _kcc_lex_scan_char@\\
\mbox{}\verb@    b .Llex_done@\\
\mbox{}\verb@@\\
\mbox{}\verb@.Llex_go_punct:@\\
\mbox{}\verb@    bl _kcc_lex_scan_punct@\\
\mbox{}\verb@    b .Llex_done@\\
\mbox{}\verb@@\\
\mbox{}\verb@.Llex_done:@\\
\mbox{}\verb@    ldp x21, x22, [sp, #32]@\\
\mbox{}\verb@    ldp x19, x20, [sp, #16]@\\
\mbox{}\verb@    ldp x29, x30, [sp], #48@\\
\mbox{}\verb@    ret@\\
\mbox{}\verb@@{\NWsep}
\end{list}
\vspace{-1.5ex}
\footnotesize
\begin{list}{}{\setlength{\itemsep}{-\parsep}\setlength{\itemindent}{-\leftmargin}}
\item \NWtxtMacroRefIn\ \NWlink{weft?}{?}.

\item{}
\end{list}
\end{minipage}\vspace{4ex}
\end{flushleft}
\section{Location Recording}

Record the current line and column into the token globals.  Column is
computed as \texttt{cursor $-$ line\_start $+$ 1}.

\begin{flushleft} \small
\begin{minipage}{\linewidth}\label{scrap59}\raggedright\small
\NWtarget{weft?}{} $\langle\,${\itshape Lexer: record location}\nobreak\ {\footnotesize {?}}$\,\rangle\equiv$
\vspace{-1ex}
\begin{list}{}{} \item
\mbox{}\verb@@\\
\mbox{}\verb@// -- kcc_lex_record_loc ---------------------------------------------@\\
\mbox{}\verb@// Internal: records current line/col into tok globals.@\\
\mbox{}\verb@.p2align 2@\\
\mbox{}\verb@_kcc_lex_record_loc:@\\
\mbox{}\verb@    stp x29, x30, [sp, #-16]!@\\
\mbox{}\verb@    mov x29, sp@\\
\mbox{}\verb@@\\
\mbox{}\verb@    adrp x0, _kcc_lex_line@{\tt @}\verb@PAGE@\\
\mbox{}\verb@    add  x0, x0, _kcc_lex_line@{\tt @}\verb@PAGEOFF@\\
\mbox{}\verb@    ldr  x0, [x0]@\\
\mbox{}\verb@@\\
\mbox{}\verb@    adrp x1, _kcc_tok_line@{\tt @}\verb@PAGE@\\
\mbox{}\verb@    add  x1, x1, _kcc_tok_line@{\tt @}\verb@PAGEOFF@\\
\mbox{}\verb@    str  x0, [x1]@\\
\mbox{}\verb@@\\
\mbox{}\verb@    // col = cursor - line_start + 1@\\
\mbox{}\verb@    adrp x0, _kcc_lex_cursor@{\tt @}\verb@PAGE@\\
\mbox{}\verb@    add  x0, x0, _kcc_lex_cursor@{\tt @}\verb@PAGEOFF@\\
\mbox{}\verb@    ldr  x0, [x0]@\\
\mbox{}\verb@@\\
\mbox{}\verb@    adrp x1, _kcc_lex_line_start@{\tt @}\verb@PAGE@\\
\mbox{}\verb@    add  x1, x1, _kcc_lex_line_start@{\tt @}\verb@PAGEOFF@\\
\mbox{}\verb@    ldr  x1, [x1]@\\
\mbox{}\verb@@\\
\mbox{}\verb@    sub x0, x0, x1@\\
\mbox{}\verb@    add x0, x0, #1@\\
\mbox{}\verb@@\\
\mbox{}\verb@    adrp x1, _kcc_tok_col@{\tt @}\verb@PAGE@\\
\mbox{}\verb@    add  x1, x1, _kcc_tok_col@{\tt @}\verb@PAGEOFF@\\
\mbox{}\verb@    str  x0, [x1]@\\
\mbox{}\verb@@\\
\mbox{}\verb@    ldp x29, x30, [sp], #16@\\
\mbox{}\verb@    ret@\\
\mbox{}\verb@@{\NWsep}
\end{list}
\vspace{-1.5ex}
\footnotesize
\begin{list}{}{\setlength{\itemsep}{-\parsep}\setlength{\itemindent}{-\leftmargin}}
\item \NWtxtMacroRefIn\ \NWlink{weft?}{?}.

\item{}
\end{list}
\end{minipage}\vspace{4ex}
\end{flushleft}
\section{Identifier Scanner}

An identifier starts with a letter or underscore and continues with
letters, digits, or underscores (\stdref{6.4.2.1}).  After scanning,
we look up the identifier text in the keyword hash table; if found,
the token type is changed to the corresponding keyword code.

\begin{flushleft} \small
\begin{minipage}{\linewidth}\label{scrap60}\raggedright\small
\NWtarget{weft?}{} $\langle\,${\itshape Lexer: scan identifier}\nobreak\ {\footnotesize {?}}$\,\rangle\equiv$
\vspace{-1ex}
\begin{list}{}{} \item
\mbox{}\verb@@\\
\mbox{}\verb@// -- kcc_lex_scan_ident ---------------------------------------------@\\
\mbox{}\verb@// Internal: scans an identifier or keyword.@\\
\mbox{}\verb@// Cursor is at the first identifier character.@\\
\mbox{}\verb@.p2align 2@\\
\mbox{}\verb@_kcc_lex_scan_ident:@\\
\mbox{}\verb@    stp x29, x30, [sp, #-48]!@\\
\mbox{}\verb@    mov x29, sp@\\
\mbox{}\verb@    stp x19, x20, [sp, #16]@\\
\mbox{}\verb@    stp x21, x22, [sp, #32]@\\
\mbox{}\verb@@\\
\mbox{}\verb@    // Load cursor and end@\\
\mbox{}\verb@    adrp x19, _kcc_lex_cursor@{\tt @}\verb@PAGE@\\
\mbox{}\verb@    add  x19, x19, _kcc_lex_cursor@{\tt @}\verb@PAGEOFF@\\
\mbox{}\verb@    ldr  x20, [x19]             // cursor@\\
\mbox{}\verb@@\\
\mbox{}\verb@    adrp x1, _kcc_lex_end@{\tt @}\verb@PAGE@\\
\mbox{}\verb@    add  x1, x1, _kcc_lex_end@{\tt @}\verb@PAGEOFF@\\
\mbox{}\verb@    ldr  x21, [x1]              // end@\\
\mbox{}\verb@@\\
\mbox{}\verb@    mov x22, x20                // x22 = token start@\\
\mbox{}\verb@@\\
\mbox{}\verb@    // Advance past identifier characters@\\
\mbox{}\verb@.Lident_loop:@\\
\mbox{}\verb@    add x20, x20, #1@\\
\mbox{}\verb@    cmp x20, x21@\\
\mbox{}\verb@    b.hs .Lident_end@\\
\mbox{}\verb@    ldrb w0, [x20]@\\
\mbox{}\verb@    // Check: a-z, A-Z, 0-9, _@\\
\mbox{}\verb@    cmp w0, #'_'@\\
\mbox{}\verb@    b.eq .Lident_loop@\\
\mbox{}\verb@    cmp w0, #'a'@\\
\mbox{}\verb@    b.lo .Lident_not_lower@\\
\mbox{}\verb@    cmp w0, #'z'@\\
\mbox{}\verb@    b.ls .Lident_loop@\\
\mbox{}\verb@.Lident_not_lower:@\\
\mbox{}\verb@    cmp w0, #'A'@\\
\mbox{}\verb@    b.lo .Lident_not_upper@\\
\mbox{}\verb@    cmp w0, #'Z'@\\
\mbox{}\verb@    b.ls .Lident_loop@\\
\mbox{}\verb@.Lident_not_upper:@\\
\mbox{}\verb@    cmp w0, #'0'@\\
\mbox{}\verb@    b.lo .Lident_end@\\
\mbox{}\verb@    cmp w0, #'9'@\\
\mbox{}\verb@    b.ls .Lident_loop@\\
\mbox{}\verb@@\\
\mbox{}\verb@.Lident_end:@\\
\mbox{}\verb@    // Store cursor@\\
\mbox{}\verb@    str x20, [x19]@\\
\mbox{}\verb@@\\
\mbox{}\verb@    // Compute length@\\
\mbox{}\verb@    sub x2, x20, x22            // x2 = length@\\
\mbox{}\verb@@\\
\mbox{}\verb@    // Store token: start and length@\\
\mbox{}\verb@    adrp x1, _kcc_tok_start@{\tt @}\verb@PAGE@\\
\mbox{}\verb@    add  x1, x1, _kcc_tok_start@{\tt @}\verb@PAGEOFF@\\
\mbox{}\verb@    str  x22, [x1]@\\
\mbox{}\verb@@\\
\mbox{}\verb@    adrp x1, _kcc_tok_len@{\tt @}\verb@PAGE@\\
\mbox{}\verb@    add  x1, x1, _kcc_tok_len@{\tt @}\verb@PAGEOFF@\\
\mbox{}\verb@    str  x2, [x1]@\\
\mbox{}\verb@@\\
\mbox{}\verb@    // Default to identifier@\\
\mbox{}\verb@    mov x0, #TOK_IDENT@\\
\mbox{}\verb@    adrp x1, _kcc_tok_type@{\tt @}\verb@PAGE@\\
\mbox{}\verb@    add  x1, x1, _kcc_tok_type@{\tt @}\verb@PAGEOFF@\\
\mbox{}\verb@    str  x0, [x1]@\\
\mbox{}\verb@@\\
\mbox{}\verb@    // Check keyword table@\\
\mbox{}\verb@    adrp x0, _kcc_kw_table@{\tt @}\verb@PAGE@\\
\mbox{}\verb@    add  x0, x0, _kcc_kw_table@{\tt @}\verb@PAGEOFF@\\
\mbox{}\verb@    ldr  x0, [x0]               // keyword hash table@\\
\mbox{}\verb@    cbz  x0, .Lident_done       // not initialized yet@\\
\mbox{}\verb@@\\
\mbox{}\verb@    mov  x1, x22                 // key = token text@\\
\mbox{}\verb@    sub  x2, x20, x22           // keylen = token length@\\
\mbox{}\verb@    bl _kcc_ht_get@\\
\mbox{}\verb@    cbz x0, .Lident_done        // not a keyword@\\
\mbox{}\verb@@\\
\mbox{}\verb@    // It's a keyword: update token type@\\
\mbox{}\verb@    adrp x1, _kcc_tok_type@{\tt @}\verb@PAGE@\\
\mbox{}\verb@    add  x1, x1, _kcc_tok_type@{\tt @}\verb@PAGEOFF@\\
\mbox{}\verb@    str  x0, [x1]@\\
\mbox{}\verb@@\\
\mbox{}\verb@.Lident_done:@\\
\mbox{}\verb@    ldp x21, x22, [sp, #32]@\\
\mbox{}\verb@    ldp x19, x20, [sp, #16]@\\
\mbox{}\verb@    ldp x29, x30, [sp], #48@\\
\mbox{}\verb@    ret@\\
\mbox{}\verb@@{\NWsep}
\end{list}
\vspace{-1.5ex}
\footnotesize
\begin{list}{}{\setlength{\itemsep}{-\parsep}\setlength{\itemindent}{-\leftmargin}}
\item \NWtxtMacroRefIn\ \NWlink{weft?}{?}.

\item{}
\end{list}
\end{minipage}\vspace{4ex}
\end{flushleft}
\section{Number Scanner}

A preprocessing number (\stdref{6.4.8}) is surprisingly broad: it
starts with a digit (or \texttt{.}~followed by a digit) and continues
with digits, letters, underscores, dots, and \texttt{e/E/p/P} signs.
We scan the full pp-number here; value parsing is deferred to Phase~5.

\begin{flushleft} \small
\begin{minipage}{\linewidth}\label{scrap61}\raggedright\small
\NWtarget{weft?}{} $\langle\,${\itshape Lexer: scan number}\nobreak\ {\footnotesize {?}}$\,\rangle\equiv$
\vspace{-1ex}
\begin{list}{}{} \item
\mbox{}\verb@@\\
\mbox{}\verb@// -- kcc_lex_scan_number --------------------------------------------@\\
\mbox{}\verb@// Internal: scans a preprocessing number.@\\
\mbox{}\verb@.p2align 2@\\
\mbox{}\verb@_kcc_lex_scan_number:@\\
\mbox{}\verb@    stp x29, x30, [sp, #-48]!@\\
\mbox{}\verb@    mov x29, sp@\\
\mbox{}\verb@    stp x19, x20, [sp, #16]@\\
\mbox{}\verb@    stp x21, x22, [sp, #32]@\\
\mbox{}\verb@@\\
\mbox{}\verb@    adrp x19, _kcc_lex_cursor@{\tt @}\verb@PAGE@\\
\mbox{}\verb@    add  x19, x19, _kcc_lex_cursor@{\tt @}\verb@PAGEOFF@\\
\mbox{}\verb@    ldr  x20, [x19]@\\
\mbox{}\verb@@\\
\mbox{}\verb@    adrp x1, _kcc_lex_end@{\tt @}\verb@PAGE@\\
\mbox{}\verb@    add  x1, x1, _kcc_lex_end@{\tt @}\verb@PAGEOFF@\\
\mbox{}\verb@    ldr  x21, [x1]@\\
\mbox{}\verb@@\\
\mbox{}\verb@    mov x22, x20                // token start@\\
\mbox{}\verb@@\\
\mbox{}\verb@.Lnum_loop:@\\
\mbox{}\verb@    add x20, x20, #1@\\
\mbox{}\verb@    cmp x20, x21@\\
\mbox{}\verb@    b.hs .Lnum_end@\\
\mbox{}\verb@    ldrb w0, [x20]@\\
\mbox{}\verb@@\\
\mbox{}\verb@    // digits@\\
\mbox{}\verb@    cmp w0, #'0'@\\
\mbox{}\verb@    b.lo .Lnum_check_alpha@\\
\mbox{}\verb@    cmp w0, #'9'@\\
\mbox{}\verb@    b.ls .Lnum_loop@\\
\mbox{}\verb@@\\
\mbox{}\verb@.Lnum_check_alpha:@\\
\mbox{}\verb@    // letters (part of pp-number: suffixes, hex digits)@\\
\mbox{}\verb@    cmp w0, #'a'@\\
\mbox{}\verb@    b.lo .Lnum_check_upper@\\
\mbox{}\verb@    cmp w0, #'z'@\\
\mbox{}\verb@    b.ls .Lnum_check_exp@\\
\mbox{}\verb@    b .Lnum_check_dot@\\
\mbox{}\verb@@\\
\mbox{}\verb@.Lnum_check_upper:@\\
\mbox{}\verb@    cmp w0, #'A'@\\
\mbox{}\verb@    b.lo .Lnum_check_dot@\\
\mbox{}\verb@    cmp w0, #'Z'@\\
\mbox{}\verb@    b.ls .Lnum_check_exp@\\
\mbox{}\verb@    b .Lnum_check_dot@\\
\mbox{}\verb@@\\
\mbox{}\verb@.Lnum_check_exp:@\\
\mbox{}\verb@    // If e/E/p/P, next char might be +/-@\\
\mbox{}\verb@    cmp w0, #'e'@\\
\mbox{}\verb@    b.eq .Lnum_maybe_sign@\\
\mbox{}\verb@    cmp w0, #'E'@\\
\mbox{}\verb@    b.eq .Lnum_maybe_sign@\\
\mbox{}\verb@    cmp w0, #'p'@\\
\mbox{}\verb@    b.eq .Lnum_maybe_sign@\\
\mbox{}\verb@    cmp w0, #'P'@\\
\mbox{}\verb@    b.eq .Lnum_maybe_sign@\\
\mbox{}\verb@    b .Lnum_loop               // other letter, still pp-number@\\
\mbox{}\verb@@\\
\mbox{}\verb@.Lnum_maybe_sign:@\\
\mbox{}\verb@    // Peek at next char for +/-@\\
\mbox{}\verb@    add x1, x20, #1@\\
\mbox{}\verb@    cmp x1, x21@\\
\mbox{}\verb@    b.hs .Lnum_loop@\\
\mbox{}\verb@    ldrb w2, [x1]@\\
\mbox{}\verb@    cmp w2, #'+'@\\
\mbox{}\verb@    b.eq .Lnum_skip_sign@\\
\mbox{}\verb@    cmp w2, #'-'@\\
\mbox{}\verb@    b.eq .Lnum_skip_sign@\\
\mbox{}\verb@    b .Lnum_loop@\\
\mbox{}\verb@@\\
\mbox{}\verb@.Lnum_skip_sign:@\\
\mbox{}\verb@    add x20, x20, #1           // consume the letter@\\
\mbox{}\verb@    b .Lnum_loop                // sign consumed in next iteration@\\
\mbox{}\verb@@\\
\mbox{}\verb@.Lnum_check_dot:@\\
\mbox{}\verb@    cmp w0, #'.'@\\
\mbox{}\verb@    b.eq .Lnum_loop@\\
\mbox{}\verb@    cmp w0, #'_'@\\
\mbox{}\verb@    b.eq .Lnum_loop@\\
\mbox{}\verb@@\\
\mbox{}\verb@.Lnum_end:@\\
\mbox{}\verb@    str x20, [x19]@\\
\mbox{}\verb@@\\
\mbox{}\verb@    // Store token@\\
\mbox{}\verb@    sub x2, x20, x22@\\
\mbox{}\verb@    adrp x1, _kcc_tok_start@{\tt @}\verb@PAGE@\\
\mbox{}\verb@    add  x1, x1, _kcc_tok_start@{\tt @}\verb@PAGEOFF@\\
\mbox{}\verb@    str  x22, [x1]@\\
\mbox{}\verb@@\\
\mbox{}\verb@    adrp x1, _kcc_tok_len@{\tt @}\verb@PAGE@\\
\mbox{}\verb@    add  x1, x1, _kcc_tok_len@{\tt @}\verb@PAGEOFF@\\
\mbox{}\verb@    str  x2, [x1]@\\
\mbox{}\verb@@\\
\mbox{}\verb@    mov x0, #TOK_NUMBER@\\
\mbox{}\verb@    adrp x1, _kcc_tok_type@{\tt @}\verb@PAGE@\\
\mbox{}\verb@    add  x1, x1, _kcc_tok_type@{\tt @}\verb@PAGEOFF@\\
\mbox{}\verb@    str  x0, [x1]@\\
\mbox{}\verb@@\\
\mbox{}\verb@    ldp x21, x22, [sp, #32]@\\
\mbox{}\verb@    ldp x19, x20, [sp, #16]@\\
\mbox{}\verb@    ldp x29, x30, [sp], #48@\\
\mbox{}\verb@    ret@\\
\mbox{}\verb@@{\NWsep}
\end{list}
\vspace{-1.5ex}
\footnotesize
\begin{list}{}{\setlength{\itemsep}{-\parsep}\setlength{\itemindent}{-\leftmargin}}
\item \NWtxtMacroRefIn\ \NWlink{weft?}{?}.

\item{}
\end{list}
\end{minipage}\vspace{4ex}
\end{flushleft}
\section{String Literal Scanner}

A string literal starts and ends with \texttt{"} (\stdref{6.4.5}).
Within the string, backslash escapes are recognized (the actual
escape processing is deferred to Phase~5).  An unescaped newline
or end-of-file before the closing quote is an error.

\begin{flushleft} \small
\begin{minipage}{\linewidth}\label{scrap62}\raggedright\small
\NWtarget{weft?}{} $\langle\,${\itshape Lexer: scan string}\nobreak\ {\footnotesize {?}}$\,\rangle\equiv$
\vspace{-1ex}
\begin{list}{}{} \item
\mbox{}\verb@@\\
\mbox{}\verb@// -- kcc_lex_scan_string --------------------------------------------@\\
\mbox{}\verb@// Internal: scans a string literal (opening '"' at cursor).@\\
\mbox{}\verb@.p2align 2@\\
\mbox{}\verb@_kcc_lex_scan_string:@\\
\mbox{}\verb@    stp x29, x30, [sp, #-48]!@\\
\mbox{}\verb@    mov x29, sp@\\
\mbox{}\verb@    stp x19, x20, [sp, #16]@\\
\mbox{}\verb@    stp x21, x22, [sp, #32]@\\
\mbox{}\verb@@\\
\mbox{}\verb@    adrp x19, _kcc_lex_cursor@{\tt @}\verb@PAGE@\\
\mbox{}\verb@    add  x19, x19, _kcc_lex_cursor@{\tt @}\verb@PAGEOFF@\\
\mbox{}\verb@    ldr  x20, [x19]@\\
\mbox{}\verb@@\\
\mbox{}\verb@    adrp x1, _kcc_lex_end@{\tt @}\verb@PAGE@\\
\mbox{}\verb@    add  x1, x1, _kcc_lex_end@{\tt @}\verb@PAGEOFF@\\
\mbox{}\verb@    ldr  x21, [x1]@\\
\mbox{}\verb@@\\
\mbox{}\verb@    mov x22, x20                // token start (the opening ")@\\
\mbox{}\verb@    add x20, x20, #1           // skip opening "@\\
\mbox{}\verb@@\\
\mbox{}\verb@.Lstr_loop:@\\
\mbox{}\verb@    cmp x20, x21@\\
\mbox{}\verb@    b.hs .Lstr_unterminated@\\
\mbox{}\verb@    ldrb w0, [x20]@\\
\mbox{}\verb@    cmp w0, #'"'@\\
\mbox{}\verb@    b.eq .Lstr_end@\\
\mbox{}\verb@    cmp w0, #'\\'@\\
\mbox{}\verb@    b.eq .Lstr_escape@\\
\mbox{}\verb@    cmp w0, #0x0A               // newline@\\
\mbox{}\verb@    b.eq .Lstr_unterminated@\\
\mbox{}\verb@    add x20, x20, #1@\\
\mbox{}\verb@    b .Lstr_loop@\\
\mbox{}\verb@@\\
\mbox{}\verb@.Lstr_escape:@\\
\mbox{}\verb@    add x20, x20, #2           // skip backslash + escaped char@\\
\mbox{}\verb@    b .Lstr_loop@\\
\mbox{}\verb@@\\
\mbox{}\verb@.Lstr_end:@\\
\mbox{}\verb@    add x20, x20, #1           // skip closing "@\\
\mbox{}\verb@    str x20, [x19]@\\
\mbox{}\verb@@\\
\mbox{}\verb@    sub x2, x20, x22@\\
\mbox{}\verb@    adrp x1, _kcc_tok_start@{\tt @}\verb@PAGE@\\
\mbox{}\verb@    add  x1, x1, _kcc_tok_start@{\tt @}\verb@PAGEOFF@\\
\mbox{}\verb@    str  x22, [x1]@\\
\mbox{}\verb@@\\
\mbox{}\verb@    adrp x1, _kcc_tok_len@{\tt @}\verb@PAGE@\\
\mbox{}\verb@    add  x1, x1, _kcc_tok_len@{\tt @}\verb@PAGEOFF@\\
\mbox{}\verb@    str  x2, [x1]@\\
\mbox{}\verb@@\\
\mbox{}\verb@    mov x0, #TOK_STRING_LIT@\\
\mbox{}\verb@    adrp x1, _kcc_tok_type@{\tt @}\verb@PAGE@\\
\mbox{}\verb@    add  x1, x1, _kcc_tok_type@{\tt @}\verb@PAGEOFF@\\
\mbox{}\verb@    str  x0, [x1]@\\
\mbox{}\verb@@\\
\mbox{}\verb@    ldp x21, x22, [sp, #32]@\\
\mbox{}\verb@    ldp x19, x20, [sp, #16]@\\
\mbox{}\verb@    ldp x29, x30, [sp], #48@\\
\mbox{}\verb@    ret@\\
\mbox{}\verb@@\\
\mbox{}\verb@.Lstr_unterminated:@\\
\mbox{}\verb@    // Error: unterminated string — report via diagnostic engine@\\
\mbox{}\verb@    str x20, [x19]@\\
\mbox{}\verb@    adrp x0, _kcc_tok_line@{\tt @}\verb@PAGE@\\
\mbox{}\verb@    add  x0, x0, _kcc_tok_line@{\tt @}\verb@PAGEOFF@\\
\mbox{}\verb@    ldr  x0, [x0]                // line (recorded at opening quote)@\\
\mbox{}\verb@    adrp x1, _kcc_tok_col@{\tt @}\verb@PAGE@\\
\mbox{}\verb@    add  x1, x1, _kcc_tok_col@{\tt @}\verb@PAGEOFF@\\
\mbox{}\verb@    ldr  x1, [x1]                // col (recorded at opening quote)@\\
\mbox{}\verb@    adrp x2, _kcc_err_unterm_string@{\tt @}\verb@PAGE@\\
\mbox{}\verb@    add  x2, x2, _kcc_err_unterm_string@{\tt @}\verb@PAGEOFF@\\
\mbox{}\verb@    bl _kcc_diag_error@\\
\mbox{}\verb@@\\
\mbox{}\verb@    sub x2, x20, x22@\\
\mbox{}\verb@    adrp x1, _kcc_tok_start@{\tt @}\verb@PAGE@\\
\mbox{}\verb@    add  x1, x1, _kcc_tok_start@{\tt @}\verb@PAGEOFF@\\
\mbox{}\verb@    str  x22, [x1]@\\
\mbox{}\verb@@\\
\mbox{}\verb@    adrp x1, _kcc_tok_len@{\tt @}\verb@PAGE@\\
\mbox{}\verb@    add  x1, x1, _kcc_tok_len@{\tt @}\verb@PAGEOFF@\\
\mbox{}\verb@    str  x2, [x1]@\\
\mbox{}\verb@@\\
\mbox{}\verb@    mov x0, #TOK_STRING_LIT@\\
\mbox{}\verb@    adrp x1, _kcc_tok_type@{\tt @}\verb@PAGE@\\
\mbox{}\verb@    add  x1, x1, _kcc_tok_type@{\tt @}\verb@PAGEOFF@\\
\mbox{}\verb@    str  x0, [x1]@\\
\mbox{}\verb@@\\
\mbox{}\verb@    ldp x21, x22, [sp, #32]@\\
\mbox{}\verb@    ldp x19, x20, [sp, #16]@\\
\mbox{}\verb@    ldp x29, x30, [sp], #48@\\
\mbox{}\verb@    ret@\\
\mbox{}\verb@@{\NWsep}
\end{list}
\vspace{-1.5ex}
\footnotesize
\begin{list}{}{\setlength{\itemsep}{-\parsep}\setlength{\itemindent}{-\leftmargin}}
\item \NWtxtMacroRefIn\ \NWlink{weft?}{?}.

\item{}
\end{list}
\end{minipage}\vspace{4ex}
\end{flushleft}
\section{Character Literal Scanner}

Character literals are similar to strings but delimited by single
quotes (\stdref{6.4.4.4}).

\begin{flushleft} \small
\begin{minipage}{\linewidth}\label{scrap63}\raggedright\small
\NWtarget{weft?}{} $\langle\,${\itshape Lexer: scan char}\nobreak\ {\footnotesize {?}}$\,\rangle\equiv$
\vspace{-1ex}
\begin{list}{}{} \item
\mbox{}\verb@@\\
\mbox{}\verb@// -- kcc_lex_scan_char ----------------------------------------------@\\
\mbox{}\verb@// Internal: scans a character constant (opening '\'' at cursor).@\\
\mbox{}\verb@.p2align 2@\\
\mbox{}\verb@_kcc_lex_scan_char:@\\
\mbox{}\verb@    stp x29, x30, [sp, #-48]!@\\
\mbox{}\verb@    mov x29, sp@\\
\mbox{}\verb@    stp x19, x20, [sp, #16]@\\
\mbox{}\verb@    stp x21, x22, [sp, #32]@\\
\mbox{}\verb@@\\
\mbox{}\verb@    adrp x19, _kcc_lex_cursor@{\tt @}\verb@PAGE@\\
\mbox{}\verb@    add  x19, x19, _kcc_lex_cursor@{\tt @}\verb@PAGEOFF@\\
\mbox{}\verb@    ldr  x20, [x19]@\\
\mbox{}\verb@@\\
\mbox{}\verb@    adrp x1, _kcc_lex_end@{\tt @}\verb@PAGE@\\
\mbox{}\verb@    add  x1, x1, _kcc_lex_end@{\tt @}\verb@PAGEOFF@\\
\mbox{}\verb@    ldr  x21, [x1]@\\
\mbox{}\verb@@\\
\mbox{}\verb@    mov x22, x20                // token start@\\
\mbox{}\verb@    add x20, x20, #1           // skip opening '@\\
\mbox{}\verb@@\\
\mbox{}\verb@.Lchr_loop:@\\
\mbox{}\verb@    cmp x20, x21@\\
\mbox{}\verb@    b.hs .Lchr_unterminated@\\
\mbox{}\verb@    ldrb w0, [x20]@\\
\mbox{}\verb@    cmp w0, #'\''@\\
\mbox{}\verb@    b.eq .Lchr_end@\\
\mbox{}\verb@    cmp w0, #'\\'@\\
\mbox{}\verb@    b.eq .Lchr_escape@\\
\mbox{}\verb@    cmp w0, #0x0A@\\
\mbox{}\verb@    b.eq .Lchr_unterminated@\\
\mbox{}\verb@    add x20, x20, #1@\\
\mbox{}\verb@    b .Lchr_loop@\\
\mbox{}\verb@@\\
\mbox{}\verb@.Lchr_escape:@\\
\mbox{}\verb@    add x20, x20, #2@\\
\mbox{}\verb@    b .Lchr_loop@\\
\mbox{}\verb@@\\
\mbox{}\verb@.Lchr_end:@\\
\mbox{}\verb@    add x20, x20, #1           // skip closing '@\\
\mbox{}\verb@    str x20, [x19]@\\
\mbox{}\verb@@\\
\mbox{}\verb@    sub x2, x20, x22@\\
\mbox{}\verb@    adrp x1, _kcc_tok_start@{\tt @}\verb@PAGE@\\
\mbox{}\verb@    add  x1, x1, _kcc_tok_start@{\tt @}\verb@PAGEOFF@\\
\mbox{}\verb@    str  x22, [x1]@\\
\mbox{}\verb@@\\
\mbox{}\verb@    adrp x1, _kcc_tok_len@{\tt @}\verb@PAGE@\\
\mbox{}\verb@    add  x1, x1, _kcc_tok_len@{\tt @}\verb@PAGEOFF@\\
\mbox{}\verb@    str  x2, [x1]@\\
\mbox{}\verb@@\\
\mbox{}\verb@    mov x0, #TOK_CHAR_LIT@\\
\mbox{}\verb@    adrp x1, _kcc_tok_type@{\tt @}\verb@PAGE@\\
\mbox{}\verb@    add  x1, x1, _kcc_tok_type@{\tt @}\verb@PAGEOFF@\\
\mbox{}\verb@    str  x0, [x1]@\\
\mbox{}\verb@@\\
\mbox{}\verb@    ldp x21, x22, [sp, #32]@\\
\mbox{}\verb@    ldp x19, x20, [sp, #16]@\\
\mbox{}\verb@    ldp x29, x30, [sp], #48@\\
\mbox{}\verb@    ret@\\
\mbox{}\verb@@\\
\mbox{}\verb@.Lchr_unterminated:@\\
\mbox{}\verb@    str x20, [x19]@\\
\mbox{}\verb@    adrp x0, _kcc_tok_line@{\tt @}\verb@PAGE@\\
\mbox{}\verb@    add  x0, x0, _kcc_tok_line@{\tt @}\verb@PAGEOFF@\\
\mbox{}\verb@    ldr  x0, [x0]                // line (recorded at opening quote)@\\
\mbox{}\verb@    adrp x1, _kcc_tok_col@{\tt @}\verb@PAGE@\\
\mbox{}\verb@    add  x1, x1, _kcc_tok_col@{\tt @}\verb@PAGEOFF@\\
\mbox{}\verb@    ldr  x1, [x1]                // col (recorded at opening quote)@\\
\mbox{}\verb@    adrp x2, _kcc_err_unterm_char@{\tt @}\verb@PAGE@\\
\mbox{}\verb@    add  x2, x2, _kcc_err_unterm_char@{\tt @}\verb@PAGEOFF@\\
\mbox{}\verb@    bl _kcc_diag_error@\\
\mbox{}\verb@@\\
\mbox{}\verb@    sub x2, x20, x22@\\
\mbox{}\verb@    adrp x1, _kcc_tok_start@{\tt @}\verb@PAGE@\\
\mbox{}\verb@    add  x1, x1, _kcc_tok_start@{\tt @}\verb@PAGEOFF@\\
\mbox{}\verb@    str  x22, [x1]@\\
\mbox{}\verb@@\\
\mbox{}\verb@    adrp x1, _kcc_tok_len@{\tt @}\verb@PAGE@\\
\mbox{}\verb@    add  x1, x1, _kcc_tok_len@{\tt @}\verb@PAGEOFF@\\
\mbox{}\verb@    str  x2, [x1]@\\
\mbox{}\verb@@\\
\mbox{}\verb@    mov x0, #TOK_CHAR_LIT@\\
\mbox{}\verb@    adrp x1, _kcc_tok_type@{\tt @}\verb@PAGE@\\
\mbox{}\verb@    add  x1, x1, _kcc_tok_type@{\tt @}\verb@PAGEOFF@\\
\mbox{}\verb@    str  x0, [x1]@\\
\mbox{}\verb@@\\
\mbox{}\verb@    ldp x21, x22, [sp, #32]@\\
\mbox{}\verb@    ldp x19, x20, [sp, #16]@\\
\mbox{}\verb@    ldp x29, x30, [sp], #48@\\
\mbox{}\verb@    ret@\\
\mbox{}\verb@@{\NWsep}
\end{list}
\vspace{-1.5ex}
\footnotesize
\begin{list}{}{\setlength{\itemsep}{-\parsep}\setlength{\itemindent}{-\leftmargin}}
\item \NWtxtMacroRefIn\ \NWlink{weft?}{?}.

\item{}
\end{list}
\end{minipage}\vspace{4ex}
\end{flushleft}
\section{Punctuator Scanner}

The punctuator scanner implements longest-match for all C17 punctuators
(\stdref{6.4.6}).  Multi-character punctuators like \texttt{>>=} or
\texttt{...} are handled by peeking ahead.

\begin{flushleft} \small
\begin{minipage}{\linewidth}\label{scrap64}\raggedright\small
\NWtarget{weft?}{} $\langle\,${\itshape Lexer: scan punctuator}\nobreak\ {\footnotesize {?}}$\,\rangle\equiv$
\vspace{-1ex}
\begin{list}{}{} \item
\mbox{}\verb@@\\
\mbox{}\verb@// -- kcc_lex_scan_punct ---------------------------------------------@\\
\mbox{}\verb@// Internal: scans a punctuator.@\\
\mbox{}\verb@.p2align 2@\\
\mbox{}\verb@_kcc_lex_scan_punct:@\\
\mbox{}\verb@    stp x29, x30, [sp, #-48]!@\\
\mbox{}\verb@    mov x29, sp@\\
\mbox{}\verb@    stp x19, x20, [sp, #16]@\\
\mbox{}\verb@    stp x21, x22, [sp, #32]@\\
\mbox{}\verb@@\\
\mbox{}\verb@    adrp x19, _kcc_lex_cursor@{\tt @}\verb@PAGE@\\
\mbox{}\verb@    add  x19, x19, _kcc_lex_cursor@{\tt @}\verb@PAGEOFF@\\
\mbox{}\verb@    ldr  x20, [x19]@\\
\mbox{}\verb@@\\
\mbox{}\verb@    adrp x1, _kcc_lex_end@{\tt @}\verb@PAGE@\\
\mbox{}\verb@    add  x1, x1, _kcc_lex_end@{\tt @}\verb@PAGEOFF@\\
\mbox{}\verb@    ldr  x21, [x1]@\\
\mbox{}\verb@@\\
\mbox{}\verb@    mov x22, x20                // token start@\\
\mbox{}\verb@    ldrb w0, [x20]@\\
\mbox{}\verb@@\\
\mbox{}\verb@    // Single-character dispatch@\\
\mbox{}\verb@    cmp w0, #'('@\\
\mbox{}\verb@    b.eq .Lpunct_1_lparen@\\
\mbox{}\verb@    cmp w0, #')'@\\
\mbox{}\verb@    b.eq .Lpunct_1_rparen@\\
\mbox{}\verb@    cmp w0, #'{'@\\
\mbox{}\verb@    b.eq .Lpunct_1_lbrace@\\
\mbox{}\verb@    cmp w0, #'}'@\\
\mbox{}\verb@    b.eq .Lpunct_1_rbrace@\\
\mbox{}\verb@    cmp w0, #'['@\\
\mbox{}\verb@    b.eq .Lpunct_1_lbracket@\\
\mbox{}\verb@    cmp w0, #']'@\\
\mbox{}\verb@    b.eq .Lpunct_1_rbracket@\\
\mbox{}\verb@    cmp w0, #';'@\\
\mbox{}\verb@    b.eq .Lpunct_1_semi@\\
\mbox{}\verb@    cmp w0, #','@\\
\mbox{}\verb@    b.eq .Lpunct_1_comma@\\
\mbox{}\verb@    cmp w0, #'~'@\\
\mbox{}\verb@    b.eq .Lpunct_1_tilde@\\
\mbox{}\verb@    cmp w0, #'?'@\\
\mbox{}\verb@    b.eq .Lpunct_1_question@\\
\mbox{}\verb@    cmp w0, #':'@\\
\mbox{}\verb@    b.eq .Lpunct_1_colon@\\
\mbox{}\verb@@\\
\mbox{}\verb@    // Two-or-more character punctuators@\\
\mbox{}\verb@    cmp w0, #'+'@\\
\mbox{}\verb@    b.eq .Lpunct_plus@\\
\mbox{}\verb@    cmp w0, #'-'@\\
\mbox{}\verb@    b.eq .Lpunct_minus@\\
\mbox{}\verb@    cmp w0, #'*'@\\
\mbox{}\verb@    b.eq .Lpunct_star@\\
\mbox{}\verb@    cmp w0, #'/'@\\
\mbox{}\verb@    b.eq .Lpunct_slash@\\
\mbox{}\verb@    cmp w0, #'%'@\\
\mbox{}\verb@    b.eq .Lpunct_percent@\\
\mbox{}\verb@    cmp w0, #'<'@\\
\mbox{}\verb@    b.eq .Lpunct_lt@\\
\mbox{}\verb@    cmp w0, #'>'@\\
\mbox{}\verb@    b.eq .Lpunct_gt@\\
\mbox{}\verb@    cmp w0, #'='@\\
\mbox{}\verb@    b.eq .Lpunct_eq@\\
\mbox{}\verb@    cmp w0, #'!'@\\
\mbox{}\verb@    b.eq .Lpunct_bang@\\
\mbox{}\verb@    cmp w0, #'&'@\\
\mbox{}\verb@    b.eq .Lpunct_amp@\\
\mbox{}\verb@    cmp w0, #'|'@\\
\mbox{}\verb@    b.eq .Lpunct_pipe@\\
\mbox{}\verb@    cmp w0, #'^'@\\
\mbox{}\verb@    b.eq .Lpunct_caret@\\
\mbox{}\verb@    cmp w0, #'#'@\\
\mbox{}\verb@    b.eq .Lpunct_hash@\\
\mbox{}\verb@    cmp w0, #'.'@\\
\mbox{}\verb@    b.eq .Lpunct_dot@\\
\mbox{}\verb@@\\
\mbox{}\verb@    // Unrecognized character@\\
\mbox{}\verb@    add x20, x20, #1@\\
\mbox{}\verb@    str x20, [x19]@\\
\mbox{}\verb@    mov x0, #TOK_UNKNOWN@\\
\mbox{}\verb@    b .Lpunct_store@\\
\mbox{}\verb@@\\
\mbox{}\verb@// -- Simple single-character punctuators ----------------------------@\\
\mbox{}\verb@.Lpunct_1_lparen:@\\
\mbox{}\verb@    mov x0, #TOK_LPAREN@\\
\mbox{}\verb@    b .Lpunct_single@\\
\mbox{}\verb@.Lpunct_1_rparen:@\\
\mbox{}\verb@    mov x0, #TOK_RPAREN@\\
\mbox{}\verb@    b .Lpunct_single@\\
\mbox{}\verb@.Lpunct_1_lbrace:@\\
\mbox{}\verb@    mov x0, #TOK_LBRACE@\\
\mbox{}\verb@    b .Lpunct_single@\\
\mbox{}\verb@.Lpunct_1_rbrace:@\\
\mbox{}\verb@    mov x0, #TOK_RBRACE@\\
\mbox{}\verb@    b .Lpunct_single@\\
\mbox{}\verb@.Lpunct_1_lbracket:@\\
\mbox{}\verb@    mov x0, #TOK_LBRACKET@\\
\mbox{}\verb@    b .Lpunct_single@\\
\mbox{}\verb@.Lpunct_1_rbracket:@\\
\mbox{}\verb@    mov x0, #TOK_RBRACKET@\\
\mbox{}\verb@    b .Lpunct_single@\\
\mbox{}\verb@.Lpunct_1_semi:@\\
\mbox{}\verb@    mov x0, #TOK_SEMI@\\
\mbox{}\verb@    b .Lpunct_single@\\
\mbox{}\verb@.Lpunct_1_comma:@\\
\mbox{}\verb@    mov x0, #TOK_COMMA@\\
\mbox{}\verb@    b .Lpunct_single@\\
\mbox{}\verb@.Lpunct_1_tilde:@\\
\mbox{}\verb@    mov x0, #TOK_TILDE@\\
\mbox{}\verb@    b .Lpunct_single@\\
\mbox{}\verb@.Lpunct_1_question:@\\
\mbox{}\verb@    mov x0, #TOK_QUESTION@\\
\mbox{}\verb@    b .Lpunct_single@\\
\mbox{}\verb@.Lpunct_1_colon:@\\
\mbox{}\verb@    mov x0, #TOK_COLON@\\
\mbox{}\verb@    b .Lpunct_single@\\
\mbox{}\verb@@\\
\mbox{}\verb@.Lpunct_single:@\\
\mbox{}\verb@    add x20, x20, #1@\\
\mbox{}\verb@    str x20, [x19]@\\
\mbox{}\verb@    b .Lpunct_store@\\
\mbox{}\verb@@\\
\mbox{}\verb@// -- Multi-character punctuators ------------------------------------@\\
\mbox{}\verb@.Lpunct_plus:                     // + ++ +=@\\
\mbox{}\verb@    add x1, x20, #1@\\
\mbox{}\verb@    cmp x1, x21@\\
\mbox{}\verb@    b.hs .Lpunct_plus_single@\\
\mbox{}\verb@    ldrb w2, [x1]@\\
\mbox{}\verb@    cmp w2, #'+'@\\
\mbox{}\verb@    b.eq .Lpunct_inc@\\
\mbox{}\verb@    cmp w2, #'='@\\
\mbox{}\verb@    b.eq .Lpunct_add_assign@\\
\mbox{}\verb@.Lpunct_plus_single:@\\
\mbox{}\verb@    mov x0, #TOK_PLUS@\\
\mbox{}\verb@    b .Lpunct_single@\\
\mbox{}\verb@.Lpunct_inc:@\\
\mbox{}\verb@    mov x0, #TOK_INC@\\
\mbox{}\verb@    b .Lpunct_double@\\
\mbox{}\verb@.Lpunct_add_assign:@\\
\mbox{}\verb@    mov x0, #TOK_ADD_ASSIGN@\\
\mbox{}\verb@    b .Lpunct_double@\\
\mbox{}\verb@@\\
\mbox{}\verb@.Lpunct_minus:                    // - -- -= ->@\\
\mbox{}\verb@    add x1, x20, #1@\\
\mbox{}\verb@    cmp x1, x21@\\
\mbox{}\verb@    b.hs .Lpunct_minus_single@\\
\mbox{}\verb@    ldrb w2, [x1]@\\
\mbox{}\verb@    cmp w2, #'-'@\\
\mbox{}\verb@    b.eq .Lpunct_dec@\\
\mbox{}\verb@    cmp w2, #'='@\\
\mbox{}\verb@    b.eq .Lpunct_sub_assign@\\
\mbox{}\verb@    cmp w2, #'>'@\\
\mbox{}\verb@    b.eq .Lpunct_arrow@\\
\mbox{}\verb@.Lpunct_minus_single:@\\
\mbox{}\verb@    mov x0, #TOK_MINUS@\\
\mbox{}\verb@    b .Lpunct_single@\\
\mbox{}\verb@.Lpunct_dec:@\\
\mbox{}\verb@    mov x0, #TOK_DEC@\\
\mbox{}\verb@    b .Lpunct_double@\\
\mbox{}\verb@.Lpunct_sub_assign:@\\
\mbox{}\verb@    mov x0, #TOK_SUB_ASSIGN@\\
\mbox{}\verb@    b .Lpunct_double@\\
\mbox{}\verb@.Lpunct_arrow:@\\
\mbox{}\verb@    mov x0, #TOK_ARROW@\\
\mbox{}\verb@    b .Lpunct_double@\\
\mbox{}\verb@@\\
\mbox{}\verb@.Lpunct_star:                     // * *=@\\
\mbox{}\verb@    add x1, x20, #1@\\
\mbox{}\verb@    cmp x1, x21@\\
\mbox{}\verb@    b.hs .Lpunct_star_single@\\
\mbox{}\verb@    ldrb w2, [x1]@\\
\mbox{}\verb@    cmp w2, #'='@\\
\mbox{}\verb@    b.eq .Lpunct_mul_assign@\\
\mbox{}\verb@.Lpunct_star_single:@\\
\mbox{}\verb@    mov x0, #TOK_STAR@\\
\mbox{}\verb@    b .Lpunct_single@\\
\mbox{}\verb@.Lpunct_mul_assign:@\\
\mbox{}\verb@    mov x0, #TOK_MUL_ASSIGN@\\
\mbox{}\verb@    b .Lpunct_double@\\
\mbox{}\verb@@\\
\mbox{}\verb@.Lpunct_slash:                    // /  /= (comments already handled)@\\
\mbox{}\verb@    add x1, x20, #1@\\
\mbox{}\verb@    cmp x1, x21@\\
\mbox{}\verb@    b.hs .Lpunct_slash_single@\\
\mbox{}\verb@    ldrb w2, [x1]@\\
\mbox{}\verb@    cmp w2, #'='@\\
\mbox{}\verb@    b.eq .Lpunct_div_assign@\\
\mbox{}\verb@.Lpunct_slash_single:@\\
\mbox{}\verb@    mov x0, #TOK_SLASH@\\
\mbox{}\verb@    b .Lpunct_single@\\
\mbox{}\verb@.Lpunct_div_assign:@\\
\mbox{}\verb@    mov x0, #TOK_DIV_ASSIGN@\\
\mbox{}\verb@    b .Lpunct_double@\\
\mbox{}\verb@@\\
\mbox{}\verb@.Lpunct_percent:                  // % %=@\\
\mbox{}\verb@    add x1, x20, #1@\\
\mbox{}\verb@    cmp x1, x21@\\
\mbox{}\verb@    b.hs .Lpunct_pct_single@\\
\mbox{}\verb@    ldrb w2, [x1]@\\
\mbox{}\verb@    cmp w2, #'='@\\
\mbox{}\verb@    b.eq .Lpunct_mod_assign@\\
\mbox{}\verb@.Lpunct_pct_single:@\\
\mbox{}\verb@    mov x0, #TOK_PERCENT@\\
\mbox{}\verb@    b .Lpunct_single@\\
\mbox{}\verb@.Lpunct_mod_assign:@\\
\mbox{}\verb@    mov x0, #TOK_MOD_ASSIGN@\\
\mbox{}\verb@    b .Lpunct_double@\\
\mbox{}\verb@@\\
\mbox{}\verb@.Lpunct_lt:                       // < << <= <<=@\\
\mbox{}\verb@    add x1, x20, #1@\\
\mbox{}\verb@    cmp x1, x21@\\
\mbox{}\verb@    b.hs .Lpunct_lt_single@\\
\mbox{}\verb@    ldrb w2, [x1]@\\
\mbox{}\verb@    cmp w2, #'<'@\\
\mbox{}\verb@    b.eq .Lpunct_shl_check@\\
\mbox{}\verb@    cmp w2, #'='@\\
\mbox{}\verb@    b.eq .Lpunct_le@\\
\mbox{}\verb@.Lpunct_lt_single:@\\
\mbox{}\verb@    mov x0, #TOK_LT@\\
\mbox{}\verb@    b .Lpunct_single@\\
\mbox{}\verb@.Lpunct_le:@\\
\mbox{}\verb@    mov x0, #TOK_LE@\\
\mbox{}\verb@    b .Lpunct_double@\\
\mbox{}\verb@.Lpunct_shl_check:               // << or <<=@\\
\mbox{}\verb@    add x2, x20, #2@\\
\mbox{}\verb@    cmp x2, x21@\\
\mbox{}\verb@    b.hs .Lpunct_shl@\\
\mbox{}\verb@    ldrb w3, [x2]@\\
\mbox{}\verb@    cmp w3, #'='@\\
\mbox{}\verb@    b.eq .Lpunct_shl_assign@\\
\mbox{}\verb@.Lpunct_shl:@\\
\mbox{}\verb@    mov x0, #TOK_SHL@\\
\mbox{}\verb@    b .Lpunct_double@\\
\mbox{}\verb@.Lpunct_shl_assign:@\\
\mbox{}\verb@    mov x0, #TOK_SHL_ASSIGN@\\
\mbox{}\verb@    b .Lpunct_triple@\\
\mbox{}\verb@@\\
\mbox{}\verb@.Lpunct_gt:                       // > >> >= >>=@\\
\mbox{}\verb@    add x1, x20, #1@\\
\mbox{}\verb@    cmp x1, x21@\\
\mbox{}\verb@    b.hs .Lpunct_gt_single@\\
\mbox{}\verb@    ldrb w2, [x1]@\\
\mbox{}\verb@    cmp w2, #'>'@\\
\mbox{}\verb@    b.eq .Lpunct_shr_check@\\
\mbox{}\verb@    cmp w2, #'='@\\
\mbox{}\verb@    b.eq .Lpunct_ge@\\
\mbox{}\verb@.Lpunct_gt_single:@\\
\mbox{}\verb@    mov x0, #TOK_GT@\\
\mbox{}\verb@    b .Lpunct_single@\\
\mbox{}\verb@.Lpunct_ge:@\\
\mbox{}\verb@    mov x0, #TOK_GE@\\
\mbox{}\verb@    b .Lpunct_double@\\
\mbox{}\verb@.Lpunct_shr_check:@\\
\mbox{}\verb@    add x2, x20, #2@\\
\mbox{}\verb@    cmp x2, x21@\\
\mbox{}\verb@    b.hs .Lpunct_shr@\\
\mbox{}\verb@    ldrb w3, [x2]@\\
\mbox{}\verb@    cmp w3, #'='@\\
\mbox{}\verb@    b.eq .Lpunct_shr_assign@\\
\mbox{}\verb@.Lpunct_shr:@\\
\mbox{}\verb@    mov x0, #TOK_SHR@\\
\mbox{}\verb@    b .Lpunct_double@\\
\mbox{}\verb@.Lpunct_shr_assign:@\\
\mbox{}\verb@    mov x0, #TOK_SHR_ASSIGN@\\
\mbox{}\verb@    b .Lpunct_triple@\\
\mbox{}\verb@@\\
\mbox{}\verb@.Lpunct_eq:                       // = ==@\\
\mbox{}\verb@    add x1, x20, #1@\\
\mbox{}\verb@    cmp x1, x21@\\
\mbox{}\verb@    b.hs .Lpunct_assign@\\
\mbox{}\verb@    ldrb w2, [x1]@\\
\mbox{}\verb@    cmp w2, #'='@\\
\mbox{}\verb@    b.eq .Lpunct_eq_eq@\\
\mbox{}\verb@.Lpunct_assign:@\\
\mbox{}\verb@    mov x0, #TOK_ASSIGN@\\
\mbox{}\verb@    b .Lpunct_single@\\
\mbox{}\verb@.Lpunct_eq_eq:@\\
\mbox{}\verb@    mov x0, #TOK_EQ@\\
\mbox{}\verb@    b .Lpunct_double@\\
\mbox{}\verb@@\\
\mbox{}\verb@.Lpunct_bang:                     // ! !=@\\
\mbox{}\verb@    add x1, x20, #1@\\
\mbox{}\verb@    cmp x1, x21@\\
\mbox{}\verb@    b.hs .Lpunct_bang_single@\\
\mbox{}\verb@    ldrb w2, [x1]@\\
\mbox{}\verb@    cmp w2, #'='@\\
\mbox{}\verb@    b.eq .Lpunct_ne@\\
\mbox{}\verb@.Lpunct_bang_single:@\\
\mbox{}\verb@    mov x0, #TOK_BANG@\\
\mbox{}\verb@    b .Lpunct_single@\\
\mbox{}\verb@.Lpunct_ne:@\\
\mbox{}\verb@    mov x0, #TOK_NE@\\
\mbox{}\verb@    b .Lpunct_double@\\
\mbox{}\verb@@\\
\mbox{}\verb@.Lpunct_amp:                      // & && &=@\\
\mbox{}\verb@    add x1, x20, #1@\\
\mbox{}\verb@    cmp x1, x21@\\
\mbox{}\verb@    b.hs .Lpunct_amp_single@\\
\mbox{}\verb@    ldrb w2, [x1]@\\
\mbox{}\verb@    cmp w2, #'&'@\\
\mbox{}\verb@    b.eq .Lpunct_and@\\
\mbox{}\verb@    cmp w2, #'='@\\
\mbox{}\verb@    b.eq .Lpunct_and_assign@\\
\mbox{}\verb@.Lpunct_amp_single:@\\
\mbox{}\verb@    mov x0, #TOK_AMP@\\
\mbox{}\verb@    b .Lpunct_single@\\
\mbox{}\verb@.Lpunct_and:@\\
\mbox{}\verb@    mov x0, #TOK_AND@\\
\mbox{}\verb@    b .Lpunct_double@\\
\mbox{}\verb@.Lpunct_and_assign:@\\
\mbox{}\verb@    mov x0, #TOK_AND_ASSIGN@\\
\mbox{}\verb@    b .Lpunct_double@\\
\mbox{}\verb@@\\
\mbox{}\verb@.Lpunct_pipe:                     // | || |=@\\
\mbox{}\verb@    add x1, x20, #1@\\
\mbox{}\verb@    cmp x1, x21@\\
\mbox{}\verb@    b.hs .Lpunct_pipe_single@\\
\mbox{}\verb@    ldrb w2, [x1]@\\
\mbox{}\verb@    cmp w2, #'|'@\\
\mbox{}\verb@    b.eq .Lpunct_or@\\
\mbox{}\verb@    cmp w2, #'='@\\
\mbox{}\verb@    b.eq .Lpunct_or_assign@\\
\mbox{}\verb@.Lpunct_pipe_single:@\\
\mbox{}\verb@    mov x0, #TOK_PIPE@\\
\mbox{}\verb@    b .Lpunct_single@\\
\mbox{}\verb@.Lpunct_or:@\\
\mbox{}\verb@    mov x0, #TOK_OR@\\
\mbox{}\verb@    b .Lpunct_double@\\
\mbox{}\verb@.Lpunct_or_assign:@\\
\mbox{}\verb@    mov x0, #TOK_OR_ASSIGN@\\
\mbox{}\verb@    b .Lpunct_double@\\
\mbox{}\verb@@\\
\mbox{}\verb@.Lpunct_caret:                    // ^ ^=@\\
\mbox{}\verb@    add x1, x20, #1@\\
\mbox{}\verb@    cmp x1, x21@\\
\mbox{}\verb@    b.hs .Lpunct_caret_single@\\
\mbox{}\verb@    ldrb w2, [x1]@\\
\mbox{}\verb@    cmp w2, #'='@\\
\mbox{}\verb@    b.eq .Lpunct_xor_assign@\\
\mbox{}\verb@.Lpunct_caret_single:@\\
\mbox{}\verb@    mov x0, #TOK_CARET@\\
\mbox{}\verb@    b .Lpunct_single@\\
\mbox{}\verb@.Lpunct_xor_assign:@\\
\mbox{}\verb@    mov x0, #TOK_XOR_ASSIGN@\\
\mbox{}\verb@    b .Lpunct_double@\\
\mbox{}\verb@@\\
\mbox{}\verb@.Lpunct_hash:                     // # ##@\\
\mbox{}\verb@    add x1, x20, #1@\\
\mbox{}\verb@    cmp x1, x21@\\
\mbox{}\verb@    b.hs .Lpunct_hash_single@\\
\mbox{}\verb@    ldrb w2, [x1]@\\
\mbox{}\verb@    cmp w2, #'#'@\\
\mbox{}\verb@    b.eq .Lpunct_dhash@\\
\mbox{}\verb@.Lpunct_hash_single:@\\
\mbox{}\verb@    mov x0, #TOK_HASH@\\
\mbox{}\verb@    b .Lpunct_single@\\
\mbox{}\verb@.Lpunct_dhash:@\\
\mbox{}\verb@    mov x0, #TOK_DHASH@\\
\mbox{}\verb@    b .Lpunct_double@\\
\mbox{}\verb@@\\
\mbox{}\verb@.Lpunct_dot:                      // . ...@\\
\mbox{}\verb@    add x1, x20, #1@\\
\mbox{}\verb@    cmp x1, x21@\\
\mbox{}\verb@    b.hs .Lpunct_dot_single@\\
\mbox{}\verb@    ldrb w2, [x1]@\\
\mbox{}\verb@    cmp w2, #'.'@\\
\mbox{}\verb@    b.ne .Lpunct_dot_single@\\
\mbox{}\verb@    add x2, x20, #2@\\
\mbox{}\verb@    cmp x2, x21@\\
\mbox{}\verb@    b.hs .Lpunct_dot_single@\\
\mbox{}\verb@    ldrb w3, [x2]@\\
\mbox{}\verb@    cmp w3, #'.'@\\
\mbox{}\verb@    b.ne .Lpunct_dot_single@\\
\mbox{}\verb@    mov x0, #TOK_ELLIPSIS@\\
\mbox{}\verb@    b .Lpunct_triple@\\
\mbox{}\verb@.Lpunct_dot_single:@\\
\mbox{}\verb@    mov x0, #TOK_DOT@\\
\mbox{}\verb@    b .Lpunct_single@\\
\mbox{}\verb@@\\
\mbox{}\verb@// -- Length helpers --------------------------------------------------@\\
\mbox{}\verb@.Lpunct_double:@\\
\mbox{}\verb@    add x20, x20, #2@\\
\mbox{}\verb@    str x20, [x19]@\\
\mbox{}\verb@    b .Lpunct_store@\\
\mbox{}\verb@@\\
\mbox{}\verb@.Lpunct_triple:@\\
\mbox{}\verb@    add x20, x20, #3@\\
\mbox{}\verb@    str x20, [x19]@\\
\mbox{}\verb@    b .Lpunct_store@\\
\mbox{}\verb@@\\
\mbox{}\verb@.Lpunct_store:@\\
\mbox{}\verb@    // Store token type, start, length@\\
\mbox{}\verb@    adrp x1, _kcc_tok_type@{\tt @}\verb@PAGE@\\
\mbox{}\verb@    add  x1, x1, _kcc_tok_type@{\tt @}\verb@PAGEOFF@\\
\mbox{}\verb@    str  x0, [x1]@\\
\mbox{}\verb@@\\
\mbox{}\verb@    // Compute length from start@\\
\mbox{}\verb@    ldr  x20, [x19]@\\
\mbox{}\verb@    sub  x2, x20, x22@\\
\mbox{}\verb@@\\
\mbox{}\verb@    adrp x1, _kcc_tok_start@{\tt @}\verb@PAGE@\\
\mbox{}\verb@    add  x1, x1, _kcc_tok_start@{\tt @}\verb@PAGEOFF@\\
\mbox{}\verb@    str  x22, [x1]@\\
\mbox{}\verb@@\\
\mbox{}\verb@    adrp x1, _kcc_tok_len@{\tt @}\verb@PAGE@\\
\mbox{}\verb@    add  x1, x1, _kcc_tok_len@{\tt @}\verb@PAGEOFF@\\
\mbox{}\verb@    str  x2, [x1]@\\
\mbox{}\verb@@\\
\mbox{}\verb@    ldp x21, x22, [sp, #32]@\\
\mbox{}\verb@    ldp x19, x20, [sp, #16]@\\
\mbox{}\verb@    ldp x29, x30, [sp], #48@\\
\mbox{}\verb@    ret@\\
\mbox{}\verb@@{\NWsep}
\end{list}
\vspace{-1.5ex}
\footnotesize
\begin{list}{}{\setlength{\itemsep}{-\parsep}\setlength{\itemindent}{-\leftmargin}}
\item \NWtxtMacroRefIn\ \NWlink{weft?}{?}.

\item{}
\end{list}
\end{minipage}\vspace{4ex}
\end{flushleft}
\section{Token Dump}

The \texttt{--dump-tokens} mode prints each token on a separate line
in the format:
\begin{verbatim}
<line>:<col> <type> '<text>'
\end{verbatim}

\texttt{kcc\_lex\_dump\_token} prints one token;
\texttt{kcc\_lex\_dump\_all} loops through the entire source.

\begin{flushleft} \small
\begin{minipage}{\linewidth}\label{scrap65}\raggedright\small
\NWtarget{weft?}{} $\langle\,${\itshape Lexer: dump token}\nobreak\ {\footnotesize {?}}$\,\rangle\equiv$
\vspace{-1ex}
\begin{list}{}{} \item
\mbox{}\verb@@\\
\mbox{}\verb@// -- kcc_lex_dump_token ---------------------------------------------@\\
\mbox{}\verb@// void kcc_lex_dump_token(void)@\\
\mbox{}\verb@// Prints the current token in dump format.@\\
\mbox{}\verb@.globl _kcc_lex_dump_token@\\
\mbox{}\verb@.p2align 2@\\
\mbox{}\verb@_kcc_lex_dump_token:@\\
\mbox{}\verb@    stp x29, x30, [sp, #-16]!@\\
\mbox{}\verb@    mov x29, sp@\\
\mbox{}\verb@@\\
\mbox{}\verb@    // Print line number@\\
\mbox{}\verb@    adrp x0, _kcc_tok_line@{\tt @}\verb@PAGE@\\
\mbox{}\verb@    add  x0, x0, _kcc_tok_line@{\tt @}\verb@PAGEOFF@\\
\mbox{}\verb@    ldr  x0, [x0]@\\
\mbox{}\verb@    bl _kcc_io_putu@\\
\mbox{}\verb@@\\
\mbox{}\verb@    // Print ':'@\\
\mbox{}\verb@    mov x0, #':'@\\
\mbox{}\verb@    bl _kcc_io_putc@\\
\mbox{}\verb@@\\
\mbox{}\verb@    // Print column@\\
\mbox{}\verb@    adrp x0, _kcc_tok_col@{\tt @}\verb@PAGE@\\
\mbox{}\verb@    add  x0, x0, _kcc_tok_col@{\tt @}\verb@PAGEOFF@\\
\mbox{}\verb@    ldr  x0, [x0]@\\
\mbox{}\verb@    bl _kcc_io_putu@\\
\mbox{}\verb@@\\
\mbox{}\verb@    // Print ' '@\\
\mbox{}\verb@    mov x0, #' '@\\
\mbox{}\verb@    bl _kcc_io_putc@\\
\mbox{}\verb@@\\
\mbox{}\verb@    // Print token type name@\\
\mbox{}\verb@    adrp x0, _kcc_tok_type@{\tt @}\verb@PAGE@\\
\mbox{}\verb@    add  x0, x0, _kcc_tok_type@{\tt @}\verb@PAGEOFF@\\
\mbox{}\verb@    ldr  x0, [x0]@\\
\mbox{}\verb@    bl _kcc_lex_type_name        // returns pointer to name string@\\
\mbox{}\verb@    bl _kcc_io_puts@\\
\mbox{}\verb@@\\
\mbox{}\verb@    // Print " '"@\\
\mbox{}\verb@    mov x0, #' '@\\
\mbox{}\verb@    bl _kcc_io_putc@\\
\mbox{}\verb@    mov x0, #'\''@\\
\mbox{}\verb@    bl _kcc_io_putc@\\
\mbox{}\verb@@\\
\mbox{}\verb@    // Print token text@\\
\mbox{}\verb@    adrp x0, _kcc_tok_start@{\tt @}\verb@PAGE@\\
\mbox{}\verb@    add  x0, x0, _kcc_tok_start@{\tt @}\verb@PAGEOFF@\\
\mbox{}\verb@    ldr  x0, [x0]@\\
\mbox{}\verb@@\\
\mbox{}\verb@    adrp x1, _kcc_tok_len@{\tt @}\verb@PAGE@\\
\mbox{}\verb@    add  x1, x1, _kcc_tok_len@{\tt @}\verb@PAGEOFF@\\
\mbox{}\verb@    ldr  x1, [x1]@\\
\mbox{}\verb@@\\
\mbox{}\verb@    bl _kcc_io_write@\\
\mbox{}\verb@@\\
\mbox{}\verb@    // Print "'\n"@\\
\mbox{}\verb@    mov x0, #'\''@\\
\mbox{}\verb@    bl _kcc_io_putc@\\
\mbox{}\verb@    mov x0, #0x0A@\\
\mbox{}\verb@    bl _kcc_io_putc@\\
\mbox{}\verb@@\\
\mbox{}\verb@    ldp x29, x30, [sp], #16@\\
\mbox{}\verb@    ret@\\
\mbox{}\verb@@{\NWsep}
\end{list}
\vspace{-1.5ex}
\footnotesize
\begin{list}{}{\setlength{\itemsep}{-\parsep}\setlength{\itemindent}{-\leftmargin}}
\item \NWtxtMacroRefIn\ \NWlink{weft?}{?}.

\item{}
\end{list}
\end{minipage}\vspace{4ex}
\end{flushleft}
The dump-all loop repeatedly calls \texttt{kcc\_lex\_next} and
\texttt{kcc\_lex\_dump\_token} until EOF.

\begin{flushleft} \small
\begin{minipage}{\linewidth}\label{scrap66}\raggedright\small
\NWtarget{weft?}{} $\langle\,${\itshape Lexer: dump all}\nobreak\ {\footnotesize {?}}$\,\rangle\equiv$
\vspace{-1ex}
\begin{list}{}{} \item
\mbox{}\verb@@\\
\mbox{}\verb@// -- kcc_lex_dump_all -----------------------------------------------@\\
\mbox{}\verb@// void kcc_lex_dump_all(void)@\\
\mbox{}\verb@// Tokenizes entire source and dumps each token.@\\
\mbox{}\verb@.globl _kcc_lex_dump_all@\\
\mbox{}\verb@.p2align 2@\\
\mbox{}\verb@_kcc_lex_dump_all:@\\
\mbox{}\verb@    stp x29, x30, [sp, #-16]!@\\
\mbox{}\verb@    mov x29, sp@\\
\mbox{}\verb@@\\
\mbox{}\verb@.Ldump_loop:@\\
\mbox{}\verb@    bl _kcc_lex_next@\\
\mbox{}\verb@@\\
\mbox{}\verb@    // Check for EOF@\\
\mbox{}\verb@    adrp x0, _kcc_tok_type@{\tt @}\verb@PAGE@\\
\mbox{}\verb@    add  x0, x0, _kcc_tok_type@{\tt @}\verb@PAGEOFF@\\
\mbox{}\verb@    ldr  x0, [x0]@\\
\mbox{}\verb@    cbz  x0, .Ldump_done        // TOK_EOF = 0@\\
\mbox{}\verb@@\\
\mbox{}\verb@    bl _kcc_lex_dump_token@\\
\mbox{}\verb@    b .Ldump_loop@\\
\mbox{}\verb@@\\
\mbox{}\verb@.Ldump_done:@\\
\mbox{}\verb@    ldp x29, x30, [sp], #16@\\
\mbox{}\verb@    ret@\\
\mbox{}\verb@@{\NWsep}
\end{list}
\vspace{-1.5ex}
\footnotesize
\begin{list}{}{\setlength{\itemsep}{-\parsep}\setlength{\itemindent}{-\leftmargin}}
\item \NWtxtMacroRefIn\ \NWlink{weft?}{?}.

\item{}
\end{list}
\end{minipage}\vspace{4ex}
\end{flushleft}
\section{Token Type Names}

A table mapping token type codes to human-readable names for
the \texttt{--dump-tokens} output.  \texttt{kcc\_lex\_type\_name}
returns a pointer to the name string for a given token type.

\begin{flushleft} \small
\begin{minipage}{\linewidth}\label{scrap67}\raggedright\small
\NWtarget{weft?}{} $\langle\,${\itshape Lexer: token names}\nobreak\ {\footnotesize {?}}$\,\rangle\equiv$
\vspace{-1ex}
\begin{list}{}{} \item
\mbox{}\verb@@\\
\mbox{}\verb@.text@\\
\mbox{}\verb@// -- kcc_lex_type_name ----------------------------------------------@\\
\mbox{}\verb@// const char *kcc_lex_type_name(uint64_t type)@\\
\mbox{}\verb@//   x0 = token type code@\\
\mbox{}\verb@// Returns pointer to a null-terminated type name string.@\\
\mbox{}\verb@.globl _kcc_lex_type_name@\\
\mbox{}\verb@.p2align 2@\\
\mbox{}\verb@_kcc_lex_type_name:@\\
\mbox{}\verb@    stp x29, x30, [sp, #-16]!@\\
\mbox{}\verb@    mov x29, sp@\\
\mbox{}\verb@@\\
\mbox{}\verb@    // Bounds check@\\
\mbox{}\verb@    cmp x0, #108@\\
\mbox{}\verb@    b.hs .Ltype_unknown_name@\\
\mbox{}\verb@@\\
\mbox{}\verb@    // Load from pointer table@\\
\mbox{}\verb@    adrp x1, _kcc_tok_names@{\tt @}\verb@PAGE@\\
\mbox{}\verb@    add  x1, x1, _kcc_tok_names@{\tt @}\verb@PAGEOFF@\\
\mbox{}\verb@    ldr  x0, [x1, x0, lsl #3]  // names[type]@\\
\mbox{}\verb@    cbnz x0, .Ltype_name_done@\\
\mbox{}\verb@@\\
\mbox{}\verb@.Ltype_unknown_name:@\\
\mbox{}\verb@    adrp x0, _kcc_tn_unknown@{\tt @}\verb@PAGE@\\
\mbox{}\verb@    add  x0, x0, _kcc_tn_unknown@{\tt @}\verb@PAGEOFF@\\
\mbox{}\verb@@\\
\mbox{}\verb@.Ltype_name_done:@\\
\mbox{}\verb@    ldp x29, x30, [sp], #16@\\
\mbox{}\verb@    ret@\\
\mbox{}\verb@@\\
\mbox{}\verb@// -- Token name strings ---------------------------------------------@\\
\mbox{}\verb@.section __TEXT,__cstring,cstring_literals@\\
\mbox{}\verb@_kcc_tn_eof:          .asciz "eof"@\\
\mbox{}\verb@_kcc_tn_identifier:   .asciz "identifier"@\\
\mbox{}\verb@_kcc_tn_number:       .asciz "number"@\\
\mbox{}\verb@_kcc_tn_char_lit:     .asciz "char_constant"@\\
\mbox{}\verb@_kcc_tn_string_lit:   .asciz "string_literal"@\\
\mbox{}\verb@_kcc_tn_unknown:      .asciz "unknown"@\\
\mbox{}\verb@// Keywords@\\
\mbox{}\verb@_kcc_tn_auto:         .asciz "auto"@\\
\mbox{}\verb@_kcc_tn_break:        .asciz "break"@\\
\mbox{}\verb@_kcc_tn_case:         .asciz "case"@\\
\mbox{}\verb@_kcc_tn_char:         .asciz "char"@\\
\mbox{}\verb@_kcc_tn_const:        .asciz "const"@\\
\mbox{}\verb@_kcc_tn_continue:     .asciz "continue"@\\
\mbox{}\verb@_kcc_tn_default:      .asciz "default"@\\
\mbox{}\verb@_kcc_tn_do:           .asciz "do"@\\
\mbox{}\verb@_kcc_tn_double:       .asciz "double"@\\
\mbox{}\verb@_kcc_tn_else:         .asciz "else"@\\
\mbox{}\verb@_kcc_tn_enum:         .asciz "enum"@\\
\mbox{}\verb@_kcc_tn_extern:       .asciz "extern"@\\
\mbox{}\verb@_kcc_tn_float:        .asciz "float"@\\
\mbox{}\verb@_kcc_tn_for:          .asciz "for"@\\
\mbox{}\verb@_kcc_tn_goto:         .asciz "goto"@\\
\mbox{}\verb@_kcc_tn_if:           .asciz "if"@\\
\mbox{}\verb@_kcc_tn_inline:       .asciz "inline"@\\
\mbox{}\verb@_kcc_tn_int:          .asciz "int"@\\
\mbox{}\verb@_kcc_tn_long:         .asciz "long"@\\
\mbox{}\verb@_kcc_tn_register:     .asciz "register"@\\
\mbox{}\verb@_kcc_tn_restrict:     .asciz "restrict"@\\
\mbox{}\verb@_kcc_tn_return:       .asciz "return"@\\
\mbox{}\verb@_kcc_tn_short:        .asciz "short"@\\
\mbox{}\verb@_kcc_tn_signed:       .asciz "signed"@\\
\mbox{}\verb@_kcc_tn_sizeof:       .asciz "sizeof"@\\
\mbox{}\verb@_kcc_tn_static:       .asciz "static"@\\
\mbox{}\verb@_kcc_tn_struct:       .asciz "struct"@\\
\mbox{}\verb@_kcc_tn_switch:       .asciz "switch"@\\
\mbox{}\verb@_kcc_tn_typedef:      .asciz "typedef"@\\
\mbox{}\verb@_kcc_tn_union:        .asciz "union"@\\
\mbox{}\verb@_kcc_tn_unsigned:     .asciz "unsigned"@\\
\mbox{}\verb@_kcc_tn_void:         .asciz "void"@\\
\mbox{}\verb@_kcc_tn_volatile:     .asciz "volatile"@\\
\mbox{}\verb@_kcc_tn_while:        .asciz "while"@\\
\mbox{}\verb@_kcc_tn_alignas:      .asciz "_Alignas"@\\
\mbox{}\verb@_kcc_tn_alignof:      .asciz "_Alignof"@\\
\mbox{}\verb@_kcc_tn_atomic:       .asciz "_Atomic"@\\
\mbox{}\verb@_kcc_tn_bool:         .asciz "_Bool"@\\
\mbox{}\verb@_kcc_tn_complex:      .asciz "_Complex"@\\
\mbox{}\verb@_kcc_tn_generic:      .asciz "_Generic"@\\
\mbox{}\verb@_kcc_tn_imaginary:    .asciz "_Imaginary"@\\
\mbox{}\verb@_kcc_tn_noreturn:     .asciz "_Noreturn"@\\
\mbox{}\verb@_kcc_tn_static_assert: .asciz "_Static_assert"@\\
\mbox{}\verb@_kcc_tn_thread_local: .asciz "_Thread_local"@\\
\mbox{}\verb@// Punctuators@\\
\mbox{}\verb@_kcc_tn_lparen:       .asciz "l_paren"@\\
\mbox{}\verb@_kcc_tn_rparen:       .asciz "r_paren"@\\
\mbox{}\verb@_kcc_tn_lbrace:       .asciz "l_brace"@\\
\mbox{}\verb@_kcc_tn_rbrace:       .asciz "r_brace"@\\
\mbox{}\verb@_kcc_tn_lbracket:     .asciz "l_bracket"@\\
\mbox{}\verb@_kcc_tn_rbracket:     .asciz "r_bracket"@\\
\mbox{}\verb@_kcc_tn_semi:         .asciz "semi"@\\
\mbox{}\verb@_kcc_tn_comma:        .asciz "comma"@\\
\mbox{}\verb@_kcc_tn_dot:          .asciz "dot"@\\
\mbox{}\verb@_kcc_tn_arrow:        .asciz "arrow"@\\
\mbox{}\verb@_kcc_tn_plusplus:     .asciz "plusplus"@\\
\mbox{}\verb@_kcc_tn_minusminus:  .asciz "minusminus"@\\
\mbox{}\verb@_kcc_tn_amp:          .asciz "amp"@\\
\mbox{}\verb@_kcc_tn_star:         .asciz "star"@\\
\mbox{}\verb@_kcc_tn_plus:         .asciz "plus"@\\
\mbox{}\verb@_kcc_tn_minus:        .asciz "minus"@\\
\mbox{}\verb@_kcc_tn_tilde:        .asciz "tilde"@\\
\mbox{}\verb@_kcc_tn_exclaim:      .asciz "exclaim"@\\
\mbox{}\verb@_kcc_tn_slash:        .asciz "slash"@\\
\mbox{}\verb@_kcc_tn_percent:      .asciz "percent"@\\
\mbox{}\verb@_kcc_tn_lessless:     .asciz "lessless"@\\
\mbox{}\verb@_kcc_tn_greatergreater: .asciz "greatergreater"@\\
\mbox{}\verb@_kcc_tn_less:         .asciz "less"@\\
\mbox{}\verb@_kcc_tn_greater:      .asciz "greater"@\\
\mbox{}\verb@_kcc_tn_lessequal:    .asciz "lessequal"@\\
\mbox{}\verb@_kcc_tn_greaterequal: .asciz "greaterequal"@\\
\mbox{}\verb@_kcc_tn_equalequal:   .asciz "equalequal"@\\
\mbox{}\verb@_kcc_tn_exclaimequal: .asciz "exclaimequal"@\\
\mbox{}\verb@_kcc_tn_caret:        .asciz "caret"@\\
\mbox{}\verb@_kcc_tn_pipe:         .asciz "pipe"@\\
\mbox{}\verb@_kcc_tn_ampamp:       .asciz "ampamp"@\\
\mbox{}\verb@_kcc_tn_pipepipe:     .asciz "pipepipe"@\\
\mbox{}\verb@_kcc_tn_question:     .asciz "question"@\\
\mbox{}\verb@_kcc_tn_colon:        .asciz "colon"@\\
\mbox{}\verb@_kcc_tn_equal:        .asciz "equal"@\\
\mbox{}\verb@_kcc_tn_starequal:    .asciz "starequal"@\\
\mbox{}\verb@_kcc_tn_slashequal:   .asciz "slashequal"@\\
\mbox{}\verb@_kcc_tn_percentequal: .asciz "percentequal"@\\
\mbox{}\verb@_kcc_tn_plusequal:    .asciz "plusequal"@\\
\mbox{}\verb@_kcc_tn_minusequal:   .asciz "minusequal"@\\
\mbox{}\verb@_kcc_tn_lesslessequal: .asciz "lesslessequal"@\\
\mbox{}\verb@_kcc_tn_greatergreaterequal: .asciz "greatergreaterequal"@\\
\mbox{}\verb@_kcc_tn_ampequal:     .asciz "ampequal"@\\
\mbox{}\verb@_kcc_tn_caretequal:   .asciz "caretequal"@\\
\mbox{}\verb@_kcc_tn_pipeequal:    .asciz "pipeequal"@\\
\mbox{}\verb@_kcc_tn_hash:         .asciz "hash"@\\
\mbox{}\verb@_kcc_tn_hashhash:     .asciz "hashhash"@\\
\mbox{}\verb@_kcc_tn_ellipsis:     .asciz "ellipsis"@\\
\mbox{}\verb@@\\
\mbox{}\verb@// -- Token name pointer table ---------------------------------------@\\
\mbox{}\verb@// Indexed by token type code.  Entries 6-9, 54-59 are NULL (unused).@\\
\mbox{}\verb@.data@\\
\mbox{}\verb@.p2align 3@\\
\mbox{}\verb@_kcc_tok_names:@\\
\mbox{}\verb@    .quad _kcc_tn_eof             // 0  TOK_EOF@\\
\mbox{}\verb@    .quad _kcc_tn_identifier      // 1  TOK_IDENT@\\
\mbox{}\verb@    .quad _kcc_tn_number          // 2  TOK_NUMBER@\\
\mbox{}\verb@    .quad _kcc_tn_char_lit        // 3  TOK_CHAR_LIT@\\
\mbox{}\verb@    .quad _kcc_tn_string_lit      // 4  TOK_STRING_LIT@\\
\mbox{}\verb@    .quad _kcc_tn_unknown         // 5  TOK_UNKNOWN@\\
\mbox{}\verb@    .quad 0, 0, 0, 0             // 6-9  reserved@\\
\mbox{}\verb@    // Keywords 10-53@\\
\mbox{}\verb@    .quad _kcc_tn_auto            // 10@\\
\mbox{}\verb@    .quad _kcc_tn_break           // 11@\\
\mbox{}\verb@    .quad _kcc_tn_case            // 12@\\
\mbox{}\verb@    .quad _kcc_tn_char            // 13@\\
\mbox{}\verb@    .quad _kcc_tn_const           // 14@\\
\mbox{}\verb@    .quad _kcc_tn_continue        // 15@\\
\mbox{}\verb@    .quad _kcc_tn_default         // 16@\\
\mbox{}\verb@    .quad _kcc_tn_do              // 17@\\
\mbox{}\verb@    .quad _kcc_tn_double          // 18@\\
\mbox{}\verb@    .quad _kcc_tn_else            // 19@\\
\mbox{}\verb@    .quad _kcc_tn_enum            // 20@\\
\mbox{}\verb@    .quad _kcc_tn_extern          // 21@\\
\mbox{}\verb@    .quad _kcc_tn_float           // 22@\\
\mbox{}\verb@    .quad _kcc_tn_for             // 23@\\
\mbox{}\verb@    .quad _kcc_tn_goto            // 24@\\
\mbox{}\verb@    .quad _kcc_tn_if              // 25@\\
\mbox{}\verb@    .quad _kcc_tn_inline          // 26@\\
\mbox{}\verb@    .quad _kcc_tn_int             // 27@\\
\mbox{}\verb@    .quad _kcc_tn_long            // 28@\\
\mbox{}\verb@    .quad _kcc_tn_register        // 29@\\
\mbox{}\verb@    .quad _kcc_tn_restrict        // 30@\\
\mbox{}\verb@    .quad _kcc_tn_return          // 31@\\
\mbox{}\verb@    .quad _kcc_tn_short           // 32@\\
\mbox{}\verb@    .quad _kcc_tn_signed          // 33@\\
\mbox{}\verb@    .quad _kcc_tn_sizeof          // 34@\\
\mbox{}\verb@    .quad _kcc_tn_static          // 35@\\
\mbox{}\verb@    .quad _kcc_tn_struct          // 36@\\
\mbox{}\verb@    .quad _kcc_tn_switch          // 37@\\
\mbox{}\verb@    .quad _kcc_tn_typedef         // 38@\\
\mbox{}\verb@    .quad _kcc_tn_union           // 39@\\
\mbox{}\verb@    .quad _kcc_tn_unsigned        // 40@\\
\mbox{}\verb@    .quad _kcc_tn_void            // 41@\\
\mbox{}\verb@    .quad _kcc_tn_volatile        // 42@\\
\mbox{}\verb@    .quad _kcc_tn_while           // 43@\\
\mbox{}\verb@    .quad _kcc_tn_alignas         // 44@\\
\mbox{}\verb@    .quad _kcc_tn_alignof         // 45@\\
\mbox{}\verb@    .quad _kcc_tn_atomic          // 46@\\
\mbox{}\verb@    .quad _kcc_tn_bool            // 47@\\
\mbox{}\verb@    .quad _kcc_tn_complex         // 48@\\
\mbox{}\verb@    .quad _kcc_tn_generic         // 49@\\
\mbox{}\verb@    .quad _kcc_tn_imaginary       // 50@\\
\mbox{}\verb@    .quad _kcc_tn_noreturn        // 51@\\
\mbox{}\verb@    .quad _kcc_tn_static_assert   // 52@\\
\mbox{}\verb@    .quad _kcc_tn_thread_local    // 53@\\
\mbox{}\verb@    .quad 0, 0, 0, 0, 0, 0       // 54-59  reserved@\\
\mbox{}\verb@    // Punctuators 60-107@\\
\mbox{}\verb@    .quad _kcc_tn_lparen          // 60@\\
\mbox{}\verb@    .quad _kcc_tn_rparen          // 61@\\
\mbox{}\verb@    .quad _kcc_tn_lbrace          // 62@\\
\mbox{}\verb@    .quad _kcc_tn_rbrace          // 63@\\
\mbox{}\verb@    .quad _kcc_tn_lbracket        // 64@\\
\mbox{}\verb@    .quad _kcc_tn_rbracket        // 65@\\
\mbox{}\verb@    .quad _kcc_tn_semi            // 66@\\
\mbox{}\verb@    .quad _kcc_tn_comma           // 67@\\
\mbox{}\verb@    .quad _kcc_tn_dot             // 68@\\
\mbox{}\verb@    .quad _kcc_tn_arrow           // 69@\\
\mbox{}\verb@    .quad _kcc_tn_plusplus        // 70@\\
\mbox{}\verb@    .quad _kcc_tn_minusminus     // 71@\\
\mbox{}\verb@    .quad _kcc_tn_amp             // 72@\\
\mbox{}\verb@    .quad _kcc_tn_star            // 73@\\
\mbox{}\verb@    .quad _kcc_tn_plus            // 74@\\
\mbox{}\verb@    .quad _kcc_tn_minus           // 75@\\
\mbox{}\verb@    .quad _kcc_tn_tilde           // 76@\\
\mbox{}\verb@    .quad _kcc_tn_exclaim         // 77@\\
\mbox{}\verb@    .quad _kcc_tn_slash           // 78@\\
\mbox{}\verb@    .quad _kcc_tn_percent         // 79@\\
\mbox{}\verb@    .quad _kcc_tn_lessless        // 80@\\
\mbox{}\verb@    .quad _kcc_tn_greatergreater  // 81@\\
\mbox{}\verb@    .quad _kcc_tn_less            // 82@\\
\mbox{}\verb@    .quad _kcc_tn_greater         // 83@\\
\mbox{}\verb@    .quad _kcc_tn_lessequal       // 84@\\
\mbox{}\verb@    .quad _kcc_tn_greaterequal    // 85@\\
\mbox{}\verb@    .quad _kcc_tn_equalequal      // 86@\\
\mbox{}\verb@    .quad _kcc_tn_exclaimequal    // 87@\\
\mbox{}\verb@    .quad _kcc_tn_caret           // 88@\\
\mbox{}\verb@    .quad _kcc_tn_pipe            // 89@\\
\mbox{}\verb@    .quad _kcc_tn_ampamp          // 90@\\
\mbox{}\verb@    .quad _kcc_tn_pipepipe        // 91@\\
\mbox{}\verb@    .quad _kcc_tn_question        // 92@\\
\mbox{}\verb@    .quad _kcc_tn_colon           // 93@\\
\mbox{}\verb@    .quad _kcc_tn_equal           // 94@\\
\mbox{}\verb@    .quad _kcc_tn_starequal       // 95@\\
\mbox{}\verb@    .quad _kcc_tn_slashequal      // 96@\\
\mbox{}\verb@    .quad _kcc_tn_percentequal    // 97@\\
\mbox{}\verb@    .quad _kcc_tn_plusequal       // 98@\\
\mbox{}\verb@    .quad _kcc_tn_minusequal      // 99@\\
\mbox{}\verb@    .quad _kcc_tn_lesslessequal   // 100@\\
\mbox{}\verb@    .quad _kcc_tn_greatergreaterequal // 101@\\
\mbox{}\verb@    .quad _kcc_tn_ampequal        // 102@\\
\mbox{}\verb@    .quad _kcc_tn_caretequal      // 103@\\
\mbox{}\verb@    .quad _kcc_tn_pipeequal       // 104@\\
\mbox{}\verb@    .quad _kcc_tn_hash            // 105@\\
\mbox{}\verb@    .quad _kcc_tn_hashhash        // 106@\\
\mbox{}\verb@    .quad _kcc_tn_ellipsis        // 107@\\
\mbox{}\verb@@{\NWsep}
\end{list}
\vspace{-1.5ex}
\footnotesize
\begin{list}{}{\setlength{\itemsep}{-\parsep}\setlength{\itemindent}{-\leftmargin}}
\item \NWtxtMacroRefIn\ \NWlink{weft?}{?}.

\item{}
\end{list}
\end{minipage}\vspace{4ex}
\end{flushleft}
\section{Keyword Table Initialization}

At startup, we populate a hash table with all 44 C17 keywords.
Each entry maps the keyword text to its token type code.

\begin{flushleft} \small
\begin{minipage}{\linewidth}\label{scrap68}\raggedright\small
\NWtarget{weft?}{} $\langle\,${\itshape Lexer: keyword data}\nobreak\ {\footnotesize {?}}$\,\rangle\equiv$
\vspace{-1ex}
\begin{list}{}{} \item
\mbox{}\verb@@\\
\mbox{}\verb@// -- Keyword strings ------------------------------------------------@\\
\mbox{}\verb@.section __TEXT,__cstring,cstring_literals@\\
\mbox{}\verb@_kcc_kw_s_auto:       .asciz "auto"@\\
\mbox{}\verb@_kcc_kw_s_break:      .asciz "break"@\\
\mbox{}\verb@_kcc_kw_s_case:       .asciz "case"@\\
\mbox{}\verb@_kcc_kw_s_char:       .asciz "char"@\\
\mbox{}\verb@_kcc_kw_s_const:      .asciz "const"@\\
\mbox{}\verb@_kcc_kw_s_continue:   .asciz "continue"@\\
\mbox{}\verb@_kcc_kw_s_default:    .asciz "default"@\\
\mbox{}\verb@_kcc_kw_s_do:         .asciz "do"@\\
\mbox{}\verb@_kcc_kw_s_double:     .asciz "double"@\\
\mbox{}\verb@_kcc_kw_s_else:       .asciz "else"@\\
\mbox{}\verb@_kcc_kw_s_enum:       .asciz "enum"@\\
\mbox{}\verb@_kcc_kw_s_extern:     .asciz "extern"@\\
\mbox{}\verb@_kcc_kw_s_float:      .asciz "float"@\\
\mbox{}\verb@_kcc_kw_s_for:        .asciz "for"@\\
\mbox{}\verb@_kcc_kw_s_goto:       .asciz "goto"@\\
\mbox{}\verb@_kcc_kw_s_if:         .asciz "if"@\\
\mbox{}\verb@_kcc_kw_s_inline:     .asciz "inline"@\\
\mbox{}\verb@_kcc_kw_s_int:        .asciz "int"@\\
\mbox{}\verb@_kcc_kw_s_long:       .asciz "long"@\\
\mbox{}\verb@_kcc_kw_s_register:   .asciz "register"@\\
\mbox{}\verb@_kcc_kw_s_restrict:   .asciz "restrict"@\\
\mbox{}\verb@_kcc_kw_s_return:     .asciz "return"@\\
\mbox{}\verb@_kcc_kw_s_short:      .asciz "short"@\\
\mbox{}\verb@_kcc_kw_s_signed:     .asciz "signed"@\\
\mbox{}\verb@_kcc_kw_s_sizeof:     .asciz "sizeof"@\\
\mbox{}\verb@_kcc_kw_s_static:     .asciz "static"@\\
\mbox{}\verb@_kcc_kw_s_struct:     .asciz "struct"@\\
\mbox{}\verb@_kcc_kw_s_switch:     .asciz "switch"@\\
\mbox{}\verb@_kcc_kw_s_typedef:    .asciz "typedef"@\\
\mbox{}\verb@_kcc_kw_s_union:      .asciz "union"@\\
\mbox{}\verb@_kcc_kw_s_unsigned:   .asciz "unsigned"@\\
\mbox{}\verb@_kcc_kw_s_void:       .asciz "void"@\\
\mbox{}\verb@_kcc_kw_s_volatile:   .asciz "volatile"@\\
\mbox{}\verb@_kcc_kw_s_while:      .asciz "while"@\\
\mbox{}\verb@_kcc_kw_s_alignas:    .asciz "_Alignas"@\\
\mbox{}\verb@_kcc_kw_s_alignof:    .asciz "_Alignof"@\\
\mbox{}\verb@_kcc_kw_s_atomic:     .asciz "_Atomic"@\\
\mbox{}\verb@_kcc_kw_s_bool:       .asciz "_Bool"@\\
\mbox{}\verb@_kcc_kw_s_complex:    .asciz "_Complex"@\\
\mbox{}\verb@_kcc_kw_s_generic:    .asciz "_Generic"@\\
\mbox{}\verb@_kcc_kw_s_imaginary:  .asciz "_Imaginary"@\\
\mbox{}\verb@_kcc_kw_s_noreturn:   .asciz "_Noreturn"@\\
\mbox{}\verb@_kcc_kw_s_sassert:    .asciz "_Static_assert"@\\
\mbox{}\verb@_kcc_kw_s_tlocal:     .asciz "_Thread_local"@\\
\mbox{}\verb@@\\
\mbox{}\verb@// -- Keyword registration table -------------------------------------@\\
\mbox{}\verb@// Each entry: (string_ptr, string_len, token_type) = 24 bytes@\\
\mbox{}\verb@.data@\\
\mbox{}\verb@.p2align 3@\\
\mbox{}\verb@_kcc_kw_reg_table:@\\
\mbox{}\verb@    .quad _kcc_kw_s_auto,       4,  10@\\
\mbox{}\verb@    .quad _kcc_kw_s_break,      5,  11@\\
\mbox{}\verb@    .quad _kcc_kw_s_case,       4,  12@\\
\mbox{}\verb@    .quad _kcc_kw_s_char,       4,  13@\\
\mbox{}\verb@    .quad _kcc_kw_s_const,      5,  14@\\
\mbox{}\verb@    .quad _kcc_kw_s_continue,   8,  15@\\
\mbox{}\verb@    .quad _kcc_kw_s_default,    7,  16@\\
\mbox{}\verb@    .quad _kcc_kw_s_do,         2,  17@\\
\mbox{}\verb@    .quad _kcc_kw_s_double,     6,  18@\\
\mbox{}\verb@    .quad _kcc_kw_s_else,       4,  19@\\
\mbox{}\verb@    .quad _kcc_kw_s_enum,       4,  20@\\
\mbox{}\verb@    .quad _kcc_kw_s_extern,     6,  21@\\
\mbox{}\verb@    .quad _kcc_kw_s_float,      5,  22@\\
\mbox{}\verb@    .quad _kcc_kw_s_for,        3,  23@\\
\mbox{}\verb@    .quad _kcc_kw_s_goto,       4,  24@\\
\mbox{}\verb@    .quad _kcc_kw_s_if,         2,  25@\\
\mbox{}\verb@    .quad _kcc_kw_s_inline,     6,  26@\\
\mbox{}\verb@    .quad _kcc_kw_s_int,        3,  27@\\
\mbox{}\verb@    .quad _kcc_kw_s_long,       4,  28@\\
\mbox{}\verb@    .quad _kcc_kw_s_register,   8,  29@\\
\mbox{}\verb@    .quad _kcc_kw_s_restrict,   8,  30@\\
\mbox{}\verb@    .quad _kcc_kw_s_return,     6,  31@\\
\mbox{}\verb@    .quad _kcc_kw_s_short,      5,  32@\\
\mbox{}\verb@    .quad _kcc_kw_s_signed,     6,  33@\\
\mbox{}\verb@    .quad _kcc_kw_s_sizeof,     6,  34@\\
\mbox{}\verb@    .quad _kcc_kw_s_static,     6,  35@\\
\mbox{}\verb@    .quad _kcc_kw_s_struct,     6,  36@\\
\mbox{}\verb@    .quad _kcc_kw_s_switch,     6,  37@\\
\mbox{}\verb@    .quad _kcc_kw_s_typedef,    7,  38@\\
\mbox{}\verb@    .quad _kcc_kw_s_union,      5,  39@\\
\mbox{}\verb@    .quad _kcc_kw_s_unsigned,   8,  40@\\
\mbox{}\verb@    .quad _kcc_kw_s_void,       4,  41@\\
\mbox{}\verb@    .quad _kcc_kw_s_volatile,   8,  42@\\
\mbox{}\verb@    .quad _kcc_kw_s_while,      5,  43@\\
\mbox{}\verb@    .quad _kcc_kw_s_alignas,    8,  44@\\
\mbox{}\verb@    .quad _kcc_kw_s_alignof,    8,  45@\\
\mbox{}\verb@    .quad _kcc_kw_s_atomic,     7,  46@\\
\mbox{}\verb@    .quad _kcc_kw_s_bool,       5,  47@\\
\mbox{}\verb@    .quad _kcc_kw_s_complex,    8,  48@\\
\mbox{}\verb@    .quad _kcc_kw_s_generic,    8,  49@\\
\mbox{}\verb@    .quad _kcc_kw_s_imaginary,  10, 50@\\
\mbox{}\verb@    .quad _kcc_kw_s_noreturn,   9,  51@\\
\mbox{}\verb@    .quad _kcc_kw_s_sassert,    14, 52@\\
\mbox{}\verb@    .quad _kcc_kw_s_tlocal,     13, 53@\\
\mbox{}\verb@@\\
\mbox{}\verb@.equ KW_COUNT, 44@\\
\mbox{}\verb@.equ KW_ENTRY_BYTES, 24       // 3 quads per entry@\\
\mbox{}\verb@@{\NWsep}
\end{list}
\vspace{-1.5ex}
\footnotesize
\begin{list}{}{\setlength{\itemsep}{-\parsep}\setlength{\itemindent}{-\leftmargin}}
\item \NWtxtMacroRefIn\ \NWlink{weft?}{?}.

\item{}
\end{list}
\end{minipage}\vspace{4ex}
\end{flushleft}
The initialization function creates a hash table and populates it
by looping through the registration table.

\begin{flushleft} \small
\begin{minipage}{\linewidth}\label{scrap69}\raggedright\small
\NWtarget{weft?}{} $\langle\,${\itshape Lexer: keyword init}\nobreak\ {\footnotesize {?}}$\,\rangle\equiv$
\vspace{-1ex}
\begin{list}{}{} \item
\mbox{}\verb@@\\
\mbox{}\verb@.text@\\
\mbox{}\verb@// -- kcc_lex_kw_init ------------------------------------------------@\\
\mbox{}\verb@// void kcc_lex_kw_init(void)@\\
\mbox{}\verb@// Builds the keyword hash table from the registration table.@\\
\mbox{}\verb@.globl _kcc_lex_kw_init@\\
\mbox{}\verb@.p2align 2@\\
\mbox{}\verb@_kcc_lex_kw_init:@\\
\mbox{}\verb@    stp x29, x30, [sp, #-48]!@\\
\mbox{}\verb@    mov x29, sp@\\
\mbox{}\verb@    stp x19, x20, [sp, #16]@\\
\mbox{}\verb@    stp x21, x22, [sp, #32]@\\
\mbox{}\verb@@\\
\mbox{}\verb@    // Create hash table with capacity 128 (power of 2, >2x keywords)@\\
\mbox{}\verb@    mov x0, #128@\\
\mbox{}\verb@    bl _kcc_ht_init@\\
\mbox{}\verb@    mov x19, x0                  // x19 = hash table@\\
\mbox{}\verb@@\\
\mbox{}\verb@    // Store in global@\\
\mbox{}\verb@    adrp x1, _kcc_kw_table@{\tt @}\verb@PAGE@\\
\mbox{}\verb@    add  x1, x1, _kcc_kw_table@{\tt @}\verb@PAGEOFF@\\
\mbox{}\verb@    str  x19, [x1]@\\
\mbox{}\verb@@\\
\mbox{}\verb@    // Loop through registration table@\\
\mbox{}\verb@    adrp x20, _kcc_kw_reg_table@{\tt @}\verb@PAGE@\\
\mbox{}\verb@    add  x20, x20, _kcc_kw_reg_table@{\tt @}\verb@PAGEOFF@\\
\mbox{}\verb@    mov  x21, #0                 // counter@\\
\mbox{}\verb@    mov  x22, #KW_COUNT@\\
\mbox{}\verb@@\\
\mbox{}\verb@.Lkw_init_loop:@\\
\mbox{}\verb@    cmp x21, x22@\\
\mbox{}\verb@    b.hs .Lkw_init_done@\\
\mbox{}\verb@@\\
\mbox{}\verb@    // Load entry: [string_ptr, string_len, token_type]@\\
\mbox{}\verb@    ldr x1, [x20, #0]           // key = string ptr@\\
\mbox{}\verb@    ldr x2, [x20, #8]           // keylen = string length@\\
\mbox{}\verb@    ldr x3, [x20, #16]          // value = token type@\\
\mbox{}\verb@@\\
\mbox{}\verb@    mov x0, x19                  // ht@\\
\mbox{}\verb@    bl _kcc_ht_put@\\
\mbox{}\verb@@\\
\mbox{}\verb@    add x20, x20, #KW_ENTRY_BYTES@\\
\mbox{}\verb@    add x21, x21, #1@\\
\mbox{}\verb@    b .Lkw_init_loop@\\
\mbox{}\verb@@\\
\mbox{}\verb@.Lkw_init_done:@\\
\mbox{}\verb@    ldp x21, x22, [sp, #32]@\\
\mbox{}\verb@    ldp x19, x20, [sp, #16]@\\
\mbox{}\verb@    ldp x29, x30, [sp], #48@\\
\mbox{}\verb@    ret@\\
\mbox{}\verb@@{\NWsep}
\end{list}
\vspace{-1.5ex}
\footnotesize
\begin{list}{}{\setlength{\itemsep}{-\parsep}\setlength{\itemindent}{-\leftmargin}}
\item \NWtxtMacroRefIn\ \NWlink{weft?}{?}.

\item{}
\end{list}
\end{minipage}\vspace{4ex}
\end{flushleft}
\section{Lexer Global State}

The lexer maintains its state in global variables: the cursor position,
end of buffer, current line/column tracking, and the current token.

\begin{flushleft} \small
\begin{minipage}{\linewidth}\label{scrap70}\raggedright\small
\NWtarget{weft?}{} $\langle\,${\itshape Lexer: data}\nobreak\ {\footnotesize {?}}$\,\rangle\equiv$
\vspace{-1ex}
\begin{list}{}{} \item
\mbox{}\verb@@\\
\mbox{}\verb@// -- Lexer global state ---------------------------------------------@\\
\mbox{}\verb@.data@\\
\mbox{}\verb@.p2align 3@\\
\mbox{}\verb@_kcc_lex_cursor:     .quad 0  // current position in source buffer@\\
\mbox{}\verb@_kcc_lex_end:        .quad 0  // one past end of source buffer@\\
\mbox{}\verb@_kcc_lex_line:       .quad 1  // current line number (1-based)@\\
\mbox{}\verb@_kcc_lex_line_start: .quad 0  // pointer to start of current line@\\
\mbox{}\verb@@\\
\mbox{}\verb@// Current token (exported for test harness access)@\\
\mbox{}\verb@.globl _kcc_tok_type@\\
\mbox{}\verb@_kcc_tok_type:       .quad 0  // token type code@\\
\mbox{}\verb@.globl _kcc_tok_start@\\
\mbox{}\verb@_kcc_tok_start:      .quad 0  // pointer to token text in source@\\
\mbox{}\verb@.globl _kcc_tok_len@\\
\mbox{}\verb@_kcc_tok_len:        .quad 0  // token text length@\\
\mbox{}\verb@.globl _kcc_tok_line@\\
\mbox{}\verb@_kcc_tok_line:       .quad 0  // line where token starts@\\
\mbox{}\verb@.globl _kcc_tok_col@\\
\mbox{}\verb@_kcc_tok_col:        .quad 0  // column where token starts@\\
\mbox{}\verb@@\\
\mbox{}\verb@// Keyword hash table@\\
\mbox{}\verb@.globl _kcc_kw_table@\\
\mbox{}\verb@_kcc_kw_table:       .quad 0  // pointer to keyword hash table@\\
\mbox{}\verb@@\\
\mbox{}\verb@// Error messages (bare text — diagnostic engine adds file:line:col prefix)@\\
\mbox{}\verb@.section __TEXT,__cstring,cstring_literals@\\
\mbox{}\verb@_kcc_err_unterm_string:  .asciz "unterminated string literal"@\\
\mbox{}\verb@_kcc_err_unterm_char:    .asciz "unterminated character constant"@\\
\mbox{}\verb@_kcc_err_unterm_comment: .asciz "unterminated block comment"@\\
\mbox{}\verb@@{\NWsep}
\end{list}
\vspace{-1.5ex}
\footnotesize
\begin{list}{}{\setlength{\itemsep}{-\parsep}\setlength{\itemindent}{-\leftmargin}}
\item \NWtxtMacroRefIn\ \NWlink{weft?}{?}.

\item{}
\end{list}
\end{minipage}\vspace{4ex}
\end{flushleft}
% ═══════════════════════════════════════════════════════════════════
% Part IV: Driver
% ═══════════════════════════════════════════════════════════════════
\part{The Driver}

% 43-main.weft — Entry Point and Argument Parsing
%
% This is where kcc begins execution.  The _main function receives
% argc and argv from the macOS kernel, initializes the compiler's
% subsystems, and orchestrates the compilation pipeline.
%
% Phase 1 adds --dump-tokens: open a source file, tokenize it,
% and print every token with its source location.

\chapter{The Driver}
\label{ch:driver}

\section{Entry Point}

On macOS, the kernel passes control to \texttt{\_main} with:
\begin{itemize}
\item \reg{x0} = \texttt{argc} (argument count)
\item \reg{x1} = \texttt{argv} (pointer to argument vector)
\end{itemize}

We save these immediately in callee-saved registers, then initialize
the arena and I/O subsystems before doing any real work.

\begin{flushleft} \small
\begin{minipage}{\linewidth}\label{scrap71}\raggedright\small
\NWtarget{weft?}{} $\langle\,${\itshape Main: entry point}\nobreak\ {\footnotesize {?}}$\,\rangle\equiv$
\vspace{-1ex}
\begin{list}{}{} \item
\mbox{}\verb@@\\
\mbox{}\verb@// -- _main ----------------------------------------------------------@\\
\mbox{}\verb@// Entry point for kcc.@\\
\mbox{}\verb@// x0 = argc, x1 = argv@\\
\mbox{}\verb@.text@\\
\mbox{}\verb@.globl _main@\\
\mbox{}\verb@.p2align 2@\\
\mbox{}\verb@_main:@\\
\mbox{}\verb@    stp x29, x30, [sp, #-48]!@\\
\mbox{}\verb@    mov x29, sp@\\
\mbox{}\verb@    stp x19, x20, [sp, #16]@\\
\mbox{}\verb@    stp x21, x22, [sp, #32]@\\
\mbox{}\verb@@\\
\mbox{}\verb@    // Save argc and argv@\\
\mbox{}\verb@    mov x19, x0                   // x19 = argc@\\
\mbox{}\verb@    mov x20, x1                   // x20 = argv@\\
\mbox{}\verb@@\\
\mbox{}\verb@    // Initialize subsystems@\\
\mbox{}\verb@    bl _kcc_arena_init@\\
\mbox{}\verb@    bl _kcc_lex_kw_init@\\
\mbox{}\verb@@\\
\mbox{}\verb@    // Print banner@\\
\mbox{}\verb@    adrp x0, _kcc_banner@{\tt @}\verb@PAGE@\\
\mbox{}\verb@    add  x0, x0, _kcc_banner@{\tt @}\verb@PAGEOFF@\\
\mbox{}\verb@    bl _kcc_io_puts@\\
\mbox{}\verb@@\\
\mbox{}\verb@    // Check for arguments@\\
\mbox{}\verb@    cmp x19, #2@\\
\mbox{}\verb@    b.lt .Lmain_no_args@\\
\mbox{}\verb@@\\
\mbox{}\verb@    // Check if argv[1] is "--dump-tokens"@\\
\mbox{}\verb@    ldr x0, [x20, #8]            // argv[1]@\\
\mbox{}\verb@    adrp x1, _kcc_opt_dump_tokens@{\tt @}\verb@PAGE@\\
\mbox{}\verb@    add  x1, x1, _kcc_opt_dump_tokens@{\tt @}\verb@PAGEOFF@\\
\mbox{}\verb@    bl _kcc_strcmp@\\
\mbox{}\verb@    cbz x0, .Lmain_dump_tokens@\\
\mbox{}\verb@@\\
\mbox{}\verb@    // Not a flag: treat as source file (legacy Phase 0 behavior)@\\
\mbox{}\verb@    adrp x0, _kcc_msg_source@{\tt @}\verb@PAGE@\\
\mbox{}\verb@    add  x0, x0, _kcc_msg_source@{\tt @}\verb@PAGEOFF@\\
\mbox{}\verb@    bl _kcc_io_puts@\\
\mbox{}\verb@    ldr x0, [x20, #8]            // argv[1]@\\
\mbox{}\verb@    bl _kcc_io_puts@\\
\mbox{}\verb@    adrp x0, _kcc_msg_newline@{\tt @}\verb@PAGE@\\
\mbox{}\verb@    add  x0, x0, _kcc_msg_newline@{\tt @}\verb@PAGEOFF@\\
\mbox{}\verb@    bl _kcc_io_puts@\\
\mbox{}\verb@    b .Lmain_exit@\\
\mbox{}\verb@@\\
\mbox{}\verb@.Lmain_dump_tokens:@\\
\mbox{}\verb@    // --dump-tokens requires a source file argument@\\
\mbox{}\verb@    cmp x19, #3@\\
\mbox{}\verb@    b.lt .Lmain_dump_usage@\\
\mbox{}\verb@@\\
\mbox{}\verb@    // Open and mmap the source file@\\
\mbox{}\verb@    ldr x0, [x20, #16]           // argv[2] = source file path@\\
\mbox{}\verb@    bl _kcc_source_open@\\
\mbox{}\verb@@\\
\mbox{}\verb@    // Translation Phase 1: trigraph replacement (§5.1.1.2 para 1)@\\
\mbox{}\verb@    bl _kcc_source_encode@\\
\mbox{}\verb@@\\
\mbox{}\verb@    // Translation Phase 2: line splicing (§5.1.1.2 para 2)@\\
\mbox{}\verb@    bl _kcc_source_splice@\\
\mbox{}\verb@@\\
\mbox{}\verb@    // Initialize lexer with the translated source@\\
\mbox{}\verb@    adrp x0, _kcc_src_base@{\tt @}\verb@PAGE@\\
\mbox{}\verb@    add  x0, x0, _kcc_src_base@{\tt @}\verb@PAGEOFF@\\
\mbox{}\verb@    ldr  x0, [x0]@\\
\mbox{}\verb@@\\
\mbox{}\verb@    adrp x1, _kcc_src_size@{\tt @}\verb@PAGE@\\
\mbox{}\verb@    add  x1, x1, _kcc_src_size@{\tt @}\verb@PAGEOFF@\\
\mbox{}\verb@    ldr  x1, [x1]@\\
\mbox{}\verb@@\\
\mbox{}\verb@    bl _kcc_lex_init@\\
\mbox{}\verb@@\\
\mbox{}\verb@    // Dump all tokens@\\
\mbox{}\verb@    bl _kcc_lex_dump_all@\\
\mbox{}\verb@@\\
\mbox{}\verb@    // Check if any errors occurred during lexing@\\
\mbox{}\verb@    bl _kcc_diag_had_errors@\\
\mbox{}\verb@    cbnz x0, .Lmain_error_exit@\\
\mbox{}\verb@    b .Lmain_exit@\\
\mbox{}\verb@@\\
\mbox{}\verb@.Lmain_dump_usage:@\\
\mbox{}\verb@    adrp x0, _kcc_err_dump_usage@{\tt @}\verb@PAGE@\\
\mbox{}\verb@    add  x0, x0, _kcc_err_dump_usage@{\tt @}\verb@PAGEOFF@\\
\mbox{}\verb@    bl _kcc_io_err@\\
\mbox{}\verb@    mov x0, #1@\\
\mbox{}\verb@    bl _kcc_sys_exit@\\
\mbox{}\verb@@\\
\mbox{}\verb@.Lmain_no_args:@\\
\mbox{}\verb@    adrp x0, _kcc_msg_usage@{\tt @}\verb@PAGE@\\
\mbox{}\verb@    add  x0, x0, _kcc_msg_usage@{\tt @}\verb@PAGEOFF@\\
\mbox{}\verb@    bl _kcc_io_puts@\\
\mbox{}\verb@    b .Lmain_exit@\\
\mbox{}\verb@@\\
\mbox{}\verb@.Lmain_error_exit:@\\
\mbox{}\verb@    bl _kcc_io_flush@\\
\mbox{}\verb@    mov x0, #1@\\
\mbox{}\verb@    bl _kcc_sys_exit@\\
\mbox{}\verb@@\\
\mbox{}\verb@.Lmain_exit:@\\
\mbox{}\verb@    // Flush buffered output@\\
\mbox{}\verb@    bl _kcc_io_flush@\\
\mbox{}\verb@@\\
\mbox{}\verb@    // Exit cleanly@\\
\mbox{}\verb@    mov x0, #0@\\
\mbox{}\verb@    bl _kcc_sys_exit@\\
\mbox{}\verb@@\\
\mbox{}\verb@    // Unreachable@\\
\mbox{}\verb@    brk #1@\\
\mbox{}\verb@@{\NWsep}
\end{list}
\vspace{-1.5ex}
\footnotesize
\begin{list}{}{\setlength{\itemsep}{-\parsep}\setlength{\itemindent}{-\leftmargin}}
\item \NWtxtMacroRefIn\ \NWlink{weft?}{?}.

\item{}
\end{list}
\end{minipage}\vspace{4ex}
\end{flushleft}
\section{String Constants}

The driver's string literals.  These live in the read-only data section.
Each is null-terminated for use with \texttt{kcc\_io\_puts}.

\begin{flushleft} \small
\begin{minipage}{\linewidth}\label{scrap72}\raggedright\small
\NWtarget{weft?}{} $\langle\,${\itshape Main: data}\nobreak\ {\footnotesize {?}}$\,\rangle\equiv$
\vspace{-1ex}
\begin{list}{}{} \item
\mbox{}\verb@@\\
\mbox{}\verb@// -- Driver string constants ----------------------------------------@\\
\mbox{}\verb@.section __TEXT,__cstring,cstring_literals@\\
\mbox{}\verb@_kcc_banner:@\\
\mbox{}\verb@    .asciz "kcc v0.1.0 - a formal C17 compiler for ARM64\n"@\\
\mbox{}\verb@_kcc_msg_source:@\\
\mbox{}\verb@    .asciz "source: "@\\
\mbox{}\verb@_kcc_msg_usage:@\\
\mbox{}\verb@    .asciz "usage: kcc [--dump-tokens] <source.c>\n"@\\
\mbox{}\verb@_kcc_msg_newline:@\\
\mbox{}\verb@    .asciz "\n"@\\
\mbox{}\verb@_kcc_opt_dump_tokens:@\\
\mbox{}\verb@    .asciz "--dump-tokens"@\\
\mbox{}\verb@_kcc_err_dump_usage:@\\
\mbox{}\verb@    .asciz "kcc: error: --dump-tokens requires a source file\n"@\\
\mbox{}\verb@@{\NWsep}
\end{list}
\vspace{-1.5ex}
\footnotesize
\begin{list}{}{\setlength{\itemsep}{-\parsep}\setlength{\itemindent}{-\leftmargin}}
\item \NWtxtMacroRefIn\ \NWlink{weft?}{?}.

\item{}
\end{list}
\end{minipage}\vspace{4ex}
\end{flushleft}
% ═══════════════════════════════════════════════════════════════════
% Part V: Verification
% ═══════════════════════════════════════════════════════════════════
\part{Verification}

% 44-test-harness.weft — Test Runner
%
% A separate test binary (tmp/test_kcc) that exercises the compiler's
% subsystems.  Each test function returns 0 for pass and non-zero for
% failure.  The runner iterates through a registration table, calls
% each test, and reports pass/fail results.
%
% The test runner is linked against all library modules (syscalls,
% arena, strings, io, hashtab, vec, strtab, source, lexer,
% lexer_tables, diag) but uses its own _main instead of the driver's.

\chapter{Test Harness}
\label{ch:test-harness}

\section{Runner Entry Point}

The test runner's \texttt{\_main} initializes the arena and keyword
table, then iterates through the test registration table.  Each
entry contains a test name string and a function pointer.  The table
is terminated by a null name sentinel.

For each test, the runner prints the test name, calls the function,
and prints \texttt{PASS} or \texttt{FAIL} based on the return value.
After all tests, a summary line reports the pass/fail counts and the
process exits with status~0 (all passed) or~1 (any failed).

\begin{flushleft} \small
\begin{minipage}{\linewidth}\label{scrap73}\raggedright\small
\NWtarget{weft?}{} $\langle\,${\itshape Test: runner entry}\nobreak\ {\footnotesize {?}}$\,\rangle\equiv$
\vspace{-1ex}
\begin{list}{}{} \item
\mbox{}\verb@@\\
\mbox{}\verb@// -- test_kcc _main -----------------------------------------------------@\\
\mbox{}\verb@// Entry point for the test binary.@\\
\mbox{}\verb@.text@\\
\mbox{}\verb@.globl _main@\\
\mbox{}\verb@.p2align 2@\\
\mbox{}\verb@_main:@\\
\mbox{}\verb@    stp x29, x30, [sp, #-48]!@\\
\mbox{}\verb@    mov x29, sp@\\
\mbox{}\verb@    stp x19, x20, [sp, #16]@\\
\mbox{}\verb@    stp x21, x22, [sp, #32]@\\
\mbox{}\verb@@\\
\mbox{}\verb@    // Initialize subsystems@\\
\mbox{}\verb@    bl _kcc_arena_init@\\
\mbox{}\verb@    bl _kcc_lex_kw_init@\\
\mbox{}\verb@@\\
\mbox{}\verb@    // Print header@\\
\mbox{}\verb@    adrp x0, _test_header@{\tt @}\verb@PAGE@\\
\mbox{}\verb@    add  x0, x0, _test_header@{\tt @}\verb@PAGEOFF@\\
\mbox{}\verb@    bl _kcc_io_puts@\\
\mbox{}\verb@    bl _kcc_io_flush@\\
\mbox{}\verb@@\\
\mbox{}\verb@    mov x19, #0                   // x19 = pass count@\\
\mbox{}\verb@    mov x20, #0                   // x20 = fail count@\\
\mbox{}\verb@@\\
\mbox{}\verb@    // Load test table pointer@\\
\mbox{}\verb@    adrp x21, _kcc_test_table@{\tt @}\verb@PAGE@\\
\mbox{}\verb@    add  x21, x21, _kcc_test_table@{\tt @}\verb@PAGEOFF@\\
\mbox{}\verb@@\\
\mbox{}\verb@.Ltest_loop:@\\
\mbox{}\verb@    ldr x0, [x21]                // test name@\\
\mbox{}\verb@    cbz x0, .Ltest_summary       // null sentinel → done@\\
\mbox{}\verb@    ldr x22, [x21, #8]           // test function pointer@\\
\mbox{}\verb@@\\
\mbox{}\verb@    // Print "  name... "@\\
\mbox{}\verb@    adrp x0, _test_prefix@{\tt @}\verb@PAGE@\\
\mbox{}\verb@    add  x0, x0, _test_prefix@{\tt @}\verb@PAGEOFF@\\
\mbox{}\verb@    bl _kcc_io_puts@\\
\mbox{}\verb@    ldr x0, [x21]                // test name@\\
\mbox{}\verb@    bl _kcc_io_puts@\\
\mbox{}\verb@    adrp x0, _test_dots@{\tt @}\verb@PAGE@\\
\mbox{}\verb@    add  x0, x0, _test_dots@{\tt @}\verb@PAGEOFF@\\
\mbox{}\verb@    bl _kcc_io_puts@\\
\mbox{}\verb@    bl _kcc_io_flush@\\
\mbox{}\verb@@\\
\mbox{}\verb@    // Reset diagnostic counters before each test@\\
\mbox{}\verb@    bl _kcc_diag_reset@\\
\mbox{}\verb@@\\
\mbox{}\verb@    // Call test function@\\
\mbox{}\verb@    blr x22@\\
\mbox{}\verb@@\\
\mbox{}\verb@    // Check result: 0 = pass, non-zero = fail@\\
\mbox{}\verb@    cbnz x0, .Ltest_fail@\\
\mbox{}\verb@@\\
\mbox{}\verb@    // PASS@\\
\mbox{}\verb@    adrp x0, _test_pass@{\tt @}\verb@PAGE@\\
\mbox{}\verb@    add  x0, x0, _test_pass@{\tt @}\verb@PAGEOFF@\\
\mbox{}\verb@    bl _kcc_io_puts@\\
\mbox{}\verb@    bl _kcc_io_flush@\\
\mbox{}\verb@    add x19, x19, #1@\\
\mbox{}\verb@    b .Ltest_next@\\
\mbox{}\verb@@\\
\mbox{}\verb@.Ltest_fail:@\\
\mbox{}\verb@    adrp x0, _test_fail@{\tt @}\verb@PAGE@\\
\mbox{}\verb@    add  x0, x0, _test_fail@{\tt @}\verb@PAGEOFF@\\
\mbox{}\verb@    bl _kcc_io_puts@\\
\mbox{}\verb@    bl _kcc_io_flush@\\
\mbox{}\verb@    add x20, x20, #1@\\
\mbox{}\verb@@\\
\mbox{}\verb@.Ltest_next:@\\
\mbox{}\verb@    add x21, x21, #16            // advance to next table entry@\\
\mbox{}\verb@    b .Ltest_loop@\\
\mbox{}\verb@@\\
\mbox{}\verb@.Ltest_summary:@\\
\mbox{}\verb@    // Print summary: "\nResults: N passed, M failed\n"@\\
\mbox{}\verb@    adrp x0, _test_results@{\tt @}\verb@PAGE@\\
\mbox{}\verb@    add  x0, x0, _test_results@{\tt @}\verb@PAGEOFF@\\
\mbox{}\verb@    bl _kcc_io_puts@\\
\mbox{}\verb@    mov x0, x19@\\
\mbox{}\verb@    bl _kcc_io_putu@\\
\mbox{}\verb@    adrp x0, _test_passed_str@{\tt @}\verb@PAGE@\\
\mbox{}\verb@    add  x0, x0, _test_passed_str@{\tt @}\verb@PAGEOFF@\\
\mbox{}\verb@    bl _kcc_io_puts@\\
\mbox{}\verb@    mov x0, x20@\\
\mbox{}\verb@    bl _kcc_io_putu@\\
\mbox{}\verb@    adrp x0, _test_failed_str@{\tt @}\verb@PAGE@\\
\mbox{}\verb@    add  x0, x0, _test_failed_str@{\tt @}\verb@PAGEOFF@\\
\mbox{}\verb@    bl _kcc_io_puts@\\
\mbox{}\verb@@\\
\mbox{}\verb@    bl _kcc_io_flush@\\
\mbox{}\verb@@\\
\mbox{}\verb@    // Exit: 0 if all passed, 1 if any failed@\\
\mbox{}\verb@    cmp x20, #0@\\
\mbox{}\verb@    b.ne .Ltest_exit_fail@\\
\mbox{}\verb@    mov x0, #0@\\
\mbox{}\verb@    bl _kcc_sys_exit@\\
\mbox{}\verb@.Ltest_exit_fail:@\\
\mbox{}\verb@    mov x0, #1@\\
\mbox{}\verb@    bl _kcc_sys_exit@\\
\mbox{}\verb@@{\NWsep}
\end{list}
\vspace{-1.5ex}
\footnotesize
\begin{list}{}{\setlength{\itemsep}{-\parsep}\setlength{\itemindent}{-\leftmargin}}
\item \NWtxtMacroRefIn\ \NWlink{weft?}{?}.

\item{}
\end{list}
\end{minipage}\vspace{4ex}
\end{flushleft}
\section{Test Assertion Helpers}

Two helper functions simplify test writing:

\texttt{kcc\_test\_lex\_setup} initializes the lexer with a test
source string, applying translation phases~1 and~2.  It also sets
\texttt{\_kcc\_src\_path} to \texttt{<test>} for diagnostic output.

\texttt{kcc\_test\_expect\_tok} calls \texttt{lex\_next} and checks
whether the resulting token type matches the expected value.  Returns
0 on match, 1 on mismatch.

\begin{flushleft} \small
\begin{minipage}{\linewidth}\label{scrap74}\raggedright\small
\NWtarget{weft?}{} $\langle\,${\itshape Test: assert helpers}\nobreak\ {\footnotesize {?}}$\,\rangle\equiv$
\vspace{-1ex}
\begin{list}{}{} \item
\mbox{}\verb@@\\
\mbox{}\verb@// -- kcc_test_lex_setup -------------------------------------------------@\\
\mbox{}\verb@// void kcc_test_lex_setup(const char *src, size_t len)@\\
\mbox{}\verb@//   x0 = pointer to test source string@\\
\mbox{}\verb@//   x1 = length in bytes (excluding null terminator)@\\
\mbox{}\verb@// Sets up the lexer with translation phases 1 and 2.@\\
\mbox{}\verb@.globl _kcc_test_lex_setup@\\
\mbox{}\verb@.p2align 2@\\
\mbox{}\verb@_kcc_test_lex_setup:@\\
\mbox{}\verb@    stp x29, x30, [sp, #-32]!@\\
\mbox{}\verb@    mov x29, sp@\\
\mbox{}\verb@    stp x19, x20, [sp, #16]@\\
\mbox{}\verb@@\\
\mbox{}\verb@    mov x19, x0                   // src pointer@\\
\mbox{}\verb@    mov x20, x1                   // src length@\\
\mbox{}\verb@@\\
\mbox{}\verb@    // Set source path to "<test>"@\\
\mbox{}\verb@    adrp x2, _kcc_test_path@{\tt @}\verb@PAGE@\\
\mbox{}\verb@    add  x2, x2, _kcc_test_path@{\tt @}\verb@PAGEOFF@\\
\mbox{}\verb@    adrp x3, _kcc_src_path@{\tt @}\verb@PAGE@\\
\mbox{}\verb@    add  x3, x3, _kcc_src_path@{\tt @}\verb@PAGEOFF@\\
\mbox{}\verb@    str  x2, [x3]@\\
\mbox{}\verb@@\\
\mbox{}\verb@    // Set src_base and src_size for translation phases@\\
\mbox{}\verb@    adrp x3, _kcc_src_base@{\tt @}\verb@PAGE@\\
\mbox{}\verb@    add  x3, x3, _kcc_src_base@{\tt @}\verb@PAGEOFF@\\
\mbox{}\verb@    str  x19, [x3]@\\
\mbox{}\verb@@\\
\mbox{}\verb@    adrp x3, _kcc_src_size@{\tt @}\verb@PAGE@\\
\mbox{}\verb@    add  x3, x3, _kcc_src_size@{\tt @}\verb@PAGEOFF@\\
\mbox{}\verb@    str  x20, [x3]@\\
\mbox{}\verb@@\\
\mbox{}\verb@    // Apply translation phase 1 (trigraphs → arena copy)@\\
\mbox{}\verb@    bl _kcc_source_encode@\\
\mbox{}\verb@@\\
\mbox{}\verb@    // Apply translation phase 2 (line splicing, in-place)@\\
\mbox{}\verb@    bl _kcc_source_splice@\\
\mbox{}\verb@@\\
\mbox{}\verb@    // Initialize lexer with the translated buffer@\\
\mbox{}\verb@    adrp x0, _kcc_src_base@{\tt @}\verb@PAGE@\\
\mbox{}\verb@    add  x0, x0, _kcc_src_base@{\tt @}\verb@PAGEOFF@\\
\mbox{}\verb@    ldr  x0, [x0]@\\
\mbox{}\verb@    adrp x1, _kcc_src_size@{\tt @}\verb@PAGE@\\
\mbox{}\verb@    add  x1, x1, _kcc_src_size@{\tt @}\verb@PAGEOFF@\\
\mbox{}\verb@    ldr  x1, [x1]@\\
\mbox{}\verb@    bl _kcc_lex_init@\\
\mbox{}\verb@@\\
\mbox{}\verb@    ldp x19, x20, [sp, #16]@\\
\mbox{}\verb@    ldp x29, x30, [sp], #32@\\
\mbox{}\verb@    ret@\\
\mbox{}\verb@@\\
\mbox{}\verb@// -- kcc_test_expect_tok ------------------------------------------------@\\
\mbox{}\verb@// int kcc_test_expect_tok(int expected_type)@\\
\mbox{}\verb@//   x0 = expected token type code@\\
\mbox{}\verb@// Calls lex_next, checks tok_type against expected.@\\
\mbox{}\verb@// Returns 0 if match, 1 if mismatch.@\\
\mbox{}\verb@.globl _kcc_test_expect_tok@\\
\mbox{}\verb@.p2align 2@\\
\mbox{}\verb@_kcc_test_expect_tok:@\\
\mbox{}\verb@    stp x29, x30, [sp, #-32]!@\\
\mbox{}\verb@    mov x29, sp@\\
\mbox{}\verb@    str x19, [sp, #16]@\\
\mbox{}\verb@@\\
\mbox{}\verb@    mov x19, x0                   // save expected type@\\
\mbox{}\verb@@\\
\mbox{}\verb@    bl _kcc_lex_next@\\
\mbox{}\verb@@\\
\mbox{}\verb@    adrp x0, _kcc_tok_type@{\tt @}\verb@PAGE@\\
\mbox{}\verb@    add  x0, x0, _kcc_tok_type@{\tt @}\verb@PAGEOFF@\\
\mbox{}\verb@    ldr  x0, [x0]@\\
\mbox{}\verb@@\\
\mbox{}\verb@    cmp x0, x19@\\
\mbox{}\verb@    b.eq .Lexpect_match@\\
\mbox{}\verb@    mov x0, #1                   // FAIL@\\
\mbox{}\verb@    b .Lexpect_done@\\
\mbox{}\verb@.Lexpect_match:@\\
\mbox{}\verb@    mov x0, #0                   // PASS@\\
\mbox{}\verb@.Lexpect_done:@\\
\mbox{}\verb@    ldr x19, [sp, #16]@\\
\mbox{}\verb@    ldp x29, x30, [sp], #32@\\
\mbox{}\verb@    ret@\\
\mbox{}\verb@@{\NWsep}
\end{list}
\vspace{-1.5ex}
\footnotesize
\begin{list}{}{\setlength{\itemsep}{-\parsep}\setlength{\itemindent}{-\leftmargin}}
\item \NWtxtMacroRefIn\ \NWlink{weft?}{?}.

\item{}
\end{list}
\end{minipage}\vspace{4ex}
\end{flushleft}
\section{Runner Data}

String constants for the test runner's output formatting.

\begin{flushleft} \small
\begin{minipage}{\linewidth}\label{scrap75}\raggedright\small
\NWtarget{weft?}{} $\langle\,${\itshape Test: runner data}\nobreak\ {\footnotesize {?}}$\,\rangle\equiv$
\vspace{-1ex}
\begin{list}{}{} \item
\mbox{}\verb@@\\
\mbox{}\verb@// -- Test runner string constants ---------------------------------------@\\
\mbox{}\verb@.section __TEXT,__cstring,cstring_literals@\\
\mbox{}\verb@_test_header:     .asciz "kcc test suite\n==============\n"@\\
\mbox{}\verb@_test_prefix:     .asciz "  "@\\
\mbox{}\verb@_test_dots:       .asciz "... "@\\
\mbox{}\verb@_test_pass:       .asciz "PASS\n"@\\
\mbox{}\verb@_test_fail:       .asciz "FAIL\n"@\\
\mbox{}\verb@_test_results:    .asciz "\nResults: "@\\
\mbox{}\verb@_test_passed_str: .asciz " passed, "@\\
\mbox{}\verb@_test_failed_str: .asciz " failed\n"@\\
\mbox{}\verb@_kcc_test_path:   .asciz "<test>"@\\
\mbox{}\verb@@{\NWsep}
\end{list}
\vspace{-1.5ex}
\footnotesize
\begin{list}{}{\setlength{\itemsep}{-\parsep}\setlength{\itemindent}{-\leftmargin}}
\item \NWtxtMacroRefIn\ \NWlink{weft?}{?}.

\item{}
\end{list}
\end{minipage}\vspace{4ex}
\end{flushleft}
% 45-conformance-tests.weft — Lexer Conformance Tests
%
% Each test function sets up a small source string, runs the lexer,
% and verifies the resulting token stream.  Tests return 0 (pass)
% or 1 (fail).  All tests are registered in a global table read by
% the test runner (Chapter~\ref{ch:test-harness}).

\chapter{Conformance Tests}
\label{ch:conformance-tests}

\section{Keyword Recognition}

Tests that representative keywords from all categories are correctly
identified, and that identifiers resembling keyword prefixes are
\emph{not} misidentified.

\begin{flushleft} \small
\begin{minipage}{\linewidth}\label{scrap76}\raggedright\small
\NWtarget{weft?}{} $\langle\,${\itshape Test: lex keywords}\nobreak\ {\footnotesize {?}}$\,\rangle\equiv$
\vspace{-1ex}
\begin{list}{}{} \item
\mbox{}\verb@@\\
\mbox{}\verb@// -- test_lex_keywords --------------------------------------------------@\\
\mbox{}\verb@// Source: "auto int return _Static_assert foo"@\\
\mbox{}\verb@// Expected: AUTO, INT, RETURN, _STATIC_ASSERT, IDENT, EOF@\\
\mbox{}\verb@.text@\\
\mbox{}\verb@.globl _test_lex_keywords@\\
\mbox{}\verb@.p2align 2@\\
\mbox{}\verb@_test_lex_keywords:@\\
\mbox{}\verb@    stp x29, x30, [sp, #-16]!@\\
\mbox{}\verb@    mov x29, sp@\\
\mbox{}\verb@@\\
\mbox{}\verb@    adrp x0, _test_src_kw@{\tt @}\verb@PAGE@\\
\mbox{}\verb@    add  x0, x0, _test_src_kw@{\tt @}\verb@PAGEOFF@\\
\mbox{}\verb@    mov  x1, #TEST_SRC_KW_LEN@\\
\mbox{}\verb@    bl _kcc_test_lex_setup@\\
\mbox{}\verb@@\\
\mbox{}\verb@    mov x0, #TOK_AUTO@\\
\mbox{}\verb@    bl _kcc_test_expect_tok@\\
\mbox{}\verb@    cbnz x0, .Ltkw_fail@\\
\mbox{}\verb@@\\
\mbox{}\verb@    mov x0, #TOK_INT@\\
\mbox{}\verb@    bl _kcc_test_expect_tok@\\
\mbox{}\verb@    cbnz x0, .Ltkw_fail@\\
\mbox{}\verb@@\\
\mbox{}\verb@    mov x0, #TOK_RETURN@\\
\mbox{}\verb@    bl _kcc_test_expect_tok@\\
\mbox{}\verb@    cbnz x0, .Ltkw_fail@\\
\mbox{}\verb@@\\
\mbox{}\verb@    mov x0, #TOK__STATIC_ASSERT@\\
\mbox{}\verb@    bl _kcc_test_expect_tok@\\
\mbox{}\verb@    cbnz x0, .Ltkw_fail@\\
\mbox{}\verb@@\\
\mbox{}\verb@    mov x0, #TOK_IDENT@\\
\mbox{}\verb@    bl _kcc_test_expect_tok@\\
\mbox{}\verb@    cbnz x0, .Ltkw_fail@\\
\mbox{}\verb@@\\
\mbox{}\verb@    mov x0, #TOK_EOF@\\
\mbox{}\verb@    bl _kcc_test_expect_tok@\\
\mbox{}\verb@    cbnz x0, .Ltkw_fail@\\
\mbox{}\verb@@\\
\mbox{}\verb@    mov x0, #0@\\
\mbox{}\verb@    ldp x29, x30, [sp], #16@\\
\mbox{}\verb@    ret@\\
\mbox{}\verb@.Ltkw_fail:@\\
\mbox{}\verb@    mov x0, #1@\\
\mbox{}\verb@    ldp x29, x30, [sp], #16@\\
\mbox{}\verb@    ret@\\
\mbox{}\verb@@{\NWsep}
\end{list}
\vspace{-1.5ex}
\footnotesize
\begin{list}{}{\setlength{\itemsep}{-\parsep}\setlength{\itemindent}{-\leftmargin}}
\item \NWtxtMacroRefIn\ \NWlink{weft?}{?}.

\item{}
\end{list}
\end{minipage}\vspace{4ex}
\end{flushleft}
\section{Numeric Literals}

Tests that decimal, hexadecimal, floating-point, and scientific
notation numbers are all classified as \texttt{TOK\_NUMBER} (pp-number).

\begin{flushleft} \small
\begin{minipage}{\linewidth}\label{scrap77}\raggedright\small
\NWtarget{weft?}{} $\langle\,${\itshape Test: lex numbers}\nobreak\ {\footnotesize {?}}$\,\rangle\equiv$
\vspace{-1ex}
\begin{list}{}{} \item
\mbox{}\verb@@\\
\mbox{}\verb@// -- test_lex_numbers ---------------------------------------------------@\\
\mbox{}\verb@// Source: "42 0xFF 3.14 1e10 0x1p2"@\\
\mbox{}\verb@// Expected: NUMBER, NUMBER, NUMBER, NUMBER, NUMBER, EOF@\\
\mbox{}\verb@.text@\\
\mbox{}\verb@.globl _test_lex_numbers@\\
\mbox{}\verb@.p2align 2@\\
\mbox{}\verb@_test_lex_numbers:@\\
\mbox{}\verb@    stp x29, x30, [sp, #-16]!@\\
\mbox{}\verb@    mov x29, sp@\\
\mbox{}\verb@@\\
\mbox{}\verb@    adrp x0, _test_src_num@{\tt @}\verb@PAGE@\\
\mbox{}\verb@    add  x0, x0, _test_src_num@{\tt @}\verb@PAGEOFF@\\
\mbox{}\verb@    mov  x1, #TEST_SRC_NUM_LEN@\\
\mbox{}\verb@    bl _kcc_test_lex_setup@\\
\mbox{}\verb@@\\
\mbox{}\verb@    mov x0, #TOK_NUMBER@\\
\mbox{}\verb@    bl _kcc_test_expect_tok@\\
\mbox{}\verb@    cbnz x0, .Ltnum_fail@\\
\mbox{}\verb@@\\
\mbox{}\verb@    mov x0, #TOK_NUMBER@\\
\mbox{}\verb@    bl _kcc_test_expect_tok@\\
\mbox{}\verb@    cbnz x0, .Ltnum_fail@\\
\mbox{}\verb@@\\
\mbox{}\verb@    mov x0, #TOK_NUMBER@\\
\mbox{}\verb@    bl _kcc_test_expect_tok@\\
\mbox{}\verb@    cbnz x0, .Ltnum_fail@\\
\mbox{}\verb@@\\
\mbox{}\verb@    mov x0, #TOK_NUMBER@\\
\mbox{}\verb@    bl _kcc_test_expect_tok@\\
\mbox{}\verb@    cbnz x0, .Ltnum_fail@\\
\mbox{}\verb@@\\
\mbox{}\verb@    mov x0, #TOK_NUMBER@\\
\mbox{}\verb@    bl _kcc_test_expect_tok@\\
\mbox{}\verb@    cbnz x0, .Ltnum_fail@\\
\mbox{}\verb@@\\
\mbox{}\verb@    mov x0, #TOK_EOF@\\
\mbox{}\verb@    bl _kcc_test_expect_tok@\\
\mbox{}\verb@    cbnz x0, .Ltnum_fail@\\
\mbox{}\verb@@\\
\mbox{}\verb@    mov x0, #0@\\
\mbox{}\verb@    ldp x29, x30, [sp], #16@\\
\mbox{}\verb@    ret@\\
\mbox{}\verb@.Ltnum_fail:@\\
\mbox{}\verb@    mov x0, #1@\\
\mbox{}\verb@    ldp x29, x30, [sp], #16@\\
\mbox{}\verb@    ret@\\
\mbox{}\verb@@{\NWsep}
\end{list}
\vspace{-1.5ex}
\footnotesize
\begin{list}{}{\setlength{\itemsep}{-\parsep}\setlength{\itemindent}{-\leftmargin}}
\item \NWtxtMacroRefIn\ \NWlink{weft?}{?}.

\item{}
\end{list}
\end{minipage}\vspace{4ex}
\end{flushleft}
\section{String and Character Literals}

Tests that string literals and character constants are correctly
classified.

\begin{flushleft} \small
\begin{minipage}{\linewidth}\label{scrap78}\raggedright\small
\NWtarget{weft?}{} $\langle\,${\itshape Test: lex strings}\nobreak\ {\footnotesize {?}}$\,\rangle\equiv$
\vspace{-1ex}
\begin{list}{}{} \item
\mbox{}\verb@@\\
\mbox{}\verb@// -- test_lex_strings ---------------------------------------------------@\\
\mbox{}\verb@// Source: "hello" 'x'@\\
\mbox{}\verb@// (The outer quotes in the .asciz are assembler string delimiters;@\\
\mbox{}\verb@//  the inner \" produce the actual " character in the test source.)@\\
\mbox{}\verb@// Expected: STRING_LIT, CHAR_LIT, EOF@\\
\mbox{}\verb@.text@\\
\mbox{}\verb@.globl _test_lex_strings@\\
\mbox{}\verb@.p2align 2@\\
\mbox{}\verb@_test_lex_strings:@\\
\mbox{}\verb@    stp x29, x30, [sp, #-16]!@\\
\mbox{}\verb@    mov x29, sp@\\
\mbox{}\verb@@\\
\mbox{}\verb@    adrp x0, _test_src_str@{\tt @}\verb@PAGE@\\
\mbox{}\verb@    add  x0, x0, _test_src_str@{\tt @}\verb@PAGEOFF@\\
\mbox{}\verb@    mov  x1, #TEST_SRC_STR_LEN@\\
\mbox{}\verb@    bl _kcc_test_lex_setup@\\
\mbox{}\verb@@\\
\mbox{}\verb@    mov x0, #TOK_STRING_LIT@\\
\mbox{}\verb@    bl _kcc_test_expect_tok@\\
\mbox{}\verb@    cbnz x0, .Ltstr_fail@\\
\mbox{}\verb@@\\
\mbox{}\verb@    mov x0, #TOK_CHAR_LIT@\\
\mbox{}\verb@    bl _kcc_test_expect_tok@\\
\mbox{}\verb@    cbnz x0, .Ltstr_fail@\\
\mbox{}\verb@@\\
\mbox{}\verb@    mov x0, #TOK_EOF@\\
\mbox{}\verb@    bl _kcc_test_expect_tok@\\
\mbox{}\verb@    cbnz x0, .Ltstr_fail@\\
\mbox{}\verb@@\\
\mbox{}\verb@    mov x0, #0@\\
\mbox{}\verb@    ldp x29, x30, [sp], #16@\\
\mbox{}\verb@    ret@\\
\mbox{}\verb@.Ltstr_fail:@\\
\mbox{}\verb@    mov x0, #1@\\
\mbox{}\verb@    ldp x29, x30, [sp], #16@\\
\mbox{}\verb@    ret@\\
\mbox{}\verb@@{\NWsep}
\end{list}
\vspace{-1.5ex}
\footnotesize
\begin{list}{}{\setlength{\itemsep}{-\parsep}\setlength{\itemindent}{-\leftmargin}}
\item \NWtxtMacroRefIn\ \NWlink{weft?}{?}.

\item{}
\end{list}
\end{minipage}\vspace{4ex}
\end{flushleft}
\section{Punctuators}

Tests that single-character and multi-character punctuators are
correctly recognized.

\begin{flushleft} \small
\begin{minipage}{\linewidth}\label{scrap79}\raggedright\small
\NWtarget{weft?}{} $\langle\,${\itshape Test: lex punctuators}\nobreak\ {\footnotesize {?}}$\,\rangle\equiv$
\vspace{-1ex}
\begin{list}{}{} \item
\mbox{}\verb@@\\
\mbox{}\verb@// -- test_lex_punctuators -----------------------------------------------@\\
\mbox{}\verb@// Source: "(){};,+-*/"@\\
\mbox{}\verb@// Expected: LPAREN RPAREN LBRACE RBRACE SEMI COMMA PLUS MINUS STAR SLASH EOF@\\
\mbox{}\verb@.text@\\
\mbox{}\verb@.globl _test_lex_punctuators@\\
\mbox{}\verb@.p2align 2@\\
\mbox{}\verb@_test_lex_punctuators:@\\
\mbox{}\verb@    stp x29, x30, [sp, #-16]!@\\
\mbox{}\verb@    mov x29, sp@\\
\mbox{}\verb@@\\
\mbox{}\verb@    adrp x0, _test_src_punct@{\tt @}\verb@PAGE@\\
\mbox{}\verb@    add  x0, x0, _test_src_punct@{\tt @}\verb@PAGEOFF@\\
\mbox{}\verb@    mov  x1, #TEST_SRC_PUNCT_LEN@\\
\mbox{}\verb@    bl _kcc_test_lex_setup@\\
\mbox{}\verb@@\\
\mbox{}\verb@    mov x0, #TOK_LPAREN@\\
\mbox{}\verb@    bl _kcc_test_expect_tok@\\
\mbox{}\verb@    cbnz x0, .Ltpunct_fail@\\
\mbox{}\verb@@\\
\mbox{}\verb@    mov x0, #TOK_RPAREN@\\
\mbox{}\verb@    bl _kcc_test_expect_tok@\\
\mbox{}\verb@    cbnz x0, .Ltpunct_fail@\\
\mbox{}\verb@@\\
\mbox{}\verb@    mov x0, #TOK_LBRACE@\\
\mbox{}\verb@    bl _kcc_test_expect_tok@\\
\mbox{}\verb@    cbnz x0, .Ltpunct_fail@\\
\mbox{}\verb@@\\
\mbox{}\verb@    mov x0, #TOK_RBRACE@\\
\mbox{}\verb@    bl _kcc_test_expect_tok@\\
\mbox{}\verb@    cbnz x0, .Ltpunct_fail@\\
\mbox{}\verb@@\\
\mbox{}\verb@    mov x0, #TOK_SEMI@\\
\mbox{}\verb@    bl _kcc_test_expect_tok@\\
\mbox{}\verb@    cbnz x0, .Ltpunct_fail@\\
\mbox{}\verb@@\\
\mbox{}\verb@    mov x0, #TOK_COMMA@\\
\mbox{}\verb@    bl _kcc_test_expect_tok@\\
\mbox{}\verb@    cbnz x0, .Ltpunct_fail@\\
\mbox{}\verb@@\\
\mbox{}\verb@    mov x0, #TOK_PLUS@\\
\mbox{}\verb@    bl _kcc_test_expect_tok@\\
\mbox{}\verb@    cbnz x0, .Ltpunct_fail@\\
\mbox{}\verb@@\\
\mbox{}\verb@    mov x0, #TOK_MINUS@\\
\mbox{}\verb@    bl _kcc_test_expect_tok@\\
\mbox{}\verb@    cbnz x0, .Ltpunct_fail@\\
\mbox{}\verb@@\\
\mbox{}\verb@    mov x0, #TOK_STAR@\\
\mbox{}\verb@    bl _kcc_test_expect_tok@\\
\mbox{}\verb@    cbnz x0, .Ltpunct_fail@\\
\mbox{}\verb@@\\
\mbox{}\verb@    mov x0, #TOK_SLASH@\\
\mbox{}\verb@    bl _kcc_test_expect_tok@\\
\mbox{}\verb@    cbnz x0, .Ltpunct_fail@\\
\mbox{}\verb@@\\
\mbox{}\verb@    mov x0, #TOK_EOF@\\
\mbox{}\verb@    bl _kcc_test_expect_tok@\\
\mbox{}\verb@    cbnz x0, .Ltpunct_fail@\\
\mbox{}\verb@@\\
\mbox{}\verb@    mov x0, #0@\\
\mbox{}\verb@    ldp x29, x30, [sp], #16@\\
\mbox{}\verb@    ret@\\
\mbox{}\verb@.Ltpunct_fail:@\\
\mbox{}\verb@    mov x0, #1@\\
\mbox{}\verb@    ldp x29, x30, [sp], #16@\\
\mbox{}\verb@    ret@\\
\mbox{}\verb@@{\NWsep}
\end{list}
\vspace{-1.5ex}
\footnotesize
\begin{list}{}{\setlength{\itemsep}{-\parsep}\setlength{\itemindent}{-\leftmargin}}
\item \NWtxtMacroRefIn\ \NWlink{weft?}{?}.

\item{}
\end{list}
\end{minipage}\vspace{4ex}
\end{flushleft}
\section{Comments}

Tests that both line comments (\texttt{//}) and block comments
(\texttt{/* */}) are correctly skipped by the lexer.

\begin{flushleft} \small
\begin{minipage}{\linewidth}\label{scrap80}\raggedright\small
\NWtarget{weft?}{} $\langle\,${\itshape Test: lex comments}\nobreak\ {\footnotesize {?}}$\,\rangle\equiv$
\vspace{-1ex}
\begin{list}{}{} \item
\mbox{}\verb@@\\
\mbox{}\verb@// -- test_lex_comments --------------------------------------------------@\\
\mbox{}\verb@// Source: "x /* comment */ y // line\nz"@\\
\mbox{}\verb@// Expected: IDENT(x) IDENT(y) IDENT(z) EOF@\\
\mbox{}\verb@.text@\\
\mbox{}\verb@.globl _test_lex_comments@\\
\mbox{}\verb@.p2align 2@\\
\mbox{}\verb@_test_lex_comments:@\\
\mbox{}\verb@    stp x29, x30, [sp, #-16]!@\\
\mbox{}\verb@    mov x29, sp@\\
\mbox{}\verb@@\\
\mbox{}\verb@    adrp x0, _test_src_comments@{\tt @}\verb@PAGE@\\
\mbox{}\verb@    add  x0, x0, _test_src_comments@{\tt @}\verb@PAGEOFF@\\
\mbox{}\verb@    mov  x1, #TEST_SRC_COMMENTS_LEN@\\
\mbox{}\verb@    bl _kcc_test_lex_setup@\\
\mbox{}\verb@@\\
\mbox{}\verb@    mov x0, #TOK_IDENT@\\
\mbox{}\verb@    bl _kcc_test_expect_tok@\\
\mbox{}\verb@    cbnz x0, .Ltcmt_fail@\\
\mbox{}\verb@@\\
\mbox{}\verb@    mov x0, #TOK_IDENT@\\
\mbox{}\verb@    bl _kcc_test_expect_tok@\\
\mbox{}\verb@    cbnz x0, .Ltcmt_fail@\\
\mbox{}\verb@@\\
\mbox{}\verb@    mov x0, #TOK_IDENT@\\
\mbox{}\verb@    bl _kcc_test_expect_tok@\\
\mbox{}\verb@    cbnz x0, .Ltcmt_fail@\\
\mbox{}\verb@@\\
\mbox{}\verb@    mov x0, #TOK_EOF@\\
\mbox{}\verb@    bl _kcc_test_expect_tok@\\
\mbox{}\verb@    cbnz x0, .Ltcmt_fail@\\
\mbox{}\verb@@\\
\mbox{}\verb@    mov x0, #0@\\
\mbox{}\verb@    ldp x29, x30, [sp], #16@\\
\mbox{}\verb@    ret@\\
\mbox{}\verb@.Ltcmt_fail:@\\
\mbox{}\verb@    mov x0, #1@\\
\mbox{}\verb@    ldp x29, x30, [sp], #16@\\
\mbox{}\verb@    ret@\\
\mbox{}\verb@@{\NWsep}
\end{list}
\vspace{-1.5ex}
\footnotesize
\begin{list}{}{\setlength{\itemsep}{-\parsep}\setlength{\itemindent}{-\leftmargin}}
\item \NWtxtMacroRefIn\ \NWlink{weft?}{?}.

\item{}
\end{list}
\end{minipage}\vspace{4ex}
\end{flushleft}
\section{Trigraph Replacement}

Tests that trigraph sequences are correctly replaced before
tokenization.  The source \texttt{??=} becomes \texttt{\#},
which the lexer should tokenize as \texttt{TOK\_HASH}.

\begin{flushleft} \small
\begin{minipage}{\linewidth}\label{scrap81}\raggedright\small
\NWtarget{weft?}{} $\langle\,${\itshape Test: lex trigraphs}\nobreak\ {\footnotesize {?}}$\,\rangle\equiv$
\vspace{-1ex}
\begin{list}{}{} \item
\mbox{}\verb@@\\
\mbox{}\verb@// -- test_lex_trigraphs -------------------------------------------------@\\
\mbox{}\verb@// Source: "??= x ??( ??) ??< ??>"@\\
\mbox{}\verb@// After trigraph replacement: "# x [ ] { }"@\\
\mbox{}\verb@// Expected: HASH IDENT LBRACKET RBRACKET LBRACE RBRACE EOF@\\
\mbox{}\verb@.text@\\
\mbox{}\verb@.globl _test_lex_trigraphs@\\
\mbox{}\verb@.p2align 2@\\
\mbox{}\verb@_test_lex_trigraphs:@\\
\mbox{}\verb@    stp x29, x30, [sp, #-16]!@\\
\mbox{}\verb@    mov x29, sp@\\
\mbox{}\verb@@\\
\mbox{}\verb@    adrp x0, _test_src_trigraph@{\tt @}\verb@PAGE@\\
\mbox{}\verb@    add  x0, x0, _test_src_trigraph@{\tt @}\verb@PAGEOFF@\\
\mbox{}\verb@    mov  x1, #TEST_SRC_TRIGRAPH_LEN@\\
\mbox{}\verb@    bl _kcc_test_lex_setup@\\
\mbox{}\verb@@\\
\mbox{}\verb@    mov x0, #TOK_HASH@\\
\mbox{}\verb@    bl _kcc_test_expect_tok@\\
\mbox{}\verb@    cbnz x0, .Lttrig_fail@\\
\mbox{}\verb@@\\
\mbox{}\verb@    mov x0, #TOK_IDENT@\\
\mbox{}\verb@    bl _kcc_test_expect_tok@\\
\mbox{}\verb@    cbnz x0, .Lttrig_fail@\\
\mbox{}\verb@@\\
\mbox{}\verb@    mov x0, #TOK_LBRACKET@\\
\mbox{}\verb@    bl _kcc_test_expect_tok@\\
\mbox{}\verb@    cbnz x0, .Lttrig_fail@\\
\mbox{}\verb@@\\
\mbox{}\verb@    mov x0, #TOK_RBRACKET@\\
\mbox{}\verb@    bl _kcc_test_expect_tok@\\
\mbox{}\verb@    cbnz x0, .Lttrig_fail@\\
\mbox{}\verb@@\\
\mbox{}\verb@    mov x0, #TOK_LBRACE@\\
\mbox{}\verb@    bl _kcc_test_expect_tok@\\
\mbox{}\verb@    cbnz x0, .Lttrig_fail@\\
\mbox{}\verb@@\\
\mbox{}\verb@    mov x0, #TOK_RBRACE@\\
\mbox{}\verb@    bl _kcc_test_expect_tok@\\
\mbox{}\verb@    cbnz x0, .Lttrig_fail@\\
\mbox{}\verb@@\\
\mbox{}\verb@    mov x0, #TOK_EOF@\\
\mbox{}\verb@    bl _kcc_test_expect_tok@\\
\mbox{}\verb@    cbnz x0, .Lttrig_fail@\\
\mbox{}\verb@@\\
\mbox{}\verb@    mov x0, #0@\\
\mbox{}\verb@    ldp x29, x30, [sp], #16@\\
\mbox{}\verb@    ret@\\
\mbox{}\verb@.Lttrig_fail:@\\
\mbox{}\verb@    mov x0, #1@\\
\mbox{}\verb@    ldp x29, x30, [sp], #16@\\
\mbox{}\verb@    ret@\\
\mbox{}\verb@@{\NWsep}
\end{list}
\vspace{-1.5ex}
\footnotesize
\begin{list}{}{\setlength{\itemsep}{-\parsep}\setlength{\itemindent}{-\leftmargin}}
\item \NWtxtMacroRefIn\ \NWlink{weft?}{?}.

\item{}
\end{list}
\end{minipage}\vspace{4ex}
\end{flushleft}
\section{Line Splicing}

Tests that backslash-newline sequences are removed before tokenization.
The source \texttt{in\textbackslash\textbackslash{}n\_t x} (with a physical newline
after the backslash) becomes \texttt{int x} after splicing, so
\texttt{int} is recognized as a keyword.

\begin{flushleft} \small
\begin{minipage}{\linewidth}\label{scrap82}\raggedright\small
\NWtarget{weft?}{} $\langle\,${\itshape Test: lex line splice}\nobreak\ {\footnotesize {?}}$\,\rangle\equiv$
\vspace{-1ex}
\begin{list}{}{} \item
\mbox{}\verb@@\\
\mbox{}\verb@// -- test_lex_line_splice -----------------------------------------------@\\
\mbox{}\verb@// Source: "in\<newline>t x" (stored as bytes: i n \ 0x0A t space x)@\\
\mbox{}\verb@// After line splicing: "int x"@\\
\mbox{}\verb@// Expected: INT IDENT(x) EOF@\\
\mbox{}\verb@.text@\\
\mbox{}\verb@.globl _test_lex_line_splice@\\
\mbox{}\verb@.p2align 2@\\
\mbox{}\verb@_test_lex_line_splice:@\\
\mbox{}\verb@    stp x29, x30, [sp, #-16]!@\\
\mbox{}\verb@    mov x29, sp@\\
\mbox{}\verb@@\\
\mbox{}\verb@    adrp x0, _test_src_splice@{\tt @}\verb@PAGE@\\
\mbox{}\verb@    add  x0, x0, _test_src_splice@{\tt @}\verb@PAGEOFF@\\
\mbox{}\verb@    mov  x1, #TEST_SRC_SPLICE_LEN@\\
\mbox{}\verb@    bl _kcc_test_lex_setup@\\
\mbox{}\verb@@\\
\mbox{}\verb@    mov x0, #TOK_INT@\\
\mbox{}\verb@    bl _kcc_test_expect_tok@\\
\mbox{}\verb@    cbnz x0, .Ltspl_fail@\\
\mbox{}\verb@@\\
\mbox{}\verb@    mov x0, #TOK_IDENT@\\
\mbox{}\verb@    bl _kcc_test_expect_tok@\\
\mbox{}\verb@    cbnz x0, .Ltspl_fail@\\
\mbox{}\verb@@\\
\mbox{}\verb@    mov x0, #TOK_EOF@\\
\mbox{}\verb@    bl _kcc_test_expect_tok@\\
\mbox{}\verb@    cbnz x0, .Ltspl_fail@\\
\mbox{}\verb@@\\
\mbox{}\verb@    mov x0, #0@\\
\mbox{}\verb@    ldp x29, x30, [sp], #16@\\
\mbox{}\verb@    ret@\\
\mbox{}\verb@.Ltspl_fail:@\\
\mbox{}\verb@    mov x0, #1@\\
\mbox{}\verb@    ldp x29, x30, [sp], #16@\\
\mbox{}\verb@    ret@\\
\mbox{}\verb@@{\NWsep}
\end{list}
\vspace{-1.5ex}
\footnotesize
\begin{list}{}{\setlength{\itemsep}{-\parsep}\setlength{\itemindent}{-\leftmargin}}
\item \NWtxtMacroRefIn\ \NWlink{weft?}{?}.

\item{}
\end{list}
\end{minipage}\vspace{4ex}
\end{flushleft}
\section{Test Registration Table}

The test table is a flat array of (name, function-pointer) pairs,
terminated by a null sentinel.  The test runner in
Chapter~\ref{ch:test-harness} iterates through this table.

\begin{flushleft} \small
\begin{minipage}{\linewidth}\label{scrap83}\raggedright\small
\NWtarget{weft?}{} $\langle\,${\itshape Test: registration}\nobreak\ {\footnotesize {?}}$\,\rangle\equiv$
\vspace{-1ex}
\begin{list}{}{} \item
\mbox{}\verb@@\\
\mbox{}\verb@// -- Test registration table --------------------------------------------@\\
\mbox{}\verb@.data@\\
\mbox{}\verb@.p2align 3@\\
\mbox{}\verb@.globl _kcc_test_table@\\
\mbox{}\verb@_kcc_test_table:@\\
\mbox{}\verb@    .quad _test_name_keywords@\\
\mbox{}\verb@    .quad _test_lex_keywords@\\
\mbox{}\verb@    .quad _test_name_numbers@\\
\mbox{}\verb@    .quad _test_lex_numbers@\\
\mbox{}\verb@    .quad _test_name_strings@\\
\mbox{}\verb@    .quad _test_lex_strings@\\
\mbox{}\verb@    .quad _test_name_punctuators@\\
\mbox{}\verb@    .quad _test_lex_punctuators@\\
\mbox{}\verb@    .quad _test_name_comments@\\
\mbox{}\verb@    .quad _test_lex_comments@\\
\mbox{}\verb@    .quad _test_name_trigraphs@\\
\mbox{}\verb@    .quad _test_lex_trigraphs@\\
\mbox{}\verb@    .quad _test_name_line_splice@\\
\mbox{}\verb@    .quad _test_lex_line_splice@\\
\mbox{}\verb@    .quad 0                       // sentinel: null name@\\
\mbox{}\verb@    .quad 0@\\
\mbox{}\verb@@{\NWsep}
\end{list}
\vspace{-1.5ex}
\footnotesize
\begin{list}{}{\setlength{\itemsep}{-\parsep}\setlength{\itemindent}{-\leftmargin}}
\item \NWtxtMacroRefIn\ \NWlink{weft?}{?}.

\item{}
\end{list}
\end{minipage}\vspace{4ex}
\end{flushleft}
\section{Test Source Data}

Source strings and test names used by the conformance tests above.
Each string has a companion \texttt{.equ} that computes its length
(excluding the null terminator) for passing to \texttt{kcc\_lex\_init}.

\begin{flushleft} \small
\begin{minipage}{\linewidth}\label{scrap84}\raggedright\small
\NWtarget{weft?}{} $\langle\,${\itshape Test: lex data}\nobreak\ {\footnotesize {?}}$\,\rangle\equiv$
\vspace{-1ex}
\begin{list}{}{} \item
\mbox{}\verb@@\\
\mbox{}\verb@// -- Test source strings ------------------------------------------------@\\
\mbox{}\verb@.section __TEXT,__cstring,cstring_literals@\\
\mbox{}\verb@@\\
\mbox{}\verb@_test_src_kw:       .asciz "auto int return _Static_assert foo"@\\
\mbox{}\verb@_test_src_kw_end:@\\
\mbox{}\verb@.equ TEST_SRC_KW_LEN, _test_src_kw_end - _test_src_kw - 1@\\
\mbox{}\verb@@\\
\mbox{}\verb@_test_src_num:      .asciz "42 0xFF 3.14 1e10 0x1p2"@\\
\mbox{}\verb@_test_src_num_end:@\\
\mbox{}\verb@.equ TEST_SRC_NUM_LEN, _test_src_num_end - _test_src_num - 1@\\
\mbox{}\verb@@\\
\mbox{}\verb@_test_src_str:@\\
\mbox{}\verb@    .byte '"'@\\
\mbox{}\verb@    .ascii "hello"@\\
\mbox{}\verb@    .byte '"'@\\
\mbox{}\verb@    .ascii " "@\\
\mbox{}\verb@    .byte '\''@\\
\mbox{}\verb@    .ascii "x"@\\
\mbox{}\verb@    .byte '\''@\\
\mbox{}\verb@    .byte 0@\\
\mbox{}\verb@_test_src_str_end:@\\
\mbox{}\verb@.equ TEST_SRC_STR_LEN, _test_src_str_end - _test_src_str - 1@\\
\mbox{}\verb@@\\
\mbox{}\verb@_test_src_punct:    .asciz "(){};,+-*/"@\\
\mbox{}\verb@_test_src_punct_end:@\\
\mbox{}\verb@.equ TEST_SRC_PUNCT_LEN, _test_src_punct_end - _test_src_punct - 1@\\
\mbox{}\verb@@\\
\mbox{}\verb@_test_src_comments:@\\
\mbox{}\verb@    .ascii "x /* comment */ y // line"@\\
\mbox{}\verb@    .byte 0x0A@\\
\mbox{}\verb@    .ascii "z"@\\
\mbox{}\verb@    .byte 0@\\
\mbox{}\verb@_test_src_comments_end:@\\
\mbox{}\verb@.equ TEST_SRC_COMMENTS_LEN, _test_src_comments_end - _test_src_comments - 1@\\
\mbox{}\verb@@\\
\mbox{}\verb@_test_src_trigraph: .asciz "??= x ??( ??) ??< ??>"@\\
\mbox{}\verb@_test_src_trigraph_end:@\\
\mbox{}\verb@.equ TEST_SRC_TRIGRAPH_LEN, _test_src_trigraph_end - _test_src_trigraph - 1@\\
\mbox{}\verb@@\\
\mbox{}\verb@_test_src_splice:@\\
\mbox{}\verb@    .ascii "in"@\\
\mbox{}\verb@    .byte '\\'@\\
\mbox{}\verb@    .byte 0x0A@\\
\mbox{}\verb@    .ascii "t x"@\\
\mbox{}\verb@    .byte 0@\\
\mbox{}\verb@_test_src_splice_end:@\\
\mbox{}\verb@.equ TEST_SRC_SPLICE_LEN, _test_src_splice_end - _test_src_splice - 1@\\
\mbox{}\verb@@\\
\mbox{}\verb@// Test name strings@\\
\mbox{}\verb@_test_name_keywords:     .asciz "lex_keywords"@\\
\mbox{}\verb@_test_name_numbers:      .asciz "lex_numbers"@\\
\mbox{}\verb@_test_name_strings:      .asciz "lex_strings"@\\
\mbox{}\verb@_test_name_punctuators:  .asciz "lex_punctuators"@\\
\mbox{}\verb@_test_name_comments:     .asciz "lex_comments"@\\
\mbox{}\verb@_test_name_trigraphs:    .asciz "lex_trigraphs"@\\
\mbox{}\verb@_test_name_line_splice:  .asciz "lex_line_splice"@\\
\mbox{}\verb@@{\NWsep}
\end{list}
\vspace{-1.5ex}
\footnotesize
\begin{list}{}{\setlength{\itemsep}{-\parsep}\setlength{\itemindent}{-\leftmargin}}
\item \NWtxtMacroRefIn\ \NWlink{weft?}{?}.

\item{}
\end{list}
\end{minipage}\vspace{4ex}
\end{flushleft}
% ═══════════════════════════════════════════════════════════════════
% Output Routing
% ═══════════════════════════════════════════════════════════════════
% assembly.weft — Output Routing
%
% This file contains all output directives for kcc.  It is the single
% point that maps logical code fragments to physical assembly files.
%
% When adding a new .s file to kcc, add its output directive HERE.
% Fragment definitions go in the appropriate chapter .weft file.

\chapter*{Output Routing}
\addcontentsline{toc}{chapter}{Output Routing}

The following output directives route code fragments to their
physical assembly files under \texttt{tmp/src/}.

\bigskip

\noindent\textbf{main.s} --- the compiler entry point:

\begin{flushleft} \small
\begin{minipage}{\linewidth}\label{scrap85}\raggedright\small
\NWtarget{weft?}{} \verb@"tmp/src/main.s"@\nobreak\ {\footnotesize {?}}$\equiv$
\vspace{-1ex}
\begin{list}{}{} \item
\mbox{}\verb@@\\
\mbox{}\verb@@\hbox{$\langle\,${\itshape ARM64: common macros}\nobreak\ {\footnotesize \NWlink{weft14}{14}}$\,\rangle$}\verb@@\\
\mbox{}\verb@@\hbox{$\langle\,${\itshape Main: entry point}\nobreak\ {\footnotesize \NWlink{weft?}{?}}$\,\rangle$}\verb@@\\
\mbox{}\verb@@\hbox{$\langle\,${\itshape Main: data}\nobreak\ {\footnotesize \NWlink{weft?}{?}}$\,\rangle$}\verb@@\\
\mbox{}\verb@@{\NWsep}
\end{list}
\vspace{-1.5ex}
\footnotesize
\begin{list}{}{\setlength{\itemsep}{-\parsep}\setlength{\itemindent}{-\leftmargin}}

\item{}
\end{list}
\end{minipage}\vspace{4ex}
\end{flushleft}
\noindent\textbf{syscalls.s} --- macOS system call wrappers:

\begin{flushleft} \small
\begin{minipage}{\linewidth}\label{scrap86}\raggedright\small
\NWtarget{weft?}{} \verb@"tmp/src/syscalls.s"@\nobreak\ {\footnotesize {?}}$\equiv$
\vspace{-1ex}
\begin{list}{}{} \item
\mbox{}\verb@@\\
\mbox{}\verb@@\hbox{$\langle\,${\itshape ARM64: common macros}\nobreak\ {\footnotesize \NWlink{weft14}{14}}$\,\rangle$}\verb@@\\
\mbox{}\verb@@\hbox{$\langle\,${\itshape Syscall: constants}\nobreak\ {\footnotesize \NWlink{weft15a}{15a}}$\,\rangle$}\verb@@\\
\mbox{}\verb@@\hbox{$\langle\,${\itshape Syscall: exit}\nobreak\ {\footnotesize \NWlink{weft15b}{15b}}$\,\rangle$}\verb@@\\
\mbox{}\verb@@\hbox{$\langle\,${\itshape Syscall: write}\nobreak\ {\footnotesize \NWlink{weft16a}{16a}}$\,\rangle$}\verb@@\\
\mbox{}\verb@@\hbox{$\langle\,${\itshape Syscall: read}\nobreak\ {\footnotesize \NWlink{weft16b}{16b}}$\,\rangle$}\verb@@\\
\mbox{}\verb@@\hbox{$\langle\,${\itshape Syscall: open}\nobreak\ {\footnotesize \NWlink{weft17a}{17a}}$\,\rangle$}\verb@@\\
\mbox{}\verb@@\hbox{$\langle\,${\itshape Syscall: close}\nobreak\ {\footnotesize \NWlink{weft17b}{17b}}$\,\rangle$}\verb@@\\
\mbox{}\verb@@\hbox{$\langle\,${\itshape Syscall: mmap}\nobreak\ {\footnotesize \NWlink{weft18}{18}}$\,\rangle$}\verb@@\\
\mbox{}\verb@@\hbox{$\langle\,${\itshape Syscall: lseek}\nobreak\ {\footnotesize \NWlink{weft19}{19}}$\,\rangle$}\verb@@\\
\mbox{}\verb@@{\NWsep}
\end{list}
\vspace{-1.5ex}
\footnotesize
\begin{list}{}{\setlength{\itemsep}{-\parsep}\setlength{\itemindent}{-\leftmargin}}

\item{}
\end{list}
\end{minipage}\vspace{4ex}
\end{flushleft}
\noindent\textbf{arena.s} --- arena memory allocator:

\begin{flushleft} \small
\begin{minipage}{\linewidth}\label{scrap87}\raggedright\small
\NWtarget{weft?}{} \verb@"tmp/src/arena.s"@\nobreak\ {\footnotesize {?}}$\equiv$
\vspace{-1ex}
\begin{list}{}{} \item
\mbox{}\verb@@\\
\mbox{}\verb@@\hbox{$\langle\,${\itshape ARM64: common macros}\nobreak\ {\footnotesize \NWlink{weft14}{14}}$\,\rangle$}\verb@@\\
\mbox{}\verb@@\hbox{$\langle\,${\itshape Syscall: constants}\nobreak\ {\footnotesize \NWlink{weft15a}{15a}}$\,\rangle$}\verb@@\\
\mbox{}\verb@@\hbox{$\langle\,${\itshape Arena: constants}\nobreak\ {\footnotesize \NWlink{weft20a}{20a}}$\,\rangle$}\verb@@\\
\mbox{}\verb@@\hbox{$\langle\,${\itshape Arena: init}\nobreak\ {\footnotesize \NWlink{weft20b}{20b}}$\,\rangle$}\verb@@\\
\mbox{}\verb@@\hbox{$\langle\,${\itshape Arena: alloc fast path}\nobreak\ {\footnotesize \NWlink{weft23}{23}}$\,\rangle$}\verb@@\\
\mbox{}\verb@@\hbox{$\langle\,${\itshape Arena: alloc slow path}\nobreak\ {\footnotesize \NWlink{weft24a}{24a}}$\,\rangle$}\verb@@\\
\mbox{}\verb@@\hbox{$\langle\,${\itshape Arena: data}\nobreak\ {\footnotesize \NWlink{weft25}{25}}$\,\rangle$}\verb@@\\
\mbox{}\verb@@{\NWsep}
\end{list}
\vspace{-1.5ex}
\footnotesize
\begin{list}{}{\setlength{\itemsep}{-\parsep}\setlength{\itemindent}{-\leftmargin}}

\item{}
\end{list}
\end{minipage}\vspace{4ex}
\end{flushleft}
\noindent\textbf{strings.s} --- string and memory operations:

\begin{flushleft} \small
\begin{minipage}{\linewidth}\label{scrap88}\raggedright\small
\NWtarget{weft?}{} \verb@"tmp/src/strings.s"@\nobreak\ {\footnotesize {?}}$\equiv$
\vspace{-1ex}
\begin{list}{}{} \item
\mbox{}\verb@@\\
\mbox{}\verb@@\hbox{$\langle\,${\itshape ARM64: common macros}\nobreak\ {\footnotesize \NWlink{weft14}{14}}$\,\rangle$}\verb@@\\
\mbox{}\verb@@\hbox{$\langle\,${\itshape String: strlen}\nobreak\ {\footnotesize \NWlink{weft40a}{40a}}$\,\rangle$}\verb@@\\
\mbox{}\verb@@\hbox{$\langle\,${\itshape String: memcpy}\nobreak\ {\footnotesize \NWlink{weft40b}{40b}}$\,\rangle$}\verb@@\\
\mbox{}\verb@@\hbox{$\langle\,${\itshape String: memcmp}\nobreak\ {\footnotesize \NWlink{weft?}{?}}$\,\rangle$}\verb@@\\
\mbox{}\verb@@\hbox{$\langle\,${\itshape String: memset}\nobreak\ {\footnotesize \NWlink{weft?}{?}}$\,\rangle$}\verb@@\\
\mbox{}\verb@@\hbox{$\langle\,${\itshape String: strcmp}\nobreak\ {\footnotesize \NWlink{weft?}{?}}$\,\rangle$}\verb@@\\
\mbox{}\verb@@{\NWsep}
\end{list}
\vspace{-1.5ex}
\footnotesize
\begin{list}{}{\setlength{\itemsep}{-\parsep}\setlength{\itemindent}{-\leftmargin}}

\item{}
\end{list}
\end{minipage}\vspace{4ex}
\end{flushleft}
\noindent\textbf{io.s} --- buffered file I/O:

\begin{flushleft} \small
\begin{minipage}{\linewidth}\label{scrap89}\raggedright\small
\NWtarget{weft?}{} \verb@"tmp/src/io.s"@\nobreak\ {\footnotesize {?}}$\equiv$
\vspace{-1ex}
\begin{list}{}{} \item
\mbox{}\verb@@\\
\mbox{}\verb@@\hbox{$\langle\,${\itshape ARM64: common macros}\nobreak\ {\footnotesize \NWlink{weft14}{14}}$\,\rangle$}\verb@@\\
\mbox{}\verb@@\hbox{$\langle\,${\itshape IO: flush}\nobreak\ {\footnotesize \NWlink{weft?}{?}}$\,\rangle$}\verb@@\\
\mbox{}\verb@@\hbox{$\langle\,${\itshape IO: putc}\nobreak\ {\footnotesize \NWlink{weft?}{?}}$\,\rangle$}\verb@@\\
\mbox{}\verb@@\hbox{$\langle\,${\itshape IO: puts}\nobreak\ {\footnotesize \NWlink{weft?}{?}}$\,\rangle$}\verb@@\\
\mbox{}\verb@@\hbox{$\langle\,${\itshape IO: write}\nobreak\ {\footnotesize \NWlink{weft?}{?}}$\,\rangle$}\verb@@\\
\mbox{}\verb@@\hbox{$\langle\,${\itshape IO: putu}\nobreak\ {\footnotesize \NWlink{weft?}{?}}$\,\rangle$}\verb@@\\
\mbox{}\verb@@\hbox{$\langle\,${\itshape IO: err}\nobreak\ {\footnotesize \NWlink{weft?}{?}}$\,\rangle$}\verb@@\\
\mbox{}\verb@@\hbox{$\langle\,${\itshape IO: data}\nobreak\ {\footnotesize \NWlink{weft?}{?}}$\,\rangle$}\verb@@\\
\mbox{}\verb@@{\NWsep}
\end{list}
\vspace{-1.5ex}
\footnotesize
\begin{list}{}{\setlength{\itemsep}{-\parsep}\setlength{\itemindent}{-\leftmargin}}

\item{}
\end{list}
\end{minipage}\vspace{4ex}
\end{flushleft}
\noindent\textbf{hashtab.s} --- hash table (FNV-1a, open addressing):

\begin{flushleft} \small
\begin{minipage}{\linewidth}\label{scrap90}\raggedright\small
\NWtarget{weft?}{} \verb@"tmp/src/hashtab.s"@\nobreak\ {\footnotesize {?}}$\equiv$
\vspace{-1ex}
\begin{list}{}{} \item
\mbox{}\verb@@\\
\mbox{}\verb@@\hbox{$\langle\,${\itshape ARM64: common macros}\nobreak\ {\footnotesize \NWlink{weft14}{14}}$\,\rangle$}\verb@@\\
\mbox{}\verb@@\hbox{$\langle\,${\itshape Hash: constants}\nobreak\ {\footnotesize \NWlink{weft27a}{27a}}$\,\rangle$}\verb@@\\
\mbox{}\verb@@\hbox{$\langle\,${\itshape Hash: FNV-1a}\nobreak\ {\footnotesize \NWlink{weft26}{26}}$\,\rangle$}\verb@@\\
\mbox{}\verb@@\hbox{$\langle\,${\itshape Hash: init}\nobreak\ {\footnotesize \NWlink{weft27b}{27b}}$\,\rangle$}\verb@@\\
\mbox{}\verb@@\hbox{$\langle\,${\itshape Hash: get}\nobreak\ {\footnotesize \NWlink{weft29}{29}}$\,\rangle$}\verb@@\\
\mbox{}\verb@@\hbox{$\langle\,${\itshape Hash: put}\nobreak\ {\footnotesize \NWlink{weft30}{30}}$\,\rangle$}\verb@@\\
\mbox{}\verb@@{\NWsep}
\end{list}
\vspace{-1.5ex}
\footnotesize
\begin{list}{}{\setlength{\itemsep}{-\parsep}\setlength{\itemindent}{-\leftmargin}}

\item{}
\end{list}
\end{minipage}\vspace{4ex}
\end{flushleft}
\noindent\textbf{vec.s} --- dynamic vector:

\begin{flushleft} \small
\begin{minipage}{\linewidth}\label{scrap91}\raggedright\small
\NWtarget{weft?}{} \verb@"tmp/src/vec.s"@\nobreak\ {\footnotesize {?}}$\equiv$
\vspace{-1ex}
\begin{list}{}{} \item
\mbox{}\verb@@\\
\mbox{}\verb@@\hbox{$\langle\,${\itshape ARM64: common macros}\nobreak\ {\footnotesize \NWlink{weft14}{14}}$\,\rangle$}\verb@@\\
\mbox{}\verb@@\hbox{$\langle\,${\itshape Vec: constants}\nobreak\ {\footnotesize \NWlink{weft31}{31}}$\,\rangle$}\verb@@\\
\mbox{}\verb@@\hbox{$\langle\,${\itshape Vec: init}\nobreak\ {\footnotesize \NWlink{weft32}{32}}$\,\rangle$}\verb@@\\
\mbox{}\verb@@\hbox{$\langle\,${\itshape Vec: push}\nobreak\ {\footnotesize \NWlink{weft36}{36}}$\,\rangle$}\verb@@\\
\mbox{}\verb@@\hbox{$\langle\,${\itshape Vec: get}\nobreak\ {\footnotesize \NWlink{weft37}{37}}$\,\rangle$}\verb@@\\
\mbox{}\verb@@{\NWsep}
\end{list}
\vspace{-1.5ex}
\footnotesize
\begin{list}{}{\setlength{\itemsep}{-\parsep}\setlength{\itemindent}{-\leftmargin}}

\item{}
\end{list}
\end{minipage}\vspace{4ex}
\end{flushleft}
\noindent\textbf{strtab.s} --- string interning:

\begin{flushleft} \small
\begin{minipage}{\linewidth}\label{scrap92}\raggedright\small
\NWtarget{weft?}{} \verb@"tmp/src/strtab.s"@\nobreak\ {\footnotesize {?}}$\equiv$
\vspace{-1ex}
\begin{list}{}{} \item
\mbox{}\verb@@\\
\mbox{}\verb@@\hbox{$\langle\,${\itshape ARM64: common macros}\nobreak\ {\footnotesize \NWlink{weft14}{14}}$\,\rangle$}\verb@@\\
\mbox{}\verb@@\hbox{$\langle\,${\itshape Strtab: init}\nobreak\ {\footnotesize \NWlink{weft39a}{39a}}$\,\rangle$}\verb@@\\
\mbox{}\verb@@\hbox{$\langle\,${\itshape Strtab: intern}\nobreak\ {\footnotesize \NWlink{weft39b}{39b}}$\,\rangle$}\verb@@\\
\mbox{}\verb@@\hbox{$\langle\,${\itshape Strtab: data}\nobreak\ {\footnotesize \NWlink{weft39c}{39c}}$\,\rangle$}\verb@@\\
\mbox{}\verb@@{\NWsep}
\end{list}
\vspace{-1.5ex}
\footnotesize
\begin{list}{}{\setlength{\itemsep}{-\parsep}\setlength{\itemindent}{-\leftmargin}}

\item{}
\end{list}
\end{minipage}\vspace{4ex}
\end{flushleft}
\noindent\textbf{diag.s} --- diagnostic engine:

\begin{flushleft} \small
\begin{minipage}{\linewidth}\label{scrap93}\raggedright\small
\NWtarget{weft?}{} \verb@"tmp/src/diag.s"@\nobreak\ {\footnotesize {?}}$\equiv$
\vspace{-1ex}
\begin{list}{}{} \item
\mbox{}\verb@@\\
\mbox{}\verb@@\hbox{$\langle\,${\itshape ARM64: common macros}\nobreak\ {\footnotesize \NWlink{weft14}{14}}$\,\rangle$}\verb@@\\
\mbox{}\verb@@\hbox{$\langle\,${\itshape Diag: constants}\nobreak\ {\footnotesize \NWlink{weft?}{?}}$\,\rangle$}\verb@@\\
\mbox{}\verb@@\hbox{$\langle\,${\itshape Diag: putu stderr}\nobreak\ {\footnotesize \NWlink{weft?}{?}}$\,\rangle$}\verb@@\\
\mbox{}\verb@@\hbox{$\langle\,${\itshape Diag: emit}\nobreak\ {\footnotesize \NWlink{weft?}{?}}$\,\rangle$}\verb@@\\
\mbox{}\verb@@\hbox{$\langle\,${\itshape Diag: error}\nobreak\ {\footnotesize \NWlink{weft?}{?}}$\,\rangle$}\verb@@\\
\mbox{}\verb@@\hbox{$\langle\,${\itshape Diag: warning}\nobreak\ {\footnotesize \NWlink{weft?}{?}}$\,\rangle$}\verb@@\\
\mbox{}\verb@@\hbox{$\langle\,${\itshape Diag: note}\nobreak\ {\footnotesize \NWlink{weft?}{?}}$\,\rangle$}\verb@@\\
\mbox{}\verb@@\hbox{$\langle\,${\itshape Diag: had errors}\nobreak\ {\footnotesize \NWlink{weft?}{?}}$\,\rangle$}\verb@@\\
\mbox{}\verb@@\hbox{$\langle\,${\itshape Diag: reset}\nobreak\ {\footnotesize \NWlink{weft?}{?}}$\,\rangle$}\verb@@\\
\mbox{}\verb@@\hbox{$\langle\,${\itshape Diag: data}\nobreak\ {\footnotesize \NWlink{weft?}{?}}$\,\rangle$}\verb@@\\
\mbox{}\verb@@{\NWsep}
\end{list}
\vspace{-1.5ex}
\footnotesize
\begin{list}{}{\setlength{\itemsep}{-\parsep}\setlength{\itemindent}{-\leftmargin}}

\item{}
\end{list}
\end{minipage}\vspace{4ex}
\end{flushleft}
\noindent\textbf{source.s} --- source file management:

\begin{flushleft} \small
\begin{minipage}{\linewidth}\label{scrap94}\raggedright\small
\NWtarget{weft?}{} \verb@"tmp/src/source.s"@\nobreak\ {\footnotesize {?}}$\equiv$
\vspace{-1ex}
\begin{list}{}{} \item
\mbox{}\verb@@\\
\mbox{}\verb@@\hbox{$\langle\,${\itshape ARM64: common macros}\nobreak\ {\footnotesize \NWlink{weft14}{14}}$\,\rangle$}\verb@@\\
\mbox{}\verb@@\hbox{$\langle\,${\itshape Syscall: constants}\nobreak\ {\footnotesize \NWlink{weft15a}{15a}}$\,\rangle$}\verb@@\\
\mbox{}\verb@@\hbox{$\langle\,${\itshape Source: constants}\nobreak\ {\footnotesize \NWlink{weft?}{?}}$\,\rangle$}\verb@@\\
\mbox{}\verb@@\hbox{$\langle\,${\itshape Source: open file}\nobreak\ {\footnotesize \NWlink{weft?}{?}}$\,\rangle$}\verb@@\\
\mbox{}\verb@@\hbox{$\langle\,${\itshape Encoding: translate}\nobreak\ {\footnotesize \NWlink{weft?}{?}}$\,\rangle$}\verb@@\\
\mbox{}\verb@@\hbox{$\langle\,${\itshape Splice: translate}\nobreak\ {\footnotesize \NWlink{weft?}{?}}$\,\rangle$}\verb@@\\
\mbox{}\verb@@\hbox{$\langle\,${\itshape Source: data}\nobreak\ {\footnotesize \NWlink{weft?}{?}}$\,\rangle$}\verb@@\\
\mbox{}\verb@@{\NWsep}
\end{list}
\vspace{-1.5ex}
\footnotesize
\begin{list}{}{\setlength{\itemsep}{-\parsep}\setlength{\itemindent}{-\leftmargin}}

\item{}
\end{list}
\end{minipage}\vspace{4ex}
\end{flushleft}
\noindent\textbf{lexer.s} --- C17 tokenizer:

\begin{flushleft} \small
\begin{minipage}{\linewidth}\label{scrap95}\raggedright\small
\NWtarget{weft?}{} \verb@"tmp/src/lexer.s"@\nobreak\ {\footnotesize {?}}$\equiv$
\vspace{-1ex}
\begin{list}{}{} \item
\mbox{}\verb@@\\
\mbox{}\verb@@\hbox{$\langle\,${\itshape ARM64: common macros}\nobreak\ {\footnotesize \NWlink{weft14}{14}}$\,\rangle$}\verb@@\\
\mbox{}\verb@@\hbox{$\langle\,${\itshape Lexer: token constants}\nobreak\ {\footnotesize \NWlink{weft?}{?}}$\,\rangle$}\verb@@\\
\mbox{}\verb@@\hbox{$\langle\,${\itshape Lexer: init}\nobreak\ {\footnotesize \NWlink{weft?}{?}}$\,\rangle$}\verb@@\\
\mbox{}\verb@@\hbox{$\langle\,${\itshape Lexer: record location}\nobreak\ {\footnotesize \NWlink{weft?}{?}}$\,\rangle$}\verb@@\\
\mbox{}\verb@@\hbox{$\langle\,${\itshape Lexer: next token}\nobreak\ {\footnotesize \NWlink{weft?}{?}}$\,\rangle$}\verb@@\\
\mbox{}\verb@@\hbox{$\langle\,${\itshape Lexer: scan identifier}\nobreak\ {\footnotesize \NWlink{weft?}{?}}$\,\rangle$}\verb@@\\
\mbox{}\verb@@\hbox{$\langle\,${\itshape Lexer: scan number}\nobreak\ {\footnotesize \NWlink{weft?}{?}}$\,\rangle$}\verb@@\\
\mbox{}\verb@@\hbox{$\langle\,${\itshape Lexer: scan string}\nobreak\ {\footnotesize \NWlink{weft?}{?}}$\,\rangle$}\verb@@\\
\mbox{}\verb@@\hbox{$\langle\,${\itshape Lexer: scan char}\nobreak\ {\footnotesize \NWlink{weft?}{?}}$\,\rangle$}\verb@@\\
\mbox{}\verb@@\hbox{$\langle\,${\itshape Lexer: scan punctuator}\nobreak\ {\footnotesize \NWlink{weft?}{?}}$\,\rangle$}\verb@@\\
\mbox{}\verb@@\hbox{$\langle\,${\itshape Lexer: dump token}\nobreak\ {\footnotesize \NWlink{weft?}{?}}$\,\rangle$}\verb@@\\
\mbox{}\verb@@\hbox{$\langle\,${\itshape Lexer: dump all}\nobreak\ {\footnotesize \NWlink{weft?}{?}}$\,\rangle$}\verb@@\\
\mbox{}\verb@@\hbox{$\langle\,${\itshape Lexer: data}\nobreak\ {\footnotesize \NWlink{weft?}{?}}$\,\rangle$}\verb@@\\
\mbox{}\verb@@{\NWsep}
\end{list}
\vspace{-1.5ex}
\footnotesize
\begin{list}{}{\setlength{\itemsep}{-\parsep}\setlength{\itemindent}{-\leftmargin}}

\item{}
\end{list}
\end{minipage}\vspace{4ex}
\end{flushleft}
\noindent\textbf{lexer\_tables.s} --- keyword table and token names:

\begin{flushleft} \small
\begin{minipage}{\linewidth}\label{scrap96}\raggedright\small
\NWtarget{weft?}{} \verb@"tmp/src/lexer_tables.s"@\nobreak\ {\footnotesize {?}}$\equiv$
\vspace{-1ex}
\begin{list}{}{} \item
\mbox{}\verb@@\\
\mbox{}\verb@@\hbox{$\langle\,${\itshape ARM64: common macros}\nobreak\ {\footnotesize \NWlink{weft14}{14}}$\,\rangle$}\verb@@\\
\mbox{}\verb@@\hbox{$\langle\,${\itshape Lexer: keyword data}\nobreak\ {\footnotesize \NWlink{weft?}{?}}$\,\rangle$}\verb@@\\
\mbox{}\verb@@\hbox{$\langle\,${\itshape Lexer: keyword init}\nobreak\ {\footnotesize \NWlink{weft?}{?}}$\,\rangle$}\verb@@\\
\mbox{}\verb@@\hbox{$\langle\,${\itshape Lexer: token names}\nobreak\ {\footnotesize \NWlink{weft?}{?}}$\,\rangle$}\verb@@\\
\mbox{}\verb@@{\NWsep}
\end{list}
\vspace{-1.5ex}
\footnotesize
\begin{list}{}{\setlength{\itemsep}{-\parsep}\setlength{\itemindent}{-\leftmargin}}

\item{}
\end{list}
\end{minipage}\vspace{4ex}
\end{flushleft}
\bigskip
\noindent\textbf{Test binaries} --- test runner and conformance tests
(assembled separately; linked into \texttt{tmp/test\_kcc}):

\medskip
\noindent\textbf{test\_runner.s} --- test harness entry point:

\begin{flushleft} \small
\begin{minipage}{\linewidth}\label{scrap97}\raggedright\small
\NWtarget{weft?}{} \verb@"tmp/tests/test_runner.s"@\nobreak\ {\footnotesize {?}}$\equiv$
\vspace{-1ex}
\begin{list}{}{} \item
\mbox{}\verb@@\\
\mbox{}\verb@@\hbox{$\langle\,${\itshape ARM64: common macros}\nobreak\ {\footnotesize \NWlink{weft14}{14}}$\,\rangle$}\verb@@\\
\mbox{}\verb@@\hbox{$\langle\,${\itshape Test: runner entry}\nobreak\ {\footnotesize \NWlink{weft?}{?}}$\,\rangle$}\verb@@\\
\mbox{}\verb@@\hbox{$\langle\,${\itshape Test: assert helpers}\nobreak\ {\footnotesize \NWlink{weft?}{?}}$\,\rangle$}\verb@@\\
\mbox{}\verb@@\hbox{$\langle\,${\itshape Test: runner data}\nobreak\ {\footnotesize \NWlink{weft?}{?}}$\,\rangle$}\verb@@\\
\mbox{}\verb@@{\NWsep}
\end{list}
\vspace{-1.5ex}
\footnotesize
\begin{list}{}{\setlength{\itemsep}{-\parsep}\setlength{\itemindent}{-\leftmargin}}

\item{}
\end{list}
\end{minipage}\vspace{4ex}
\end{flushleft}
\noindent\textbf{test\_lexer.s} --- lexer conformance tests:

\begin{flushleft} \small
\begin{minipage}{\linewidth}\label{scrap98}\raggedright\small
\NWtarget{weft?}{} \verb@"tmp/tests/test_lexer.s"@\nobreak\ {\footnotesize {?}}$\equiv$
\vspace{-1ex}
\begin{list}{}{} \item
\mbox{}\verb@@\\
\mbox{}\verb@@\hbox{$\langle\,${\itshape ARM64: common macros}\nobreak\ {\footnotesize \NWlink{weft14}{14}}$\,\rangle$}\verb@@\\
\mbox{}\verb@@\hbox{$\langle\,${\itshape Lexer: token constants}\nobreak\ {\footnotesize \NWlink{weft?}{?}}$\,\rangle$}\verb@@\\
\mbox{}\verb@@\hbox{$\langle\,${\itshape Test: lex keywords}\nobreak\ {\footnotesize \NWlink{weft?}{?}}$\,\rangle$}\verb@@\\
\mbox{}\verb@@\hbox{$\langle\,${\itshape Test: lex numbers}\nobreak\ {\footnotesize \NWlink{weft?}{?}}$\,\rangle$}\verb@@\\
\mbox{}\verb@@\hbox{$\langle\,${\itshape Test: lex strings}\nobreak\ {\footnotesize \NWlink{weft?}{?}}$\,\rangle$}\verb@@\\
\mbox{}\verb@@\hbox{$\langle\,${\itshape Test: lex punctuators}\nobreak\ {\footnotesize \NWlink{weft?}{?}}$\,\rangle$}\verb@@\\
\mbox{}\verb@@\hbox{$\langle\,${\itshape Test: lex comments}\nobreak\ {\footnotesize \NWlink{weft?}{?}}$\,\rangle$}\verb@@\\
\mbox{}\verb@@\hbox{$\langle\,${\itshape Test: lex trigraphs}\nobreak\ {\footnotesize \NWlink{weft?}{?}}$\,\rangle$}\verb@@\\
\mbox{}\verb@@\hbox{$\langle\,${\itshape Test: lex line splice}\nobreak\ {\footnotesize \NWlink{weft?}{?}}$\,\rangle$}\verb@@\\
\mbox{}\verb@@\hbox{$\langle\,${\itshape Test: registration}\nobreak\ {\footnotesize \NWlink{weft?}{?}}$\,\rangle$}\verb@@\\
\mbox{}\verb@@\hbox{$\langle\,${\itshape Test: lex data}\nobreak\ {\footnotesize \NWlink{weft?}{?}}$\,\rangle$}\verb@@\\
\mbox{}\verb@@{\NWsep}
\end{list}
\vspace{-1.5ex}
\footnotesize
\begin{list}{}{\setlength{\itemsep}{-\parsep}\setlength{\itemindent}{-\leftmargin}}

\item{}
\end{list}
\end{minipage}\vspace{4ex}
\end{flushleft}
% ═══════════════════════════════════════════════════════════════════
% Appendices
% ═══════════════════════════════════════════════════════════════════
\appendix

% C-roadmap.weft — Complete phase roadmap for kcc
%
% This appendix documents every phase required to produce a
% fully conforming C17 compiler from bare ARM64 assembly.

\chapter{Development Roadmap}
\label{ch:roadmap}

This appendix enumerates the complete sequence of implementation
phases that take \kcc{} from a bootable skeleton (Phase~0) to a
self-hosting, conforming C17 compiler (Phase~8).  Each phase lists
the \texttt{.weft} files to be created or extended, the new tangled
output files, and the acceptance criteria that must hold before the
phase is considered complete.

% ────────────────────────────────────────────────────────────────
\section{Phase 0 --- Project Skeleton}
\label{sec:phase0}

\textbf{Goal.}  Prove the literate infrastructure works end to end:
\texttt{just build} produces a native ARM64 binary that boots,
prints a banner, and exits cleanly.  No C source is compiled yet.

\textbf{weft files.}
\begin{itemize}
\item \texttt{kcc.weft} --- master file, LaTeX preamble, includes.
\item \texttt{literate/assembly.weft} --- output routing (all
      output directives live here).
\item \texttt{literate/foundations/01-normative-corpus.weft} ---
      standards and citation conventions (prose only).
\item \texttt{literate/foundations/02-conformance-model.weft} ---
      conformance obligations (prose only).
\item \texttt{literate/machine/07-arm64-conventions.weft} ---
      register map, Apple ABI, \texttt{FUNC}/\texttt{ENDFUNC}/\texttt{SYSCALL} macros.
\item \texttt{literate/machine/08-syscalls.weft} ---
      \texttt{exit}, \texttt{write}, \texttt{read}, \texttt{open},
      \texttt{close}, \texttt{mmap} wrappers.
\item \texttt{literate/machine/09-memory-management.weft} ---
      arena allocator (bump pointer over \texttt{mmap} chunks).
\item \texttt{literate/machine/11-string-operations.weft} ---
      \texttt{strlen}, \texttt{memcpy}, \texttt{memcmp}, \texttt{memset}.
\item \texttt{literate/machine/12-io.weft} ---
      buffered file I/O (4\,KiB static buffer, flush/putc/puts).
\item \texttt{literate/driver/43-main.weft} ---
      \texttt{\_main} entry point, argument parsing, banner.
\end{itemize}

\textbf{Tangled output.}
\texttt{main.s}, \texttt{syscalls.s}, \texttt{arena.s},
\texttt{strings.s}, \texttt{io.s}.

\textbf{Acceptance.}
\begin{enumerate}
\item \texttt{just tangle} succeeds with no errors.
\item \texttt{just build} assembles and links with no errors.
\item \texttt{tmp/kcc} prints its banner and exits 0.
\item \texttt{tmp/kcc hello.c} echoes the source file name.
\end{enumerate}

\textbf{Status.} Complete.

% ────────────────────────────────────────────────────────────────
\section{Phase 1 --- Data Structures and Lexer Foundation}
\label{sec:phase1}

\textbf{Goal.}  Read a C source file from disk and tokenize it into
a stream of C17 preprocessing tokens (\stdref{6.4}).

\textbf{weft files.}
\begin{itemize}
\item \texttt{literate/machine/10-data-structures.weft} ---
      hash table (open addressing, FNV-1a) and dynamic vector,
      both arena-allocated.
\item \texttt{literate/machine/14-source-manager.weft} ---
      \texttt{mmap}-based file reading, source location tracking
      (file/line/column as a packed 64-bit value).
\item \texttt{literate/frontend/17-lexer.weft} ---
      character classification tables, tokenizer for all C17
      token classes: identifiers, keywords, integer/floating
      constants, string/char literals, punctuators.
\end{itemize}

\textbf{New tangled files.}
\texttt{hashtab.s}, \texttt{vec.s}, \texttt{strtab.s},
\texttt{source.s}, \texttt{lexer.s}, \texttt{lexer\_tables.s}.

\textbf{Acceptance.}
\begin{enumerate}
\item \texttt{tmp/kcc --dump-tokens hello.c} prints every token
      with its source location.
\item All 44 C17 keywords are recognized.
\item Integer, floating, character, and string literals are
      correctly classified (value parsing deferred to Phase~5).
\item Unterminated strings and unrecognized characters produce
      diagnostics on stderr.
\end{enumerate}

\textbf{Status.} Complete.

% ────────────────────────────────────────────────────────────────
\section{Phase 2 --- Complete Lexer and Diagnostics}
\label{sec:phase2}

\textbf{Goal.}  Finish the tokenizer to full C17 conformance and
build the diagnostic engine required by \stdref{5.1.1.3}.

\textbf{weft files.}
\begin{itemize}
\item \texttt{literate/machine/13-diagnostics.weft} ---
      diagnostic engine: error/warning/note classification,
      source-location printing (``file:line:col: error: ...''),
      error-count tracking, fatal-error limit.
\item \texttt{literate/frontend/15-source-encoding.weft} ---
      Phase~1 of translation: trigraph replacement, source
      character set mapping (\stdref{5.1.1.2} phase 1).
\item \texttt{literate/frontend/16-line-splicing.weft} ---
      Phase~2 of translation: backslash-newline splicing
      (\stdref{5.1.1.2} phase 2).
\item \texttt{literate/verification/44-test-harness.weft} ---
      test runner: discover tests, run each, report pass/fail.
\item \texttt{literate/verification/45-conformance-tests.weft} ---
      initial lexer conformance tests.
\end{itemize}

\textbf{New tangled files.}
\texttt{diag.s}, \texttt{test\_runner.s}, \texttt{test\_lexer.s}.

\textbf{Acceptance.}
\begin{enumerate}
\item Diagnostics follow the format
      \texttt{file:line:col: error: message}.
\item Trigraphs and line splicing are handled before tokenization.
\item \texttt{just test} runs the test suite and reports results.
\item The lexer correctly tokenizes the GCC torture test suite's
      lexer-visible subset (no preprocessing yet).
\end{enumerate}

\textbf{Status.} Complete.

% ────────────────────────────────────────────────────────────────
\section{Phase 3 --- Preprocessor}
\label{sec:phase3}

\textbf{Goal.}  Implement the full C17 preprocessor
(\stdref{6.10}), covering translation phases 3 and 4.

\textbf{weft files.}
\begin{itemize}
\item \texttt{literate/frontend/18-preprocessor.weft} ---
      macro definition (\texttt{\#define}), macro expansion
      (object-like and function-like with stringification and
      token pasting), conditional compilation
      (\texttt{\#if}/\texttt{\#ifdef}/\texttt{\#elif}/\texttt{\#else}/\texttt{\#endif}),
      file inclusion (\texttt{\#include}), line control
      (\texttt{\#line}), error directive (\texttt{\#error}),
      pragma handling (\texttt{\#pragma}, including
      \texttt{\_Pragma} operator), null directive.
\item \texttt{literate/frontend/19-string-literals.weft} ---
      Phase~5: escape sequence processing, execution charset
      mapping.
\item \texttt{literate/frontend/20-concatenation.weft} ---
      Phase~6: adjacent string literal concatenation.
\end{itemize}

\textbf{New tangled files.}
\texttt{preprocessor.s}.

\textbf{Acceptance.}
\begin{enumerate}
\item \texttt{tmp/kcc -E file.c} outputs the preprocessed token
      stream (translation phases 1--6 complete).
\item Recursive and self-referential macros terminate correctly
      per \stdref{6.10.3.4}.
\item \texttt{\#include} searches \texttt{-I} directories and
      the system include path.
\item String concatenation across macro boundaries works.
\item Predefined macros (\texttt{\_\_FILE\_\_},
      \texttt{\_\_LINE\_\_}, \texttt{\_\_DATE\_\_},
      \texttt{\_\_TIME\_\_}, \texttt{\_\_STDC\_\_},
      \texttt{\_\_STDC\_VERSION\_\_}) are set correctly.
\end{enumerate}

% ────────────────────────────────────────────────────────────────
\section{Phase 4 --- Parser}
\label{sec:phase4}

\textbf{Goal.}  Build a recursive-descent parser that consumes
the preprocessed token stream and produces a typed abstract syntax
tree (AST) for every C17 construct.

\textbf{weft files.}
\begin{itemize}
\item \texttt{literate/frontend/21-parser.weft} ---
      parser core: recursive descent infrastructure, AST node
      allocation (arena), AST node types, error recovery
      (synchronization tokens).
\item \texttt{literate/frontend/22-declarations.weft} ---
      declaration parsing: declaration specifiers, declarators
      (including abstract and function), the typedef hack
      (lexer feedback), \texttt{struct}/\texttt{union}/\texttt{enum}
      definitions, \texttt{\_Static\_assert}.
\item \texttt{literate/frontend/23-expressions.weft} ---
      expression parsing: precedence climbing for all C17
      operators (15 precedence levels), postfix operators
      (call, subscript, member access, increment/decrement),
      unary operators, cast expressions, compound literals,
      \texttt{\_Generic} selection, \texttt{sizeof}/\texttt{\_Alignof}.
\item \texttt{literate/frontend/24-statements.weft} ---
      statement parsing: compound, \texttt{if}/\texttt{else},
      \texttt{switch}/\texttt{case}/\texttt{default},
      \texttt{while}, \texttt{do}/\texttt{while}, \texttt{for},
      \texttt{goto}, \texttt{continue}, \texttt{break},
      \texttt{return}, labeled statements, expression statements.
\end{itemize}

\textbf{New tangled files.}
\texttt{parser.s}, \texttt{declarations.s},
\texttt{expressions.s}, \texttt{statements.s}.

\textbf{Acceptance.}
\begin{enumerate}
\item \texttt{tmp/kcc --dump-ast file.c} prints the AST in a
      readable tree format.
\item All syntactically valid C17 programs parse without error.
\item Syntax errors produce useful diagnostics with the
      offending token highlighted.
\item The typedef ambiguity is resolved correctly (e.g.,
      \texttt{T * x;} parsed as declaration when \texttt{T} is
      a typedef name).
\item \texttt{sizeof(int)} and \texttt{sizeof x} both parse.
\end{enumerate}

% ────────────────────────────────────────────────────────────────
\section{Phase 5 --- Type System and Semantic Analysis}
\label{sec:phase5}

\textbf{Goal.}  Implement the C17 type system and semantic
analysis pass that decorates the AST with type information and
checks all constraints.

\textbf{weft files.}
\begin{itemize}
\item \texttt{literate/semantics/25-symbol-table.weft} ---
      four namespaces (\stdref{6.2.3}): ordinary identifiers,
      tags, labels, struct/union members.  Scope management:
      file, function, block, function-prototype scopes.
      Linkage resolution (external, internal, none).
\item \texttt{literate/semantics/26-type-system.weft} ---
      internal type representation: basic types, qualified types,
      pointer, array (fixed, variable-length, incomplete),
      function types, struct/union/enum types.
      Type compatibility (\stdref{6.2.7}) and composite type
      construction.
\item \texttt{literate/semantics/27-type-conversions.weft} ---
      integer promotions (\stdref{6.3.1.1}), usual arithmetic
      conversions (\stdref{6.3.1.8}), implicit conversions,
      explicit casts, pointer conversions.
\item \texttt{literate/semantics/28-type-checking.weft} ---
      constraint checking for every expression operator per
      \stdref{6.5}: operand type requirements, result types,
      lvalue-to-rvalue conversion, array-to-pointer decay,
      function-to-pointer decay.
\item \texttt{literate/semantics/29-initializers.weft} ---
      initializer processing (\stdref{6.7.9}): designated
      initializers, brace elision, string initializers for
      arrays, recursive struct/union/array initialization.
\item \texttt{literate/semantics/30-constant-expressions.weft} ---
      compile-time evaluation (\stdref{6.6}): integer constant
      expressions, arithmetic constant expressions, address
      constants.
\end{itemize}

\textbf{New tangled files.}
\texttt{symtab.s}, \texttt{types.s}, \texttt{conversions.s},
\texttt{sema.s}.

\textbf{Acceptance.}
\begin{enumerate}
\item Type errors produce clear diagnostics citing the
      relevant Standard clause.
\item \texttt{tmp/kcc --dump-types file.c} prints every
      expression with its resolved type.
\item Implicit conversions (integer promotions, UAC, pointer
      decay) are inserted automatically.
\item Struct member access chains resolve correctly through
      anonymous structs/unions.
\item \texttt{enum} constants are typed as \texttt{int} per
      \stdref{6.7.2.2}.
\item Initializer lists with designators and brace elision
      produce the correct initialized object.
\end{enumerate}

% ────────────────────────────────────────────────────────────────
\section{Phase 6 --- Intermediate Representation and Optimization}
\label{sec:phase6}

\textbf{Goal.}  Lower the typed AST into a linear intermediate
representation (IR) with virtual registers, and implement
essential optimization passes.

\textbf{weft files.}
\begin{itemize}
\item \texttt{literate/ir/31-ir-design.weft} ---
      IR design: linear IR with virtual registers, basic blocks,
      control-flow graph, SSA-like properties (each virtual
      register assigned once).  IR instruction set: arithmetic,
      comparison, branch, call, load, store, alloca, phi, cast.
\item \texttt{literate/ir/32-ir-builder.weft} ---
      AST-to-IR lowering: expression evaluation, short-circuit
      boolean operators, \texttt{switch} lowering (jump table or
      if-chain), loop lowering, struct passing/returning.
\item \texttt{literate/ir/33-ir-validation.weft} ---
      IR well-formedness checks: every use dominated by its
      definition, consistent types, correct basic-block
      termination, reachability analysis.
\item \texttt{literate/optimization/34-ir-passes.weft} ---
      pass infrastructure: pass manager, iteration control,
      analysis caching.
\item \texttt{literate/optimization/35-peephole.weft} ---
      algebraic simplifications: constant folding, strength
      reduction, identity elimination (\texttt{x+0}, \texttt{x*1}).
\item \texttt{literate/optimization/36-dce.weft} ---
      dead code elimination: remove instructions whose results
      are never used, remove unreachable basic blocks.
\end{itemize}

\textbf{New tangled files.}
\texttt{ir.s}, \texttt{opt.s}.

\textbf{Acceptance.}
\begin{enumerate}
\item \texttt{tmp/kcc --dump-ir file.c} prints the IR in a
      readable textual format.
\item IR validation passes with no errors for all valid programs.
\item Constant expressions evaluate at compile time.
\item Dead stores and unused computations are eliminated.
\item The IR correctly models all C17 control flow, including
      \texttt{goto}, \texttt{switch} fall-through, and
      \texttt{setjmp}/\texttt{longjmp} (treated as opaque calls).
\end{enumerate}

% ────────────────────────────────────────────────────────────────
\section{Phase 7 --- ARM64 Backend}
\label{sec:phase7}

\textbf{Goal.}  Generate native ARM64 machine code and emit
Mach-O object files directly --- no external assembler needed.
This is the self-referential core: ARM64 assembly generating
ARM64 machine code.

\textbf{weft files.}
\begin{itemize}
\item \texttt{literate/backend/37-abi-aapcs64.weft} ---
      AAPCS64 calling convention implementation: parameter
      classification (stages A, B, C), return value handling,
      HFA/HVA detection, Apple-specific deltas (small struct
      passing, variadic arguments).
\item \texttt{literate/backend/38-isel.weft} ---
      instruction selection: pattern matching from IR
      instructions to ARM64 instruction sequences.  Covering
      arithmetic, comparisons, branches, loads/stores (with
      addressing mode selection), function calls, returns.
\item \texttt{literate/backend/39-regalloc.weft} ---
      register allocation: linear scan allocator, spill code
      generation, callee-saved register assignment, frame
      layout computation.
\item \texttt{literate/backend/40-emit-arm64.weft} ---
      ARM64 instruction encoding: binary encoding of every
      instruction format (data processing, loads/stores,
      branches, system), relocation handling.
\item \texttt{literate/backend/41-emit-macho.weft} ---
      Mach-O object file emission: header, load commands
      (\texttt{LC\_SEGMENT\_64}, \texttt{LC\_SYMTAB},
      \texttt{LC\_DYSYMTAB}), sections (\texttt{\_\_TEXT},
      \texttt{\_\_DATA}, \texttt{\_\_BSS}), symbol table,
      string table, relocation entries.
\item \texttt{literate/driver/42-driver.weft} ---
      full compilation pipeline orchestration: phase sequencing,
      multi-file compilation, object file output.
\end{itemize}

\textbf{New tangled files.}
\texttt{abi.s}, \texttt{isel.s}, \texttt{regalloc.s},
\texttt{emit.s}, \texttt{driver.s}.

\textbf{Acceptance.}
\begin{enumerate}
\item \texttt{tmp/kcc -c file.c -o file.o} produces a valid
      Mach-O object file.
\item \texttt{ld file.o -o file -lSystem ...} links the object
      file into a working executable.
\item Generated code passes the C17 conformance test suite
      (integer arithmetic, control flow, function calls,
      structs, pointers, arrays).
\item Calling convention is correct: \kcc{}-compiled code can
      call system libraries and vice versa.
\item The object file is accepted by Apple's \texttt{ld} linker
      without warnings.
\end{enumerate}

% ────────────────────────────────────────────────────────────────
\section{Phase 8 --- Self-Hosting}
\label{sec:phase8}

\textbf{Goal.}  Demonstrate that \kcc{} can compile a
non-trivial C program --- ideally a C subset of itself or an
equivalent program --- and produce correct output.  This is the
ultimate proof that the compiler is correct and complete.

\textbf{Activities.}
\begin{itemize}
\item Write a C version of a significant \kcc{} subsystem
      (e.g., the lexer or the arena allocator) and compile it
      with \kcc{}.
\item Differential testing: compile the same C source with both
      \kcc{} and \texttt{clang}, then compare the behavior of
      the two executables on a large test corpus.
\item Bootstrap test: if feasible, write a minimal C compiler in
      C (a ``stage-1'' compiler) and compile it with \kcc{}.
      Then use the stage-1 compiler to compile itself.  Compare
      the two stage-1 binaries.
\item Performance characterization: measure compilation speed
      and generated code quality against \texttt{clang -O0} on
      a representative benchmark suite.
\end{itemize}

\textbf{Acceptance.}
\begin{enumerate}
\item At least one non-trivial C program (${}>1000$ lines)
      compiles and runs correctly under \kcc{}.
\item Differential testing against \texttt{clang} reveals no
      behavioral discrepancies on conforming programs.
\item The \stdref{J.3} implementation-defined behavior document
      is complete (Appendix~A).
\item The UB policy matrix is complete (Appendix~B).
\item \texttt{just weave} produces a typeset document that
      serves as a comprehensive textbook on both compiler
      construction and ARM64 systems programming.
\end{enumerate}

% ────────────────────────────────────────────────────────────────
\section{Summary}
\label{sec:roadmap-summary}

\begin{center}
\begin{tabular}{clll}
\hline
Phase & Milestone & Key Output & \texttt{.weft} Files \\
\hline
0 & Skeleton boots        & \texttt{tmp/kcc} runs        & 01, 02, 07--09, 11, 12, 43 \\
1 & Lexer                 & \texttt{--dump-tokens}        & 10, 14, 17 \\
2 & Diagnostics + tests   & \texttt{just test}            & 13, 15, 16, 44, 45 \\
3 & Preprocessor          & \texttt{-E}                   & 18, 19, 20 \\
4 & Parser                & \texttt{--dump-ast}           & 21--24 \\
5 & Type system           & \texttt{--dump-types}         & 25--30 \\
6 & IR + optimization     & \texttt{--dump-ir}            & 31--36 \\
7 & Backend               & \texttt{-c file.c -o file.o}  & 37--42 \\
8 & Self-hosting          & Differential testing          & --- \\
\hline
\end{tabular}
\end{center}

Total: 45 \texttt{.weft} source files, approximately 25 tangled
assembly files, producing a single native \texttt{tmp/kcc} binary
that compiles C17 to ARM64 Mach-O --- with no dependency on any
existing C compiler.

% ═══════════════════════════════════════════════════════════════════
% Indexes
% ═══════════════════════════════════════════════════════════════════

{\small\begin{list}{}{\setlength{\itemsep}{-\parsep}\setlength{\itemindent}{-\leftmargin}}
\item \verb@"tmp/src/arena.s"@ {\footnotesize {\NWtxtDefBy} \NWlink{weft?}{?}.}
\item \verb@"tmp/src/diag.s"@ {\footnotesize {\NWtxtDefBy} \NWlink{weft?}{?}.}
\item \verb@"tmp/src/hashtab.s"@ {\footnotesize {\NWtxtDefBy} \NWlink{weft?}{?}.}
\item \verb@"tmp/src/io.s"@ {\footnotesize {\NWtxtDefBy} \NWlink{weft?}{?}.}
\item \verb@"tmp/src/lexer.s"@ {\footnotesize {\NWtxtDefBy} \NWlink{weft?}{?}.}
\item \verb@"tmp/src/lexer_tables.s"@ {\footnotesize {\NWtxtDefBy} \NWlink{weft?}{?}.}
\item \verb@"tmp/src/main.s"@ {\footnotesize {\NWtxtDefBy} \NWlink{weft?}{?}.}
\item \verb@"tmp/src/source.s"@ {\footnotesize {\NWtxtDefBy} \NWlink{weft?}{?}.}
\item \verb@"tmp/src/strings.s"@ {\footnotesize {\NWtxtDefBy} \NWlink{weft?}{?}.}
\item \verb@"tmp/src/strtab.s"@ {\footnotesize {\NWtxtDefBy} \NWlink{weft?}{?}.}
\item \verb@"tmp/src/syscalls.s"@ {\footnotesize {\NWtxtDefBy} \NWlink{weft?}{?}.}
\item \verb@"tmp/src/vec.s"@ {\footnotesize {\NWtxtDefBy} \NWlink{weft?}{?}.}
\item \verb@"tmp/tests/test_lexer.s"@ {\footnotesize {\NWtxtDefBy} \NWlink{weft?}{?}.}
\item \verb@"tmp/tests/test_runner.s"@ {\footnotesize {\NWtxtDefBy} \NWlink{weft?}{?}.}
\end{list}}

{\small\begin{list}{}{\setlength{\itemsep}{-\parsep}\setlength{\itemindent}{-\leftmargin}}
\item $\langle\,$Arena: alloc fast path\nobreak\ {\footnotesize \NWlink{weft23}{23}}$\,\rangle$ {\footnotesize {\NWtxtRefIn} \NWlink{weft?}{?}.}
\item $\langle\,$Arena: alloc link new chunk\nobreak\ {\footnotesize \NWlink{weft24b}{24b}}$\,\rangle$ {\footnotesize {\NWtxtRefIn} \NWlink{weft24a}{24a}.}
\item $\langle\,$Arena: alloc slow path\nobreak\ {\footnotesize \NWlink{weft24a}{24a}}$\,\rangle$ {\footnotesize {\NWtxtRefIn} \NWlink{weft?}{?}.}
\item $\langle\,$Arena: constants\nobreak\ {\footnotesize \NWlink{weft20a}{20a}}$\,\rangle$ {\footnotesize {\NWtxtRefIn} \NWlink{weft?}{?}.}
\item $\langle\,$Arena: data\nobreak\ {\footnotesize \NWlink{weft25}{25}}$\,\rangle$ {\footnotesize {\NWtxtRefIn} \NWlink{weft?}{?}.}
\item $\langle\,$Arena: init\nobreak\ {\footnotesize \NWlink{weft20b}{20b}}$\,\rangle$ {\footnotesize {\NWtxtRefIn} \NWlink{weft?}{?}.}
\item $\langle\,$Arena: init chunk header\nobreak\ {\footnotesize \NWlink{weft22}{22}}$\,\rangle$ {\footnotesize {\NWtxtRefIn} \NWlink{weft20b}{20b}.}
\item $\langle\,$Arena: init mmap call\nobreak\ {\footnotesize \NWlink{weft21}{21}}$\,\rangle$ {\footnotesize {\NWtxtRefIn} \NWlink{weft20b}{20b}.}
\item $\langle\,$ARM64: common macros\nobreak\ {\footnotesize \NWlink{weft14}{14}}$\,\rangle$ {\footnotesize {\NWtxtRefIn} \NWlink{weft?}{?}\NWlink{weft?}{, ?}\NWlink{weft?}{, ?}\NWlink{weft?}{, ?}\NWlink{weft?}{, ?}\NWlink{weft?}{, ?}\NWlink{weft?}{, ?}\NWlink{weft?}{, ?}\NWlink{weft?}{, ?}\NWlink{weft?}{, ?}\NWlink{weft?}{, ?}\NWlink{weft?}{, ?}\NWlink{weft?}{, ?}\NWlink{weft?}{, ?}.
}
\item $\langle\,$Diag: constants\nobreak\ {\footnotesize \NWlink{weft?}{?}}$\,\rangle$ {\footnotesize {\NWtxtRefIn} \NWlink{weft?}{?}.}
\item $\langle\,$Diag: data\nobreak\ {\footnotesize \NWlink{weft?}{?}}$\,\rangle$ {\footnotesize {\NWtxtRefIn} \NWlink{weft?}{?}.}
\item $\langle\,$Diag: emit\nobreak\ {\footnotesize \NWlink{weft?}{?}}$\,\rangle$ {\footnotesize {\NWtxtRefIn} \NWlink{weft?}{?}.}
\item $\langle\,$Diag: error\nobreak\ {\footnotesize \NWlink{weft?}{?}}$\,\rangle$ {\footnotesize {\NWtxtRefIn} \NWlink{weft?}{?}.}
\item $\langle\,$Diag: had errors\nobreak\ {\footnotesize \NWlink{weft?}{?}}$\,\rangle$ {\footnotesize {\NWtxtRefIn} \NWlink{weft?}{?}.}
\item $\langle\,$Diag: note\nobreak\ {\footnotesize \NWlink{weft?}{?}}$\,\rangle$ {\footnotesize {\NWtxtRefIn} \NWlink{weft?}{?}.}
\item $\langle\,$Diag: putu stderr\nobreak\ {\footnotesize \NWlink{weft?}{?}}$\,\rangle$ {\footnotesize {\NWtxtRefIn} \NWlink{weft?}{?}.}
\item $\langle\,$Diag: reset\nobreak\ {\footnotesize \NWlink{weft?}{?}}$\,\rangle$ {\footnotesize {\NWtxtRefIn} \NWlink{weft?}{?}.}
\item $\langle\,$Diag: warning\nobreak\ {\footnotesize \NWlink{weft?}{?}}$\,\rangle$ {\footnotesize {\NWtxtRefIn} \NWlink{weft?}{?}.}
\item $\langle\,$Encoding: translate\nobreak\ {\footnotesize \NWlink{weft?}{?}}$\,\rangle$ {\footnotesize {\NWtxtRefIn} \NWlink{weft?}{?}.}
\item $\langle\,$Hash: constants\nobreak\ {\footnotesize \NWlink{weft27a}{27a}}$\,\rangle$ {\footnotesize {\NWtxtRefIn} \NWlink{weft?}{?}.}
\item $\langle\,$Hash: FNV-1a\nobreak\ {\footnotesize \NWlink{weft26}{26}}$\,\rangle$ {\footnotesize {\NWtxtRefIn} \NWlink{weft?}{?}.}
\item $\langle\,$Hash: get\nobreak\ {\footnotesize \NWlink{weft29}{29}}$\,\rangle$ {\footnotesize {\NWtxtRefIn} \NWlink{weft?}{?}.}
\item $\langle\,$Hash: init\nobreak\ {\footnotesize \NWlink{weft27b}{27b}}$\,\rangle$ {\footnotesize {\NWtxtRefIn} \NWlink{weft?}{?}.}
\item $\langle\,$Hash: put\nobreak\ {\footnotesize \NWlink{weft30}{30}}$\,\rangle$ {\footnotesize {\NWtxtRefIn} \NWlink{weft?}{?}.}
\item $\langle\,$IO: data\nobreak\ {\footnotesize \NWlink{weft?}{?}}$\,\rangle$ {\footnotesize {\NWtxtRefIn} \NWlink{weft?}{?}.}
\item $\langle\,$IO: err\nobreak\ {\footnotesize \NWlink{weft?}{?}}$\,\rangle$ {\footnotesize {\NWtxtRefIn} \NWlink{weft?}{?}.}
\item $\langle\,$IO: flush\nobreak\ {\footnotesize \NWlink{weft?}{?}}$\,\rangle$ {\footnotesize {\NWtxtRefIn} \NWlink{weft?}{?}.}
\item $\langle\,$IO: putc\nobreak\ {\footnotesize \NWlink{weft?}{?}}$\,\rangle$ {\footnotesize {\NWtxtRefIn} \NWlink{weft?}{?}.}
\item $\langle\,$IO: puts\nobreak\ {\footnotesize \NWlink{weft?}{?}}$\,\rangle$ {\footnotesize {\NWtxtRefIn} \NWlink{weft?}{?}.}
\item $\langle\,$IO: putu\nobreak\ {\footnotesize \NWlink{weft?}{?}}$\,\rangle$ {\footnotesize {\NWtxtRefIn} \NWlink{weft?}{?}.}
\item $\langle\,$IO: write\nobreak\ {\footnotesize \NWlink{weft?}{?}}$\,\rangle$ {\footnotesize {\NWtxtRefIn} \NWlink{weft?}{?}.}
\item $\langle\,$Lexer: data\nobreak\ {\footnotesize \NWlink{weft?}{?}}$\,\rangle$ {\footnotesize {\NWtxtRefIn} \NWlink{weft?}{?}.}
\item $\langle\,$Lexer: dump all\nobreak\ {\footnotesize \NWlink{weft?}{?}}$\,\rangle$ {\footnotesize {\NWtxtRefIn} \NWlink{weft?}{?}.}
\item $\langle\,$Lexer: dump token\nobreak\ {\footnotesize \NWlink{weft?}{?}}$\,\rangle$ {\footnotesize {\NWtxtRefIn} \NWlink{weft?}{?}.}
\item $\langle\,$Lexer: init\nobreak\ {\footnotesize \NWlink{weft?}{?}}$\,\rangle$ {\footnotesize {\NWtxtRefIn} \NWlink{weft?}{?}.}
\item $\langle\,$Lexer: keyword data\nobreak\ {\footnotesize \NWlink{weft?}{?}}$\,\rangle$ {\footnotesize {\NWtxtRefIn} \NWlink{weft?}{?}.}
\item $\langle\,$Lexer: keyword init\nobreak\ {\footnotesize \NWlink{weft?}{?}}$\,\rangle$ {\footnotesize {\NWtxtRefIn} \NWlink{weft?}{?}.}
\item $\langle\,$Lexer: next token\nobreak\ {\footnotesize \NWlink{weft?}{?}}$\,\rangle$ {\footnotesize {\NWtxtRefIn} \NWlink{weft?}{?}.}
\item $\langle\,$Lexer: record location\nobreak\ {\footnotesize \NWlink{weft?}{?}}$\,\rangle$ {\footnotesize {\NWtxtRefIn} \NWlink{weft?}{?}.}
\item $\langle\,$Lexer: scan char\nobreak\ {\footnotesize \NWlink{weft?}{?}}$\,\rangle$ {\footnotesize {\NWtxtRefIn} \NWlink{weft?}{?}.}
\item $\langle\,$Lexer: scan identifier\nobreak\ {\footnotesize \NWlink{weft?}{?}}$\,\rangle$ {\footnotesize {\NWtxtRefIn} \NWlink{weft?}{?}.}
\item $\langle\,$Lexer: scan number\nobreak\ {\footnotesize \NWlink{weft?}{?}}$\,\rangle$ {\footnotesize {\NWtxtRefIn} \NWlink{weft?}{?}.}
\item $\langle\,$Lexer: scan punctuator\nobreak\ {\footnotesize \NWlink{weft?}{?}}$\,\rangle$ {\footnotesize {\NWtxtRefIn} \NWlink{weft?}{?}.}
\item $\langle\,$Lexer: scan string\nobreak\ {\footnotesize \NWlink{weft?}{?}}$\,\rangle$ {\footnotesize {\NWtxtRefIn} \NWlink{weft?}{?}.}
\item $\langle\,$Lexer: token constants\nobreak\ {\footnotesize \NWlink{weft?}{?}}$\,\rangle$ {\footnotesize {\NWtxtRefIn} \NWlink{weft?}{?}\NWlink{weft?}{, ?}.
}
\item $\langle\,$Lexer: token names\nobreak\ {\footnotesize \NWlink{weft?}{?}}$\,\rangle$ {\footnotesize {\NWtxtRefIn} \NWlink{weft?}{?}.}
\item $\langle\,$Main: data\nobreak\ {\footnotesize \NWlink{weft?}{?}}$\,\rangle$ {\footnotesize {\NWtxtRefIn} \NWlink{weft?}{?}.}
\item $\langle\,$Main: entry point\nobreak\ {\footnotesize \NWlink{weft?}{?}}$\,\rangle$ {\footnotesize {\NWtxtRefIn} \NWlink{weft?}{?}.}
\item $\langle\,$Source: constants\nobreak\ {\footnotesize \NWlink{weft?}{?}}$\,\rangle$ {\footnotesize {\NWtxtRefIn} \NWlink{weft?}{?}.}
\item $\langle\,$Source: data\nobreak\ {\footnotesize \NWlink{weft?}{?}}$\,\rangle$ {\footnotesize {\NWtxtRefIn} \NWlink{weft?}{?}.}
\item $\langle\,$Source: open file\nobreak\ {\footnotesize \NWlink{weft?}{?}}$\,\rangle$ {\footnotesize {\NWtxtRefIn} \NWlink{weft?}{?}.}
\item $\langle\,$Splice: translate\nobreak\ {\footnotesize \NWlink{weft?}{?}}$\,\rangle$ {\footnotesize {\NWtxtRefIn} \NWlink{weft?}{?}.}
\item $\langle\,$String: memcmp\nobreak\ {\footnotesize \NWlink{weft?}{?}}$\,\rangle$ {\footnotesize {\NWtxtRefIn} \NWlink{weft?}{?}.}
\item $\langle\,$String: memcpy\nobreak\ {\footnotesize \NWlink{weft40b}{40b}}$\,\rangle$ {\footnotesize {\NWtxtRefIn} \NWlink{weft?}{?}.}
\item $\langle\,$String: memset\nobreak\ {\footnotesize \NWlink{weft?}{?}}$\,\rangle$ {\footnotesize {\NWtxtRefIn} \NWlink{weft?}{?}.}
\item $\langle\,$String: strcmp\nobreak\ {\footnotesize \NWlink{weft?}{?}}$\,\rangle$ {\footnotesize {\NWtxtRefIn} \NWlink{weft?}{?}.}
\item $\langle\,$String: strlen\nobreak\ {\footnotesize \NWlink{weft40a}{40a}}$\,\rangle$ {\footnotesize {\NWtxtRefIn} \NWlink{weft?}{?}.}
\item $\langle\,$Strtab: data\nobreak\ {\footnotesize \NWlink{weft39c}{39c}}$\,\rangle$ {\footnotesize {\NWtxtRefIn} \NWlink{weft?}{?}.}
\item $\langle\,$Strtab: init\nobreak\ {\footnotesize \NWlink{weft39a}{39a}}$\,\rangle$ {\footnotesize {\NWtxtRefIn} \NWlink{weft?}{?}.}
\item $\langle\,$Strtab: intern\nobreak\ {\footnotesize \NWlink{weft39b}{39b}}$\,\rangle$ {\footnotesize {\NWtxtRefIn} \NWlink{weft?}{?}.}
\item $\langle\,$Syscall: close\nobreak\ {\footnotesize \NWlink{weft17b}{17b}}$\,\rangle$ {\footnotesize {\NWtxtRefIn} \NWlink{weft?}{?}.}
\item $\langle\,$Syscall: constants\nobreak\ {\footnotesize \NWlink{weft15a}{15a}}$\,\rangle$ {\footnotesize {\NWtxtRefIn} \NWlink{weft?}{?}\NWlink{weft?}{, ?}\NWlink{weft?}{, ?}.
}
\item $\langle\,$Syscall: exit\nobreak\ {\footnotesize \NWlink{weft15b}{15b}}$\,\rangle$ {\footnotesize {\NWtxtRefIn} \NWlink{weft?}{?}.}
\item $\langle\,$Syscall: lseek\nobreak\ {\footnotesize \NWlink{weft19}{19}}$\,\rangle$ {\footnotesize {\NWtxtRefIn} \NWlink{weft?}{?}.}
\item $\langle\,$Syscall: mmap\nobreak\ {\footnotesize \NWlink{weft18}{18}}$\,\rangle$ {\footnotesize {\NWtxtRefIn} \NWlink{weft?}{?}.}
\item $\langle\,$Syscall: open\nobreak\ {\footnotesize \NWlink{weft17a}{17a}}$\,\rangle$ {\footnotesize {\NWtxtRefIn} \NWlink{weft?}{?}.}
\item $\langle\,$Syscall: read\nobreak\ {\footnotesize \NWlink{weft16b}{16b}}$\,\rangle$ {\footnotesize {\NWtxtRefIn} \NWlink{weft?}{?}.}
\item $\langle\,$Syscall: write\nobreak\ {\footnotesize \NWlink{weft16a}{16a}}$\,\rangle$ {\footnotesize {\NWtxtRefIn} \NWlink{weft?}{?}.}
\item $\langle\,$Test: assert helpers\nobreak\ {\footnotesize \NWlink{weft?}{?}}$\,\rangle$ {\footnotesize {\NWtxtRefIn} \NWlink{weft?}{?}.}
\item $\langle\,$Test: lex comments\nobreak\ {\footnotesize \NWlink{weft?}{?}}$\,\rangle$ {\footnotesize {\NWtxtRefIn} \NWlink{weft?}{?}.}
\item $\langle\,$Test: lex data\nobreak\ {\footnotesize \NWlink{weft?}{?}}$\,\rangle$ {\footnotesize {\NWtxtRefIn} \NWlink{weft?}{?}.}
\item $\langle\,$Test: lex keywords\nobreak\ {\footnotesize \NWlink{weft?}{?}}$\,\rangle$ {\footnotesize {\NWtxtRefIn} \NWlink{weft?}{?}.}
\item $\langle\,$Test: lex line splice\nobreak\ {\footnotesize \NWlink{weft?}{?}}$\,\rangle$ {\footnotesize {\NWtxtRefIn} \NWlink{weft?}{?}.}
\item $\langle\,$Test: lex numbers\nobreak\ {\footnotesize \NWlink{weft?}{?}}$\,\rangle$ {\footnotesize {\NWtxtRefIn} \NWlink{weft?}{?}.}
\item $\langle\,$Test: lex punctuators\nobreak\ {\footnotesize \NWlink{weft?}{?}}$\,\rangle$ {\footnotesize {\NWtxtRefIn} \NWlink{weft?}{?}.}
\item $\langle\,$Test: lex strings\nobreak\ {\footnotesize \NWlink{weft?}{?}}$\,\rangle$ {\footnotesize {\NWtxtRefIn} \NWlink{weft?}{?}.}
\item $\langle\,$Test: lex trigraphs\nobreak\ {\footnotesize \NWlink{weft?}{?}}$\,\rangle$ {\footnotesize {\NWtxtRefIn} \NWlink{weft?}{?}.}
\item $\langle\,$Test: registration\nobreak\ {\footnotesize \NWlink{weft?}{?}}$\,\rangle$ {\footnotesize {\NWtxtRefIn} \NWlink{weft?}{?}.}
\item $\langle\,$Test: runner data\nobreak\ {\footnotesize \NWlink{weft?}{?}}$\,\rangle$ {\footnotesize {\NWtxtRefIn} \NWlink{weft?}{?}.}
\item $\langle\,$Test: runner entry\nobreak\ {\footnotesize \NWlink{weft?}{?}}$\,\rangle$ {\footnotesize {\NWtxtRefIn} \NWlink{weft?}{?}.}
\item $\langle\,$Vec: constants\nobreak\ {\footnotesize \NWlink{weft31}{31}}$\,\rangle$ {\footnotesize {\NWtxtRefIn} \NWlink{weft?}{?}.}
\item $\langle\,$Vec: get\nobreak\ {\footnotesize \NWlink{weft37}{37}}$\,\rangle$ {\footnotesize {\NWtxtRefIn} \NWlink{weft?}{?}.}
\item $\langle\,$Vec: init\nobreak\ {\footnotesize \NWlink{weft32}{32}}$\,\rangle$ {\footnotesize {\NWtxtRefIn} \NWlink{weft?}{?}.}
\item $\langle\,$Vec: push\nobreak\ {\footnotesize \NWlink{weft36}{36}}$\,\rangle$ {\footnotesize {\NWtxtRefIn} \NWlink{weft?}{?}.}
\end{list}}


\end{document}
